\chapter{逐篇核读档案与全量候选索引}\label{ch:catalog}

本章前半部分是已完成正文与关键图表核对的文献档案；后半部分是尚未完成核读或待进一步筛选的候选目录。待核目录的题名、年份与链接是检索信息，不等于对其方法、发现或相关性的确认。

\section{如何反查本地原文}

文献键是稳定检索标识。各档案提供本地PDF链接、DOI或版本入口和页码。PDF阅读器可能限制相对文件链接，可直接在项目的 papers 文件夹打开对应文件；完整路径及全部状态也在 paper\_index.html 和 paper\_index.csv 中。

\section{Ottaviani1993：1993年}\label{paper:Ottaviani1993}

\textbf{Nonlinear collisionless magnetic reconnection}\citep{Ottaviani1993}\

M. Ottaviani;  F. Porcelli.

\textsf{核读范围：}JET-P(93)66 PDF1-16全部，扫描件OCR后核关键公式与图；\textsf{图表：}图1-4，关键式11和图3视觉核验.

\href{run:papers/Ottaviani1993.pdf}{打开本地PDF：\texttt{papers/Ottaviani1993.pdf}}\quad\href{https://doi.org/10.1103/PhysRevLett.71.3802}{DOI／版本入口}

\paragraph{在报告中的研究脉络}\hyperref[story:onset]{启动、爆发与磁岛}。关联表示本报告的阅读线索，不自动表示这些论文互相引用。

\paragraph{实际模型}无耗散、不可压、二维周期冷两场惯性模型，忽略Larmor、密度和压强梯度，F=psi-de\textsuperscript{2}lap psi被输运。

\paragraph{计算／实验细节}psi0=cos x, Lx=pi, epsilon=Lx/Ly=.5, de/(2Lx)=.04，Delta'de约2.03，仅m=1不稳定；伪谱1024×64 Fourier分量；t约125小尺度不可解析而终止。

\paragraph{核实的发现}早期非线性de<w<<Lx内对流胞保持宏观结构，电流子层按delta\textasciitilde{}de exp(-位移/de)缩窄，流速层仍约de；幅度方程A''约A+c2 A\textasciicircum{}4给准爆发增长。

\paragraph{限制与未完成问题}是早期非线性与大Delta'结果；到宏观岛宽时增长变慢的另一计算仅文字报告；不能把有限时形式发散当成物理无限电流；无收敛扫描且细层未解析后主动停止。读取的是投稿前JET版，不冒充期刊最终版逐字核对。

\paragraph{本报告判断}重要新现象是多尺度电流／速度结构与非线性加速，排除了该序列中必然Sweet–Parker式减速；不提供任意参数的最终饱和／热化预测。

\section{Schep1994：1994年}\label{paper:Schep1994}

\textbf{Generalized two-fluid theory of nonlinear magnetic structures}\citep{Schep1994}\

T. J. Schep;  F. Pegoraro;  B. N. Kuvshinov.

\textsf{核读范围：}PDF1-11全部（正文2843-2852）；\textsf{图表：}式1-18视觉复核；图1、2及全文理论.

\href{run:papers/Schep1994.pdf}{打开本地PDF：\texttt{papers/Schep1994.pdf}}\quad\href{https://doi.org/10.1063/1.870523}{DOI／版本入口}

\paragraph{在报告中的研究脉络}\hyperref[story:onset]{启动、爆发与磁岛}。关联表示本报告的阅读线索，不自动表示这些论文互相引用。

\paragraph{实际模型}低beta强导引场、等温电子，平行离子流忽略，保留电子惯性和电子陀螺黏性抵消；离子响应采用全FLR非线性构造。

\paragraph{计算／实验细节}频率低于离子回旋／磁声、高于离子声响应；slab, yz周期；定常传播解背景指数密度及线性磁剪切；本文无时变重联模拟网格。

\paragraph{核实的发现}构造能量与非正则Hamilton结构，二维有两族电子Casimir，冻结组合是Psi\_e±de sqrt(beta\_e) ln(n/n0)，不是Psi或Psi\_e；有限惯性去掉电流本身的奇点，但梯度仍可能很大，广义通量X点分离。

\paragraph{限制与未完成问题}第2845页明确全FLR离子方程是ansatz，并非从Vlasov严格推导；第2849页未求解完整非线性本征问题；第2851页指出相关传播尺度Landau共振使流体模型失效；机构库封面年份错误，正文1994。

\paragraph{本报告判断}适合严谨说明约化模型的构造条件与守恒；不能把全FLR构造称为完整动力学或已完成数值重联验证。

\section{Biskamp1997：1997年}\label{paper:Biskamp1997}

\textbf{Two-fluid theory of collisionless magnetic reconnection}\citep{Biskamp1997}\

Dieter Biskamp;  E. Schwarz;  J. F. Drake.

\textsf{核读范围：}1-8（全部）；\textsf{图表：}图1-5及表I-III；表I-II视觉核验.

\href{run:papers/Biskamp1997.pdf}{打开本地PDF：\texttt{papers/Biskamp1997.pdf}}\quad\href{https://doi.org/10.1063/1.872211}{DOI／版本入口}

\paragraph{在报告中的研究脉络}\hyperref[story:fast]{局地快率与必要条件}。关联表示本报告的阅读线索，不自动表示这些论文互相引用。

\paragraph{实际模型}准中性、标量压强、均匀密度、两流体惯性；电子黏性与离子黏性正则化；忽略电阻；常密度使标量电子压强不进入感应。

\paragraph{计算／实验细节}2D周期2pi方盒；伪谱2/3去混叠；N=256至2048；用六阶扩散替代四阶以集中小尺度；EMHD de=0.1至0.01；离子扫描di=0.4至0.05，di/de约6.6-13.3。

\paragraph{核实的发现}电子惯性层长宽随de缩小；表I速率对de与超黏性弱敏感；表II在Alfven单位E*=Edi，di下降时E*由0.68到0.34；电子KH湍流主要改变有限耗散率，而非增加已快的重联率。

\paragraph{限制与未完成问题}未确定di趋零渐近；恒密度强制会产生局地能量源汇；低beta模型未纳入平行电子压强效应，不能推广成强导引场都慢；FLR仅粗略处理。

\paragraph{本报告判断}把拓扑改变尺度与吞吐尺度区分开；不是严格无耗散模型，也不证明普适0.1或无限系统尺寸无关。

\section{L020：2000年}\label{paper:L020}

\textbf{The onset of turbulence in collisionless magnetic reconnection}\citep{L020}\

B. N. Rogers;  J. F. Drake;  Michael A. Shay.

\textsf{核读范围：}GRL 27, 3157–3160（2000），物理页1–4全部，含模型、二次不稳定性、消融、PIC对照与参考文献；\textsf{图表：}图1–4与三次色散式(1)，物理页2–3直接视觉核验；扫描文本的符号以原图为准.

\href{run:papers/L020.pdf}{打开本地PDF：\texttt{papers/L020.pdf}}\quad\href{https://doi.org/10.1029/2000gl000038}{DOI／版本入口}

\paragraph{在报告中的研究脉络}\hyperref[story:three]{三维传播与有限X线}。关联表示本报告的阅读线索，不自动表示这些论文互相引用。

\paragraph{实际模型}三维等温双流体，Hall与电子惯性、无显式电阻；承接早期F3D模型。标量压力缺离子FLR、梯度B漂移共振和Landau阻尼；PIC只另检验初始一维片的二维LHDI。

\paragraph{计算／实验细节}双周期片初始w0=di，质量比277.8、de/di=.06，512×256×64、10×5×1.25di，全周期，初始总压1.5B0\textsuperscript{2}；Ti+Te=B0\textsuperscript{2}/(4πn∞)。比较de=.2di、Bg=2B0、w0=.5di等；无绝对时间步和完整收敛表。

\paragraph{核实的发现}非线性重联先形成陡峭梯度，约t=4离子回旋时间出现电子剪切流不稳定性，约6.5再出现LHDI；不是初始片线性阶段就同样不稳定。电子电流层由约de变为2.5de，仍为电子尺度；三维平均率反而比二维低约25\%。ESF惯性速度关联支撑X线平均电场，空间积分近零，体现异常黏性而非异常电阻。LHDI削弱分离线面内电子电流/Hall场30–40\%，密度梯度趋边缘稳定；去掉直接粒子关联通量影响很小，抑制Bz异常输运则继续陡化并崩溃。

\paragraph{限制与未完成问题}Lz仅1.25di，能容纳电子短波但排除更长三维结构；de约3格。二维PIC对照验证增长量级，不能替代非线性三维动理学验证。模拟磁声激波在真实动理学中可能被Landau阻尼；有限beta梯度B稳定作用等被省略。没有证明湍流必然加快重联。

\paragraph{本报告判断}直接纠正湍流等于异常电阻、等于加速的简化图像；应分别测动量通量、磁场输运、粒子通量与总率，并记录不稳定性启动顺序。可与2026十矩三维研究形成长期证据链。

\section{Birn2001：2001年}\label{paper:Birn2001}

\textbf{Geospace Environmental Modeling (GEM) Magnetic Reconnection Challenge}\citep{Birn2001}\

Joachim Birn;  J. F. Drake;  Michael A. Shay;  B. N. Rogers;  R. E. Denton;  M. Hesse;  M. M. Kuznetsova;  Z. W. Ma;  Amitava Bhattacharjee;  A. Otto;  Philip L. Pritchett.

\textsf{核读范围：}1-5（全部）；\textsf{图表：}图1；PDF第3-4页数值与说明.

\href{run:papers/Birn2001.pdf}{打开本地PDF：\texttt{papers/Birn2001.pdf}}\quad\href{https://doi.org/10.1029/1999ja900449}{DOI／版本入口}

\paragraph{在报告中的研究脉络}\hyperref[story:fast]{局地快率与必要条件}。关联表示本报告的阅读线索，不自动表示这些论文互相引用。

\paragraph{实际模型}GEM跨代码比较；有限种子Harris片，二维；Hall组的电子惯性／耗散并不统一。

\paragraph{计算／实验细节}a=0.5di, n\_b/n\_0=0.2, T\_e/T\_i=0.2, mi/me=25, Lx=25.6di, Lz=12.8di；网格和无质量电子代码的耗散留给参赛者；psi0=0.1。

\paragraph{核实的发现}图1的Hall／混合／PIC通量曲线接近，峰值率约0.24；但粒子模型电子层和离子喷流宽度可比对应流体模型大约两倍。

\paragraph{限制与未完成问题}有限种子刻意跳过线性触发；过量扩散降低速率；原文末段明确限定二维，三维次级不稳定性和亚暴触发未解决。

\paragraph{本报告判断}适合作为非线性通量率基准，不能据此声称所有扩散区细节或自然触发均可由Hall MHD替代。

\section{F766：2001年}\label{paper:F766}

\textbf{3D EMHD reconnection in a laboratory plasma}\citep{F766}\

R. L. Stenzel;  J. Manuel Urrutia;  M. C. Griskey;  K. D. Strohmaier.

\textsf{核读范围：}EPS53卷553–560，2001年全文8页；库的2014为平台登记错误，不能用作发表年份；\textsf{图表：}图1–9，直接核验页5磁通/发光及页7双涡合并FRC演化.

\href{run:papers/F766.pdf}{打开本地PDF：\texttt{papers/F766.pdf}}\quad\href{https://doi.org/10.1186/bf03353269}{DOI／版本入口}

\paragraph{在报告中的研究脉络}\hyperref[story:diagnostic]{实验与可区分诊断}。关联表示本报告的阅读线索，不自动表示这些论文互相引用。

\paragraph{实际模型}UCLA实验三维EMHD参数区的无持续驱动磁涡松弛；先由回路天线建立反向偶极场再断电，三维磁探针跨重复放电扫描并由curl B求电流，Langmuir探针及光电二极管诊断。

\paragraph{计算／实验细节}氩等离子体柱直径1m长2.5m、ne≤1e12cm\textsuperscript{-3}、Te≤3eV、B0约5–10G、中性压.3mTorr；空间步.75cm、时间10ns，天线脉冲约5μs、下降.75μs。双10cm天线相距25cm，de约.75cm、非绝热电子区约5cm、电子碰撞自由程≥20cm。

\paragraph{核实的发现}关断后势场cusp零点变为不规则spiral零点，约5μs闭合磁通变开放；双涡经退化零点/零线合并为寿命约单涡两倍的FRC。无观测到细de电流片、whistler湍流或明显异常电阻；回路磁通变化约5V，Eθ.1V/cm、J2A/cm\textsuperscript{2}推σ20Ω\textsuperscript{-1}cm\textsuperscript{-1}与Spitzer28接近。电子温度约倍增、发光主要发生关断后；早期离子密度脉冲更可能是外加场驱动。

\paragraph{限制与未完成问题}黏性/压力张量主导来自轨道尺度比de/碰撞尺度的推断，未直接测完整Pe散度预算；导电率接近经典值并不独立确认破冻结项。空间分辨约de，不能排除更薄瞬态结构；跨放电重构依赖可重复性。文中广义涡量加号错误，正确B*=B-de\textsuperscript{2}\ensuremath{\nabla}\textsuperscript{2}B；零点存在并不自动证明拓扑变化，射流动能也非必然不可逆热化。

\paragraph{本报告判断}作为EMHD三维零点演化和流体模拟预测的独立实验检验；报告保留无薄片/无湍流这个负结果及诊断分辨限制，而不将其写成普适禁止。

\section{L029：2001年}\label{paper:L029}

\textbf{Hybrid and Hall-MHD simulations of collisionless reconnection: Effects of plasma pressure tensor}\citep{L029}\

Lin Yin;  D. Winske;  S. Peter Gary;  Joachim Birn.

\textsf{核读范围：}LA-UR-01-2484，封面及ISSS-6论文正文1–4页（物理页1–5）全部；这是正式会议短文，不是同作者另两篇期刊文章；\textsf{图表：}图1–2及电子压力演化式；物理页3和5直接视觉核验.

\href{run:papers/L029.pdf}{打开本地PDF：\texttt{papers/L029.pdf}}\quad\href{https://digital.library.unt.edu/ark:/67531/metadc718188/}{DOI／版本入口}

\paragraph{在报告中的研究脉络}\hyperref[story:fast]{局地快率与必要条件}。关联表示本报告的阅读线索，不自动表示这些论文互相引用。

\paragraph{实际模型}无电子惯性、无显式电阻的Hall流体与动理学离子混合对照，电子非对角压力进入欧姆定律。文中印出的压力方程为输运/形变加Ωe/τ各向同性松弛，没有显式磁旋转项，不能不加区别称完整Vlasov十矩方程。另用粒子估计离子非对角压力修正Hall动量。

\paragraph{计算／实验细节}Harris片α=.4di、10×10di和128\textsuperscript{2}；另一组α=.8di、20×20di和256\textsuperscript{2}。上游beta=.2，Ti/Te=5，y周期；x边界横向速度/场为零。交错有限差分/leapfrog；Hall磁场RK4，每流体步5–10磁场子步，磁场六阶超黏性和流体平滑；τ及绝对步长未列。

\paragraph{核实的发现}电子压力在X点、Hall项在其周围支撑Ez；两种模型整体率和构型近似，却在离子面外速度明显不同：标量离子压力使vz一直峰于X点，混合模型将动量搬到两侧。粒子Hall用Piyz及Pixz修正后恢复分裂速度峰，主要改善来自Piyz。电子压力还通过面内电场拖曳离子，影响出入层动力学。

\paragraph{限制与未完成问题}会议论文只有两幅结果图，平均率改善没有给定量误差或收敛；没有完整电子旋转项的所印方程与现代十矩定义不同。粒子压力纠正依赖离子分布，已不是闭合纯流体。MHD嵌入只说明接口程序，没有展示其全球结果，不能据此称全球耦合已验证。

\paragraph{本报告判断}重要早期反例：总重联率接近，并不认证离子动量输运正确；应把Piyz和vz空间结构加入Hall模型验收。后续十矩及混合模型的必要性可由本篇的直接对照说明。

\section{L204：2001年}\label{paper:L204}

\textbf{Hall effects on field-aligned current generation in three-dimensional magnetic reconnection}\citep{L204}\

N. Yokokawa;  M. Fujimoto;  Y. Yamade;  Tohru Mukai.

\textsf{核读范围：}Earth Planets Space 53, 501–507 (2001)全部7页；出版年份不是数据库误收的2014；\textsf{图表：}图1–8及初始化/数值段，图7–8直接视觉核验，比较固定Ry与固定Ry/D两条扫描路线.

\href{run:papers/L204.pdf}{打开本地PDF：\texttt{papers/L204.pdf}}\quad\href{https://doi.org/10.1186/bf03353262}{DOI／版本入口}

\paragraph{在报告中的研究脉络}\hyperref[story:three]{三维传播与有限X线}。关联表示本报告的阅读线索，不自动表示这些论文互相引用。

\paragraph{实际模型}三维Hall MHD，冷无质量电子、离子压力和焦耳/黏性加热；人为局域增强电阻触发，因此不研究电子破冻结机制。

\paragraph{计算／实验细节}D=1.5di磁尾主例，Bx=tanh(z/D)，n=1+2sech\textsuperscript{2}(z/D)，离子β=0.5；η=D[4exp(-x\textsuperscript{2}/Rx\textsuperscript{2}-y\textsuperscript{2}/Ry\textsuperscript{2}-z\textsuperscript{2}/Rz\textsuperscript{2})+1]/Rm，Rm=200，Rx=4D, Ry=10di, Rz=2D。盒80D×10Ry×20D，格0.4D×0.1Ry×0.2D，四阶差分/RK4，Re=200。x, z对称边界，y周期；厚度至10di、Rm100–400扫描。

\paragraph{核实的发现}Hall电子先晨向后昏向偏转产生纬向分层的强场向电流，传播到远离X线处，区别于MHD喷流分叉产生的晨昏电流对。D=1.5例峰上行FAC约0.7初始片电流，换算约7nA/m\textsuperscript{2}和总1.5×10\textasciicircum{}5A。固定Ry增加D时MHD分量增强；固定Ry/D=6.7则Hall分量到D=10仍占优，MHD/Hall峰比约0.6。说明纵横比与三维程度不能被片厚单变量替代。

\paragraph{限制与未完成问题}Ry/D>7且D到数十di仍重要是由有限Rm扫描和Hall峰\ensuremath{\propto}√Rm外推，非全部直接计算。无偶极场、电离层耦合、电子加速动力学；对称边界抑制跨尾KH。模拟没有重现传统亚暴电流楔，不可宣称已解释极光。Ry<7di喷流不足的结果未作图；网格随D变粗且无完整收敛表。

\paragraph{本报告判断}重要发现是Hall改变大尺度电流闭合形态，非仅局部四极磁场；实验/空间联系应把模型预测FAC和实际电子能谱、极光辐射分开。

\section{Rogers2001：2001年}\label{paper:Rogers2001}

\textbf{Role of Dispersive Waves in Collisionless Magnetic Reconnection}\citep{Rogers2001}\

B. N. Rogers;  R. E. Denton;  J. F. Drake;  M. A. Shay.

\textsf{核读范围：}1-4全部；\textsf{图表：}图1-4；图3、4视觉核验.

\href{run:papers/Rogers2001.pdf}{打开本地PDF：\texttt{papers/Rogers2001.pdf}}\quad\href{https://doi.org/10.1103/physrevlett.87.195004}{DOI／版本入口}

\paragraph{在报告中的研究脉络}\hyperref[story:fast]{局地快率与必要条件}。关联表示本报告的阅读线索，不自动表示这些论文互相引用。

\paragraph{实际模型}二维等温双流体，电子惯性，周期双电流片；EMHD和约化KAW机制推导及完整线性两流体色散关系。

\paragraph{计算／实验细节}主要扫描mi/me=82.6, Lx×Ly=20di×10di, w0=di，Bg/Bx=3和30；另用质量比300；早期图1盒10×5di、质量比300；正文未报告网格及耗散实现。

\paragraph{核实的发现}同时允许whistler和KAW时率约.12；仅一分支约.06；强导引冷等离子体无二次色散波时率显著低、长片。控制量mu\_k=(me/mi)(c\_A\textsuperscript{2}+c\_s\textsuperscript{2})/c\_Ak\textsuperscript{2}与beta\_k/2=c\_s\textsuperscript{2}/c\_Ak\textsuperscript{2}，而非单独Bg/Bx。

\paragraph{限制与未完成问题}线性波向非线性外推时By随时间空间变化，图2用Bx上游作方便替代，作者明确承认；不是任意闭合、对等离子体或动理学的必要条件定理；无网格收敛展示。

\paragraph{本报告判断}证据支持其等温流体模型的波模／几何关联，后来的压力机制及动理学反例需要单独核验。

\section{Shay2001：2001年}\label{paper:Shay2001}

\textbf{Alfvénic collisionless magnetic reconnection and the Hall term}\citep{Shay2001}\

Michael A. Shay;  J. F. Drake;  B. N. Rogers;  R. E. Denton.

\textsf{核读范围：}期刊3759-3772全部14页；\textsf{图表：}图1-13正文图注；模型页、图7-9和11-12视觉核验.

\href{run:papers/Shay2001.pdf}{打开本地PDF：\texttt{papers/Shay2001.pdf}}\quad\href{https://doi.org/10.1029/1999ja001007}{DOI／版本入口}

\paragraph{在报告中的研究脉络}\hyperref[story:closure]{压力、热流与闭合}。关联表示本报告的阅读线索，不自动表示这些论文互相引用。

\paragraph{实际模型}GEM Hall/惯性标量压强双流体、混合粒子离子、全PIC对比；流体gamma5/3，常de惯性反演。

\paragraph{计算／实验细节}25.6×12.8di，x周期z导体，背景密度.2；有限岛种子跳过线性阶段。质量比25主比较512×256；混合质量比9-100最高1024×512。磁方程六阶谱耗散，密度/压强二阶防负扩散；混合电子子步为离子1/50-1/100，PIC c/VA25和场6子步。

\paragraph{核实的发现}Hall、混合、PIC总磁通率相近，但离子动理学喷流约宽两倍、峰流速约一半；常de使流体电子层比PIC约窄两倍。电子质量变化时内层宽度、峰喷流随de与电子Alfven速度变化，Bx\textsuperscript{2}/L近常数；me0时层缩到网格而总率仍接近。分离线密度耗尽：上游.2，Hall最低.03、混合.1、PIC.15；对应Hall电流与磁压增强。

\paragraph{限制与未完成问题}小盒反平行、有限种子及质量比9-100不能证明任意系统、启动、导引场和闭合都无关。无质量电子断线来自六阶/网格耗散；没有严格理想Hall断线证据。均匀小电阻不足的论证使用电子层宽长比近固定1/5等几何假设，不能排除后来的其他电阻Hall分支。密度扩散防负数也可能影响空穴细节；其系数与完整离散流程本文未列。

\paragraph{本报告判断}总吞吐相近不保证密度、离子喷流或电子扩散层正确；原文Hall最小模型结论应限制到已测试的问题。

\section{L035：2002年}\label{paper:L035}

\textbf{Simulations of current sheet thinning and reconnection}\citep{L035}\

Lin Yin;  D. Winske.

\textsf{核读范围：}正式JGR 2002JA009507全文11页，含图页及参考文献；\textsf{图表：}图1–9、式1–14，图9嵌入/全Hall磁通重合直接视觉核验；核对式3与式9–11.

\href{run:papers/L035.pdf}{打开本地PDF：\texttt{papers/L035.pdf}}\quad\href{https://doi.org/10.1029/2002ja009507}{DOI／版本入口}

\paragraph{在报告中的研究脉络}\hyperref[story:fast]{局地快率与必要条件}。关联表示本报告的阅读线索，不自动表示这些论文互相引用。

\paragraph{实际模型}二维混合粒子离子/流体电子与Hall标量离子模型对照，电子无质量但演化经验各向同性化压力张量；式3不含完整Vlasov磁旋转项，非未经简化的十矩方程。无显式常电阻，四阶磁超扩散。

\paragraph{计算／实验细节}小盒10×10di, 128\textsuperscript{2}，初始a=2di，βi=0.2, Ti/Te=5，混合20背景+平均20片离子/格；边界驱动最高入流0.1VA后关闭，Hall驱动至30、混合至40。嵌入测试60×20di、512×256、a=4di，有限法向B约0.05B0；Hall内区x5.5–14.5, y0–50。交错网格/leapfrog，Hall场RK4每流体步5–10子步；过渡erf权重。

\paragraph{核实的发现}压缩生成厚等离子体片中的电子承载薄电流层。混合模型离子出平面速度从中心向片边及X点两侧分散，Hall标量Pi保持中心峰；离子运动使片变宽，需更久驱动才能起重联。测试粒子在Hall场中恢复混合形态，支持非对角Pi输运解释。嵌入Hall/MHD测试给平滑场与重联磁通时间曲线几乎重合，较2001会议文已有实际验证。

\paragraph{限制与未完成问题}薄化阶段Pixz因片边粒子少而噪声大，没有反馈实现，作用仍是推测；Piyz反馈验证指向既往论文。嵌入外域尚与内域同格，未展示节约算量/AMR；权重梯度\ensuremath{\nabla}α×E'人为生磁场，作者用平滑处理，不是严格守恒耦合。电子闭合时间和超扩散系数等复现信息不全；无完整局部非理想预算，模型能研究离子尺度但不能声称解决全部真实电子物理。

\paragraph{本报告判断}把其作为‘外部压缩—电子薄层—模型间离子输运差别’和早期多模型嵌入证据，区分已计算结果与拟议改进。

\section{Shay2003：2003年}\label{paper:Shay2003}

\textbf{Inherently three dimensional magnetic reconnection: A mechanism for bursty bulk flows?}\citep{Shay2003}\

M. A. Shay;  J. F. Drake;  M. Swisdak;  W. Dorland;  B. N. Rogers.

\textsf{核读范围：}期刊4页全部；\textsf{图表：}图1-3与正文全部，图像视觉核查.

\href{run:papers/Shay2003.pdf}{打开本地PDF：\texttt{papers/Shay2003.pdf}}\quad\href{https://doi.org/10.1029/2002gl016267}{DOI／版本入口}

\paragraph{在报告中的研究脉络}\hyperref[story:three]{三维传播与有限X线}。关联表示本报告的阅读线索，不自动表示这些论文互相引用。

\paragraph{实际模型}三维等温F3D Hall+电子惯性，全部方程加四阶耗散。

\paragraph{计算／实验细节}原文x流出、z流入、y电流；51.2×128×51.2di，网格.1×2×.1di，质量比25；半厚.5/1/2，n1-1.5。随机磁场与离子电流扰动最大.2；孤立X线试验用sech\textsuperscript{2}(y/10)种子，|Jy|>2.5端点；y网格做两值检查。

\paragraph{核实的发现}较厚2di电流片产生约10di有限长X线且可整体向电子方向移动，较薄片多个X线逐渐连接为全局重联。一端速度随初始片厚倒数，另一端可静止、反向扩展或同向平移；后者伴三维慢模样压力/磁压反相结构。

\paragraph{限制与未完成问题}小噪声时厚片在有限时间不启动，不是严格线性稳定证明。耗散系数与完整时间算法此文未列；电流方向粗网格抑制微不稳定性。磁尾1-4Re和BBF只量级类比，缺平衡法向场及动理学；2021改变盒子长宽比已显示固定长度不是仅由片厚决定。

\paragraph{本报告判断}早期发现介观三维局域X线及合并，不应把有限盒中的孤立结构升格为普适亚暴触发模型。

\section{L049：2004年}\label{paper:L049}

\textbf{Developments in the theory of collisionless reconnection in magnetic configurations with a strong guide field}\citep{L049}\

Francesco Pegoraro;  D. Borgogno;  F. Califano;  Del Sarto,  D.;  E. Yu. Echkina;  D. Grasso;  T. V. Liseykina;  Francesco Porcelli.

\textsf{核读范围：}1-11全部；\textsf{图表：}图1-5，PDF第5-7、9页视觉复核.

\href{run:papers/L049.pdf}{打开本地PDF：\texttt{papers/L049.pdf}}\quad\href{https://doi.org/10.5194/npg-11-567-2004}{DOI／版本入口}

\paragraph{在报告中的研究脉络}\hyperref[story:onset]{启动、爆发与磁岛}。关联表示本报告的阅读线索，不自动表示这些论文互相引用。

\paragraph{实际模型}作者团队强导引场约化双流体及漂移动理学结果综述；二维Hamilton两场；暖电子模型需保留抵消抗磁漂移惯性的压力张量修正。

\paragraph{计算／实验细节}psi0=-L/[2cosh\textsuperscript{2}(x/L)]，Lx=2Ly=4pi L，Nx=2048, Ny=1024，de=.3L, rho\_s/de=0-1.5；只一个线性不稳定模，随机扰动。

\paragraph{核实的发现}暖电子G±不断卷绕细丝，粗粒化磁岛可饱和而总能量守恒；冷电子极限出现二次KH、单极和偶极涡，岛形脉动与流体／电子动能交换；rho\_s/de接近1时饱和前未见KH。漂移动理学按不同正则动量薄片发生不同方向卷绕。

\paragraph{限制与未完成问题}第573页明确暖电子细丝无限变薄导致流体近似不一致；综述未给完整算法细节；未证明三维、任意尺度或热化过程。

\paragraph{本报告判断}可以学习磁拓扑与广义拓扑不相同以及宏观饱和不等于微观静止；不能把细丝当成已被解析的真实耗散。

\section{L057：2004年}\label{paper:L057}

\textbf{Scaling of forced magnetic reconnection in the Hall-magnetohydrodynamic Taylor problem}\citep{L057}\

Richard Fitzpatrick.

\textsf{核读范围：}期刊937–946全文；物理页1封面、2–11正文/参考文献全部；\textsf{图表：}图1–13；图10–13与耗散分解视觉核验.

\href{run:papers/L057.pdf}{打开本地PDF：\texttt{papers/L057.pdf}}\quad\href{https://doi.org/10.1063/1.1640378}{DOI／版本入口}

\paragraph{在报告中的研究脉络}\hyperref[story:global]{层长、驱动与全局完成}。关联表示本报告的阅读线索，不自动表示这些论文互相引用。

\paragraph{实际模型}二维不可压零导引场Hall MHD，无电子惯性，含均匀电阻/离子黏性及数值用超电阻；Taylor导体墙位移驱动。

\paragraph{计算／实验细节}x±1导体墙、y周期L=8，psi0=-x\textsuperscript{2}/2；墙扰动Xi=Xi0[1-(1+t/tau)exp(-t/tau)]，tau1。256×128、x集中网格、dt1e-3，二阶空间/时间，Harned–Mikic型半隐式。线性Xi0=1e-4；非线性常Xi0=.01、nu=2.5e-9，电阻约1e-6–1e-3，按图另扫。

\paragraph{核实的发现}J诊断定义为R/eta，不能直接当物理重联率。线性R\textasciitilde{}di*eta\textasciicircum{}(1/3)*Xi0，Hall转折di\textasciitilde{}eta\textasciicircum{}(1/3)；非线性峰值约R\textasciitilde{}di\textasciicircum{}(3/2)*Xi0\textsuperscript{2}，电阻/超电阻依赖较弱。图10中低eta时超电阻贡献主导，不能称纯电阻快重联。Hall电流局域化、离子与细层解耦。

\paragraph{限制与未完成问题}指数来自有限扫描拟合，di接近系统大小时失效；没有独立网格/步长收敛数据。无电子惯性使极低耗散极限不可无限外推，短波正则化必要；峰值非稳态/全程平均。与自发薄片的上游场选择不同。

\paragraph{本报告判断}适合学习驱动协议、线性与非线性、R和R/eta的区分；不能把所有快Hall重联概括成同一di指数。

\section{L058：2004年}\label{paper:L058}

\textbf{Scaling of forced magnetic reconnection in the Hall-magnetohydrodynamical Taylor problem with arbitrary guide field}\citep{L058}\

Richard Fitzpatrick.

\textsf{核读范围：}期刊3961–3968全文；物理页1封面、2–9正文/参考文献全部；\textsf{图表：}图1–7，式13–24；物理页4、6、8视觉复核.

\href{run:papers/L058.pdf}{打开本地PDF：\texttt{papers/L058.pdf}}\quad\href{https://doi.org/10.1063/1.1768956}{DOI／版本入口}

\paragraph{在报告中的研究脉络}\hyperref[story:global]{层长、驱动与全局完成}。关联表示本报告的阅读线索，不自动表示这些论文互相引用。

\paragraph{实际模型}消去快压缩波的冷离子四场模型，保留小压缩率与大导引场乘积；无电子惯性。c\_beta=sqrt(beta\_F/(1+beta\_F)), d\_beta=c\_beta di，beta\_F=Gamma P0/Bg\textsuperscript{2}。

\paragraph{计算／实验细节}Taylor同L057几何；256×128、dt1e-3、二阶半隐式，超电阻2.5e-9，mu=kappa1e-7。beta\_F从1e-2到1e2；比较固定d\_beta，不是固定di。典型eta1e-6, Xi0=.01, d\_beta=.5。

\paragraph{核实的发现}固定d\_beta时跨四个数量级beta峰值和初始爆发相似；初始约8个Alfven时间后骤增，峰值约电阻MHD30倍，但首波仅长到最终岛宽约25\%，之后慢而不稳并趋Rutherford增长。拟合Rmax\textasciitilde{}d\_beta\textasciicircum{}(3/2)*Xi0\textsuperscript{2}；强导引场下d\_beta趋声回旋尺度。

\paragraph{限制与未完成问题}“任意导引场”仍受快压缩波分离、冷离子、忽略密度变化等排序约束，非完整任意beta双流体。beta\_F与本报告beta常规定义差Gamma/2。对Shay模型的反驳依赖难以定义的上游边界及原文线性场强说法，不足独立推翻R\textasciitilde{}Bup\textsuperscript{2}/sqrt(n)。W/d\_beta约10案例指数拟合本身不精确。

\paragraph{本报告判断}突出初始爆发与长期磁通处理不等价；四场约化需从压力平衡和小压缩率逐步推导，不能直接设所有div(v)=0。

\section{Shay2004：2004年}\label{paper:Shay2004}

\textbf{The scaling of embedded collisionless reconnection}\citep{Shay2004}\

M. A. Shay;  J. F. Drake;  M. Swisdak;  B. N. Rogers.

\textsf{核读范围：}期刊2199-2213共15页全部；\textsf{图表：}表1、图1-11；图3-5、8、10及方程视觉核验.

\href{run:papers/Shay2004.pdf}{打开本地PDF：\texttt{papers/Shay2004.pdf}}\quad\href{https://doi.org/10.1063/1.1705650}{DOI／版本入口}

\paragraph{在报告中的研究脉络}\hyperref[story:global]{层长、驱动与全局完成}。关联表示本报告的阅读线索，不自动表示这些论文互相引用。

\paragraph{实际模型}二维F3D准中性等温Hall加常系数电子惯性Bprime=(1-de\textsuperscript{2}lap)B；非完整可压双流体惯性。

\paragraph{计算／实验细节}全周期正弦双片，n1-1.5, T=1, mi/me25；Lx51.2-409.6di, Ly12.8-102.4di，通常dx=.1di、de=.2di。四阶耗散所有方程mu4=5.1e-5；run7 dx=.05、mu4=5e-7。四个大算例初始岛半宽.55di；加B及离子流噪声1e-4。

\paragraph{核实的发现}局地Bd用离子Jz峰的25\%位置定义，上游/下游边界并非任意取di；下游取离子电子流速交点。大系统后期Vin约.1局地Alfven速度，进入该阶段时岛半宽约6.1-9.9di。最大算例最后3\%时段完成约90\%磁通重联、75\%岛宽增长。小Delta-prime时即使呈X形，压力背压抵消张力仍慢；大Delta-prime时电子层变薄、Hall增强伴离子层从长片缩短。

\paragraph{限制与未完成问题}无导引场、单质量比；质量比无关只是援引先前研究和预期。启动持续时间强依赖尺寸与初值，非普适0.1。de仅两格，run7同时改网格和耗散不能当作单独网格扫描。局地Bd定义含25\%阈值；迟期边界干扰后停止比较。印刷式2以Ji为粒子通量却压力项写T/n grad n，且温度文字Ti+Ti，不能照抄作一致守恒方程。

\paragraph{本报告判断}关键贡献是分开启动与快速阶段、实际测定局地上游场；宏观归一化与局地归一化不同，不能据后期结果抹去启动难题。

\section{Cassak2005：2005年}\label{paper:Cassak2005}

\textbf{Catastrophe Model for Fast Magnetic Reconnection Onset}\citep{Cassak2005}\

P. A. Cassak;  Michael A. Shay;  J. F. Drake.

\textsf{核读范围：}1-10（全部，arXiv v2）；\textsf{图表：}图1-4；图3视觉核验.

\href{run:papers/Cassak2005.pdf}{打开本地PDF：\texttt{papers/Cassak2005.pdf}}\quad\href{https://doi.org/10.1103/physrevlett.95.235002}{DOI／版本入口}

\paragraph{在报告中的研究脉络}\hyperref[story:onset]{启动、爆发与磁岛}。关联表示本报告的阅读线索，不自动表示这些论文互相引用。

\paragraph{实际模型}F3D二维双周期双Harris片、等温、均匀电阻；含电子惯性；四阶耗散加在全部方程。

\paragraph{计算／实验细节}409.6di×204.8di，dx=dy=0.1di，mi/me=25，初始w0=2di，eta4=2e-5；快慢两支启动程序不同，最终方程相同。

\paragraph{核实的发现}相同eta'=0.015有约0.06与0.01两个速率；图3慢快临界0.011-0.013，快慢临界0.020-0.0225；图4先降电阻再恢复仍保持快支，检验历史依赖。

\paragraph{限制与未完成问题}de仅两格，论文未展示独立网格收敛；强导引场rho\_s替代在该文只是外推；DIII-D温度估算低约一数量级；MRX缺片宽不能定量比较；二维未处理湍流。

\paragraph{本报告判断}双稳态证据来自同参不同历史与迟滞试验，而非仅凭两条速率曲线；实验比较不能写成普遍定量验证。

\section{E459：2005年}\label{paper:E459}

\textbf{Current layer cascade in collisionless electron-magnetohydrodynamic reconnection and electron compressibility effects}\citep{E459}\

Del Sarto,  D.;  F. Califano;  Francesco Pegoraro.

\textsf{核读范围：}HAL封面及PoP12, 012317(2005)全部8页正文，共9页；发表在线2004、卷期2005；\textsf{图表：}图1–9及式1–10；物理页6–8直接核验广义磁通/电流涡旋与能量曲线.

\href{run:papers/E459.pdf}{打开本地PDF：\texttt{papers/E459.pdf}}\quad\href{https://doi.org/10.1063/1.1829064}{DOI／版本入口}

\paragraph{在报告中的研究脉络}\hyperref[story:onset]{启动、爆发与磁岛}。关联表示本报告的阅读线索，不自动表示这些论文互相引用。

\paragraph{实际模型}冷电子惯性EMHD，静止离子、强均匀导向场，无碰撞无动理学；F=ψ-de\textsuperscript{2}\ensuremath{\nabla}\textsuperscript{2}ψ随电子流冻结，W=b-de\textsuperscript{2}\ensuremath{\nabla}\textsuperscript{2}b有磁力源。电荷分离压缩修正仅把W中的de换成λe=de√(1+Ωce\textsuperscript{2}/ωpe\textsuperscript{2})。

\paragraph{计算／实验细节}Lx=4πL, Ly=2πL，1024×512，de=.06–1L；de=.12/.3另用2048×512。ψ0=-L/(2cosh\textsuperscript{2}(x/L))，k=1的Δ'约14.3。独立射流试验1024\textsuperscript{2}、δ约de、u约.085L/τw；主比较de=.3, λe=.3/.67。

\paragraph{核实的发现}平滑磁岛内不重联的F细丝形成电流条带，条带相撞折转，随后弯曲射流/KH型次级不稳定性卷出蘑菇状涡旋，破坏规则电流层级联。提高de或电荷分离压缩延缓涡旋形成；压缩独立射流约20个whistler时间无明显改变。晚期磁场两个分量和电子动能振荡交换，伴磁岛脉动。

\paragraph{限制与未完成问题}离子质量形式上取无穷，长时间结果不能自动用于真实电子离子系统。饱和前能量/Casimir误差<1\%，晚期不可解析细尺度使总能误差达5\%/10\%；故晚期阻尼不能定量归因物理耗散。能量重分配不是热化；模型无离子异常动量交换。未完整列出时间推进与边界实现，引用旧代码。

\paragraph{本报告判断}将L313的正则磁通冻结变成实际非线性形态证据；是电子压缩抑制次级流体不稳定性的重要原始模拟。

\section{L076：2005年}\label{paper:L076}

\textbf{The scaling of forced collisionless reconnection}\citep{L076}\

Brian Sullivan;  B. N. Rogers;  Michael A. Shay.

\textsf{核读范围：}2005期刊版122312物理页1–12全部；不是2008同题论文；\textsf{图表：}表I、图1–14；模型页2、图11、14视觉核验.

\href{run:papers/L076_2005.pdf}{打开本地PDF：\texttt{papers/L076\_2005.pdf}}\quad\href{https://doi.org/10.1063/1.2146910}{DOI／版本入口}

\paragraph{在报告中的研究脉络}\hyperref[story:global]{层长、驱动与全局完成}。关联表示本报告的阅读线索，不自动表示这些论文互相引用。

\paragraph{实际模型}二维等温两流体、简化B*=(1-me/mi Laplacian)B惯性；正弦宽电流片，稳定/边界稳定的撕裂初态；离子动量内加局域体力。

\paragraph{计算／实验细节}双周期102.4di方盒，512×1024、dx=.2、dy=.1，质量比25；另质量比100用1024×2048。四阶差分、二阶trapezoidal leapfrog；密度/动量四阶耗散5.1e-5。驱动Gaussian宽wx=10、wy=2.22、tau\_f=10，F∞=.005/.03/.05/.11，另连续爬升、宽度7.07/14.14、Bg=5、低温对照。

\paragraph{核实的发现}强驱动先堆积磁通/压力，再在片厚约.7–.8局地di时转入快重联；弱驱动形成新准平衡，无快阶段。临界F∞在.005与.0075之间，仅适用于该设定。绝对率随局地Bd\textsuperscript{2}/sqrt(n)增加；较宽驱动的绝对率更大，但除以Bd\textsuperscript{2}后反而较小。强迫案例的归一出流/率比自发薄片对照约低50\%。岛宽停止长大时内部场仍能增强，磁通仍改变。

\paragraph{限制与未完成问题}准稳态X点靠数值扩散破冻结；de法向仅约2格，无完整分辨率扫描。惯性算子是模型近似，非变密度精确双流体。表I、正文交叉图号、图9/10的F∞标注存在不一致，复现须核对输入；不把驱动阈值推广为通用常数。

\paragraph{本报告判断}外部供给控制上游场、局地层处理磁通的两层关系得到明确展示；必须同时报告绝对率和局地归一率。

\section{Ren2005：2005年}\label{paper:Ren2005}

\textbf{Experimental Verification of the Hall Effect during Magnetic Reconnection in a Laboratory Plasma}\citep{Ren2005}\

Yang Ren;  Masaaki Yamada;  Stefan P Gerhardt;  Hantao Ji;  Russell M. Kulsrud;  Aleksey Kuritsyn.

\textsf{核读范围：}055003-1至4（全部4页）；\textsf{图表：}图1-4及式1；主要结果图视觉核验.

\href{run:papers/Ren2005.pdf}{打开本地PDF：\texttt{papers/Ren2005.pdf}}\quad\href{https://doi.org/10.1103/physrevlett.95.055003}{DOI／版本入口}

\paragraph{在报告中的研究脉络}\hyperref[story:diagnostic]{实验与可区分诊断}。关联表示本报告的阅读线索，不自动表示这些论文互相引用。

\paragraph{实际模型}MRX无初始净导引场的H/D实验；二维磁探针与高分辨率1D扫描。

\paragraph{计算／实验细节}71通道径向间距1.25mm，系统误差小于10G；90通道2D图轴向4cm、径向扫描1cm且多次放电平均；重联约40微秒对Alfven约1微秒。

\paragraph{核实的发现}四极Hall磁场高分辨率50-60G，对重联场约100G；粗2D图20-40G。Hall结构径向约di、轴向约5di；碰撞率指标delta/lambda\_mfp由.2到20，BHall/Brec约.6降.25。

\paragraph{限制与未完成问题}电子20km/s由电流及电子承担假设推断，并非直接流速测量；Hall电场140V/m对Erec约100V/m是上游量级估计，部分区域BR不准，未全项闭合，也不证明Hall在X点断线。误差棒为放电差异。

\paragraph{本报告判断}强证据是四极场的空间极性及碰撞率关联；与湍动同时存在，不能排除其他耗散机制。

\section{Cassak2006：2006年}\label{paper:Cassak2006}

\textbf{A Model for Spontaneous Onset of Fast Magnetic Reconnection}\citep{Cassak2006}\

P. A. Cassak;  J. F. Drake;  M. A. Shay.

\textsf{核读范围：}arXiv physics/0604001v1全部11页；\textsf{图表：}图1-2全部视觉核验.

\href{run:papers/Cassak2006.pdf}{打开本地PDF：\texttt{papers/Cassak2006.pdf}}\quad\href{https://doi.org/10.1086/505690}{DOI／版本入口}

\paragraph{在报告中的研究脉络}\hyperref[story:onset]{启动、爆发与磁岛}。关联表示本报告的阅读线索，不自动表示这些论文互相引用。

\paragraph{实际模型}无导引场F3D等温惯性Hall双流体；宽正弦片内嵌SP层，均匀电阻。

\paragraph{计算／实验细节}409.6×204.8di0，网格.1di0，质量比25，D4=2e-5；Bx=cos(2pi y/Ly)、n=1+.5sin\textsuperscript{2}、T1。扰动.004；电阻.009和.0025两组均转快；展示的.0025组从.009组t5364继承，过渡扰动至约11000消失。

\paragraph{核实的发现}固定电阻后，重联输运较强上游磁场使内嵌SP层逐渐变薄；约t18000离子电子解耦，至约19500转入开放排气。慢阶段vin约eta/delta，快阶段以.17VAup拟合；电子层降至.2di0。建立长期缓慢积累与短时加速的自发启动候选机制，无须在启动时外加降低电阻。

\paragraph{限制与未完成问题}所示运行并非从最初到最后电阻从未改过，后期自发转换与早期重启应分开说明。原文400离子时间的扩散时间估计与图中约1500的整体转换窗仅量级接近，不宜称严格上界得到数值验证。太阳11小时积累、28秒转换是选定参数估算；VTF支持仅私人通信。未含导引场、演化电阻、热反馈、3D和后来的高S磁岛竞争，未唯一解释耀斑触发。

\paragraph{本报告判断}关键新证据是强场被慢流入自动带入后引发Hall转变；宏观耀斑时间相似只能算相容，不能作机制判定。

\section{L089：2006年}\label{paper:L089}

\textbf{Physical origin of the quadrupole out-of-plane magnetic field in Hall-magnetohydrodynamic reconnection}\citep{L089}\

Dmitri Uzdensky;  Russell M. Kulsrud.

\textsf{核读范围：}arXiv astro-ph/0605309v1全部44页，正文1–32、附录A/B33–40、参考文献和图42–44均读完；\textsf{图表：}式1–69、A1–A14、B1–B18，图1–3；图1/3与电子惯性阈值页27直接视觉核验.

\href{run:papers/L089.pdf}{打开本地PDF：\texttt{papers/L089.pdf}}\quad\href{https://doi.org/10.1063/1.2209627}{DOI／版本入口}

\paragraph{在报告中的研究脉络}\hyperref[story:diagnostic]{实验与可区分诊断}。关联表示本报告的阅读线索，不自动表示这些论文互相引用。

\paragraph{实际模型}定常不可压2.5维理想EMHD，无导向场；内域离子静止均匀背景，电子质量与电阻主推导忽略，标量Pe。外缘须恢复离子电流，内电子非理想层不在解域。

\paragraph{计算／实验细节}解析无网格/扫描。给定面内磁通Ψ=-cEzt+B0δ(y\textsuperscript{2}/δ\textsuperscript{2}-x\textsuperscript{2}/L\textsuperscript{2})/2及外给Ez，δ\ensuremath{\ll}L。定义V=\ensuremath{\int}dlpol/Bpol，D=4πne/c，Q=(L/di)(vrec/VA)；没有求出L、δ、Ez的全局闭合。

\paragraph{核实的发现}电中性与不可压迫使电子沿扩张/收缩磁通管进出，给Φe=-cEzV、Bz=-DΦe=cDEzV；四极Bz近分离面呈对数增强。在严格定常理想EMHD内域，Bz沿电子流线不变、vez沿面内磁力线不变，故差动拉伸的生成源必须在离子电流尚重要的过渡区，不能放在纯EMHD内域。附录给\ensuremath{\nabla}\textsuperscript{2}Ψ=F(Ψ)一致性条件。双极面内E将离子拉入形成对向束，HallB偏转为排流。

\paragraph{限制与未完成问题}此为给定磁几何的局部结构解，不自洽预测快率；分离面对数奇性、所有理想轨道汇X点是理想化失效。估计惯性破坏区y*≈δ[(me/mi)ln(mi/me)]\textasciicircum{}1/4，不能等同通常de厚度的普适测量；未求完整正则化解。假设剖面的平坦jz与边界推出的尖锐vez有不一致，仅用磁压回馈定性修正。Pe与电势在修正电场中简并，真实离子加热需碰撞/波散射；对向束的二阶矩增加不自动是热化。

\paragraph{本报告判断}第一章可用来推导Hall四极场及说明内外区域匹配。四极场作为两流体标志的强说法仅在本文对称无导向场假设内成立；与L137给定流动例和强导向场F644比较，不能将形态当充分因果诊断。

\section{Cassak2007：2007年}\label{paper:Cassak2007}

\textbf{Scaling of asymmetric magnetic reconnection: General theory and collisional simulations}\citep{Cassak2007}\

P. A. Cassak;  M. A. Shay.

\textsf{核读范围：}期刊102114全部11页；\textsf{图表：}式1-37、表I-II及图1-7；核心图表3-7视觉核验.

\href{run:papers/Cassak2007.pdf}{打开本地PDF：\texttt{papers/Cassak2007.pdf}}\quad\href{https://doi.org/10.1063/1.2795630}{DOI／版本入口}

\paragraph{在报告中的研究脉络}\hyperref[ch:numerics]{模型与数值可信度}。关联表示本报告的阅读线索，不自动表示这些论文互相引用。

\paragraph{实际模型}二维反平行非对称守恒标度；数值验证是电阻MHD，不是Hall，作为后续Hall理论的基础。

\paragraph{计算／实验细节}F3D显式梯形leapfrog和有限差分；409.6×204.8L0双周期，dx.1；w0=2、eta.05、gamma5/3、D4=1e-4。六组B01/B02为1/1、1/2、1/3、1/4、1/5、2/4；密度均1、最低beta1；温度不对称，未加Joule热。

\paragraph{核实的发现}磁X点、流停滞点和电流峰不重合，较低Alfven速度侧的等离子体穿越X点。磁岛向弱场鼓起，但X线相对于盒子向强场漂移；片内X点仍偏弱场侧，两种说法参照不同。强场外侧出现纯扩散区，与喷流宽度不同；分片tanh比简单平移tanh更能描述剖面。流出速度、通量率与两侧B和rho的混合标度符合受测参数范围。

\paragraph{限制与未完成问题}流出密度公式假设磁管混合时体积近不变；能量标度忽略压力贡献，不是严格完整热力学解。表I使用实测局地B和层长，强侧厚度还依赖电流峰对称化重建。仅直接扫描磁场不对称且密度相同；密度差与Hall验证留给他文。均匀大电阻且无欧姆热；强场纯扩散区未必存在于真实无碰撞实验。

\paragraph{本报告判断}应把无指定耗散的守恒约束与确定层长宽所需的闭合分开；不能用此文单独定出Hall绝对率。

\section{Arnold2008：2008年}\label{paper:Arnold2008}

\textbf{A semi-implicit Hall-MHD solver using whistler wave preconditioning}\citep{Arnold2008}\

Lukas Arnold;  Jürgen Dreher;  Rainer Grauer.

\textsf{核读范围：}1-11（arXiv v1全部）；\textsf{图表：}图1-4、表1、式13-21；主要结果图表视觉核验.

\href{run:papers/Arnold2008.pdf}{打开本地PDF：\texttt{papers/Arnold2008.pdf}}\quad\href{https://doi.org/10.1016/j.cpc.2007.11.018}{DOI／版本入口}

\paragraph{在报告中的研究脉络}\hyperref[ch:numerics]{模型与数值可信度}。关联表示本报告的阅读线索，不自动表示这些论文互相引用。

\paragraph{实际模型}质量为零电子Hall MHD；仅Hall感应项Crank–Nicolson半隐式，其余RK2显式；局地3×3近似Jacobian、Richardson平滑与非线性多重网格。

\paragraph{计算／实验细节}64\textsuperscript{2}单whistler模式测试；GEM类对称子域Lx=5di, Ly=2.5di, Nx=256, Ny=128；显式dt=.0002；隐式.002, .004, .0082。

\paragraph{核实的发现}高波数误差较快衰减，低波数由多重网格修正；2-3个V周期收敛。步长最大41倍，通量曲线接近显式且减dt改善，计算时间约减80\%。

\paragraph{限制与未完成问题}是特定算例的耗时与时间步检验，不等于任何AMR重联均加速；Alfven波仍限制显式其余部分。预印本没有详列GEM电阻及所有运行参数。

\paragraph{本报告判断}局地whistler预条件比纯流体Jacobian更适合Hall刚性；隐式稳定不保证对线性增长率的时间精度。

\section{Cassak2008：2008年}\label{paper:Cassak2008}

\textbf{Scaling of asymmetric Hall magnetic reconnection}\citep{Cassak2008}\

P. A. Cassak;  Michael A. Shay.

\textsf{核读范围：}期刊PDF1-5全部；\textsf{图表：}表1、图1-3，PDF第2-3页视觉复核.

\href{run:papers/Cassak2008.pdf}{打开本地PDF：\texttt{papers/Cassak2008.pdf}}\quad\href{https://doi.org/10.1029/2008gl035268}{DOI／版本入口}

\paragraph{在报告中的研究脉络}\hyperref[story:global]{层长、驱动与全局完成}。关联表示本报告的阅读线索，不自动表示这些论文互相引用。

\paragraph{实际模型}F3D2D反平行非对称双周期Harris；绝热离子gamma5/3、冷电子；电子惯性近似采用常密度且忽略离子速度。

\paragraph{计算／实验细节}204.8di×102.4di, dx=.05di, mi/me25, w0=2di, beta\_min4；无二阶电阻黏性，所有方程四阶扩散约6.25e-6 di\textsuperscript{3} vAmax；磁噪5e-5, 速度噪.08。场比至3、密度比至4。

\paragraph{核实的发现}九组算例率与混合Alfven标度基本一致；对称基准率.0587，离子喷流.685vA、电子2.067vA；离子／电子停滞点及X点可三者分离；磁岛向弱场鼓起，片向强场移动。两种率测法差1-10\%。

\paragraph{限制与未完成问题}电子惯性未一致保留密度不对称；高beta、无导引、二维；另一离子流速定义绝对值差60\%但标度相同。作者限可可靠研究的不对称度<4；更大时电子片失真；表1ABN1行p、T、rho不相容，复现需查原始输入。

\paragraph{本报告判断}支持有限范围的反平行标度，不证明磁层顶强不对称与导引场情形普适0.1。

\section{F717：2008年}\label{paper:F717}

\textbf{The scaling of forced collisionless reconnection}\citep{F717}\

Brian P. Sullivan;  Barrett N. Rogers.

\textsf{核读范围：}arXiv0906.0334v1物理页1–10全部；对应2008年期刊，页眉2021为排版日期；\textsf{图表：}图1–6；物理页7的式17/18与图3/4视觉复核.

\href{run:papers/F717.pdf}{打开本地PDF：\texttt{papers/F717.pdf}}\quad\href{https://doi.org/10.1063/1.2992136}{DOI／版本入口}

\paragraph{在报告中的研究脉络}\hyperref[story:global]{层长、驱动与全局完成}。关联表示本报告的阅读线索，不自动表示这些论文互相引用。

\paragraph{实际模型}等温含简化电子惯性两流体；三维局域离子体力驱动，与同尺寸二维比较；无初始导引场。

\paragraph{计算／实验细节}25.6di周期立方体、128×256×128、dx=dz=.2、dy=.1、质量比25；四阶空间差分、二阶trapezoidal leapfrog，四阶密度/动量耗散5.1e-5。wx=wz=5；典型2D F∞=.15、3D .2；三维率为场线识别磁岛面积内磁通导数，另比较中心切片及除以10di的代理。

\paragraph{核实的发现}三维磁通和等离子体可沿±z绕开重联区，强迫压缩较弱、重联持续更短；该例初始强迫区磁通被重联的比例二维超过70\%、三维约10\%。磁通分布沿电子载流方向偏移。盘状出流几何给Vin/Vout≈2delta/D，二维矩形则delta/D；局地Bd归一后率仍相近。

\paragraph{限制与未完成问题}3D总磁通率除以2wz只是人为的每单位长度比较，非一般3D重联率定理。盒子小、等温、X点数值扩散，无压力张量与粒子谱；未给三维网格收敛。预印本式17a/18a把delta/D乘到Vout，和邻文及图6的Vout/VA≈1矛盾；式10的tan与稳定性论述亦不一致，tau\_f前后写法不同，不能原样作复现输入。

\paragraph{本报告判断}新现象在三维磁通绕流和电子输运导致偏斜；相近的局地率不能推出相近的全局转换效率。

\section{Hakim2008：2008年}\label{paper:Hakim2008}

\textbf{Extended MHD Modelling with the Ten-Moment Equations}\citep{Hakim2008}\

Ammar H. Hakim.

\textsf{核读范围：}期刊36-43页（全部8页）；\textsf{图表：}图1-4、式19-26、42-47；图3-4视觉核验.

\href{run:papers/Hakim2008.pdf}{打开本地PDF：\texttt{papers/Hakim2008.pdf}}\quad\href{https://doi.org/10.1007/s10894-007-9116-z}{DOI／版本入口}

\paragraph{在报告中的研究脉络}\hyperref[ch:numerics]{模型与数值可信度}。关联表示本报告的阅读线索，不自动表示这些论文互相引用。

\paragraph{实际模型}一般推导为两物种十矩全Maxwell，但实际GEM例是十矩离子、充分松弛的五矩电子；热流散度设零，离子碰撞设零。

\paragraph{计算／实验细节}GEM近似设置：Lx=8pi, Ly=4pi, 网格256×128, a=.5di, mi/me=25, nb=.2, Ti/Te=5, psi0=.1；x周期，y开边界；初始电流仅电子承担。

\paragraph{核实的发现}通量在t约10后快速增长，作者与GEM/五矩基准比较；电子层边缘强流后期激发counter-streaming流体不稳定性。论文还展示真实质量比1/1832.6的1D色散Riemann算例。

\paragraph{限制与未完成问题}不能把1D真实质量比当作2D重联质量比；正文明确数值算法未详述而引用2006论文，GEM没有独立网格收敛或压力张量对PIC细节验证。零热流不是普适物理闭合。

\paragraph{本报告判断}属于早期模型可行性验证；摘要“ten-moment”不能理解为该重联例两物种均演化完整压力张量。

\section{L119：2008年}\label{paper:L119}

\textbf{Experimental study of the Hall effect and electron diffusion region during magnetic reconnection in a laboratory plasma}\citep{L119}\

Yang Ren;  Masaaki Yamada;  Hantao Ji;  S. Dorfman;  Stefan P Gerhardt;  Russel Kulsrud.

\textsf{核读范围：}PPPL4329全部46物理页，包括封面、完整正文、图1–18与参考文献；正文印刷页1–39；\textsf{图表：}图1–18；直接核验物理页28图13欧姆预算、33图16尺度、34图17速度、37图18摩擦模型比较（34页渲染备查，速度另由原文及图17数据描述核对）.

\href{run:papers/L119.pdf}{打开本地PDF：\texttt{papers/L119.pdf}}\quad\href{https://doi.org/10.1063/1.2936269}{DOI／版本入口}

\paragraph{在报告中的研究脉络}\hyperref[story:diagnostic]{实验与可区分诊断}。关联表示本报告的阅读线索，不自动表示这些论文互相引用。

\paragraph{实际模型}MRX无外加导场实验，五列磁探针由BT导出面内电流，再用ve≈-j/en和Mach离子探针验证双尺度；另与具有MRX通量芯边界的NPIC二维全动理学模拟比较。稳态一维电子排出摩擦模型只用于机制估计。

\paragraph{计算／实验细节}径向磁测最高2.5mm、轴向3cm；H/D/He扫描密度与碰撞性，约半数放电出现可辨四极场，分析选择这些放电。模拟75×150cm、两通量芯面距40cm、质量比75、初始n2.6e19m\textsuperscript{-3}、线圈下降200初始离子回旋时、通量芯吸收粒子；用整体率与离子片厚选参数。

\paragraph{核实的发现}Hall四极幅度约上游场30–60\%，电子超离子Alfvén排出与更宽更慢离子通道并存。δBT定义为峰值到40\%速度点的半宽，在轴向速度峰处测得δBT≈8de；vout≈.11VAe（VAe用电子层边缘B），均无显著离子质量依赖。长度约40–80de且随密度降低而缩短，不支持统一L\ensuremath{\propto}de；与二维1–2de/更快电子流存在真实差别。乘积(vout/VAe)(δ/de)约.88接近1。

\paragraph{限制与未完成问题}图13E约170V/m、估另一Hall分量约50V/m、经典电阻约40V/m、所估惯性约1V/m，Pe张量及涨落预算未测全，不能宣布耗散机制解决。摩擦模型u\textsuperscript{2}+Au=1，A=νL/VAe；经典预测.3–.6，经验ν*预测.1–.3仍可能高两倍；ν*由ET/jT反推并假设各方向同样阻尼，不是独立异常碰撞测量。匹配全局率的PIC非无调参预测；速度从磁场求导、密度与离子流近似得，非直接电子谱测量。选择四极事件限制普适性，长度误差大不能据其图拟合严格反比例。

\paragraph{本报告判断}详细报告实验直接量、反演量、选样、分辨率、边界及未闭合欧姆项；第一章给摩擦模型推导，第五章说明同尺度指数并不等于电子层绝对厚度/速度已获验证。

\section{L122：2008年}\label{paper:L122}

\textbf{Hall resistive tearing mode: A variational formulation}\citep{L122}\

Bhimsen K. Shivamoggi.

\textsf{核读范围：}arXiv0801.3453v1全部8页，无图，读完变分泛函与两支渐近推导；\textsf{图表：}无图；页6–7直接视觉核验式33–44，独立对式38作弱Hall展开.

\href{run:papers/L122.pdf}{打开本地PDF：\texttt{papers/L122.pdf}}\quad\href{https://doi.org/10.1209/0295-5075/83/55002}{DOI／版本入口}

\paragraph{在报告中的研究脉络}\hyperref[story:onset]{启动、爆发与磁岛}。关联表示本报告的阅读线索，不自动表示这些论文互相引用。

\paragraph{实际模型}不可压、准中性、无电子惯性的电阻Hall四场线性撕裂，内层F≈x；Hall参数σ=di/a、ε=ηc\textsuperscript{2}τA/a\textsuperscript{2}。

\paragraph{计算／实验细节}解析Rayleigh–Ritz，高斯电流试探函数J(k)=exp(-μk\textsuperscript{2}/2)，假定μω/ε\ensuremath{\gg}1；没有网格/独立模拟，无显式Δ'外层匹配参数和连续波数扫描。

\paragraph{核实的发现}提出三层结构的变分近似，将电阻支与Hall支写为ω\textsuperscript{2}-√(εω)-σ\textsuperscript{2}/4=0，形式上给ω≈ε\textasciicircum{}1/3或σ/2及转换σ≈ε\textasciicircum{}1/3；线性b在x、y皆奇，对应四极场。

\paragraph{限制与未完成问题}不能把形式上η→0有限增长当完整理想Hall断线证明：电子非理想项仍来自电阻，外层匹配与奇异极限未解决。式31含ωJ\textsuperscript{2}/k\textsuperscript{2}，高斯J(0)=1使普通积分发散，但式33使用有限值而未说明正则化/匹配。式36–37给μ\textsuperscript{2}=ωε，电阻支上μω/ε=1而非\ensuremath{\gg}1，不能同时宣称该渐近式严格恢复该支。原文式43写+σ\textsuperscript{2}ε\textasciicircum{}1/6/6；直接展开式38应为+σ\textsuperscript{2}/(6ε\textasciicircum{}1/3)，属于所读版本代数不一致。称ε\textasciicircum{}1/3为FKR亦缺小Δ'条件，不宜替代经典常ψ的S\textasciicircum{}-3/5标度。

\paragraph{本报告判断}记录为有启发但需审查的早期近似，报告不给它与完整匹配/数值本征解同等证据权重；用作学习如何检查近似适用条件的例子。

\section{L123：2008年}\label{paper:L123}

\textbf{Influence of the Hall effect on the reconnection rate at line-tied magnetic X-points}\citep{L123}\

I. J. D. Craig;  Yuri E. Litvinenko.

\textsf{核读范围：}A\&A484, 847–850（2008），物理页1–4全部，含线性方程、拟合、讨论与参考文献；\textsf{图表：}图1–2与式(13)–(16)视觉核验；重新换算拟合的系统长度指数.

\href{run:papers/L123.pdf}{打开本地PDF：\texttt{papers/L123.pdf}}\quad\href{https://doi.org/10.1051/0004-6361:200809569}{DOI／版本入口}

\paragraph{在报告中的研究脉络}\hyperref[story:onset]{启动、爆发与磁岛}。关联表示本报告的阅读线索，不自动表示这些论文互相引用。

\paragraph{实际模型}冷、可压、线性化电阻Hall X点，忽略电子惯性和压力。平衡ψE=xy，研究小扰动在刚性线缚方盒中的振荡衰减；不是非线性开放边界Hall层。

\paragraph{计算／实验细节}[-1, 1]\textsuperscript{2}、δψ初始=.1(1-x\textsuperscript{2})(1-y\textsuperscript{2})，墙u=0、δψ=0、∂nZ=0。η1.5e-4至2e-2，di/√η=.1–.9，ν=.1η，20余例；最终拟合排除高η数据。未给格点数/步长/离散算法的可复现明细。

\paragraph{核实的发现}Hall经沿平衡场的电流梯度产生Z，再驱动轴向W，增强扰动衰减。拟合log10(Δα/α0)=-.16+1.89log10di-.31log10η；无Hall参考α0≈π\textsuperscript{2}/(2ln\textsuperscript{2}η)。作者忽略弱η依赖后提出相对衰减增强随L\textasciicircum{}-2衰减，表明边界和尺度可改变Hall收益。

\paragraph{限制与未完成问题}实际di/√η<1，只覆盖弱至过渡Hall的小扰动。α是模幅衰减率，能量指数为2α；不能直接与非线性Ez/(BVA)比较。若保持物理n、B和磁扩散率固定，di\ensuremath{\propto}L\textasciicircum{}-1且η\ensuremath{\propto}L\textasciicircum{}-1，完整拟合给L\textasciicircum{}-1.58，而非严格L\textasciicircum{}-2；后者额外丢弃η依赖。大尺度和强Hall外推缺直接验证。

\paragraph{本报告判断}结论重点是具体线缚线性问题中的尺寸依赖；不应把近似L\textasciicircum{}-2包装成已推翻所有快Hall重联的普遍定理。报告保留作者结论并显式说明拟合与简化步骤。

\section{Malyshkin2008：2008年}\label{paper:Malyshkin2008}

\textbf{Model of Hall Reconnection}\citep{Malyshkin2008}\

Leonid Malyshkin.

\textsf{核读范围：}arXiv0810.3278v2 1-5全部，含发表后说明；\textsf{图表：}图1及式11-26，PDF3-4页视觉核验.

\href{run:papers/Malyshkin2008.pdf}{打开本地PDF：\texttt{papers/Malyshkin2008.pdf}}\quad\href{https://doi.org/10.1103/physrevlett.101.225001}{DOI／版本入口}

\paragraph{在报告中的研究脉络}\hyperref[story:global]{层长、驱动与全局完成}；\hyperref[story:closure]{压力、热流与闭合}。关联表示本报告的阅读线索，不自动表示这些论文互相引用。

\paragraph{实际模型}二维不可压、标量压力、均匀电阻、无导引无电子惯性Hall MHD；准稳态局地泰勒展开。

\paragraph{计算／实验细节}没有新数值模拟；外给Bm及L\textsuperscript{2}=-2Bm/(dyy By)\_up；电阻用磁扩散单位；S=LVA/eta。

\paragraph{核实的发现}Ez/(BmVA)约sqrt(sqrt(3)/S+2di\textsuperscript{2}/L\textsuperscript{2})。Hall强时电子层delta约eta L/(VA di)、离子层Delta约di，率约di/L。注意电子层上游离子Vin约VA L/(Sdi)很小，而离子层外磁通输入VR约VA di/L；两者不能混称。

\paragraph{限制与未完成问题}L和Bm必须由全局解确定，未由本理论求出；小电阻极限电子电流发散且电子惯性/压力张量被删去。条件Ez远小于BmVA且无层内不稳定性；不能外推di/L约1。热化密度反馈仅讨论，没有求解能量方程。

\paragraph{本报告判断}给出局地Hall层的连接标度，未解决全局率选择；正文中的collisionless不能读成X点无需非理想项。

\section{Toth2008：2008年}\label{paper:Toth2008}

\textbf{Hall magnetohydrodynamics on block-adaptive grids}\citep{Toth2008}\

Gábor Tóth;  Yingjuan Ma;  Tamas I. Gombosi.

\textsf{核读范围：}1-23（作者稿全部）；\textsf{图表：}图1-9、表1-2、式14-35；主要结果图表视觉核验.

\href{run:papers/Toth2008.pdf}{打开本地PDF：\texttt{papers/Toth2008.pdf}}\quad\href{https://doi.org/10.1016/j.jcp.2008.04.010}{DOI／版本入口}

\paragraph{在报告中的研究脉络}\hyperref[ch:numerics]{模型与数值可信度}。关联表示本报告的阅读线索，不自动表示这些论文互相引用。

\paragraph{实际模型}保守Hall MHD，忽略电子压强梯度与惯性；BATSRUS三维块AMR、广义坐标、MUSCL；显式RK2或BDF2及Newton-Krylov-Schwarz；8-wave散度控制。

\paragraph{计算／实验细节}AMR面电流需三阶磁场ghost插值得二阶curl，细网格通量覆盖粗面。GEM25.6×12.8，64×128，dx=.4, dz=.1, eta=.005, psi扰动.1或.5，z零梯度开边界；并行测试约250万格。

\paragraph{核实的发现}whistler波速随1/dx，直接一阶LLF不一致；minmod显式退为一阶，光滑区MC保二阶；隐式W=0减少通量黏性。GEM小种子分裂成双X点与中间岛，较大种子保持单X点但通量率相近。三维磁尾有四极场、准周期plasmoid。

\paragraph{限制与未完成问题}GEM有显式电阻且电子层不在模型中；三维物理解读留待后续，不能当成FTE机制证实。强扩展显式到1616核好，隐式约256核后下降；32核隐式14倍、混合20倍是具体机器算例。

\paragraph{本报告判断}展示离散与AMR的一致性问题，拓扑细节可对初始种子敏感而通量率不敏感；需要分开数值验证与物理验证。

\section{Drake2009Pickup：2009年}\label{paper:Drake2009Pickup}

\textbf{Ion heating resulting from pickup in magnetic reconnection exhausts}\citep{Drake2009Pickup}\

J. F. Drake;  M. Swisdak;  T. D. Phan;  P. A. Cassak;  M. A. Shay;  S. T. Lepri;  R. P. Lin;  E. Quataert;  T. H. Zurbuchen.

\textsf{核读范围：}PDF 1–14，正文§1–6、方法、结果、讨论、附录A、参考文献全部；图1–7、表I和关键公式直接查看PDF；\textsf{图表：}PDF 3–10图1–7、表I；PDF 7式(7)–(10)，PDF 12–13附录A；图6斜率0.39，图5与表I需区分标量与径向温度.

\href{run:papers/Drake2009Pickup.pdf}{打开本地PDF：\texttt{papers/Drake2009Pickup.pdf}}\quad\href{https://doi.org/10.1029/2008JA013701}{DOI／版本入口}

\paragraph{在报告中的研究脉络}\hyperref[story:energy]{电子独有重联与热化}。关联表示本报告的阅读线索，不自动表示这些论文互相引用。

\paragraph{实际模型}反平行二维p3d PIC + 2D/1D场测试离子 + de Hoffmann–Teller参考系解析 + Wind/ACE观测；不是Hall MHD温度闭合

\paragraph{计算／实验细节}双周期Harris，Lx×Ly=204.8×102.4 dp，半宽1.5 dp，网格0.05 dp，质量比25，c/vA=15，背景密度0.2，Te=1/12、Ti=5/12 mpvA\textsuperscript{2}，低密度区100粒子/格；无系统质量比/网格收敛序列。测试轨道取下游近稳态场。

\paragraph{核实的发现}入流跨窄Hall场边界后失磁化、拾取状增能与反向束；冷入流解析ΔTi≈mi v0\textsuperscript{2}/3，预测的各向异性并不准确。Wind 22次高剪切穿越支持平方标度但系数仅预测的0.39；ACE重离子径向温度有质量标度线索。

\paragraph{限制与未完成问题}原文保留各向异性散射、远下游慢激波结构与膨胀解释；不提供电子分配，不能解释全部高能尾或强导场阈值。SWICS 12 min采样、只测Trr且场方向改变；未读取原始数据/运行代码/前序仪器方法。

\paragraph{本报告判断}补足拾取参照及定量不符；与Yoon2023对电场净功和主导机制的归因有实质差异，尚未在同一场/轨道/时间窗复算，不能宣布统一或谁已被排除。末速度恢复不能单独代替电功积分。

\section{L137：2009年}\label{paper:L137}

\textbf{ANALYTICAL DESCRIPTION OF STEADY MAGNETIC RECONNECTION IN HALL MAGNETOHYDRODYNAMICS}\citep{L137}\

Yuri E. Litvinenko.

\textsf{核读范围：}ApJ694, 1464–1470（2009）全部7页，含局域轴向场附录A；\textsf{图表：}图1几何、式(24)–(48)、附录A1–A4；页5的O点展开和层长表达直接视觉检查.

\href{run:papers/L137.pdf}{打开本地PDF：\texttt{papers/L137.pdf}}\quad\href{https://doi.org/10.1088/0004-637x/694/2/1464}{DOI／版本入口}

\paragraph{在报告中的研究脉络}\hyperref[story:global]{层长、驱动与全局完成}。关联表示本报告的阅读线索，不自动表示这些论文互相引用。

\paragraph{实际模型}稳态二维不可压无黏电阻Hall、标量压力；以横向导数远大于纵向导数作弱二维展开，预设Harris形J0和轴向Z1=b0tanh(x/l)。不是全局初边值解。

\paragraph{计算／实验细节}无新模拟；ψ、φ、W、Z按y幂展开，给定E及轴向梯度b0。半厚l=η/(E+di b0)，流入φ1=-E tanh；太阳参数仅作Hall模型失效量级估算。

\paragraph{核实的发现}Hall把平面流与四极轴向场耦合，使层厚、有限纵向尺度和流出相互联系。正常Hall极性b0>0时，低阶中心磁通竟为O点而非X点；作者和审稿人明确讨论其可能是数学伪象。四极轴向场在此预设流场下即使di=0也可存在，四极形状本身不唯一认证Hall。附录改用局域Z1=b0tanh×sech\textsuperscript{2}后，厚度改为η/E、流入不反向需E>di b0，且远处轴向速度不再发散。

\paragraph{限制与未完成问题}弱二维假设恰在中心可能失效，O点不能当作被验证的新物理岛；主文Z1远处不消失导致W无界。附录未完成二阶ψ2和新层长。E仍外给，0.1需额外采用w≈10di和Ba≈1；未从模型选出普适率。太阳参数下薄层可能小于de，使无电子惯性模型失效。

\paragraph{本报告判断}模型假设影响结论的清楚案例：只读摘要会漏掉中心拓扑和附录改变标度。宜用来教学生检查ansatz、适用域和局地解的全球匹配，不把形式解析解等同真实重联结构。

\section{L145：2009年}\label{paper:L145}

\textbf{On the physical interpretation of Malyshkin’s (2008) model of resistive Hall magnetohydrodynamic reconnection}\citep{L145}\

Dmitri Uzdensky.

\textsf{核读范围：}arXiv0901.0520v1物理页1–2全文；页眉自动日期2018不是出版年，期刊2009；\textsf{图表：}无图；式(1)–(6)全部重算，页1视觉核对.

\href{run:papers/L145.pdf}{打开本地PDF：\texttt{papers/L145.pdf}}\quad\href{https://doi.org/10.1063/1.3125819}{DOI／版本入口}

\paragraph{在报告中的研究脉络}\hyperref[story:global]{层长、驱动与全局完成}。关联表示本报告的阅读线索，不自动表示这些论文互相引用。

\paragraph{实际模型}稳态、二维、不可压、电阻Hall，电子质量忽略；外给电子层半长L、上游Bm、n和磁扩散率η。标量压力沿出流可能与磁力同阶但不改变量级论证。

\paragraph{计算／实验细节}解析论文，无新数值运行。用高斯单位，η=η'c\textsuperscript{2}/(4π)为磁扩散率，η'才是电阻率；坐标x入流、y出流、z电流。

\paragraph{核实的发现}电子出流由磁力与电阻拖曳平衡，非电子惯性控制。出流和面外方向都有同样电阻因子，消去得veyδ=diVA、δ=ηL/(diVA)、vin=diVA/L、cEz=BmVA di/L。因而速率不显含η不表示电阻没有作用，也不表示出流不由碰撞拖曳决定。Ez中的diVA不依赖离子质量。

\paragraph{限制与未完成问题}L与Bm仍是未决定的外部输入；若L随宏观尺度增长，即使η消去，率仍可能太慢。Bm不能直接替换成离子层外B0，必须匹配两层电流。固定L, Bm令η趋零使δ缩小、vey增大，终将失去忽略电子惯性的适用性。没有新模拟或独立尺寸扫描。

\paragraph{本报告判断}为已读Malyshkin2008、Simakov2009提供透明的力平衡解释与明确的未闭合问题。适合放入第一章逐步推导，并阻止把局地代数解称全局快率理论。

\section{L154：2009年}\label{paper:L154}

\textbf{The hall effect in magnetic reconnection: Hybrid versus Hall‐less hybrid simulations}\citep{L154}\

K. Malakit;  P. A. Cassak;  Michael A. Shay;  J. F. Drake.

\textsf{核读范围：}GRL36, L07107（2009）物理页1–5全文，方法、关Hall实验、分辨率和局域电阻对照均读；\textsf{图表：}图1–2全部；物理页2方程和页3时序/归一化图直接检查.

\href{run:papers/L154.pdf}{打开本地PDF：\texttt{papers/L154.pdf}}\quad\href{https://doi.org/10.1029/2009gl037538}{DOI／版本入口}

\paragraph{在报告中的研究脉络}\hyperref[story:fast]{局地快率与必要条件}。关联表示本报告的阅读线索，不自动表示这些论文互相引用。

\paragraph{实际模型}P3D二维混合，动理学离子、me=0、Te=0；Hall-less仅在推进B的EOHM删Hall项，粒子推进EION仍保留J×B以维持平均离子力。两种电场本来就不同；不能当作自洽双物种Maxwell模型。

\paragraph{计算／实验细节}204.8×102.4di、1024×512，dx=.2，100粒子/格、Ti=.5、上游n1，双Harris片宽1、Bg0、全周期、种子.3。均匀η=.015、η4=1e-3；关Hall时t300。另η=.007在dx=.2与.1作分辨对照；未给绝对时间步。

\paragraph{核实的发现}标准混合率约.03，无Hall约.015并符合实测上游SP预测；在t300关闭Hall后不到10离子回旋时间降至慢率，层延长。随后绝对率回升由上游堆积引起，除以上游B\textasciicircum{}(3/2)后与SP支相符。η=.007时半宽.52只有2.5格，出现网格岛与假快率；5格则恢复稳定慢率。额外采用局域cosh电阻，Hall-less也能快，解释与早期文献的差异。

\paragraph{限制与未完成问题}标准例电阻不小，η=0仍有η4，去掉η4则层塌到网格；不是无耗散证明。无Hall出岛对率影响小仅适用于本系统，未排除极大盒磁岛链。没有电子动理学和三维。图2d把Hall-less时间1700压到300，不能据曲线同横坐标比较启动时间。

\paragraph{本报告判断}强的模型消融和数值反例：必须同时控制均匀/局域耗散、网格和上游场；结论限定于冷无质量电子混合体系，不能外推成任何压力闭合下Hall普遍必要。

\section{Malyshkin2009：2009年}\label{paper:Malyshkin2009}

\textbf{Model of Two-Fluid Reconnection}\citep{Malyshkin2009}\

Leonid Malyshkin.

\textsf{核读范围：}arXiv0904.0660v3 1-4全部；\textsf{图表：}图1、式3-23与讨论；PDF2-3页视觉核验.

\href{run:papers/Malyshkin2009.pdf}{打开本地PDF：\texttt{papers/Malyshkin2009.pdf}}\quad\href{https://doi.org/10.1103/physrevlett.103.235004}{DOI／版本入口}

\paragraph{在报告中的研究脉络}\hyperref[story:global]{层长、驱动与全局完成}。关联表示本报告的阅读线索，不自动表示这些论文互相引用。

\paragraph{实际模型}二维准中性常密度双流体，标量压力、有限电子惯性、均匀有效电阻、准稳态、无导引。

\paragraph{计算／实验细节}无新模拟；扩展动量保留-de\textsuperscript{2}(j·grad)j；全局Bext、Lext给定，局地电子半长L与离子尺度区分。

\paragraph{核实的发现}惯性强支给delta约de、离子厚di、电子喷流约电子Alfven速度；Ez约eta Bext/de，L约VA de di/eta。离子在电子层内可很慢，在电子喷流减速区才继续加速；存在可同时容纳SP与较快支的参数窗。

\paragraph{限制与未完成问题}率显式正比电阻；collisionless命名指把eta解释为待动理学计算的有效电阻。Hall过渡支的长度未唯一求出，转变路径是假说；稳态模型不能描述迟滞过程。预测长度me\textasciicircum{}(1/2)与所引PIC的me\textasciicircum{}(3/8)不同；实验层宽约8倍差异仍未解释。

\paragraph{本报告判断}这是有残余有效耗散的稳态惯性理论，不与Andres的非稳态无电阻拓扑变化等价。

\section{Simakov2009：2009年}\label{paper:Simakov2009}

\textbf{Quantitative analytical model for magnetic reconnection in Hall magnetohydrodynamics}\citep{Simakov2009}\

Andrei N. Simakov;  Luis Chacòn.

\textsf{核读范围：}LANL2008-09-29预稿PDF1-34全部，期刊2009年PoP扩展篇；\textsf{图表：}式1-25逐页图像核对、两页35算例表及图1-4全部视觉核验.

\href{run:papers/Simakov2009.pdf}{打开本地PDF：\texttt{papers/Simakov2009.pdf}}\quad\href{https://doi.org/10.1063/1.3077269}{DOI／版本入口}

\paragraph{在报告中的研究脉络}\hyperref[story:global]{层长、驱动与全局完成}；\hyperref[story:closure]{压力、热流与闭合}。关联表示本报告的阅读线索，不自动表示这些论文互相引用。

\paragraph{实际模型}二维不可压高beta Hall MHD；无电子惯性、无导引；均匀电阻加超电阻；离散区域解析近似加非线性岛合并。

\paragraph{计算／实验细节}Fadeev平衡epsilon=.2，35算例di=0至1、eta=0至1e-3；表中etaH最高7.6e-6（正文写7.6e-7不一致）。解析电流形状系数mu=16；在峰值时测电子片厚/长、上游Bin，比较半高宽与速度峰两定义。本文未列实际网格、时间步和边界截断，数值算法指向Chacon-Knoll2003。

\paragraph{核实的发现}统一匹配SP与电子Hall极限，强Hall局地率约sqrt(2)di/w；电子层细长也可耗散系数不显含的快率，超电阻造成开放X形却仍可处于慢支；有限四极Bz本身不判快。35算例归一化率跨度约1200倍，理论比值范围.5-2.3；固定这组岛合并的强Hall区率大致随di上升。

\paragraph{限制与未完成问题}一般过渡公式是两渐近极限的插值，层长w和Bin从模拟取值，并未预测宏观驱动选择。忽略电子惯性产生过快电子喷流和大于动理学常见值的局地率；不能无条件推至di无限大。只比较峰值附近准稳态，未完成整段时史理论。预稿题名与PRL2008短文相近但不是同一篇。

\paragraph{本报告判断}是电子层几何与率的有用定量检验，不是全部参数都由初值唯一预测的全局理论。

\section{Sullivan2009：2009年}\label{paper:Sullivan2009}

\textbf{Extension of the electron dissipation region in collisionless Hall magnetohydrodynamics reconnection}\citep{Sullivan2009}\

Brian P. Sullivan;  A. Bhattacharjee;  Yi-Min Huang.

\textsf{核读范围：}arXiv0905.4765v3 1-11全部；\textsf{图表：}图1-13；图4-6、8-13视觉复核.

\href{run:papers/Sullivan2009.pdf}{打开本地PDF：\texttt{papers/Sullivan2009.pdf}}\quad\href{https://doi.org/10.1063/1.3251423}{DOI／版本入口}

\paragraph{在报告中的研究脉络}\hyperref[story:global]{层长、驱动与全局完成}。关联表示本报告的阅读线索，不自动表示这些论文互相引用。

\paragraph{实际模型}等温压缩性双流体；参考密度常系数B'惯性；超电阻断线，密度和动量另含四阶扩散。

\paragraph{计算／实验细节}周期12.8di方盒、512\textsuperscript{2}、dx.025；四岛正弦平衡、总温1、beta最小11.45、V0.1；全部初始电流由离子携带；mu4=5.1e-7；扫描eta\_H约3.9e-5至3.9e-3及mi/me25-200；另盒6.4和25.6。

\paragraph{核实的发现}无惯性扫描L\_e约eta\_H\textasciicircum{}(1/3)，delta\_e约eta\_H\textasciicircum{}(4/9)，局地归一化率随eta\_H变，而原始率和全局归一化率约.1-.2且弱依赖；惯性扫描长度可按(me/mi)\textasciicircum{}(3/8)，厚度近1/3幂，长宽比弱变化。

\paragraph{限制与未完成问题}无惯性超电阻扫描最长仍不到di，未复现动力学长约10di片及大量次级岛；最低eta层接近网格不可靠；3/8与1/3在一数量级质量窗无法区分；真实质量比15de长度和1de宽度为外推；边界、导引与惯性超电阻同时重要区未扫描。

\paragraph{本报告判断}重要发现是电子层长度也依赖断线机制，以及局地与全局归一化可能给出不同表述；不是Hall成功复现任意长电子层的证明。

\section{Cassak2010：2010年}\label{paper:Cassak2010}

\textbf{A saddle-node bifurcation model of magnetic reconnection onset}\citep{Cassak2010}\

P. A. Cassak;  Michael A. Shay;  J. F. Drake.

\textsf{核读范围：}062105-1至7（全部）；\textsf{图表：}图1-8；图像复核待补；主要结果页已视觉复核.

\href{run:papers/Cassak2010.pdf}{打开本地PDF：\texttt{papers/Cassak2010.pdf}}\quad\href{https://doi.org/10.1063/1.3435269}{DOI／版本入口}

\paragraph{在报告中的研究脉络}\hyperref[story:onset]{启动、爆发与磁岛}。关联表示本报告的阅读线索，不自动表示这些论文互相引用。

\paragraph{实际模型}以片宽为变量的鞍结分岔现象学模型；F3D双流体与无电子惯性Hall对照。

\paragraph{计算／实验细节}409.6×204.8di, dx=.1；有限质量比25的旧例D4=2e-5，无惯性新例D4=8e-5作用全部方程，初始片宽1.2di, T=1, eta=.012，磁种子.004与随机5e-6。

\paragraph{核实的发现}无惯性下相同eta仍有快慢支与迟滞；片宽相图有非单调结构，电子出流随层变薄增加，与k约1/delta的双流体波速相关。

\paragraph{限制与未完成问题}模型没有自洽给出层长；因上游场随时间变，正文明确不做片宽演化定量比较。无惯性算例保留更强超电阻，不能称无耗散。Hall-电阻平衡必然不稳定的论断后来被Huang2011在其他几何中反驳。

\paragraph{本报告判断}强证据是关闭电子惯性仍保留历史依赖；对任意电阻Hall层不稳定的外推需要保留争论。2010称已做收敛，不等于2005正文展示了独立收敛曲线。

\section{E469：2010年}\label{paper:E469}

\textbf{Nonlinear gyrofluid simulations of collisionless reconnection}\citep{E469}\

D. Grasso;  E. Tassi;  François L. Waelbroeck.

\textsf{核读范围：}HAL封面加作者稿1-16页，共17页全读；\textsf{图表：}模型式1-26及图1-6；全部图像视觉核验.

\href{run:papers/E469.pdf}{打开本地PDF：\texttt{papers/E469.pdf}}\quad\href{https://doi.org/10.1063/1.3475440}{DOI／版本入口}

\paragraph{在报告中的研究脉络}\hyperref[story:onset]{启动、爆发与磁岛}。关联表示本报告的阅读线索，不自动表示这些论文互相引用。

\paragraph{实际模型}二维强导引场Hamilton陀螺流体三场；演化离子导心密度、电子密度及电子广义磁通；Pade极化与陀螺平均。

\paragraph{计算／实验细节}[-pi, pi]\textsuperscript{2}周期，1024×128；de.2，rho\_s/rho\_i取.01/.2/.4；平衡sech\textsuperscript{2}x截断23个Fourier模。均匀导心密度，离子密度四涡扰动使电子初始扰动0；总能损耗<1\%内视为可信。

\paragraph{核实的发现}有限离子Larmor半径使静电势保留电子密度的细丝，虽陀螺平均势仍平滑；分离线附近垂直电场峰增强、平行电场出现细结构。岛内密度趋平，细丝更集中于分离线；改变能量通道，并在所测条件抑制冷等离子体中的二次KH。交换rho\_s与rho\_i时岛形/率可相近，而广义不变量和势结构不同。

\paragraph{限制与未完成问题}0.4离子半径、.01声半径一组作者自称非物理隔离极限；不等于真实热电子/离子空间等离子体。没有分布函数，电子束与加热是流体结构的暗示，不能量化非热加速。小尺度不断形成，无网格扫描与完整离散算法，能量误差<1\%不能单独证明细丝收敛。式13明确磁通由G+与G-之和恢复；结论处误写difference；文中inverse Kelvin-Helmholtz operator实际指Helmholtz逆算子。

\paragraph{本报告判断}新发现是相近总重联发展掩盖显著不同的电场/密度细结构；保守模型能量项不能自动解释为不可逆热化。

\section{E472：2010年}\label{paper:E472}

\textbf{Gyrofluid simulations of collisionless reconnection in the presence of diamagnetic effects}\citep{E472}\

E. Tassi;  François L. Waelbroeck;  D. Grasso.

\textsf{核读范围：}HAL封面及JPCS260, 012020(2010)全文9页，共10页；\textsf{图表：}表1、图1–2、式1–12；页6/8磁岛及六场结构直接视觉核验.

\href{run:papers/E472.pdf}{打开本地PDF：\texttt{papers/E472.pdf}}\quad\href{http://porto.polito.it/2380084/1/2380084.pdf}{DOI／版本入口}

\paragraph{在报告中的研究脉络}\hyperref[story:onset]{启动、爆发与磁岛}。关联表示本报告的阅读线索，不自动表示这些论文互相引用。

\paragraph{实际模型}三场Hamiltonian gyrofluid，ni纯gyro平均E×B输运，无离子平行流；ne和惯性磁通耦合，等温、低β强导向，无Landau阻尼；线性平衡密度梯度提供抗磁漂移。

\paragraph{计算／实验细节}双周期的是扰动，背景n'0x本身不周期；x\ensuremath{\in}[-π, π], y\ensuremath{\in}[-aπ, aπ], a=1.5/2，1024×128，非线性升至1024×256；de=.2，ρi=.35/.4, ρs=.2/.4，v*=-.04, -.16, -.4, -.64。线性种子10\textasciicircum{}-6，非线性10\textasciicircum{}-3。

\paragraph{核实的发现}抗磁漂移降低线性增长，v*=-.64例稳定；Ti>Te时岛沿离子抗磁方向传播，理论ω=(ω*e+ω*i)/2，γ\textsuperscript{2}=γ0\textsuperscript{2}-(ω*e-ω*i)\textsuperscript{2}/4。强梯度非线性岛更窄，分离线电子密度瓣不对称。热离子极限φ含电子细丝和离子gyrocenter大尺度两项，不能再简单φ\ensuremath{\propto}ne。

\paragraph{限制与未完成问题}表1增长率与渐近理论差15\%–38\%，频率量值也有明显偏差，不能称精确吻合；Δ'增大改善趋势但ρi/de仅约2。有限幅岛的所谓饱和未展示长期收敛。强导向假设下作者明确不直接应用磁层顶；不能与L174出流阻断准则混为一谈。没有粒子热化/平行动量反馈。

\paragraph{本报告判断}与L174比较线性抗磁稳定、岛旋转、小幅饱和和自发终止是不同问题。

\section{Katz2010：2010年}\label{paper:Katz2010}

\textbf{Laboratory Observation of Localized Onset of Magnetic Reconnection}\citep{Katz2010}\

Noam Katz;  Jan Egedal;  Will Fox;  Ari Le;  Jeff Bonde;  Arturs Vrublevskis.

\textsf{核读范围：}MIT封面+PRL4页，共5页全读；\textsf{图表：}图1-5及式1-5；核心图2-5、公式视觉核对.

\href{run:papers/Katz2010.pdf}{打开本地PDF：\texttt{papers/Katz2010.pdf}}\quad\href{https://doi.org/10.1103/physrevlett.104.255004}{DOI／版本入口}

\paragraph{在报告中的研究脉络}\hyperref[story:diagnostic]{实验与可区分诊断}。关联表示本报告的阅读线索，不自动表示这些论文互相引用。

\paragraph{实际模型}VTF强导引场环形氩实验；三维电势和离子极化电流半经验模型，非Hall模拟论文。

\paragraph{计算／实验细节}n约2e18/m\textsuperscript{3}、Te15eV、Ti1eV；主导环向场56mT在R1m。感应电场测6环向位置，电流密度8cm分辨率；浮动势作等离子体势近似。模型b'.035T/m, l0=1.7m, 椭率1.7；电势幅6V、宽.12m由实验拟合；只早期10-20μs。

\paragraph{核实的发现}轴对称外部驱动仍出现局域起始，感应电场双向绕环传播约2倍Alfven速度；加入静电场后的总电场仍局域。快重联仅在测到外区q约2或3时，本文详析q2；旋转电势允许极化电流闭合局域电流下降。半经验增长时间22μs，与实验20±6μs一致；二维先前对应振荡，三维符号改变可给指数增长。

\paragraph{限制与未完成问题}离子仅边缘磁化；浮动势近似、均匀密度、线性几何与拟合势使分离线不吻合。增长模型使用实测电场空间形状和势/感应电场共同增长关系，并非无参数第一性原理预测。约2VA不应抹成精确VA，更非验证Hall唯一驱动；局域微观冻结破坏未由本篇完整辨析。

\paragraph{本报告判断}提供全局几何、电流闭合与三维启动的实验约束；比较Hall模型时必须把感应场与总场分开。

\section{L159：2010年}\label{paper:L159}

\textbf{Comparison of Secondary Islands in Collisional Reconnection to Hall Reconnection}\citep{L159}\

L. S. Shepherd;  P. A. Cassak.

\textsf{核读范围：}arXiv1006.1883v2对应2010 PRL105, 015004，物理页1–4全部；\textsf{图表：}图1–4、三阶段判据；页3的磁岛排出时空图和多岛对照视觉核验.

\href{run:papers/L159.pdf}{打开本地PDF：\texttt{papers/L159.pdf}}\quad\href{https://doi.org/10.1103/physrevlett.105.015004}{DOI／版本入口}

\paragraph{在报告中的研究脉络}\hyperref[story:onset]{启动、爆发与磁岛}。关联表示本报告的阅读线索，不自动表示这些论文互相引用。

\paragraph{实际模型}二维电阻Hall/F3D、电子惯性、恒温、无导引场；另从相同状态关Hall和电子惯性作电阻MHD对照。

\paragraph{计算／实验细节}819.2×409.6di双周期双Harris片，初始宽12di、质量比25、η=.008、T1、上游n1，种子.016并加随机磁扰动打破排岛对称。dx=.2和.1比较，展示后者，D4=1.75e-4。Sg约1e5，LSP约200；未列绝对步长和噪声幅度。

\paragraph{核实的发现}明确分离三阶段：t约700、片半高半宽约2di时先出次级岛，t约1780、半宽约di时再进入Hall。率约.004→.005→数倍至近一数量级提高。Hall启动后短强层将原有次级岛及竞争X线排走；另一张片三岛例也显示排出。t1120关Hall/惯性后率最终约.009，符合N=2的SP×√N估计，未达到Hall支。

\paragraph{限制与未完成问题}只有该S、几何和有限岛数，不能推出任意高S磁岛链必慢于Hall、也不能推出Hall与磁岛普遍不能共存。D4减小使Hall率略增；质量比25未扫描。恒温忽略Ohmic热反馈，缺导引场/三维/Spitzer温度依赖。时空图每70时间单位采样再插值，不能由平滑斜线提取更细瞬态。冠状观测的阶段对应只为猜想。

\paragraph{本报告判断}与后续Huang2011、Shi2024等联读：此处是Hall排出现有岛，后续研究另出现Hall层持续生岛或往复转换；差别需按参数和阶段说明。

\section{L160：2010年}\label{paper:L160}

\textbf{Fast and slow two-fluid magnetic reconnection}\citep{L160}\

Leonid M Malyshkin.

\textsf{核读范围：}arXiv1005.3760v1物理页1–16全部；2010 Physica Scripta长文，扩展已读2009 PRL推导，不算独立模拟验证；\textsf{图表：}式(1)–(53)、表1和图1–2；表1及图2直接视觉核验，中心惯性/压力项逐项核对.

\href{run:papers/L160.pdf}{打开本地PDF：\texttt{papers/L160.pdf}}\quad\href{https://doi.org/10.1088/0031-8949/2010/t142/014033}{DOI／版本入口}

\paragraph{在报告中的研究脉络}\hyperref[story:global]{层长、驱动与全局完成}。关联表示本报告的阅读线索，不自动表示这些论文互相引用。

\paragraph{实际模型}准中性不可压非相对论双流体，任意Z和质量比，标量物种压力、常电阻、无黏性，二维对称准稳态。保留中心质量动量中的电子相对漂移应力及广义欧姆惯性。

\paragraph{计算／实验细节}解析论文，无计算网格。外给Bext、Lext、η；电子层δ, L与离子层Δ, Lext分开。c=4π=1；γtilde是Hall/理想项比，不是时间增长率。快支要求Lext/de\ensuremath{\ll}S≤Lext\textsuperscript{2}/(de di)。

\paragraph{核实的发现}从两种粒子动量严格重组得附加ω+\textsuperscript{2}de\textsuperscript{2}(j·\ensuremath{\nabla})j；若把它吸入共同参考系压力，压力便不再标量。对称稳态X点仍Ez=ηj0，导出δ≥约de、快支δ≈de、Δ≈di、L≈VA de di/η、Ez≈ηBext/de。电子先达到电子Alfvén喷流，离子在更远下游继续加速。质量相等的对等离子体Hall系数消失，本模型却错误给慢SP，作者归因遗漏压力各向异性。

\paragraph{限制与未完成问题}快支仍正比有效η，不是η=0惯性稳态断线；η需另由动理学决定。中间Hall支的L不能求出，把其安排在固定S阈值上作为快慢转换属于猜想。图2无分支稳定性、转换时长或迟滞计算；不能作动态证据。B在两层边缘近相等、压力曲率和薄层均为假设；长电子喷流外冻结，与PIC压力非理想区不完全相符。提出Hall-Petschek激波但未完整推导其双流体跳跃条件。

\paragraph{本报告判断}第一章可用其完整质量重组避免漏动量项；结论章应把静态可行分支与Cassak的实测迟滞分开。它解释特定电阻/惯性平衡，未统一决定实际无碰撞快率。

\section{L164：2010年}\label{paper:L164}

\textbf{Slow shocks and conduction fronts from Petschek reconnection of skewed magnetic fields: two-fluid effects}\citep{L164}\

D. W. Longcope;  S. J. Bradshaw.

\textsf{核读范围：}PDF 1–45 including all Sections 1–7 and references;  no internal appendix；\textsf{图表：}Direct PDF visual: Eqs.4–7 (p10),  Eqs.14–18 (p12),  Figs.4–7, 9–14 (pp17, 19, 20, 24, 26, 28, 30, 32, 35, 37). Remaining figures/captions read in text..

\href{run:papers/L164.pdf}{打开本地PDF：\texttt{papers/L164.pdf}}\quad\href{https://doi.org/10.1088/0004-637X/718/2/1491}{DOI／版本入口}

\paragraph{在报告中的研究脉络}\hyperref[story:energy]{电子独有重联与热化}。关联表示本报告的阅读线索，不自动表示这些论文互相引用。

\paragraph{实际模型}采用准中性、共同密度和共同速度的一维两温碰撞流体描述平行激波，再映射到外部给定的Petschek偏斜磁场重联出流。共同速度来自本问题无电流的特殊约化，并非一般Hall假设。保留离子黏性、两物种经典热导及库仑温度交换；电子惯性忽略，压力梯度电场维持电子力平衡。它不是自洽Hall扩散区模拟。

\paragraph{计算／实验细节}例子包含入流马赫数2.4、3.6、5，初始双曲正切速度层宽1–500个归一化离子黏性长度；电子和离子Prandtl数0.013、0.33，交换参数0.019。输运系数随温度的5/2次方变化，交换频率随nTe\textasciicircum{}(-3/2)变化。中心镜面、远端持续入流；稳态用射击积分，时变网格和时间步未报告。偏斜角约≤120°时几何约化较可靠。

\paragraph{核实的发现}离子黏性子激波先把出流动能转为离子内能，随后长下游区域的碰撞交换加热电子，电子传导又形成前导热波；三个区域及相关时间尺度明显分离。早期离子温度可超过最终平衡值，电子升温滞后，而远下游最终跳跃仍可与单温结果一致。初始速度层足够宽时，局地两温大小关系也可改变，不能压成固定加热比。

\paragraph{限制与未完成问题}原文发现部分强激波早期离子热流达到或超过自由流值，黏性修正甚至违反压力正定条件，故经典闭合恰在关键加热层可能失效。电子热流较小与未包含X点直接电子加热有关，不能外推到所有重联。模型不含自洽Hall场、完整动理学尾、动态磁通管截面和反平行极限，也没有进行实验Ti/Te定量验证。

\paragraph{本报告判断}这是两温流体可研究热化过程的清楚实例：它把离子流能耗散、物种交换和电子传导在空间与时间上分开，同时主动揭示经典离子闭合的失效。可据此说明“双流体并非不能研究加热”，但不能把给定出流的碰撞激波结果当作无碰撞Hall重联的普适分配律。文中与拾取加热等动理学机制的联系应作为模型边界保留。

\section{L166：2010年}\label{paper:L166}

\textbf{Three-dimensional simulations of the orientation and structure of reconnection X-lines}\citep{L166}\

Ron Schreier;  M. Swisdak;  J. F. Drake;  P. A. Cassak.

\textsf{核读范围：}arXiv1009.1650v1物理页1–12全部，含方法、结果、讨论和参考文献；\textsf{图表：}图1–6及式(1)–(7)；图4电流等值面与图6角度峰值直接视觉核验，注意图4横纵轴比例不同.

\href{run:papers/L166.pdf}{打开本地PDF：\texttt{papers/L166.pdf}}\quad\href{https://doi.org/10.1063/1.3494218}{DOI／版本入口}

\paragraph{在报告中的研究脉络}\hyperref[story:three]{三维传播与有限X线}。关联表示本报告的阅读线索，不自动表示这些论文互相引用。

\paragraph{实际模型}三维可压缩Hall F3D、电子惯性、标量等温压力；非对称双Harris电流层，全周期，比较二维和三维X线方向选择。

\paragraph{计算／实验细节}B1=√3、B2=3，ρ1=4、ρ2=1，剪切角54.7°；场分量(1, 0, √2)和(-1, 0, 2√2)，w0=0.5di，T=1，质量比25。三维盒51.2×51.2×409.6di，网格各0.2di=de，四阶扩散0.001，无显式η、ν；二阶时间、四阶空间。随机磁扰动10\textasciicircum{}-3，kz最高5；无质量比及网格收敛扫描。

\paragraph{核实的发现}二维在Sonnerup磁剪切准则禁止的边界情形仍发生重联，率约0.02。三维初有多方向X线，后来少数占优；Canny/Hough对Jxz等值面投影测得α=35.9±8.5°，与最大出流预测34.3°相容。X线长度早期即停止增长，作者归因不同共振面磁岛相互作用，区别于反平行单共振面无限延伸。

\paragraph{限制与未完成问题}最大归一化率预测31.7°也在误差内，不能区分两种选择原则。X线由电流形态代理识别，部分场向电流结构不是重联；无三维零点。周期盒离散共振面k·B=0限制可选方向；单几何且无盒长扫描，有限长度结论不能普适化。de只有一格，无法确认电子层结构；固定纵横比率公式在三维未经证明。

\paragraph{本报告判断}三维发现是方向竞争与跨共振面相互限制；验证强度是排除一个过严准则、与另一预测相容，尚非唯一识别全局最优化机制。

\section{Baalrud2011：2011年}\label{paper:Baalrud2011}

\textbf{Hall magnetohydrodynamic reconnection in the plasmoid unstable regime}\citep{Baalrud2011}\

Scott Baalrud;  Amitava Bhattacharjee;  Yi-Min Huang;  K. Germaschewski.

\textsf{核读范围：}arXiv1108.3129v1 1-8全部；\textsf{图表：}图1-9，图8-9视觉复核.

\href{run:papers/Baalrud2011.pdf}{打开本地PDF：\texttt{papers/Baalrud2011.pdf}}\quad\href{https://doi.org/10.1063/1.3633473}{DOI／版本入口}

\paragraph{在报告中的研究脉络}\hyperref[story:onset]{启动、爆发与磁岛}。关联表示本报告的阅读线索，不自动表示这些论文互相引用。

\paragraph{实际模型}不可压标量压力、无电子惯性、电阻Hall四场；SP平衡宽度随S变化，扰动电阻内层d\_eta小于di而di小于delta\_SP。

\paragraph{计算／实验细节}线性本征对照S\_L=1e8, di/L=1e-5, 3e-5, 1e-4；非线性岛合并L=1, S=1e6, di=1e-5至1e-3；算法网格引用前作未完整列出。

\paragraph{核实的发现}最快模增长率正比于(di/L)的6/13次幂与S的7/13次幂，系数约3.2（以Gamma0=2vA/L归一）；最快波数kL正比于(di/L)的1/13次幂与S的11/26次幂，系数约2.9。全局层仍长薄时扰动已出现四极Bz。非线性增长转折在di/L约2e-4，与内电阻层约1.8e-4吻合。

\paragraph{限制与未完成问题}非线性计算只确认转折，作者承认无法精确验证6/13指数；gamma>>Gamma0, k delta\_SP<<1与三层排序限制适用范围；不是非线性最终率公式。PDF自动日期2018与arxiv2011不一致，按arxiv版本记录。

\paragraph{本报告判断}Hall先影响撕裂内层再改变全局平衡是重要发现；不能把线性指数和稳定单X点视为同一结论。

\section{E512：2011年}\label{paper:E512}

\textbf{Reduced fluid-kinetic equations for low-frequency dynamics,  magnetic reconnection,  and electron heating in low-beta plasmas}\citep{E512}\

A. Zocco;  A. A. Schekochihin.

\textsf{核读范围：}arXiv1104.4622v2全部28页，含低β排序、完整KREHM推导、能量、Hermite谱及附录A/B的二维不变量与线性双层匹配；\textsf{图表：}本文无结果图；直接查看页12/15/24原版公式，核验能量符号、Hermite截止和匹配系数，非仅摘要或综述转述.

\href{run:papers/E512.pdf}{打开本地PDF：\texttt{papers/E512.pdf}}\quad\href{https://doi.org/10.1063/1.3628639}{DOI／版本入口}

\paragraph{在报告中的研究脉络}\hyperref[story:energy]{电子独有重联与热化}。关联表示本报告的阅读线索，不自动表示这些论文互相引用。

\paragraph{实际模型}KREHM是由gyrokinetics作βe\textasciitilde{}me/mi\ensuremath{\ll}1次级展开的流体—动理学模型，不是纯双流体：强导场、低频、k平行\ensuremath{\ll}k垂直、de与ρs保留，电子FLR忽略；离子gi=0初值子类，保留Γ0=I0(b)exp(-b)极化、ui平行0。电子分布扣除密度/流矩后的ge保留平行速度结构，修正Lenard–Bernstein碰撞，磁扩散率νei de\textsuperscript{2}。

\paragraph{计算／实验细节}理论及算法方案，无本篇非线性算例网格或新模拟数据。Hermite g0=g1=0、g2=δT平行/(√2T0)、m≥3流送耦合相邻阶并受νm阻尼。推荐模式数必须越过mc=[3|k平行|vthe/(2√2ν)]\textasciicircum{}(2/3)，而非固定十矩足够；半碰撞极限κ平行=vthe\textsuperscript{2}/(2ν)，vthe=√(2Te/me)。

\paragraph{核实的发现}ge=0一般不是动力学不变子空间；平行温度项交换流体/电子自由能但本模型中不独自破坏磁通冻结，破冻结仍由η或电子惯性。Wfluid点=-Q-Ohmic，He点=Q-Dcoll，Q可正可负而Dcoll≥0。慢/稳态相混合Em\textasciitilde{}m\textasciicircum{}(-1/2)exp[-(m/mc)\textasciicircum{}(3/2)]使ν→0+耗散有限；快速线性增长有mγ=(|k平行|vthe/(2√2γ))\textsuperscript{2}截止，若mγ\ensuremath{\ll}mc则Dcoll随ν消失。附录系统给无碰撞/半碰撞单双层撕裂及两尺度匹配，低β边界明确。

\paragraph{限制与未完成问题}不能用于一般无导场高β电子扩散区；无离子Landau阻尼、垂直相混合/离子黏性；极冷离子还需βe\ensuremath{\ll}τ\textsuperscript{3}以免丢失离子碰撞排序。小于de的垂直奇异梯度未保证受电子加热通道正则化，可能需额外超扩散；有限加热是条件性理论预测而非本篇已算出的三维湍流事实。原版式110–112把定义为正耗散的Dcoll写负号，与83/96矛盾，本报告依定义重算取正。

\paragraph{本报告判断}第一章新增从动理学约化到Hermite级联/熵产生的学习推导，明确不属于纯双流体；第四章提出闭合应同时预测率、层结构及不可逆加热，不能用能量转移替代热化。

\section{Huang2011：2011年}\label{paper:Huang2011}

\textbf{Onset of fast reconnection in Hall magnetohydrodynamics mediated by the plasmoid instability}\citep{Huang2011}\

Yi-Min Huang;  A. Bhattacharjee;  Brian P. Sullivan.

\textsf{核读范围：}1-25（arXiv v2全部）；\textsf{图表：}图1-11、式1-6；主要结果图视觉核验；主要结果页已视觉复核.

\href{run:papers/Huang2011.pdf}{打开本地PDF：\texttt{papers/Huang2011.pdf}}\quad\href{https://doi.org/10.1063/1.3606363}{DOI／版本入口}

\paragraph{在报告中的研究脉络}\hyperref[story:onset]{启动、爆发与磁岛}。关联表示本报告的阅读线索，不自动表示这些论文互相引用。

\paragraph{实际模型}二维等温、标量压力、无电子惯性、电阻Hall MHD；HMHD五点差分/梯形leapfrog；psi、By演化保证二维散度。

\paragraph{计算／实验细节}S=5e5，L/di=2500, 5000, 10000，对照纯电阻；T=3，导体滑移边界，仅上半域，6400乘1024非均匀网格，最小dx=1.2e-4L, dz=1.4e-5L；另在Sdi=100扫描L/di500, 1000, 2500。

\paragraph{核实的发现}2500形成持久单X点；5000经历磁岛触发Hall、排岛、再拉长到约500di、再生岛的间歇循环，率.005-.03；10000未进入Hall。主SP宽大于di时，次级薄层仍可触发。较大系统在相同Sdi下片更长更宽、开放角更小。

\paragraph{限制与未完成问题}Hall分支仍由电阻断线；此循环与无碰撞PIC连续产岛机制不等同。宽度定义为四分之一峰值全宽，不能与半高半宽直接比（.7di对应.17di）。相图边界仅统计估计，Sdi临界随系统尺寸未得到定律；对太阳日冕/磁尾Spitzer参数该中间区未必适用。

\paragraph{本报告判断}破除Hall必然单X点的说法，但不能由形状相像认定动力学机制一致。小系统欧姆项与收敛支持电阻可单独维持Hall分支，作者与Cassak2010有实质争论。

\section{L156：2011年}\label{paper:L156}

\textbf{A Discontinuous Galerkin Method for Ideal Two-Fluid Plasma Equations}\citep{L156}\

John Loverich;  Ammar Hakim;  Uri Shumlak.

\textsf{核读范围：}arXiv1003.4542v1全部40页；期刊2011年9卷240–268，预印本2010年；\textsf{图表：}图1–19及全部算法公式；直接看页12、32、34、37的CFL、磁岛、重联通量及晚期射流结构.

\href{run:papers/L156.pdf}{打开本地PDF：\texttt{papers/L156.pdf}}\quad\href{https://doi.org/10.4208/cicp.250509.210610a}{DOI／版本入口}

\paragraph{在报告中的研究脉络}\hyperref[ch:numerics]{模型与数值可信度}。关联表示本报告的阅读线索，不自动表示这些论文互相引用。

\paragraph{实际模型}电子/离子各五矩Euler方程加完整Maxwell与双曲散度清理；保留位移电流、非准中性、电子惯性，无显式电阻粘性，标量压力γ=5/3。

\paragraph{计算／实验细节}矩形DG一次/二次基给二/三阶空间精度，局部Lax通量、Gaussian积分、特征minmod/TVB限制器、SSPRK3。限制线性系数时高阶系数清零。GEM质量比25、c/VA=10、Te/Ti=.2，128×64与512×256，dt=.002/.0005，M=0或25；另做真实质量比1836的500–50000格色散激波与解析多模电子波。

\paragraph{核实的发现}粗网格二阶出现持续磁岛而三阶无岛；高分辨二/三阶通量接近，三阶一小岛随后合并，增加耗散也能消岛，因此磁岛与数值耗散不是单调关系。晚期碰撞射流有复杂结构，有限精度和硬件精度能改变磁岛运动方向及局部动量约10\%。源项振荡使稳定时间步与相位准确性不同；激波细小波列可能是真实双流体色散。

\paragraph{限制与未完成问题}无Landau阻尼及压力张量闭合；通量一致不证明X点非理想电场预算闭合。边界有限磁通导致晚期饱和。预印本有明确量纲/符号排印问题：p12三阶CFL写c/Δx，正确量纲Δx/c；能量内能应p/(γ-1)，Harris电流应按curl B/μ0重算，不能照抄式5.12。

\paragraph{本报告判断}数值章用作DG实现及可复现性案例；比较需要同时报告阶数、限制器、网格、dt与机器精度，不能用一张磁岛图断言新物理。

\section{L174：2011年}\label{paper:L174}

\textbf{Model for Incomplete Reconnection in Sawtooth Crashes}\citep{L174}\

Matthew Beidler;  P. A. Cassak.

\textsf{核读范围：}arXiv1111.0590v1，物理页1–5全部，含MAST估算与文末限制；\textsf{图表：}图1–3、式(1)–(3)，页3的速度交叉、率下降和局地结构视觉核验.

\href{run:papers/L174.pdf}{打开本地PDF：\texttt{papers/L174.pdf}}\quad\href{https://doi.org/10.1103/physrevlett.107.255002}{DOI／版本入口}

\paragraph{在报告中的研究脉络}\hyperref[story:onset]{启动、爆发与磁岛}。关联表示本报告的阅读线索，不自动表示这些论文互相引用。

\paragraph{实际模型}二维周期平板F3D，冷离子、绝热电子γe=5/3、Hall和电子惯性；高压区起初远离有理面，由重联入流自行带入。并非真实三维环形m=n=1锯齿模拟。

\paragraph{计算／实验细节}102.4×204.8di，Harris宽.5di，质量比25、dx=.05和.025对照，D4=2e-5、无显式η。压力梯度宽2，p1/p2=5/25，初始密度1、平均导引场幅5并总压平衡；磁种子.01、噪声2e-5。比较均匀5/5、25/25与弱梯度5/7；χparallel=.08测试无显著变化。

\paragraph{核实的发现}快重联自身把高压核心带向X线，约t140梯度到达、t195电子抗磁速度达到两侧最大电子喷流速度的平均，通量率下降并转负振荡；仍有自由磁能。X点抗磁漂移、电流层变宽和磁剪切减弱随后出现。均匀压力或梯度不足的对照不停止，支持局地压强梯度的自限制机制。

\paragraph{限制与未完成问题}|v*i-v*e|out>vout为此模型的判据，实际只验证冷离子，热离子/强传导/三维待检验。后续结构只猜为漂移不稳定性，环形互换模未模拟。MAST v*e≈13.5与vout≈14.2km/s采用推定电子层宽2rho\_e及间接场/梯度，离子资料缺失；Bh取Bphi(1+rs/R0)的原文表达不是通常正交投影所得的场强，数量一致性不应宣称精确几何验证。

\paragraph{本报告判断}重要新现象是快重联可以输运进能终止自身的压强梯度，不需要耗尽所有磁通。实验比较属于原理与量级相容；不把摘要关于各装置包括ITER的外推当作普适验证。

\section{E471：2012年}\label{paper:E471}

\textbf{Numerical investigation of a compressible gyrofluid model for collisionless magnetic reconnection}\citep{E471}\

Luca Comisso;  D. Grasso;  E. Tassi;  François L. Waelbroeck.

\textsf{核读范围：}HAL封面及PoP19, 042103(2012)全部12页，共13页；含gyrokinetic来源与能量分析；\textsf{图表：}图1–10、式1–42；页7/9/11直接核验离子速度宽度、密度分布与能量.

\href{run:papers/E471.pdf}{打开本地PDF：\texttt{papers/E471.pdf}}\quad\href{https://doi.org/10.1063/1.3697860}{DOI／版本入口}

\paragraph{在报告中的研究脉络}\hyperref[story:energy]{电子独有重联与热化}。关联表示本报告的阅读线索，不自动表示这些论文互相引用。

\paragraph{实际模型}四场Hamiltonian gyrofluid，演化ni, ui, ne, ue的正则动量，保留平行离子压缩及Padé离子回旋平均；低β、等温，无Landau阻尼。ni为gyrocenter密度，不是粒子密度，准中性通过极化关系保证。

\paragraph{计算／实验细节}盒[-π, π]×[-2π, 2π]双周期，伪谱1024×128；Aeq=.1sech\textsuperscript{2}x的23项Fourier近似，de=.2, di=2, ρs=.4，ρi=10\textasciicircum{}-4–.8，Δ'=59.9。质量比100，电子β=.08；不同ρi按γLt/重联通量比较。

\paragraph{核实的发现}增加ρi使离子平行速度条带变宽，数据提示宽度\ensuremath{\propto}ρi，速度及X点电流尖峰下降；离子gyrocenter密度偏向中性线，电子密度仍偏分离线。热离子磁岛较长且开角较大，高次谐波更强。I±=D±diρi ni与G±=F±deρs ne四个冻结不变量解释两种结构；给出动力学固定正则动量的不变量对应。

\paragraph{限制与未完成问题}ρi扫描尚不足确立普适宽度定律；di/de=10并非真实质量比，低β使离子平行动量贡献小，不能推广高β。总能晚期误差<1.5\%，未给全面网格/时间收敛；磁能转至密度/静电扰动能是可逆过程，不能等同温升或非热粒子加速。等温两速度形式仅闭合表示，不是实际分布。

\paragraph{本报告判断}提供比双场更丰富的流体结构与能量分配检验，第一章须明确定义gyrocenter与粒子密度。

\section{GuideExp2012：2012年}\label{paper:GuideExp2012}

\textbf{Quantitative Study of Guide-Field Effects on Hall Reconnection in a Laboratory Plasma}\citep{GuideExp2012}\

T. D. Tharp;  Masaaki Yamada;  Hantao Ji;  E. Lawrence;  S. Dorfman;  Clayton E. Myers;  Jongsoo Yoo.

\textsf{核读范围：}PPPL-4753全部8页（正文编号1-5，含封面说明）；\textsf{图表：}图1-7、式1-3；主要结果图视觉核验.

\href{run:papers/GuideExp2012.pdf}{打开本地PDF：\texttt{papers/GuideExp2012.pdf}}\quad\href{https://doi.org/10.1103/physrevlett.109.165002}{DOI／版本入口}

\paragraph{在报告中的研究脉络}\hyperref[story:diagnostic]{实验与可区分诊断}。关联表示本报告的阅读线索，不自动表示这些论文互相引用。

\paragraph{实际模型}MRX外加可控环向导引场；counterhelicity pull，磁场探针测Hall项与感应电场，中心Langmuir探针测n/Te。

\paragraph{计算／实验细节}Bg/Brec从0到约1；delta约2cm，di约5cm，rho\_s约2.5cm；Hall项取X点上游3cm、下游5cm，归一化使用重联分量Brec而非总场。

\paragraph{核实的发现}导引场改变四极场形状与幅值，重联率降低比典型局地模拟更强；磁通芯阻碍排气，环向导引场被压缩、堆积，对全局压力平衡产生影响。

\paragraph{限制与未完成问题}正文明确不能把速率降低分成局部Hall与全局边界两部分；密度仅中心一点，完整压力平衡/离子流速测量留待未来。所读是PPPL预印本，式号图号以此版本为准。

\paragraph{本报告判断}边界可显著控制观察到的导引场依赖；不可当成所有开边界空间重联的普适Bg标度。

\section{L184：2012年}\label{paper:L184}

\textbf{Observation of explosive collisionless reconnection in 3D nonlinear gyrofluid simulations}\citep{L184}\

A. Biancalani;  B. Scott.

\textsf{核读范围：}arXiv1110.4595v2全部8页；\textsf{图表：}全部方程和图1-2；模型页、两幅增长率/能量图视觉核验.

\href{run:papers/L184.pdf}{打开本地PDF：\texttt{papers/L184.pdf}}\quad\href{https://doi.org/10.1209/0295-5075/97/15005}{DOI／版本入口}

\paragraph{在报告中的研究脉络}\hyperref[story:onset]{启动、爆发与磁岛}。关联表示本报告的阅读线索，不自动表示这些论文互相引用。

\paragraph{实际模型}GEM陀螺流体代码（不是GEM重联挑战）：低频强导引场、电子惯性、温度矩、Pade FLR及Landau闭合。

\paragraph{计算／实验细节}氘质量比约3704；Δprime rho\_s=.75，de/rho\_s=.41-4.1；de/Lperp示例.04；tau\_i1另0比较，epshat18350。Nx最高2048、Ny512，第三维点数未列；x Dirichlet，y/平行周期；白噪声种子1/10。beta扫描同时缩放盒子和片宽；能量源/耗散审计及x盒加倍，种子半倍/两倍检查。

\paragraph{核实的发现}非线性增长率比线性增长率提高一个数量级：beta1.6e-3时.009对.0008（Lperp/cs归一）。加速时X点分成两个Y点，饱和时恢复；岛内次级涡度及分离线尖电流；有限离子温度本组不改定性。

\paragraph{限制与未完成问题}两幅图仅显示能量/增长率；X-Y拓扑及湍动机制主要文字报告，未在本篇给图诊断或因果剔除。beta与尺度共同变化，不能称独立beta因果扫描。未明确三维是否必需。GE负只表明净数值能量损失，不能证明数值耗散对磁拓扑和重联率贡献严格为零；仍需Ohm残差、系数/网格扫描。

\paragraph{本报告判断}提供小Δprime也可爆发加速的数值证据；尺度独立极限和高速流/三维因果仍未解决。

\section{Shepherd2012：2012年}\label{paper:Shepherd2012}

\textbf{Guide field dependence of 3‐D X‐line spreading during collisionless magnetic reconnection}\citep{Shepherd2012}\

L. S. Shepherd;  P. A. Cassak.

\textsf{核读范围：}arXiv1208.2922v1全部9页，排版日期2021非发表年；\textsf{图表：}图1-6；图3-6视觉检查，传播方向、阈值和误差核对.

\href{run:papers/Shepherd2012.pdf}{打开本地PDF：\texttt{papers/Shepherd2012.pdf}}\quad\href{https://doi.org/10.1029/2012ja017867}{DOI／版本入口}

\paragraph{在报告中的研究脉络}\hyperref[story:three]{三维传播与有限X线}。关联表示本报告的阅读线索，不自动表示这些论文互相引用。

\paragraph{实际模型}F3D三维等温Hall加电子惯性；全初始电流由电子承担。

\paragraph{计算／实验细节}51.2×25.6×256di；dx=dy=.05, dz=1（部分.5检查）；质量比25；初始半厚.4di。导引场0至3；D4xy2.5e-5，D4z随导引场用.064/.081/.097，曾加倍检查；电阻0。局域磁扰动加随机模，Jmax>3.3定义重联，Hall电场>.01作参照；只拟合斜模相互作用前的时段。

\paragraph{核实的发现}弱导引场主要向电子电流方向延伸；较强导引场约Bg/B0≥2呈双向、近导引场Alfven速度。Bg2.5测得两方向3.0和2.9VA0；与载流子速度和导引场Alfven速度的最大值模型相符。后期斜向X线相互作用，破坏简单的单X线传播图像。

\paragraph{限制与未完成问题}阈值选择导致15\%-20\%速度不确定度；弱场反方向零是观测窗口内未分辨，非严格不传播。只验证电子承担全部电流，未验证一般离子分担系数。质量比25、初始薄片，无自洽薄化。临界场估计用了形成后电子层de与Brec约.4；后续传播理论强调前方未重联片厚，须区分。

\paragraph{本报告判断}支持受限参数下传播方向的导引场转变；不能直接作为太阳耀斑或真实质量比的定量预报。

\section{F625：2013年}\label{paper:F625}

\textbf{Influence of the dissipation mechanism on collisionless magnetic reconnection in symmetric and asymmetric current layers}\citep{F625}\

N. Aunai;  Michael Hesse;  Carrie Black;  Rebekah Evans;  Maria Kuznetsova.

\textsf{核读范围：}arXiv:1303.0442v1本地12页全部核读，包括对称/非对称两组数值实验、讨论及参考文献；\textsf{图表：}图1—10、表I—II和图注全文核对；另检查物理页7、10的布局与图表对应；正文若干颜色、图号和算例字母与图注不一致，结论以参数和曲线实际标记对应，不依赖颜色描述.

\href{run:papers/F625.pdf}{打开本地PDF：\texttt{papers/F625.pdf}}\quad\href{https://doi.org/10.1063/1.4795727}{DOI／版本入口}

\paragraph{在报告中的研究脉络}\hyperref[ch:numerics]{模型与数值可信度}。关联表示本报告的阅读线索，不自动表示这些论文互相引用。

\paragraph{实际模型}二维混合模型：Boris推进动力学质子，无质量、等温、准中性流体电子；忽略位移电流。Ohm定律含Hall、标量电子压强梯度及数值/均匀电阻/超电阻三种非理想处理。电磁场二阶差分、半格错置、预测校正；x周期、y封闭导体。

\paragraph{计算／实验细节}对称片λ=0.5di，密度1、上游β=1、Te=0.25；250di×50di、2500×500为基准，部分算例同时改变域和网格。η=0.002为网格控制组，η=0.014电阻组，ν=0.0003超电阻组。非对称Bx=0.5+tanh、n=1-(S+S\textsuperscript{2})/3，两侧|B|=0.5/1.5、n=1/(1/3)，Bg=1、总温度1.5、Ti/Te=5。Δy=0.15/0.076/0.05di，Lx=100/150/300di，ν=0.0005；分析限制在相应Alfvén返回时间之前。

\paragraph{核实的发现}对称组宏观结构和速率变化较小，网格控制与超电阻率约0.065/0.06（基准单位），均匀电阻约0.055。非对称网格控制组电子层反复拉长、数值撕裂成岛、再次拉长，速率随分辨率明显变化；电阻及超电阻局域化电子层、无反复岛生成，后者率较前者高。缩小周期域可使回流加宽电流层、抑制岛或人为维持短X线，不能把这种稳态当局部物理机制。

\paragraph{限制与未完成问题}此文检验‘任意非理想项都等价’而没有证明哪种流体闭合等价于真实电子动力学；数值岛不能直接认作自然等离子体岛。无电子惯性/压力张量，超电阻系数也非由微观过程导出。部分分辨率比较同时改变域长；超电阻参数完整扫描尚未完成，作者只报告初步弱敏感性。初态非Vlasov平衡，会发出反射波。

\paragraph{本报告判断}必须纳入：Hall项存在不保证闭合无关或网格无关；与近期分辨率争议及导引场磁岛合并共同组成模型、离散、边界三个层次的对照。

\section{F672：2013年}\label{paper:F672}

\textbf{Gyro-induced acceleration of magnetic reconnection}\citep{F672}\

L. Comisso;  D. Grasso;  F. L. Waelbroeck;  D. Borgogno.

\textsf{核读范围：}arXiv1309.1753v1物理页1–40全部，含附录色散关系、图版与表I；\textsf{图表：}图1–13、表I；图6、10、11和表I视觉复核.

\href{run:papers/F672.pdf}{打开本地PDF：\texttt{papers/F672.pdf}}\quad\href{https://doi.org/10.1063/1.4821840}{DOI／版本入口}

\paragraph{在报告中的研究脉络}\hyperref[story:onset]{启动、爆发与磁岛}。关联表示本报告的阅读线索，不自动表示这些论文互相引用。

\paragraph{实际模型}二维强导引场低beta等温四场gyrofluid；离子平行流、Padé环平均、粒子准中性、电子惯性，四组Lagrangian不变量；无显式碰撞或Landau阻尼。

\paragraph{计算／实验细节}sech\textsuperscript{2}x磁通平衡A0=.1, L1，23个Fourier模，Bplane/Bg最大约.08；双周期x±pi, y±a*pi；AB3伪谱最高4096×512，微尺度滤波。非线性Ly4pi、ky.5、Delta'59.9；de由.2降至.025且di=10de, rho\_s=2de, rho\_i=4de。

\paragraph{核实的发现}热离子且尺度充分分离时出现两次非线性增长率加速，中间停滞；固定rho\_tau=.2236比较冷热离子，首次加速只在热离子出现。大k*rho\_i时phi≈-rho\_i\textsuperscript{2} ne保留密度细结构，冷离子phi=Lap\textasciicircum{}-1 ne被平滑，局地E×B流可差一个量级。第二次加速伴随出流开口扩大。表I以初始最大面内场参照的率由de=.2时.2827降至.05时.1176、.025时.1004。

\paragraph{限制与未完成问题}γ\_eff=d ln(delta psi)/dt不等于磁通率；率用初始场最大值作上游参照，非动态层边缘。固定Delta'继续de→0会进入小Delta'并使快阶段消失，不是无限理想极限0.1。不含电子FLR/粒子不稳定性；微层不断变薄，到不可解析时停止。一次x加倍网格率不变、峰值ue变不到6\%，不是所有算例全收敛。正文图2参数rho\_s=.1与图注.2冲突，待作者输入确认。

\paragraph{本报告判断}双加速的证据来自对照、场结构和不变量关系，不能简化成“热离子使重联更快”；第一个小尺度加速与第二个宏观开口必须分开。

\section{L189：2013年}\label{paper:L189}

\textbf{Experimental Study of Two-Fluid Effect during Magnetic Reconnection in the UTST Merging Experiment}\citep{L189}\

Kotaro YAMASAKI;  S. Kamio;  Koichiro Takemura;  Qinghong Cao;  T. Watanabe;  Hirotomo Itagaki;  Takuma Yamada;  Michiaki Inomoto;  Yasushi Ono.

\textsf{核读范围：}PFR2013卷8文章2401130全部4页正文、图表与诊断说明；\textsf{图表：}图1–6，直接核验页2诊断关系及页3势降/磁通/电流的时序和采样几何.

\href{run:papers/L189.pdf}{打开本地PDF：\texttt{papers/L189.pdf}}\quad\href{https://doi.org/10.1585/pfr.8.2401130}{DOI／版本入口}

\paragraph{在报告中的研究脉络}\hyperref[story:diagnostic]{实验与可区分诊断}。关联表示本报告的阅读线索，不自动表示这些论文互相引用。

\paragraph{实际模型}UTST双球形托卡马克合并实验，磁探针阵列重建ψ、jt、Et；浮动电势径向扫描与离子Doppler谱联系两流体静电势和离子排出。

\paragraph{计算／实验细节}典型Bp=.01T、Bt=.1T，δ约50mm，作者引ρi79mm/平均自由程120mm；浮动探针10mm管/1mm钨尖，光谱沿R=.1m到采样点视线积分。200–205μs势梯度出现，Er估400–800V/m，VA24km/s，C\textsuperscript{2+}流约27km/s，估τacc2.1–4.2μs与回旋周期3.9μs比较。

\paragraph{核实的发现}重联磁通变化加速、电流下降与X点附近浮动电势变陡同时出现；含X点视线观测到近Alfvén径向流。静电加速时间量级与观测相容，提示磁化电子形成电势对离子排出有作用。

\paragraph{限制与未完成问题}浮动电势只有在Te空间变化足够小时才可代替空间电势，本文未提供完整同步Te修正；势降不能单独量出非准中性或证明Hall独占因果。谱线有视线积分/流动展宽，作者明确承认温度受宏观流影响，不能由谱宽直接确定不可逆加热。正文CIII与图6HeII标注不一致，ρi79mm与C回旋半径16mm属不同口径；热化/激波为解释而非直接测量。

\paragraph{本报告判断}实验章分别列直接观测、诊断换算条件和机制推断；比较加速与热化，不能把相关时序写成已证实Hall加热定律。

\section{L190：2013年}\label{paper:L190}

\textbf{Laboratory Study of Hall Reconnection in Partially Ionized Plasmas}\citep{L190}\

Eric E. Lawrence;  Hantao Ji;  Masaaki Yamada;  Jongsoo Yoo.

\textsf{核读范围：}本地PPPL-4764技术报告16物理页全部检查；物理页3—14为12页论文正文、图表和参考文献，另含封面、声明和空白页。实际核读2012年5月作者稿，对应2013年PRL，未声称已逐字核对期刊终稿；\textsf{图表：}图1—5及图注全部核对；放大物理页12图4—5的误差棒、两种归一化和电阻率比较.

\href{run:papers/L190.pdf}{打开本地PDF：\texttt{papers/L190.pdf}}\quad\href{https://doi.org/10.1103/physrevlett.110.015001}{DOI／版本入口}

\paragraph{在报告中的研究脉络}\hyperref[ch:numerics]{模型与数值可信度}；\hyperref[story:diagnostic]{实验与可区分诊断}。关联表示本报告的阅读线索，不自动表示这些论文互相引用。

\paragraph{实际模型}MRX部分电离氦实验，测Hall四极场、离子喷流与重联电场；不是电子—离子双流体模拟。解释需要离子—中性拖曳和有限设备尺寸。

\paragraph{计算／实验细节}pull重联准稳态约40μs（约12τA）；245通道磁探针阵列覆盖20cm×16cm，径向0.6—2cm、轴向3cm；Mach/Langmuir探针及移动电流片重建径向剖面。填气5.5—45mTorr，nn≈0.17—1.4×10\textasciicircum{}15cm\textasciicircum{}-3，ne≈2×10\textasciicircum{}13cm\textasciicircum{}-3、Te≈6eV；两端ρtot/ρi约8与70。放电中中性密度仅在±50\%内与填气推断相容。

\paragraph{核实的发现}加中性粒子后离子喷流按√(ρi/ρtot)降低，接近按总质量密度定义的Alfvén速度；归一Hall四极场形状及约5cm的峰间尺度、离子喷流通道宽度没有相应扩大。因此按离子Alfvén速度归一的率下降，按总密度速度归一则约为常数且量级1。电子已受中性碰撞影响时仍可保留Hall结构，不能把Hall等同于严格无碰撞。

\paragraph{限制与未完成问题}边界自由理论预期Hall尺度扩大，最高中性密度例预期有效尺度83cm超过设备，有限尺寸是候选解释而非已证明原因。中性压力未直接测量；流剖面依赖多发重复，轴对称重建有假设；基准‘ionized’仍含显著中性粒子。归一率约1来自本装置和尺度限制，不能用作普适常数。

\paragraph{本报告判断}纳入实验章节和多组分模型适用边界，区分电子—离子双流体与带电—中性双流体，建议显式中性动量和实际边界联合验证。

\section{L191：2013年}\label{paper:L191}

\textbf{Nonlinear variational method for predicting fast collisionless magnetic reconnection}\citep{L191}\

Makoto Hirota;  P.J. Morrison;  Yasutomo Ishii;  M. Yagi;  N. Aiba.

\textsf{核读范围：}arXiv1301.3196v2全部22页，含附录A–D能量积分、流映射和对OP理论的检验；\textsf{图表：}图1–9、B1；页12图8势能指数和页14图9电阻扫描直接视觉核验.

\href{run:papers/L191.pdf}{打开本地PDF：\texttt{papers/L191.pdf}}\quad\href{https://doi.org/10.1088/0029-5515/53/6/063024}{DOI／版本入口}

\paragraph{在报告中的研究脉络}\hyperref[story:onset]{启动、爆发与磁岛}。关联表示本报告的阅读线索，不自动表示这些论文互相引用。

\paragraph{实际模型}冷离子/冷电子无耗散双场惯性RMHD，ψe=ψ-de\textsuperscript{2}\ensuremath{\nabla}\textsuperscript{2}ψ随不可压缩流冻结；Lagrangian流映射变分法加直接非线性模拟。

\paragraph{计算／实验细节}双周期ψeq=ψ0cos(2πx/Lx)，de/Lx=.01，Ly/Lx=4π，kLx=.5；y谱至200模，x均匀差分约10\textasciicircum{}4点。大Δ'序Lx\textsuperscript{3}/(8Ly\textsuperscript{2})\ensuremath{\ll}de。τ0\textasciicircum{}-1=dekτH\textasciicircum{}-1，τH\textasciicircum{}-1=α\textsuperscript{2}ψ0。

\paragraph{核实的发现}Y端电流层对应试探映射给线性U=-.776εhat\textsuperscript{2}、γ=.881/τ0；非线性U约-.439εhat\textsuperscript{3}，故εhat¨约.659εhat\textsuperscript{2}、εdot\ensuremath{\propto}ε\textasciicircum{}3/2，比指数快。数值能量支持该指数而不支持OP早期ε¨\ensuremath{\propto}ε\textasciicircum{}4；附录D指出边界积分对极窄子层任意参数敏感，不能闭合能量。弱电阻提高线性率，但τ0/τe>1时非线性趋减速。

\paragraph{限制与未完成问题}变分试探族和小中间区近似并非精确非线性解；de\ensuremath{\ll}ε\ensuremath{\ll}Lx适用，不能把数学有限时间发散当无限能量。细电流尖峰达到网格后终止。F697后来承认本篇y模数和数值黏性漏掉次级plasmoid。TFTR τ0约90μs/τe270μs仅量级对应，非完整环形锯齿/部分重联验证；加热和逃逸电子是展望。

\paragraph{本报告判断}第一章适合由能量守恒逐步推导爆发指数，第三章必须连同2017分辨率修正讲述。

\section{L193：2013年}\label{paper:L193}

\textbf{Study of the effects of guide field on Hall reconnection}\citep{L193}\

T. D. Tharp;  Masaaki Yamada;  Hantao Ji;  E. Lawrence;  S. Dorfman;  Clayton E. Myers;  Jongsoo Yoo;  Yi-Min Huang;  Amitava Bhattacharjee.

\textsf{核读范围：}作者公开期刊版物理页1–6全部；\textsf{图表：}图1–6；图3–6视觉复核、测点和归一细查.

\href{run:papers/L193.pdf}{打开本地PDF：\texttt{papers/L193.pdf}}\quad\href{https://doi.org/10.1063/1.4805244}{DOI／版本入口}

\paragraph{在报告中的研究脉络}\hyperref[story:diagnostic]{实验与可区分诊断}。关联表示本报告的阅读线索，不自动表示这些论文互相引用。

\paragraph{实际模型}MRX氘等离子体导引场实验与有限电阻Hall MHD；等温初态、总压绝热gamma5/3，Ti=Te。

\paragraph{计算／实验细节}300多个磁探针，Langmuir单点密度温度；delta约2cm、di约5cm、rho\_s约2.5cm；两种螺度运行加外导引场约±3Brec。模拟25×12.5di，出流周期/入流导体，Harris h=.5di、扰动.05、T0=.125、eta=.005(Sdi200)，峰值测量；网格/步长未列。

\paragraph{核实的发现}导引场改变Hall四极结构并减弱Hall电场；实验测上游4cm和下游8cm，局地Hall项与重联电场随Bg同降。长文新增全局约20cm导引场压缩，反螺度中局地场约外加值2倍，出流背压与率下降相关；约0.08降到Bg/Brec约1.5时0.02。平板模拟未产生这部分压缩。

\paragraph{限制与未完成问题}必须用局地增强Bg才与模拟相近；同/反螺度是不同初始和全局设置，不能拼成单参数扫描。图3各面板独立色标，不能从颜色比较幅值。测得Hall平衡在X点外，非证明Hall项支撑X点破冻结。导引场背压对速率的因果贡献未定量隔离。

\paragraph{本报告判断}与2012短文分开登记；局地Hall证据良好，但几何/边界未被验证的限制需要与实验一致性一起呈现。

\section{Loureiro2013Heating：2013年}\label{paper:Loureiro2013Heating}

\textbf{Fast collisionless reconnection and electron heating in strongly magnetized plasmas}\citep{Loureiro2013Heating}\

N. F. Loureiro;  A. A. Schekochihin;  A. Zocco.

\textsf{核读范围：}1-6（全部）；\textsf{图表：}PDF2-5页直接视觉核验图1-5、式4-6；6页全文含脚注参考文献，无文内附录。.

\href{run:papers/Loureiro2013Heating.pdf}{打开本地PDF：\texttt{papers/Loureiro2013Heating.pdf}}\quad\href{https://doi.org/10.1103/PhysRevLett.111.025002}{DOI／版本入口}

\paragraph{在报告中的研究脉络}\hyperref[story:energy]{电子独有重联与热化}。关联表示本报告的阅读线索，不自动表示这些论文互相引用。

\paragraph{实际模型}KREHM，流体变量+电子Hermite动理学；二维周期伪谱，非纯Hall/有限矩。

\paragraph{计算／实验细节}384\textsuperscript{2}；M=30-500；de=rhoi=rhos=.25a，Lx=2πa、Ly扫描，主Δ'a20；νcoll=1/(Δt M\textsuperscript{4})，νH=.25(Δx/π)\textasciicircum{}6/Δt，六阶空间超扩散。

\paragraph{核实的发现}主例重联峰t22.7τA、耗散峰29.4τA；电子相混热化沿分离线，M变化对主重联结果影响较小；大盒峰率约.2vABy/c。

\paragraph{限制与未完成问题}Landau标签含高阶矩空间超扩散及超碰撞；没有离子不可逆加热竞争。约60\%初始磁能电子热化不是Qi/Qe，且对应大盒数值未画出；忽略常规电阻但保留碰撞的形式不一致作者已说明。

\paragraph{本报告判断}模型身份、温度/能量/熵产区别及跨文献判断见独立加热专题；不把本篇单独作为普适分配定律。

\section{Stanier2013：2013年}\label{paper:Stanier2013}

\textbf{Two-fluid simulations of driven reconnection in the mega-ampere spherical tokamak}\citep{Stanier2013}\

Adam  John Stanier;  P. K. Browning;  Mykola Gordovskyy;  K. G. McClements;  M. Gryaznevich;  V. S. Lukin.

\textsf{核读范围：}arXiv1308.2855v1全部12页；\textsf{图表：}图1-12、表I与模型式1-5；算法页及图3-12视觉核验.

\href{run:papers/Stanier2013.pdf}{打开本地PDF：\texttt{papers/Stanier2013.pdf}}\quad\href{https://doi.org/10.1063/1.4830104}{DOI／版本入口}

\paragraph{在报告中的研究脉络}\hyperref[story:global]{层长、驱动与全局完成}；\hyperref[story:diagnostic]{实验与可区分诊断}。关联表示本报告的阅读线索，不自动表示这些论文互相引用。

\paragraph{实际模型}MAST合并二维平面与环形轴对称可压缩Hall MHD，电子无惯性、标量pe；pi=pe=p/2，电阻/超电阻热、离子黏性、各向异性热传导。

\paragraph{计算／实验细节}L0=1m，B0=.5T、Bp0=.1T、n0=5e18/m\textsuperscript{3}、Te=Ti10eV，di=.145m，tau0=.29μs。HiFi四阶Jacobi谱元；通常180×360单元(720×1440有效点)，拉伸最小dz=.23mm；低超电阻360×540单元。Crank-Nicolson自适应隐式，MHD dt上限.05tau0、Hall.01；eta1e-5、mu1e-3、超电阻1e-6至1e-10。R .2-2m、Z±2.2m导体壁；环形无滑移，平面滑移；两绳半径.4m、总电流268kA，无初始线圈形成过程。

\paragraph{核实的发现}电阻小到约1e-5产生磁压堆积、反弹和速率振荡；改变超电阻呈厚片、易生中心岛的中间态、薄层开放X三种情形。平面对称性使中心岛滞留并阻塞合并；环形几何破坏对称后岛可排出，合并成嵌套磁面球形托卡马克。Hall电子平行压缩形成四极样密度结构，叠加环形内侧压缩，产生沿实验诊断弦的双峰密度，外峰随合并消失。

\paragraph{限制与未完成问题}作者增强黏性和超电阻以便计算，eta热导固定、pi=pe；没有独立Te/Ti演化、FLR和电子惯性。降低超电阻同时改变局地加热和rho\_s，不能只当作几何参数扫描。低超电阻峰率双网格差.2\%只覆盖生岛前时间，未保证后期所有岛统计收敛。模拟每20tau0约5.8μs与实验.1ms采样不等，峰密度也不相同；相似剖面是定性解释而非逐时定量匹配。实验初始磁结构无法反演，只估绳宽/总电流；q>1不等于全面三维稳定性已证明。

\paragraph{本报告判断}对用户关心的新现象特别有用：装置几何能改变岛是否排出、密度诊断呈什么形状；验证必须包括几何和仪器弦位。

\section{Yoo2013：2013年}\label{paper:Yoo2013}

\textbf{Observation of Ion Acceleration and Heating during Collisionless Magnetic Reconnection in a Laboratory Plasma}\citep{Yoo2013}\

Jongsoo Yoo;  Masaaki Yamada;  Hantao Ji;  Clayton E. Myers.

\textsf{核读范围：}PRL110 215007正式版全部5页，正文1–4、参考文献4–5；PPPL作者公开PDF；无文内附录。；\textsf{图表：}图1–4；直接查看PDF2–4页，核对电势及误差棒、离子流/电功、谱线拟合与下游温度剖面。PDF特殊字符提取失真处以页面图像为准。.

\href{run:papers/Yoo2013.pdf}{打开本地PDF：\texttt{papers/Yoo2013.pdf}}\quad\href{https://doi.org/10.1103/PhysRevLett.110.215007}{DOI／版本入口}

\paragraph{在报告中的研究脉络}\hyperref[story:diagnostic]{实验与可区分诊断}。关联表示本报告的阅读线索，不自动表示这些论文互相引用。

\paragraph{实际模型}MRX近反平行、无外加导引场氦放电实验；浮动电势与电子温度换算等离子体电势，电子动量平衡积分另作一致性检验，与离子光谱/流速比较（2026-11-13纠正原摘要）。

\paragraph{计算／实验细节}从4200余次放电按重现性筛选，充气压4.5mTorr。磁探针径向最高.6cm、轴向3cm；17针浮动电势阵列最高.7cm。He II468.6nm谱用13个Gaussian处理精细结构；时间分辨5.6微秒、空间3–4cm；Mach流速用光谱校准。Te≤12eV，电子均程8–12cm，片半宽约2cm；X点经典电阻项约占重联电场10\%；2026-11-13复查：电势由浮动电势和Te换算，电子动量平衡积分是独立一致性检验。

\paragraph{核实的发现}电势井在X线附近约10V、下游约35V并变宽；500–800V/m面内场在小于di的尺度上将离子加速至约.5VA。面内场对离子电功在下游局域达20–50W/cm\textsuperscript{3}，面外电功约2–8W/cm\textsuperscript{3}。Z≈9cm后流速趋饱和而温度升高；温度横向峰比出流分布窄。电势重建中电子压力梯度不可忽略。

\paragraph{限制与未完成问题}三维和时间变化由跨放电二维扫描近似；离子温度测量有空间平均，忽略精细谱线会高估15\%–25\%。下游密度5–8×10\textasciicircum{}13cm\textsuperscript{-3}使离子均程降至2–5cm，局部碰撞与再磁化可能参与；作者也保留波粒、势阱反弹及装置边界解释，明确未解决唯一热化机制。

\paragraph{本报告判断}把L272回引的原实验单独核读，建立测得电势/加速/升温与后续随机加热解释的证据层级。支持能量进入离子的空间路径，不能直接当成某一无碰撞熵产机制已获唯一实验确认。

\section{Andres2014：2014年}\label{paper:Andres2014}

\textbf{Effects of electron inertia in collisionless magnetic reconnection}\citep{Andres2014}\

N. Andrés;  Luis Martin;  P. Dmitruk;  Daniel Gómez.

\textsf{核读范围：}arXiv1312.2437v2 1-8全部；\textsf{图表：}图1-5；式28、图4视觉核验.

\href{run:papers/Andres2014.pdf}{打开本地PDF：\texttt{papers/Andres2014.pdf}}\quad\href{https://doi.org/10.1063/1.4890021}{DOI／版本入口}

\paragraph{在报告中的研究脉络}\hyperref[story:global]{层长、驱动与全局完成}。关联表示本报告的阅读线索，不自动表示这些论文互相引用。

\paragraph{实际模型}常密度不可压、标量压力、完整电子惯性准中性双流体，非全Maxwell；伪谱、无显式电阻／黏性。

\paragraph{计算／实验细节}周期2pi×2pi，双tanh半厚2pi*.02, psi扰动.02cos x, Hall .1；网格512\textsuperscript{2}-2048\textsuperscript{2}；me/mi .015及1/1836；总能误差<1e-8, 初始零两类螺度维持1e-15。

\paragraph{核实的发现}随分辨率提高理想MHD与Hall的伪重联趋零，而有限惯性通量趋非零；真实质量比案例峰值接近.1，质量比两值总体趋势相似。

\paragraph{限制与未完成问题}未扫描Hall与种子；守恒总能并非所有局地梯度准确的充分条件。真实质量比de约.00233，2048网格dx约.00307，不能凭全局通量相近声称充分解析电子层。预印本式28右端印为(2delta-1)导致物理质量比下右端负，与图1实频不相容，报告不照抄。

\paragraph{本报告判断}最强证据是无惯性对照加网格扫描；结果只在其标量不可压模型内支持惯性作用，不排除真实动理学压力张量。

\section{Che2014Buneman：2014年}\label{paper:Che2014Buneman}

\textbf{Two-Fluid Description of Wave-Particle Interactions in Strong Buneman Turbulence}\citep{Che2014Buneman}\

H. Che.

\textsf{核读范围：}1-26（全部）；\textsf{图表：}PDF7、16-19、21页视觉核验图1、6-9和完整流体方程；所有26页含图注/参考文献读完，无附录。.

\href{run:papers/Che2014Buneman.pdf}{打开本地PDF：\texttt{papers/Che2014Buneman.pdf}}\quad\href{https://doi.org/10.1063/1.4882677}{DOI／版本入口}

\paragraph{在报告中的研究脉络}\hyperref[story:energy]{电子独有重联与热化}。关联表示本报告的阅读线索，不自动表示这些论文互相引用。

\paragraph{实际模型}3D PIC诊断与近似双流体闭合分析；未独立演化验证纯双流体重联，模拟未发展重联。

\paragraph{计算／实验细节}mi/me100, w0=.1di=de, Te=Ti=.04miVA\textsuperscript{2}；盒1×1×2di，512×512×1024；x/z周期、y导电；电子漂移9VA≈3vte、离子.9VA；|B|≈sqrt26 B0。

\paragraph{核实的发现}Buneman拖曳和相混把电子漂移能转为平行热，磁能下降伴电子垂直热增长；局地流体矩预算近似刻画PIC；离子主要协助波交换，弱热化。

\paragraph{限制与未完成问题}Dez(t)仍需动力学理论/PIC拟合；同前序Che2013运行，非独立验证；无实际重联、无实验比；式19为电子平行/垂直比。图7约化PIC熵增不等于精细Vlasov熵增，未熵粗粒化收敛验证。

\paragraph{本报告判断}模型身份、温度/能量/熵产区别及跨文献判断见独立加热专题；不把本篇单独作为普适分配定律。

\section{E475：2014年}\label{paper:E475}

\textbf{Numerical comparison between a gyrofluid and gyrokinetic model investigating collisionless magnetic reconnection}\citep{E475}\

O. Zacharias;  Luca Comisso;  D. Grasso;  Ralf Kleiber;  Matthias Borchardt;  R. Hatzky.

\textsf{核读范围：}arXiv1406.1334v1物理页1–29全部；\textsf{图表：}表I、图1–11；图8、11视觉复核，并核对正文与图上定量差别.

\href{run:papers/E475.pdf}{打开本地PDF：\texttt{papers/E475.pdf}}\quad\href{https://doi.org/10.1063/1.4882679}{DOI／版本入口}

\paragraph{在报告中的研究脉络}\hyperref[story:onset]{启动、爆发与磁岛}。关联表示本报告的阅读线索，不自动表示这些论文互相引用。

\paragraph{实际模型}等温含电子惯性的四场gyrofluid，对比EUTERPE delta-f gyrokinetic PIC（陀螺离子/漂移电子）；无密度温度梯度的强导引场平衡。

\paragraph{计算／实验细节}sech\textsuperscript{2}x平衡C=.1、剪切长度5。流体1024×128双周期伪谱/AB；EUTERPE最多1024×128，x Dirichlet、y周期，B样条有限元/RK4、最多3e7 markers、最小dt=.125，de至少16格。表I质量比1836/100、de=.1，rho\_s=.3/.2，本文beta定义少常规因子2。非线性PIC保留m=0, ±1；流体初始种子与PIC噪声启动不同。

\paragraph{核实的发现}各代码先与各自线性射击本征解核验，再比较模型；线性中间波数仍可差20\%，减小剪切幅度后差约12\%。非线性小Delta'的饱和半岛宽、振荡频率吻合较好；ky约1.6时岛宽差约30\%，ky1.9约10\%，更近边界才趋一致。rho\_s=.3降至.025时半岛宽约.04不变；升高Ti/Te显著改变振荡频率而岛宽变化较弱。

\paragraph{限制与未完成问题}只验证小Delta'非线性，不能直接给F672大Delta'双加速背书。边界不同，虽在线性局地本征模中影响小；PIC模滤波限制高次谐波。报告能量误差小于2.5\%，未换算对释放磁能的比例。正文称Ti扫描岛宽仅变5\%，但图11端点约.040→.049及.046→.053，目测约15–25\%，不能无保留转述5\%。

\paragraph{本报告判断}高质量模型对照的价值在限定可测量与参数域；流体再现岛振荡频率不等于重建了俘获粒子的完整分布。

\section{F679：2014年}\label{paper:F679}

\textbf{Bulk ion acceleration and particle heating during magnetic reconnection in a laboratory plasma}\citep{F679}\

Jongsoo Yoo;  Masaaki Yamada;  Hantao Ji;  Jonathan Jara-Almonte;  Clayton E. Myers.

\textsf{核读范围：}所记录版本全部13个PDF物理页；正文、方法、结果、讨论、图表、参考文献及适用文内附录均核读。；\textsf{图表：}直接查看PDF物理页5, 8, 9, 10, 11；Figs 5,  9–13;  Eqs (1)–(19).

\href{run:papers/F679.pdf}{打开本地PDF：\texttt{papers/F679.pdf}}\quad\href{https://doi.org/10.1063/1.4874331}{DOI／版本入口}

\paragraph{在报告中的研究脉络}\hyperref[story:energy]{电子独有重联与热化}；\hyperref[story:diagnostic]{实验与可区分诊断}。关联表示本报告的阅读线索，不自动表示这些论文互相引用。

\paragraph{实际模型}本文是磁重联实验装置的氦等离子体实验，并以二维全动理学粒子模拟进行对照。电子动量平衡用于解释面内电势，光谱与探针分别约束离子温度和流速；模拟比较有、无离子碰撞的速度分布。文中的“双流体区域”指电子与离子解耦的实验状态，并不意味着所用数值模型是有限矩双流体模型。

\paragraph{计算／实验细节}实验气压为 4.5 毫托，磁通芯间距 42 厘米，离子皮肤深度约 8 厘米，无外加导场。4200 次放电约三成通过筛选，每点平均 7 至 15 次；离子光谱的空间、时间分辨率约为 3 至 4 厘米和 5.6 微秒。粒子模拟质量比为 400，初始两物种等温，电子热速为光速的 0.183 倍；离子碰撞频率及自由程相对特征尺度与实验匹配。

\paragraph{核实的发现}面内电场约 700 伏每米，高于约 200 伏每米的重联电场，主要离子增能由面内电势驱动。有序流能小于电势能下降量，下游出现更高离子温度。无碰撞模拟的较大温度包含多束混合；加入实验量级碰撞后，分布更接近麦克斯韦形态。电子能量沉积约 340 千瓦，经典电阻贡献估计仅约 59 千瓦。

\paragraph{限制与未完成问题}实验没有完整测得电子压力张量及热通量，约 300 千瓦的热损失估计具有较大误差，因此异常电子耗散是条件性推断，不能唯一归因于所见波动。二维模拟的部分反向离子束也未被实验空间分辨率辨认。完整网格、时间步和边界细节部分依赖前序方法文献，本次未复核全部前序实现或复算原始数据。

\paragraph{本报告判断}本文把电势加速与下游分布热化联系起来，支持轨道运动、再磁化和碰撞共同作用的解释，为流体闭合提供实验约束。它没有证明只加入霍尔项就能得到正确熵产，也没有给出普适离子与电子不可逆加热比。应把物种电场功、流动能、压力功和热输运分开核算；本次依据正式版全部 13 页及关键图式作出判断。

\section{L202：2014年}\label{paper:L202}

\textbf{Fast,  purely growing collisionless reconnection as an eigenfunction problem related to but not involving linear whistler waves}\citep{L202}\

Paul M. Bellan.

\textsf{核读范围：}PoP21, 102108(2014)全部9页正文及出版封面，共10页；\textsf{图表：}图1–3与式1–21；物理页7图3本征函数、κ值直接视觉核验.

\href{run:papers/L202.pdf}{打开本地PDF：\texttt{papers/L202.pdf}}\quad\href{https://doi.org/10.1063/1.4897375}{DOI／版本入口}

\paragraph{在报告中的研究脉络}\hyperref[story:onset]{启动、爆发与磁岛}。关联表示本报告的阅读线索，不自动表示这些论文互相引用。

\paragraph{实际模型}均匀密度、离子静止、电子压强各向同性barotropic的电阻/惯性EMHD线性Harris本征问题；两条二阶ODE共四阶，偶Az/奇Bz耦合。

\paragraph{计算／实验细节}ξ=x/de，ε=de/L；N=500半域ξ\ensuremath{\in}[0, 40]，IDL两点边值+Muller找根，a(0)=-1, a(40)=s(0)=s(40)=0，并调κ使a'(0)=0。导数用一阶差分，未给收敛扫描。主要忽略(ky de)\textsuperscript{2}，ν归一化0或2，ε=.1–1。

\paragraph{核实的发现}得到纯增长、非传播的局域模，γ=√2ωwh/κ，κ为约5–18量级；电子惯性/电阻正则化中心，完全略去二者则中心电流奇异。增长时间可以与whistler频率同阶，但这不等于存在驱动重联的线性驻whistler波。MRX估γ约10\textasciicircum{}6 s\textasciicircum{}-1，理论出流零点宽1.2cm仍小于观测3.3cm。

\paragraph{限制与未完成问题}未计算非线性稳态率。以任意约10\%种子把1/γ比作总重联时间只是量级论证；MRX离子冻结/均匀n及几何近似有限。κ弱变化只是所扫区间，不能宣称me→0普适常率。文中把一般电阻MHD概括成二阶过窄：完整撕裂还需动量方程。J·E>0表示场对粒子做功，不等于不可逆热化。

\paragraph{本报告判断}用于澄清色散时间尺度与真实波、线性增长率与非线性重联率的差别，同时保留其中心正则性证明。

\section{Cassak2015：2015年}\label{paper:Cassak2015}

\textbf{Fast magnetic reconnection due to anisotropic electron pressure}\citep{Cassak2015}\

P. A. Cassak;  R. N. Baylor;  R. L. Fermo;  M. T. Beidler;  M. A. Shay;  Marc Swisdak;  J. F. Drake;  H. Karimabadi.

\textsf{核读范围：}1-5（全部）；\textsf{图表：}图1各机制消融、图2场及欧姆项；图2视觉核验.

\href{run:papers/Cassak2015.pdf}{打开本地PDF：\texttt{papers/Cassak2015.pdf}}\quad\href{https://doi.org/10.1063/1.4908545}{DOI／版本入口}

\paragraph{在报告中的研究脉络}\hyperref[story:fast]{局地快率与必要条件}；\hyperref[story:closure]{压力、热流与闭合}；\hyperref[story:diagnostic]{实验与可区分诊断}。关联表示本报告的阅读线索，不自动表示这些论文互相引用。

\paragraph{实际模型}F3D，电子CGL或离子CGL对照，独立关闭Hall与电子惯性；保留eta=0.005与D4；压力演化省略E·J加热项。

\paragraph{计算／实验细节}204.8di×102.4di，4096×2048，mi/me=25，Bg=5，n=1，所演化物种初始p\_parallel=p\_perp=5，另一种冷；D4=2.5e-5；coherent perturbation .014，velocity noise .04。

\paragraph{核实的发现}eCGL在无Hall、甚至无电子惯性对照中仍有开放排气和约0.06-0.1率；iCGL或只把电子各向异性放动量方程则约0.01长层；面外四极极性反转。

\paragraph{限制与未完成问题}CGL压强并不负责X点断线，电阻保留；仅某种各向异性状态方程有效，Egedal闭合的慢结果在文中是私人通信；PIC一致性提及但未展示图；过强或过弱导引场均转慢。

\paragraph{本报告判断}这是电子压力改变层几何的机制消融证据，不能称为任意各向异性都快或无电阻无Hall的重联证明。

\section{F644：2015年}\label{paper:F644}

\textbf{Two-fluid numerical simulations of turbulence inside Kelvin-Helmholtz vortices: Intermittency and reconnecting current sheets}\citep{F644}\

C. Rossi;  F. Califano;  Alessandro Retinò;  Luca Sorriso‐Valvo;  Pierre Henri;  Sergio Servidio;  F. Valentini;  A. Chasapis;  L. Rezeau.

\textsf{核读范围：}物理页1–10全部，含HAL与期刊封面，正式正文122303-1至8，2015年；\textsf{图表：}图1–8、表I、方程1–10；谱图3与虚拟穿越图8直接视觉核验.

\href{run:papers/F644.pdf}{打开本地PDF：\texttt{papers/F644.pdf}}\quad\href{https://doi.org/10.1063/1.4936795}{DOI／版本入口}

\paragraph{在报告中的研究脉络}\hyperref[story:energy]{电子独有重联与热化}。关联表示本报告的阅读线索，不自动表示这些论文互相引用。

\paragraph{实际模型}实际是2.5维Hall MHD，非三维；准中性无位移电流、无质量电子，双物种标量绝热γ5/3；欧姆律仅-ue×B-\ensuremath{\nabla}Pe/n，无电子惯性/显式电阻/非对角压力。

\paragraph{计算／实验细节}初始剪切uy=tanh[(x-Lx/2)/6]，B=(0, sin0.02, cos0.02)，n=1, Ti=Te=0.5；域约400×402di，4096×8192，dx≈0.1, dy≈0.05；RK3、紧致差分、y谱滤波与x六阶滤波，y周期x透明边界。t1000在60×60di局部涡内区用Hanning窗FFT与PVI，谱滤波从kdi>3主导。

\paragraph{核实的发现}KH涡破碎自发产生间歇薄片，面内谱从各向异性趋于各向同性；磁谱MHD段约-1.8、次离子段约-3.8但后者区间很窄。磁增量胖尾/平坦度随尺度增强，高PVI集中在约0.5–1di。抽选局部片有X形磁通、双侧电子入流与喷流；强导向场约5倍重联场且稍非对称，四极Hall场不明显。

\paragraph{限制与未完成问题}声称R≈0.1–0.3来自场角/流速纵横比代理而非dψ/dt或X点非理想电场；湍流非定常使稳态公式及所谓下界解释不严格，表I数值和文中比值也非逐项一致。二维标量电子方程缺显式物理断线项，磁拓扑改变需要数值滤波/耗散，未给非理想预算，不能据此证明真实电子耗散机制。无电子加热或熵产生诊断。局部60di窗口、kdi1–3短区间不支持普适谱指数。

\paragraph{本报告判断}适合展示KH湍流中多薄片/间歇性与虚拟航天器诊断；不能把其高纵横比直接当精确普适快重联率，亦不宜把流体能量变化等同不可逆动理学加热。

\section{F674：2015年}\label{paper:F674}

\textbf{Study of energy conversion and partitioning in the magnetic reconnection layer of a laboratory plasma}\citep{F674}\

Masaaki Yamada;  Jongsoo Yoo;  Jonathan Jara-Almonte;  W. Daughton;  Hantao Ji;  Russell M. Kulsrud;  Clayton E. Myers.

\textsf{核读范围：}所记录版本全部11个PDF物理页；正文、方法、结果、讨论、图表、参考文献及适用文内附录均核读。；\textsf{图表：}直接查看PDF物理页6, 7, 8, 9；Figs6–8,  TableI,  Eqs(2)–(6).

\href{run:papers/F674_published.pdf}{打开本地PDF：\texttt{papers/F674\_published.pdf}}\quad\href{https://doi.org/10.1063/1.4920960}{DOI／版本入口}

\paragraph{在报告中的研究脉络}\hyperref[story:energy]{电子独有重联与热化}。关联表示本报告的阅读线索，不自动表示这些论文互相引用。

\paragraph{实际模型}本文通过磁重联实验装置的局部测量建立有限控制体能量收支，并与装置几何下及开放边界下的全动理学粒子模拟比较。分别积分电子和离子的电场功，再区分有序流能、热储能与边界输运。它是实验能量预算和动理学对照研究，不能列为霍尔流体独立预测温度比例的数值证据。

\paragraph{计算／实验细节}本次读取 2015 年正式版 11 页，保留的 2014 年预稿按同一论文处理。开放粒子模拟采用真实质子质量比 1836，初始离子与电子温度比为 5，背景密度比为 0.23，背景电子温度比为 0.76。模拟盒为 20 乘 20 个参考离子皮肤深度，网格为 5120 乘 5120，每种粒子约一百亿个；电流片半宽为 0.5 个参考离子皮肤深度。

\paragraph{核实的发现}图 7 给出离子和电子电场沉积率分别为 0.71±0.15 与 0.35±0.06 兆瓦，约为二比一；两物种的热能变化项却分别约为 0.27±0.10 与 0.26±0.10 兆瓦。后者还包含热储能变化和边界焓流差。表一改变控制体后，离子和电子占输入磁能的份额也发生变化，说明该比例与统计区域有关。

\paragraph{限制与未完成问题}预算中的热损失具有较大误差，部分离子损失只给出上限；输出能流、边界条件和采样尺度都会影响份额。粒子模拟虽然包含真实质量比，也未覆盖任意导场、等离子体压强、三维湍流及宏观装置。本文承接 2014 年实验成果，不能把两篇文章当成相互独立的复现实验；原始数据和模拟代码未在本次复算。

\paragraph{本报告判断}常被引用的二比一是物种电场能量沉积率比，既不是离子与电子温度比，也不是两物种不可逆熵产生热比。本文证明在指定实验及控制体内可以建立有用的分物种能量预算，但没有提供普适加热闭合。应同时报告有序流动、热储能、焓流和热损失，不能只摘取一个比值作为霍尔模型的普遍预测。

\section{L213：2015年}\label{paper:L213}

\textbf{Electron-scale nested quadrupole Hall field in Cluster observations of magnetic reconnection}\citep{L213}\

Neeraj Jain;  A. Surjalal Sharma.

\textsf{核读范围：}期刊正式版719–724全部6页，正文719–723、参考文献723–724；无文内附录。旧arXiv10页版本保留，主索引改为正式版。；\textsf{图表：}图1–4，直接比较图1模拟嵌套四极、图2观测与图4虚拟穿越；核对L1–L4各分量过零次序。.

\href{run:papers/L213_published.pdf}{打开本地PDF：\texttt{papers/L213\_published.pdf}}\quad\href{https://doi.org/10.5194/angeo-33-719-2015}{DOI／版本入口}

\paragraph{在报告中的研究脉络}\hyperref[story:diagnostic]{实验与可区分诊断}。关联表示本报告的阅读线索，不自动表示这些论文互相引用。

\paragraph{实际模型}电子惯性EMHD模拟与Cluster单次电子尺度穿越对照；离子固定的早期电子动力学，不含电子动理学压力各向异性。

\paragraph{计算／实验细节}归一化长度de、时间Ωce\textasciicircum{}-1；图1比较45和145Ωce\textasciicircum{}-1，早期次级点约x=±16de。与2001-10-01磁尾SC4事件比较时取B0=10nT、de=20km。模拟片半厚约7de，观测3–5de；采用约20度量级不确定性的法向坐标。本文未完整列出网格/时间步/收敛，依赖2009年前序方法论文。

\paragraph{核实的发现}主X点出流与较弱邻近X点入流相互作用，形成不直接围绕某个X点的新外四极；原四极扩展并嵌在其中。三点同强度的对照不产生同样嵌套，因此多个X点本身不是充分条件。模拟虚拟轨道可重现观测多个磁分量的过零次序和电场双极形；较弱外极约为主Hall峰的30\%。

\paragraph{限制与未完成问题}单航天器不能唯一分离空间与时间；未穿越X点，法向场约1nT、坐标误差可达20度。观测片宽与模拟也非精确一致。固定离子模型的晚期外推没有在此直接检验；不能凭磁场形貌唯一认定机制，不能称直接拍到了完整嵌套二维结构。

\paragraph{本报告判断}与L283尘埃多物种嵌套四极形成关键对照：相似磁场图样可有不同尺度和生成机制；与F703谱线积分、L119局部场测量共同说明需要合成观测和误差传播。

\section{L214：2015年}\label{paper:L214}

\textbf{Explosive magnetic reconnection caused by an X-shaped current-vortex layer in a collisionless plasma}\citep{L214}\

Makoto Hirota;  Yuji Hattori;  P. J. Morrison.

\textsf{核读范围：}arXiv1503.06865全部26页，完整内外区匹配、能量推导及参考文献；\textsf{图表：}图1–11及式1–64；页9/12/23直接核验局域X层、能量和7/4拟合.

\href{run:papers/L214.pdf}{打开本地PDF：\texttt{papers/L214.pdf}}\quad\href{https://doi.org/10.1063/1.4921329}{DOI／版本入口}

\paragraph{在报告中的研究脉络}\hyperref[story:onset]{启动、爆发与磁岛}。关联表示本报告的阅读线索，不自动表示这些论文互相引用。

\paragraph{实际模型}冷离子、有限电子温度双场惯性RMHD；ψ±=ψ-de\textsuperscript{2}\ensuremath{\nabla}\textsuperscript{2}ψ±ρsde\ensuremath{\nabla}\textsuperscript{2}φ分别冻结。专门取de=ρs，四项正定总能守恒。

\paragraph{计算／实验细节}双周期余弦平衡，双方向Fourier谱+RK4，最小de=ρs=.005Lx用8192\textsuperscript{2}波数网格；扫描Ly/Lx及de，特征例ky/α=.171。大Δ'近似LxΔ'约16(Ly/Lx)\textsuperscript{2}。

\paragraph{核实的发现}所扫线性不稳定解分为ρsΔ'\ensuremath{\lesssim}.65早饱和与\ensuremath{\gtrsim}.65在ε>ρs后加速。非线性X电流/涡量层在局部生成，长度约Lx，不由初始线性波长单独决定；磁能主要转入离子平面流能。匹配得局域岛位移σdot=σ\textsuperscript{2}/(ly0τH)，外部测量位移ε满足σ\ensuremath{\propto}ε\textasciicircum{}3/4，从而εdot\ensuremath{\propto}ε\textasciicircum{}7/4，快于全局kink型对应解。

\paragraph{限制与未完成问题}.65阈值、ly约.5Lx和7/4依赖本平衡、de=ρs及试探几何；最优ly落在所设Δ'c=100/Lx近似域下界，不能当无约束全局最优化定理。薄层能量小不表示可先令de=0消除重联。ET是约化扰动能，不是自洽热化。无动理学/Landau阻尼，de≠ρs推广本文未验证。

\paragraph{本报告判断}新增发现是局域非线性X层与外部位移的关系，而非仅重复快重联；第一章应区分σ和ε，避免混淆2和7/4指数。

\section{L216：2015年}\label{paper:L216}

\textbf{Fast magnetic reconnection with large guide fields}\citep{L216}\

A. Stanier;  Andrei N. Simakov;  L. Chacón;  W. Daughton.

\textsf{核读范围：}物理页1–10全部；1–2封面、3–7正文、8–10图版；\textsf{图表：}图1–3；物理页5的离散模型、图1和图3视觉核对.

\href{run:papers/L216.pdf}{打开本地PDF：\texttt{papers/L216.pdf}}\quad\href{https://doi.org/10.1063/1.4905629}{DOI／版本入口}

\paragraph{在报告中的研究脉络}\hyperref[story:fast]{局地快率与必要条件}。关联表示本报告的阅读线索，不自动表示这些论文互相引用。

\paragraph{实际模型}强导引场、低beta、冷离子、等温电子约化两场；电子惯性de和声回旋半径rho\_s；电子超电阻与离子黏性有限，非零耗散支持准稳态X点电场。

\paragraph{计算／实验细节}岛合并四种两流体参数保持h=max(de, rho\_s)=.01，对比无两流体项；eta\_H扫描约1e-11至1e-8，mu/eta\_H=1e4。Harris半盒Lx/Lx0=.5–8，Lx0=100h；只改变出流方向长度，入流盒宽与初始片宽固定。

\paragraph{核实的发现}de>rho\_s、无快速色散波时仍出现对eta\_H近乎不敏感的速率。测得耗散层边缘Bx约随delta/h线性减小，delta<rho\_s后微观层长w约随delta/rho\_s缩短；外层Bxh和wv近乎不依赖eta\_H。统一标度delta\textasciitilde{}sqrt(eta\_H*wv/Bxh)/h、Ez\textasciitilde{}Bxh\textsuperscript{2}*h/wv。局地归一Harris率在16倍出流长度扫描中无明显系统下降。

\paragraph{限制与未完成问题}三条离散方程不能独立选出全部五个层参数；Bxh、wv仍需外部供给。系数有约因子2变化，厚度预测仅量级精度。尺寸独立不等于任意几何自相似放大。未证明全局层长为何被限制，不能把有限超电阻删成零后仍称稳态破冻结。

\paragraph{本报告判断}与Rogers2001构成直接反例链：快速色散波不是所测试低beta快重联的必要条件；应检查局地上游场和层长而非只比较最终0.1。

\section{L217：2015年}\label{paper:L217}

\textbf{Fluid vs. kinetic magnetic reconnection with strong guide fields}\citep{L217}\

A. Stanier;  Andrei N. Simakov;  L. Chacón;  W. Daughton.

\textsf{核读范围：}arXiv1505.03794v1物理页1–17全部；自动排版日期不作为出版年；\textsf{图表：}图1–19及式1–22；图11/12、18/19视觉复核，全部图注核读.

\href{run:papers/L217.pdf}{打开本地PDF：\texttt{papers/L217.pdf}}\quad\href{https://doi.org/10.1063/1.4932330}{DOI／版本入口}

\paragraph{在报告中的研究脉络}\hyperref[story:fast]{局地快率与必要条件}。关联表示本报告的阅读线索，不自动表示这些论文互相引用。

\paragraph{实际模型}补全L216的约化流体与VPIC对照；陀螺黏性抵消使电子惯性中的平流速度为E×B速度v，不是含抗磁修正的磁通输运速度vs。

\paragraph{计算／实验细节}PIXIE2D Newton–Krylov/物理多重网格，二阶中心差分、对流QUICK，CN或BDF2、X点集中网格；未列逐算例格点与步长。无色散波的局地对照VPIC质量比15.2、Bg17.4、vthe/c=.0625、omega\_pe/Omega\_cex=10、Ti/Te=.25；有色散波质量比380.3、Bg12、vthe/c=.308、Ti/Te=.1、每格16000粒子。尺寸扫描另用vthe/c=.25及omega\_pe/Omega\_cex=2.5并作相对论Alfven归一。

\paragraph{核实的发现}流体与PIC的速率及宏观层长接近，有色散波案例峰值相差不超过约10\%；PIC早期出现小岛/电子射流扭曲，流体未出现。无色散波的大盒更易产生二次岛，但小盒无岛也快，二次岛不是必要条件。固定de降低rho\_s与固定rho\_s增大de的双向扫描说明速率主要联系h=max(de, rho\_s)及启动期间上游场变化，不能将与beta的相关性归因于色散分支消失。

\paragraph{限制与未完成问题}尺寸率用4h上游场、入流跨盒Alfven时间窗平均，并排除末25\%磁通耗尽时段；不等于全时程一致。大盒PIC噪声改变生岛史。初始片宽、局地/初始场归一影响结论。图14与图15黏性设置不同，不能混用。层长自选机制仍未解决。

\paragraph{本报告判断}这是有参数、形态和归一控制的支持证据，同时清楚界定“率吻合”与“电子结构一致”的差别。

\section{Ng2015：2015年}\label{paper:Ng2015}

\textbf{The island coalescence problem: Scaling of reconnection in extended fluid models including higher-order moments}\citep{Ng2015}\

Jonathan Ng;  Yi-Min Huang;  Ammar Hakim;  A. Bhattacharjee;  Adam  John Stanier;  W. Daughton;  Liang Wang;  Kai Germaschewski.

\textsf{核读范围：}1-8（全部）；\textsf{图表：}图1-6文字及图注；主要结果图已视觉复核.

\href{run:papers/Ng2015.pdf}{打开本地PDF：\texttt{papers/Ng2015.pdf}}\quad\href{https://doi.org/10.1063/1.4935302}{DOI／版本入口}

\paragraph{在报告中的研究脉络}\hyperref[story:global]{层长、驱动与全局完成}。关联表示本报告的阅读线索，不自动表示这些论文互相引用。

\paragraph{实际模型}电阻MHD、Hall MHD电阻／超电阻、十矩局地松弛对照PIC；Fadeev岛合并。

\paragraph{计算／实验细节}epsilon=.4, nb=.2n0, Lx×Lz=4piλ×2piλ, x周期, z导体／反射；mi/me=25，Ti=Te，beta=1；平均窗0-1.5tA，模拟至2.5tA。Hall eta=.009或etaH=1e-4及eta=1e-5；mu=.01。

\paragraph{核实的发现}Hall区峰值／平均率对系统尺寸仅弱敏感，约λ\textasciicircum{}-0.2至\textasciicircum{}-0.25；十矩调低ki到1/(3di)改善小系统离子压力支撑，但大系统次级岛不断抛出，合并完成且缺少PIC弹回。

\paragraph{限制与未完成问题}本文比较PIC平均率约λ\textasciicircum{}-0.8；引言早期研究的\textasciicircum{}-0.5不可混用，盒子宽高比不同。十矩ki是可调参数，出现PIC中没有的次级岛；本文不是十矩全面成功的证据。

\paragraph{本报告判断}必须比较同几何、同平均窗；Hall率可高2-3倍。

\section{Numata2015Heating：2015年}\label{paper:Numata2015Heating}

\textbf{Ion and electron heating during magnetic reconnection in weakly collisional plasmas}\citep{Numata2015Heating}\

Ryusuke Numata;  N. F. Loureiro.

\textsf{核读范围：}1-18（全部）；\textsf{图表：}PDF3、5、10、11、14、17页直接视觉核验式2.1-2.5、表1、图5-7/11-12及附录A收敛图A1。.

\href{run:papers/Numata2015Heating.pdf}{打开本地PDF：\texttt{papers/Numata2015Heating.pdf}}\quad\href{https://doi.org/10.1017/S002237781400107X}{DOI／版本入口}

\paragraph{在报告中的研究脉络}\hyperref[story:energy]{电子独有重联与热化}。关联表示本报告的阅读线索，不自动表示这些论文互相引用。

\paragraph{实际模型}双物种电磁δf陀螺动理学AstroGK，2D位置2D速度；不是流体模型。

\paragraph{计算／实验细节}Lx3.2πa, Ly2.5πa, Δ'a23.2, ky a.8；me/mi=.01, Ti/Te1, rhoi/a.25, rhoe/a.025；β=.01, .03, .1, .3, 1时de/a=.25, .14, .079, .046, .025；Nx256 Ny128(β1时Nx512)，Nλ=NE64；νii=νee=νei是人为正则化选择。

\paragraph{核实的发现}低β电子相混熵产主导且晚于重联；β1时部分晚期Di/De达量级1；离子熵产集中次级岛内部，电子更多分离线/边缘；离子有平行和FLR垂直相混。

\paragraph{限制与未完成问题}瞬时Di/De不是积分普适Qi/Qe；β扫描连带de变化；人工相等碰撞率、低质量比、强导引排序；附录主要验证电子速度网格，作者承认应补离子分辨率；未实验定量匹配。

\paragraph{本报告判断}模型身份、温度/能量/熵产区别及跨文献判断见独立加热专题；不把本篇单独作为普适分配定律。

\section{Tanabe2015MAST：2015年}\label{paper:Tanabe2015MAST}

\textbf{Electron and Ion Heating Characteristics during Magnetic Reconnection in the MAST Spherical Tokamak}\citep{Tanabe2015MAST}\

H. Tanabe;  T. Yamada;  T. Watanabe;  K. Gi;  K. Kadowaki;  M. Inomoto;  R. Imazawa;  M. Gryaznevich;  C. Michael;  B. Crowley;  N. J. Conway;  R. Scannell;  J. Harrison;  I. Fitzgerald;  A. Meakins;  N. Hawkes;  K. G. McClements;  T. O’Gorman;  C. Z. Cheng;  Y. Ono;  The MAST Team.

\textsf{核读范围：}所记录版本全部5个PDF物理页；正文、方法、结果、讨论、图表、参考文献及适用文内附录均核读。；\textsf{图表：}直接查看PDF物理页2, 3, 4；All Figs1–5;  experimentgeometry, Te/ne/Titomography, timeevolution, guidefieldscan.

\href{run:papers/Tanabe2015MAST.pdf}{打开本地PDF：\texttt{papers/Tanabe2015MAST.pdf}}\quad\href{https://doi.org/10.1103/PhysRevLett.115.215004}{DOI／版本入口}

\paragraph{在报告中的研究脉络}\hyperref[story:energy]{电子独有重联与热化}；\hyperref[story:diagnostic]{实验与可区分诊断}。关联表示本报告的阅读线索，不自动表示这些论文互相引用。

\paragraph{实际模型}本文是球形托卡马克的强导场双环合并实验，以汤姆孙散射测量电子温度和密度，以离子多普勒层析重建温度空间分布，并与已有双流体模拟比较。实验同时讨论电子与离子的碰撞能量交换、各向异性热输运及封闭磁面约束。论文自身不进行独立霍尔模拟，温度结构的符合是对具体模型解释的实验约束。

\paragraph{计算／实验细节}环向导场约为 0.3 至 0.8 特斯拉，重联场约为 0.07 至 0.15 特斯拉；离子皮肤深度约 0.1 米，离子回旋半径小于 0.01 米。电子诊断分别为 300 和 130 通道，空间分辨率约 15 与 10 毫米；离子诊断为 32 弦碳离子光谱层析。二维剖面组合多个放电脉冲和测量位置，重联场依据合并后磁平衡重建估计。

\paragraph{核实的发现}电子在重联中心附近形成宽约 2 至 5 厘米的热点，温度在小于 1 毫秒内由约 10 升到 200 电子伏；离子温升分布在出流及下游停滞区。经过约 4 至 10 毫秒的两物种能量交换时间，两个温度剖面均形成三峰。更高导场使中心电子热点更尖，但所测下游离子温度基本不变，与较弱导场机制的趋势不同。

\paragraph{限制与未完成问题}原文明确指出缺乏直接磁探针测量限制了热点成因的判别。黏性、压缩阻尼与输运解释依赖模型，温度图并不是完整局地功率或熵产预算；杂质离子温度与主离子温度的对应也依赖诊断和碰撞条件。全部二维点并非单次放电同时测得；独立诊断论文、学位论文和原始脉冲数据未在本次重新分析。

\paragraph{本报告判断}这是双温流体建模线的重要实验支撑：原文直接引用 L223，并将下游温升与含碰撞黏性耗散的双流体结果联系。但热点位置相似不足以唯一验证闭合，强导场下约束增强和驻留时间延长也会改变最终温度。该文支持具体装置中的局部加热研究，没有给出普适物种分配比；正式版 5 页及全部五幅图已核读。

\section{Wang2015：2015年}\label{paper:Wang2015}

\textbf{Comparison of multi-fluid moment models with particle-in-cell simulations of collisionless magnetic reconnection}\citep{Wang2015}\

Liang Wang;  Ammar H. Hakim;  A. Bhattacharjee;  K. Germaschewski.

\textsf{核读范围：}1-15（arXiv v3全部）；\textsf{图表：}图1-10及式20-24、36；图4、8-10视觉核验.

\href{run:papers/Wang2015.pdf}{打开本地PDF：\texttt{papers/Wang2015.pdf}}\quad\href{https://doi.org/10.1063/1.4906063}{DOI／版本入口}

\paragraph{在报告中的研究脉络}\hyperref[story:closure]{压力、热流与闭合}。关联表示本报告的阅读线索，不自动表示这些论文互相引用。

\paragraph{实际模型}五矩和十矩全Maxwell，局地Hammett-Perkins松弛；k0e=k0i=1/de为targeted，10\textasciicircum{}4/de近五矩。

\paragraph{计算／实验细节}主Harris例Lx×Ly=100×50di, a=.913di, nb=.3, Ti/Te=5, mi/me=25, vA/c=1/15, 扰动.1；流体4095×2047（约8格/de），PIC5760×2880（约12格/de），每物种2.7e9粒子。另有mi/me=25, 100, 256, 400质量比扫描。

\paragraph{核实的发现}图4通量时间史明显不同；同重联磁通下targeted十矩的X点附近电子流和非对角压强较接近PIC，五矩及强松弛出现PIC没有的磁岛链。质量扫描给出电子层长约(me/mi)\textasciicircum{}(3/8)、厚约(me/mi)\textasciicircum{}(1/3)。

\paragraph{限制与未完成问题}对比结构常取不同时间（图10 PIC75、流体132.5）；Pxy峰值PIC约.024对流体.006，标量化仍偏强；下游欧姆项峰位不一致、抵消不足。真实质量比层长12-17de和厚.4-.6de是外推，不是直接计算；无完整热流张量检验。

\paragraph{本报告判断}能复现部分局地压力支撑，不能写成重联全过程或所有压力分量定量吻合；局地闭合参数仍需选择。

\section{Andres2016：2016年}\label{paper:Andres2016}

\textbf{Influence of the Hall effect and electron inertia in collisionless magnetic reconnection}\citep{Andres2016}\

N. Andrés;  P. Dmitruk;  Daniel O. Gómez.

\textsf{核读范围：}arXiv1509.01169v1 1-8全部；\textsf{图表：}式15-16与图4-6视觉复核.

\href{run:papers/Andres2016.pdf}{打开本地PDF：\texttt{papers/Andres2016.pdf}}\quad\href{https://doi.org/10.1063/1.4942418}{DOI／版本入口}

\paragraph{在报告中的研究脉络}\hyperref[story:global]{层长、驱动与全局完成}。关联表示本报告的阅读线索，不自动表示这些论文互相引用。

\paragraph{实际模型}同常密度不可压惯性Hall双流体，无显式耗散；提出准稳态局地量纲标度。

\paragraph{计算／实验细节}周期2pi方盒, 2048\textsuperscript{2}, me/mi=.015, Hall λ=.07至.16步长.01共十组，初始片厚2pi*.02、psi扰动.02cos x；另λ.1比较me/mi .015与.15。

\paragraph{核实的发现}电流拟合厚度delta约(1.3±.3)lambda\_e\textasciicircum{}(1.06±.07)；固定设置中Bin约.25-.3、半长约.45-.6变化弱；准稳态率r约(.11±.07)λ\textasciicircum{}(.98±.03)，增Hall时峰值提前。

\paragraph{限制与未完成问题}只是2.3倍参数窗；取delta与Bin近恒定时段，非严格稳态；局地改用B0=Bin, L0=Delta才写r=λ，实际原归一化系数约.11。预印本式15惯性旋度项符号、式16两等号不一致；严格对称X点若停滞且稳态，标量压力无电阻模型不能用稳态对流惯性独自支撑Ez。

\paragraph{本报告判断}数值有限窗线性趋势可记录；稳态X点解析推导需保留时间惯性或澄清局地估算位置，不能直接推广普适长度规律。

\section{E448：2016年}\label{paper:E448}

\textbf{Laboratory Observation of Resistive Electron Tearing in a Two-Fluid Reconnecting Current Sheet}\citep{E448}\

Jonathan Jara-Almonte;  Hantao Ji;  Masaaki Yamada;  Jongsoo Yoo;  W. Fox.

\textsf{核读范围：}OSTI作者稿全部5页；\textsf{图表：}图1-5；图2-5视觉核对，含误差带与统计定义.

\href{run:papers/E448.pdf}{打开本地PDF：\texttt{papers/E448.pdf}}\quad\href{https://doi.org/10.1103/physrevlett.117.095001}{DOI／版本入口}

\paragraph{在报告中的研究脉络}\hyperref[story:diagnostic]{实验与可区分诊断}。关联表示本报告的阅读线索，不自动表示这些论文互相引用。

\paragraph{实际模型}MRX碰撞氩等离子体，亚离子电流片，测得Hall场；电子惯性可忽略，Spitzer电阻破坏冻结。

\paragraph{计算／实验细节}上游n约5e19/m\textsuperscript{3}、Te3eV；片内2e20/m\textsuperscript{3}、Te4eV；de.5mm<平均自由程2mm<半厚9mm<<di20cm。约400探针，最佳径向6mm、流出方向3cm；O/X插值加Poincare-Hopf绕数、时间过境校验。200放电、约500形成事件；每放电6.4μs时间窗，Alfven时间8.8±3μs。

\paragraph{核实的发现}S不足100（主要图约10-50）即有单磁岛及链；平均岛数和片长宽比随S增大。超过3/4形成事件电场增强；平均一个岛对应全球归一化率约10\%增加；无岛时已有.1-.2快速率。长薄、强Hall、电子仍碰撞的电流片没有自动变成de宽的开放X形。

\paragraph{限制与未完成问题}未完成含实测密度剖面、压缩性、Hall四极场与有限长宽比的二维本征计算，故抵触的是不适用的单流体阈值，未唯一证明具体电子撕裂机制。径向分辨率仅约半厚的2/3，不能据此确定电子子层；岛数有大方差，电场增强是统计关联。三维结构未解决。

\paragraph{本报告判断}非常适合检验可压缩电阻Hall模拟：需复现实验平衡与驱动，不能把Hall效应与无碰撞等同。

\section{F611：2016年}\label{paper:F611}

\textbf{Hybrid simulations of magnetic reconnection with kinetic ions and fluid electron pressure anisotropy}\citep{F611}\

A. Lê;  W. Daughton;  Homayoun Karimabadi;  J. Egedal.

\textsf{核读范围：}LANL接受稿全部18个PDF物理页：第1页机构封面，稿件编号1–17页；正文1–15页、参考文献15–17页；无文内附录。；\textsf{图表：}图1–7、式1–17；直接查看电子层、各向异性、两组率曲线和Hall场剖面，特别是图7时间对齐说明。.

\href{run:papers/F611.pdf}{打开本地PDF：\texttt{papers/F611.pdf}}\quad\href{https://doi.org/10.1063/1.4943893}{DOI／版本入口}

\paragraph{在报告中的研究脉络}\hyperref[story:closure]{压力、热流与闭合}。关联表示本报告的阅读线索，不自动表示这些论文互相引用。

\paragraph{实际模型}H3D混合模型：PIC离子、无质量准中性电子，保留Hall、各向异性电子压力及电阻/超电阻。旋转对称压强状态方程在弱俘获等温与强俘获CGL极限间插值；与等温混合和全PIC比较，不是完整电子动理学或十矩。

\paragraph{计算／实验细节}force-free与Harris两种初值；40×20di域、混合1024×1024格、全PIC5120×2560格，每动理学物种约400粒/格。PIC质量比400、ωpe/Ωce=2，x周期、z导体反射；片半厚di、Ti/Te=5、βe=.1；导场比分别.3和.2。率在X线周围边长2di方盒平均；Harris混合种子更大并按初始已重联磁通作时间平移。

\paragraph{核实的发现}加入各向异性后复现跨数个di的延长电子电流层；等温电子没有同样长层。Fe=μ0(pe\ensuremath{\parallel}-pe\ensuremath{\perp})/B\textsuperscript{2}在force-free例峰值约.7，Harris长层处大于约.5。Harris的X线附近面外磁场增强至背景的约1.5–2倍，图6剖面存在双峰而非中心单峰。三种模型后期率仍接近，早期增长史不同。

\paragraph{限制与未完成问题}Harris闭合需要沿二维磁通标签追踪上游值，并额外修改归一化密度以避免中心出现非物理各向异性；不是任意几何自动适用的本地状态方程。无导场的非旋转对称压力、复杂三维非局域闭合未解决。本稿未列出η、ηH的具体取值及独立收敛扫描。PIC质量比400使部分电子仅边缘磁化，层长比混合短。

\paragraph{本报告判断}提供相近总率仍需检查层长和Hall场的直接对照。与Cassak2015并列时，应区分两种电子状态方程；作者明确这里的插值闭合本身没有产生无Hall快速支。新增推导放第一章，缺失计算输入明确留空。

\section{L210：2016年}\label{paper:L210}

\textbf{"Ideal" tearing and the transition to fast reconnection in the weakly collisional MHD and EMHD regimes}\citep{L210}\

Del Sarto,  Daniele;  Fulvia Pucci;  Anna Tenerani;  Marco Velli.

\textsf{核读范围：}arXiv1511.00035v1物理页1–15全部，含附录A与表I；PDF自动日期2018与上传年不同，期刊年另核对；\textsf{图表：}图1–3、表I、式1–22与A1–A3，图2两种归一化增长谱直接视觉核验.

\href{run:papers/L210.pdf}{打开本地PDF：\texttt{papers/L210.pdf}}\quad\href{https://doi.org/10.1002/2015ja021975}{DOI／版本入口}

\paragraph{在报告中的研究脉络}\hyperref[ch:numerics]{模型与数值可信度}。关联表示本报告的阅读线索，不自动表示这些论文互相引用。

\paragraph{实际模型}两场惯性RMHD与惯性EMHD两个极限；随后加入ρτ\textsuperscript{2}=ρs\textsuperscript{2}+ρi\textsuperscript{2}有限拉莫尔修正。本文明确没有解中间全Hall介导体制；无完整非线性模拟。

\paragraph{计算／实验细节}解析边界层加数值线性本征求解，Harris平衡及一般Δ'\ensuremath{\propto}k\textasciicircum{}-p。全局εd*=(de/L)\textsuperscript{2}，RMHD时钟L/VA，EMHD时钟L\textsuperscript{2}/(di VA)；临界长宽比需足够容纳最快模，通常L/a≥约20。图2扫εd*至10\textasciicircum{}-12，临界峰约0.39和0.37各自逆时钟。

\paragraph{核实的发现}电子惯性小仍可在足够薄片上得到宏观时间的增长。Harris惯性RMHD有a/L≈(de/L)\textasciicircum{}2/3；EMHD解析指数对εd*为5/28，数值约3/16，两者应区分。ρτ\ensuremath{\gg}de时Harris得a/L≈(de/L)\textasciicircum{}1/3(ρτ/L)\textasciicircum{}1/3，较冷极限容许更厚临界片。对次级拉长薄片重新归一化，提出解释非线性阶段增长加快的候选机制。

\paragraph{限制与未完成问题}‘理想’是宏观增长时间而非η=de=0的拓扑变化，且增长率不等于非线性归一化重联电场。EMHD解析大Δ'关系借自特定分段平衡，有限de/di尺度分离不足会破坏渐近。次级爆发解释未以本文非线性模拟验证；忽略沿层流动的稳定性影响。忽略电阻须ηm\ensuremath{\ll}γde\textsuperscript{2}，正文讨论正确而结论段反向不等号不可照抄。部分k重归一化符号和表I有印刷问题，报告自洽重推，不沿用不明缩放。

\paragraph{本报告判断}适合第一章完整教局域/全局时间与小/大Δ'匹配，突出触发阈值而非承诺普适速率；保留RMHD、EMHD和全Hall三个范围区别。

\section{L223：2016年}\label{paper:L223}

\textbf{Two-fluid and magnetohydrodynamic modelling of magnetic reconnection in the MAST spherical tokamak and the solar corona}\citep{L223}\

P. K. Browning;  S. Cardnell;  M. Evans;  F. Arese Lucini;  V. S. Lukin;  K. G. McClements;  A. Stanier.

\textsf{核读范围：}PDF 1–16: all body,  methods,  results,  discussion,  figures/captions,  references;  no internal appendix；\textsf{图表：}Direct PDF visual: Eq.1–5 and definitions (pp3–4),  Fig.3 (p7),  relaxation equations/Figs.4–8 (pp9–10, 13–14);  remaining figure captions read in full text.

\href{run:papers/L223.pdf}{打开本地PDF：\texttt{papers/L223.pdf}}\quad\href{https://doi.org/10.1088/0741-3335/58/1/014041}{DOI／版本入口}

\paragraph{在报告中的研究脉络}\hyperref[story:energy]{电子独有重联与热化}。关联表示本报告的阅读线索，不自动表示这些论文互相引用。

\paragraph{实际模型}采用二维直角和紧环径比轴对称环面Hall MHD描述同向电流磁通绳合并，保留电子漂移、电子标量压力、电阻、超电阻及热传导。前列方程仍演化总压，图3才给独立离子、电子温度的初步扩展：离子由黏性供热，电子由电阻和超电阻供热。后半篇的Taylor弛豫计算只估计磁能释放，不能混作物种加热模拟。

\paragraph{计算／实验细节}基准密度5×10\textsuperscript{18} m\textsuperscript{-3}、长度1 m、磁场0.5 T，初始Ti=Te=10 eV，归一化离子惯性长度0.145。典型电阻10\textsuperscript{-5}、离子黏性10\textsuperscript{-3}、平行电子热导0.1、垂直离子热导10\textsuperscript{-7}；图3超电阻10\textsuperscript{-8}。使用HiFi及前序模型，但本文没有报告该两温加热图专用的网格、时间步或收敛检查，完整两温方程亦未列出。

\paragraph{核实的发现}图3的初步两温环面计算产生下游离子双热峰，峰值可达keV量级，与早期MAST高驱动观测有定性联系；电子热源主要跟随超电阻涉及的电流梯度，分布不同于离子。中心或空心电子温度与岛排出的联系仍是初步解释。弛豫模型另得到几何相关磁能释放及近似极向磁场平方标度，但不能给出热量物种比例或空间分布。

\paragraph{限制与未完成问题}超电阻是自由模型参数，并直接影响电子加热位置；有限标量闭合不描述完整粒子分布或非热尾。离子热峰高于后来较低驱动实验，实验比较仍属初步，不能称为加热比例定量验证。Stanier2014学位论文、源代码和独立数据未核读，本文也未给两温热图的专门收敛证据；没有推导普适Qi/Qe。

\paragraph{本报告判断}本篇直接证明两温Hall扩展能够在指定耗散、传导和几何条件下计算离子与电子局地热图。证据应表述为“模型能做条件性热预算和实验比较”，不能表述为“Hall项本身提供不可逆热化”或“物种分配已经普遍解决”。同时必须区分前列总压方程、后续初步两温计算和独立磁能弛豫估计三种证据。

\section{L224：2016年}\label{paper:L224}

\textbf{A two-fluid study of oblique tearing modes in a force-free current sheet}\citep{L224}\

Cihan Akçay;  W. Daughton;  V. S. Lukin;  Yi-Hsin Liu.

\textsf{核读范围：}arXiv1601.07220v1全部9页，附录A1–A10及参考文献完整；\textsf{图表：}图1–4、表I、超电阻阈值式7，图1斜角峰与图2–3耗散扫描直接视觉核验.

\href{run:papers/L224.pdf}{打开本地PDF：\texttt{papers/L224.pdf}}\quad\href{https://doi.org/10.1063/1.4940945}{DOI／版本入口}

\paragraph{在报告中的研究脉络}\hyperref[story:three]{三维传播与有限X线}。关联表示本报告的阅读线索，不自动表示这些论文互相引用。

\paragraph{实际模型}HiFi全双流体准中性模型，等温电子/绝热离子、电子惯性、标量压力加简单运动黏性，无陀螺黏性；与VPIC对照。通过旋转无力平衡逐个选斜模，实际算二维线性过程，不是三维多模非线性湍流。

\paragraph{计算／实验细节}无力片λ=0.5di，域2π×2πdi，mi/me=100并部分试400；bg=1, 2.5, 4, 10，βe=.09, .025, .01, .02，Ti=Te。η=0, ηH=10\textasciicircum{}-7，扫描4×10\textasciicircum{}-8–10\textasciicircum{}-3，μ=κ=10\textasciicircum{}-4并部分扫描。HiFi三/四阶谱元非均匀格de内8–11点，BDF2；沿片周期、法向导体滑移壁。只演化扰动而冻结平衡，排除超电阻改变初始梯度。

\paragraph{核实的发现}强导向场中最快线性撕裂偏离片中心：bg10时θ≈3–4°，增长约正向模式两倍；双流体和PIC均有此趋势。两流体绝对率较PIC低25–30\%或某些角度近两倍。提高ηH使斜角峰移回θ=0，改变不稳定模式选择；ηH≤约10\textasciicircum{}-7增长趋于收敛，与ηHcrit=γde\textsuperscript{2}w\textsuperscript{2}、w=de\textsuperscript{2}Δ'/(2√π)相符。早期Eparallel主要由∂tve支撑，离开共振层超电阻贡献增加。

\paragraph{限制与未完成问题}简化两流体/动理学边界层理论均给过平的斜角谱，不能准确预报增强；附录提出平衡电流梯度j'修正但未解完整本征问题。两流体理论与动力学理论√(2/π)系数差不能解释所有角度误差。电子恒温丢失超黏热，作者称该参数下可忽略但无完整能量图。w仅约0.17–0.4de，单报de格数不能替代内层收敛；无三维线绑边界或非线性验证。印刷动量通量质量权重等不直接抄为通用方程。

\paragraph{本报告判断}高价值模型比较：可复制斜模优势，却不能视为全谱/加热完全等效；尤其应在数值章强调耗散对‘最强模式是否斜向’的定性影响。

\section{L227：2016年}\label{paper:L227}

\textbf{Magnetic reconnection and modification of the Hall physics due to cold ions at the magnetopause}\citep{L227}\

Mats André;  W. Li;  Sergio Toledo‐Redondo;  Y. V. Khotyaintsev;  Andris Vaivads;  Daniel Bruce Graham;  Cecilia Norgren;  James L. Burch;  Peter Lindqvist;  Göran T. Marklund;  R. E. Ergun;  R. B. Torbert;  Werner Magnes;  Christopher T. Russell;  Barbara L. Giles;  Thomas Earle Moore;  M. O. Chandler;  Craig J. Pollock;  D. T. Young;  L. A. Avanov;  J. Dorelli;  Daniel J. Gershman;  W. R. Paterson;  B. Lavraud;  Y. Saito.

\textsf{核读范围：}HAL insu-03669467v1共9页（封面1页+正式期刊6705–6712全部8页）；正文及参考文献。；\textsf{图表：}图1–4，关键为图3分布和图4四星Ohm预算；直接查看PDF6/7页。.

\href{run:papers/L227.pdf}{打开本地PDF：\texttt{papers/L227.pdf}}\quad\href{https://doi.org/10.1002/2016gl069665}{DOI／版本入口}

\paragraph{在报告中的研究脉络}\hyperref[story:diagnostic]{实验与可区分诊断}。关联表示本报告的阅读线索，不自动表示这些论文互相引用。

\paragraph{实际模型}2015-11-12 MMS日下点磁层顶分离面观测；热、中温、冷离子与电子的多组分广义Ohm关系，不是普通单离子双流体模拟。

\paragraph{计算／实验细节}四星间距约10km，电子/离子分布30/150ms；分离面宽70–100km、移动70–80km/s；冷离子约50eV，ρc约25km，电子约1km，中温/热离子约100/320km。正常场约10mV/m，部分短脉冲>20mV/m。

\paragraph{核实的发现}热离子已去磁时，冷离子仍与电子共同E×B漂移而部分抵消Hall电流；冷离子占一半附近Hall与冷离子漂移项相当，电子压强散度另供2–3mV/m，合计支持约10mV/m法向场。

\paragraph{限制与未完成问题}单事件的分离面正常电场预算，不是X线重联电场测量；最终E使用E·B=0重建、Ex加2mV/m校准，未采用轴向探针最终场；压力散度和curlometer为空间平均。不能从此推出普遍重联率抑制量或局部不可逆加热。

\paragraph{本报告判断}扩充两尺度图像为温度分群多尺度；单一离子标量压强闭合不能保留冷/热离子不同磁化状态。

\section{L230：2016年}\label{paper:L230}

\textbf{Role of electron inertia and reconnection dynamics in a stressedX-point collapse with a guide-field}\citep{L230}\

Graf von der Pahlen,  J.;  D. Tsiklauri.

\textsf{核读范围：}arXiv1608.01397全部11页，包括方法、二维Ohm分解、三维结构、分布函数及参考文献；\textsf{图表：}图1–8；页4参数、页6图3 Ohm平衡、页8图5三维涡管直接视觉核验.

\href{run:papers/L230.pdf}{打开本地PDF：\texttt{papers/L230.pdf}}\quad\href{https://doi.org/10.1051/0004-6361/201628071}{DOI／版本入口}

\paragraph{在报告中的研究脉络}\hyperref[story:onset]{启动、爆发与磁岛}。关联表示本报告的阅读线索，不自动表示这些论文互相引用。

\paragraph{实际模型}EPOCH相对论电磁PIC，属于两流体闭合的动理学对照，不能列为Hall MHD模拟。受压X点α=1.2，平面开放边界吸收粒子/辐射场，z周期；Bg=0–.6BP。

\paragraph{计算／实验细节}mi/me=100，Te=Ti=6×10\textasciicircum{}7K；2.5D400\textsuperscript{2}、Δ=λD、500粒子/格；3D200×200×100、Δ=2λD、200粒子/种/格；盒约4di、z长一半，至tωpe=500。正文n=10\textasciicircum{}16 m\textasciicircum{}-2为单位笔误，按体密度m\textasciicircum{}-3解释；其B0=.03207T与VAe=.1c、n=10\textasciicircum{}16不一致，后两者应给.003207T，不能无核对搬作太阳量级。

\paragraph{核实的发现}Bg增大先延迟启动；.6BP时磁岛和第二X点出现，移动X点处对流惯性ve, x∂xue, z显著支撑重联电场，短时率超过低Bg。三维出现两侧反向倾斜的电子涡管，约±16°；平面貌似逆流移动实为倾斜管沿z运动投影。高Bg电子分布形成并随后平坦化的尾部凸起。

\paragraph{限制与未完成问题}图3高Bg惯性、压强及时间惯性存在强局部抵消，不能说压强项处处消失。三维未算完整Ohm平衡或磁连接重联率，只比涡管/电流结构；Q>0不是重联证据。文中斜模条件k×B=0为笔误，应为k·B=0。四di小盒、开放边界、质量比/三维网格限制外推；导向场数值阈值依赖所用BP定义。

\paragraph{本报告判断}为流体模型需保留全对流惯性提供反例和检验目标；三维结构、非热分布不应冒充Hall流体已预测结果。

\section{F697：2017年}\label{paper:F697}

\textbf{An Explosive Scaling Law for Nonlinear Magnetic Reconnection and Its Insensitivity to Microscopic Scales}\citep{F697}\

Makoto Hirota.

\textsf{核读范围：}PFR12, 1401010(2017)全部8页，含太阳量级外推与全部结果；\textsf{图表：}图1–8及式1–23；页3/6/7直接核验冷等离子体岛链、热离子/电子及电阻扫描.

\href{run:papers/F697.pdf}{打开本地PDF：\texttt{papers/F697.pdf}}\quad\href{https://doi.org/10.1585/pfr.12.1401010}{DOI／版本入口}

\paragraph{在报告中的研究脉络}\hyperref[story:onset]{启动、爆发与磁岛}。关联表示本报告的阅读线索，不自动表示这些论文互相引用。

\paragraph{实际模型}低β强导向场简化gyrofluid，Padé极化ne-ni0=\ensuremath{\nabla}\textsuperscript{2}(1-ρi\textsuperscript{2}\ensuremath{\nabla}\textsuperscript{2})\textasciicircum{}-1φ，电子惯性及有限ρs；电阻只作用ψ-ψeq，人为维持背景。

\paragraph{计算／实验细节}双周期余弦片，Fourier谱+RK4，4096\textsuperscript{2}并以8192\textsuperscript{2}复核plasmoid；de/Lx=.005–.01、ρs或ρi同阶，Ly/Lx=5.24/5.84，η=0, 2, 4de\textsuperscript{2}γL0。

\paragraph{核实的发现}把2015局域X层εdot\ensuremath{\propto}ε\textasciicircum{}7/4和磁能释放率\ensuremath{\propto}(-δW)\textasciicircum{}4/3延伸至热离子或热电子及有限电阻例；固定非线性幅度时速率弱依赖微尺度。冷ρs=ρi=0出现间歇plasmoid，能量释放有波动且低于该拟合。作者明确承认2013工作200个y模及数值黏性漏掉次级岛，是分辨率改变现象判断的直接实例。

\paragraph{限制与未完成问题}不是全部四参数扫描或普适最快率证明；启动门槛/总时间仍依赖δin，τrec\ensuremath{\propto}δin\textasciicircum{}-3/4。磁能大部转为流能不等于加热。太阳示例取VA=2.2×10\textasciicircum{}8m/s约.73c、重联场.1T大于导向场.01T，已超出非相对论强导向约化模型的可靠排序；.3小时只能作为示意外推，非观测验证。

\paragraph{本报告判断}与L191/L214形成完整证据链，展示率的微尺度不敏感和启动时间敏感可以同时成立。

\section{F798：2017年}\label{paper:F798}

\textbf{Extended MHD modeling of tearing-driven magnetic relaxation}\citep{F798}\

Joshua Paul Sauppe;  Carl R. Sovinec.

\textsf{核读范围：}本地作者稿13页全部核读，包含模型、参数、连续松弛事件、模态分解、动量通量、限制和参考文献；引用的独立2016附录收敛材料未单独核读；\textsf{图表：}图1—16和表I全部核对；放大页7的Hall/MHD发电机竞争、页11的Maxwell与边界应力通量.

\href{run:papers/F798.pdf}{打开本地PDF：\texttt{papers/F798.pdf}}\quad\href{https://doi.org/10.1063/1.4977540}{DOI／版本入口}

\paragraph{在报告中的研究脉络}\hyperref[story:three]{三维传播与有限X线}；\hyperref[story:diagnostic]{实验与可区分诊断}。关联表示本报告的阅读线索，不自动表示这些论文互相引用。

\paragraph{实际模型}NIMROD三维圆柱RFP扩展MHD；比较普通MHD、Hall加简化电子惯性、再加一阶Braginskii离子回旋黏性。电子惯性主要限制whistler时间尺度，本研究测得它对碰撞less重联作用很小。固定离子/电子温度份额、忽略Ohmic和黏性加热，不能作完整热预算。

\paragraph{计算／实验细节}R/a=3，r-z面120×64个双五次谱元，方位Fourier m=0—5、轴向周期；r=a无滑移、法向B固定、n/T Dirichlet，散度清理和小密度扩散。S=20000，比MST约低两数量级，Pm=1、χ/ν=0.1、β轴心0.1、di/a≈0.17、Ti=Te、电子质量增大5倍。无初始轴向/角向流，力自由顺磁箍缩，加小磁流扰动；跟踪多个松弛事件。

\paragraph{核实的发现}首次准单螺旋态崩解时Hall与MHD发电机共同展平电流；后续多螺旋事件中Hall项可反向抵消部分MHD项，净效应仍展平。主m=1, n=6贡献下降、n=5/7增强与反向Hall相关。对应Lorentz力先于Reynolds响应，后者通常抵消前者；后续事件核心电流和速度改变量同向，与MST集合观测方向更一致。总磁松弛相近但动能响应不同；径向Maxwell应力输运动量占主导。

\paragraph{限制与未完成问题}不能把局部J·EH=0误用为平均Hall发电机无效；该项在平均磁场方向的相关投影可非零。净轴向/角向动量主要来自模拟壁面的黏性/回旋黏性耦合，并非已解释实验自发转动源。遗漏平衡pinch流造成约10\%到22.5\%的特定动量通量误差，作者明确承认。少数事件不足统计比较，压力、电阻剖面简化且S偏低。

\paragraph{本报告判断}纳入直接扩展MHD的新现象：Hall效应能协助也能阻碍磁松弛，不能一概表述为Hall必然加速；区分局部重联率、平均发电机与全局动量输运。

\section{L236：2017年}\label{paper:L236}

\textbf{Fast Magnetic Reconnection: “Ideal” Tearing and the Hall Effect}\citep{L236}\

Fulvia Pucci;  Marco Velli;  Anna Tenerani.

\textsf{核读范围：}arXiv1704.08793v1全部7页；\textsf{图表：}式1-21、图1-7；图1-7视觉核查.

\href{run:papers/L236.pdf}{打开本地PDF：\texttt{papers/L236.pdf}}\quad\href{https://doi.org/10.3847/1538-4357/aa7b82}{DOI／版本入口}

\paragraph{在报告中的研究脉络}\hyperref[story:onset]{启动、爆发与磁岛}。关联表示本报告的阅读线索，不自动表示这些论文互相引用。

\paragraph{实际模型}不可压缩均匀密度电阻Hall MHD，Harris反平行平衡；无电子惯性及压力。线性二维本征问题。

\paragraph{计算／实验细节}六阶两点本征方程、Lentini-Pereyra算法；主要局地Sa=1e5-1e8，di/a<1；未列完整网格/求根容差。以Ph=(di/a)Sa\textasciicircum{}(1/4)表征Hall层；数值拟合大Ph增长指数zeta=.564±.007。

\paragraph{核实的发现}Hall中间层变宽而内电阻层变薄；保留Hall对外内匹配的修改后，数值增长指数与朴素2/3估计不同。由拟合得到临界a/L约S\textasciicircum{}(-1/3)(h sqrt(S))\textasciicircum{}.29，使全局Alfven时间的最大增长为约.41。

\paragraph{限制与未完成问题}指数.29来自有限范围拟合，不是严格解析指数；要求di/a<1且全局h约S\textasciicircum{}(-1/2)至S\textasciicircum{}(-.27)。ideal指增长不随S减慢，并非电阻严格为零；未求非线性磁通率、动态片形成或三维斜模。原文beta误差传播写成3倍zeta误差不符合其式19导数；本报告只引用近似指数，不沿用该误差估计。

\paragraph{本报告判断}提供Hall改变破裂临界长宽比的线性证据；不能直接据此宣称真实动态电流片必在该宽度断裂。

\section{L238：2017年}\label{paper:L238}

\textbf{Hall-mediated magnetic reconnection and onset of plasmoid instability}\citep{L238}\

G. Vekstein;  K. Kusano.

\textsf{核读范围：}arXiv1704.06540全部21页，强迫重联线性过程、次级撕裂推导及参考文献完整；\textsf{图表：}无数据图；式1–47逐段，页14、18–19公式直接视觉核验，区分归一化di和实际di.

\href{run:papers/L238.pdf}{打开本地PDF：\texttt{papers/L238.pdf}}\quad\href{https://arxiv.org/abs/1704.06540}{DOI／版本入口}

\paragraph{在报告中的研究脉络}\hyperref[story:onset]{启动、爆发与磁岛}。关联表示本报告的阅读线索，不自动表示这些论文互相引用。

\paragraph{实际模型}不可压电阻Hall MHD、无黏性无电子惯性；撕裂稳定的无力Taylor问题，B=(0, B0sinαx, B0cosαx)，导体壁±a发生小δcosky形变，αa约1且远离π/2失稳。

\paragraph{计算／实验细节}纯解析标度，没有模拟/网格。S=aVA/ηm\ensuremath{\gg}1, d=di/a\ensuremath{\ll}1，β>S\textasciicircum{}-2/5容许不可压近似；通量用B0δ归一化。保留ε=αa和ka时两阈值d1≈S\textasciicircum{}-1/3(εka)\textasciicircum{}-1/3、d2≈S\textasciicircum{}-1/5(εka)\textasciicircum{}-1/5。

\paragraph{核实的发现}区分Hall开始影响与Hall控制整个重联的两个阈值：d<d1按MHD耗时S\textasciicircum{}3/5；d1<d<d2先Hall后回MHD，虽然片厚已远小于di；d>d2全程Hall，耗时S\textasciicircum{}1/2d\textasciicircum{}-1/2。原因是理想Hall外层随时间缩窄至电阻层。Hall可使次级撕裂增长更快却使最快波长更长、初生岛更少；在Hall电阻演化阶段另开多岛通道，区别于MHD流动停止后的Rutherford过程。

\paragraph{限制与未完成问题}条件δlayer<di不是普适充分判据，但本文结论针对小幅、受迫、特定时间序列，不能否定自发非线性Hall转换。所谓Hall快仍含S\textasciicircum{}1/2，非η无关。次级岛仅估计γ×片寿命>1，未积分时变增长、初始噪声及非线性岛竞争；无模拟验证。近边界稳定阈值的能量发散需非线性修正，不在主推导内。

\paragraph{本报告判断}重要的是双阈值、Hall到MHD的反向转换可能性，以及‘增长更快≠岛更多’；为报告的不确定性与待验证方案提供可检验标度。

\section{L241：2017年}\label{paper:L241}

\textbf{The role of guide field in magnetic reconnection driven by island coalescence}\citep{L241}\

Adam  John Stanier;  W. Daughton;  Andrei N. Simakov;  L. Chacón;  A. Lê;  H. Karimabadi;  Jonathan Ng;  A. Bhattacharjee.

\textsf{核读范围：}本地12页全部核读，包含模型、导引场扫描、系统尺寸扫描、镜模判别和参考文献；本地为arXiv:1611.05933v1作者稿，页首另有2018排版日期，不据此改动2017期刊年份；\textsf{图表：}图1—10及图注逐项核对；另放大查看物理页4的图1、页8的图6—7；未取得补充动画，动画不作为独立证据.

\href{run:papers/L241.pdf}{打开本地PDF：\texttt{papers/L241.pdf}}\quad\href{https://doi.org/10.1063/1.4976712}{DOI／版本入口}

\paragraph{在报告中的研究脉络}\hyperref[story:global]{层长、驱动与全局完成}；\hyperref[story:onset]{启动、爆发与磁岛}。关联表示本报告的阅读线索，不自动表示这些论文互相引用。

\paragraph{实际模型}二维磁岛合并；比较VPIC全动力学、动理学离子/无质量等温电子混合模型、标量压力Hall MHD和双种10矩模型。Hall MHD使用有限体积、BDF2、全隐式预条件Newton–Krylov；10矩使用局部隐式分裂与有限体积波传播，热流散度以各向同性松弛闭合。

\paragraph{计算／实验细节}Fadeev平衡ε=0.4，背景密度0.2n0，Ti/Te=1，x跨度2πλ、y跨度4πλ；x导体/粒子反射、y周期，扰动0.1B0。PIC质量比25、ωpe/Ωce=2。Hall/混合η=10\textasciicircum{}-5 μ0 di VA0，ηH=10\textasciicircum{}-4 μ0 di\textasciicircum{}3 VA0；Hall黏性0.01 mi n0 di VA0。10矩ke=1/de、ki=1/(3di)。扫描λ/di=5—25，另讨论50；Bg/Bm=0—5.66，Bm=0.353B0。速率以Bm VAm归一，平均窗口1.5全局Alfvén时间；正文没有完整列出各算例网格和粒子数，不能据本文单独复现所有算例。

\paragraph{核实的发现}λ=5di时，Bg从0增至B0（2.83Bm），PIC/混合峰值率分别增加约51\%/53\%，Hall MHD反降约13\%；强导引场下速率、喷流和堆积场相近。零导引场PIC喷流约为Hall的1/7，主要由离子非旋转对称压力应力限制。平均率随尺寸在零导引场PIC约λ\textasciicircum{}-0.8、Hall约λ\textasciicircum{}-0.25；Bg=B0时四模型指数聚集于-0.27至-0.33。强导引场PIC会生次级岛，而Hall仍为开放X形，却可有相近平均率。大尺度零导引场堆积区出现与镜模相容的斜向、随磁岛近静止的磁峰/密度凹陷。

\paragraph{限制与未完成问题}镜模进一步抑制重联的因果作用仍是猜测：色散计算用均匀双Maxwell假设，实测为非线性饱和阶段。10矩闭合长度由小尺度动力学结果校准，零导引场较大系统会产生PIC中没有的岛。Bg=B0的定量一致不表示压力分解一致；实际电流片宽度仍不同。质量比25、二维、有限尺寸扫描，不能把指数或导引场阈值当普适定律。

\paragraph{本报告判断}纳入导引场促进重联、平均率与局部峰值区分、不同模型相同速率但不同拓扑、10矩闭合失真四条主线；作为Hall MHD适用边界的直接比较证据。

\section{Ng2017：2017年}\label{paper:Ng2017}

\textbf{Simulations of anti-parallel reconnection using a nonlocal heat flux closure}\citep{Ng2017}\

Jonathan Ng;  Ammar H. Hakim;  Amitava Bhattacharjee;  Adam  John Stanier;  W. Daughton.

\textsf{核读范围：}1-5（全部）；\textsf{图表：}图1-5文字及图注；主要结果图已视觉复核.

\href{run:papers/Ng2017.pdf}{打开本地PDF：\texttt{papers/Ng2017.pdf}}\quad\href{https://doi.org/10.1063/1.4993195}{DOI／版本入口}

\paragraph{在报告中的研究脉络}\hyperref[story:global]{层长、驱动与全局完成}。关联表示本报告的阅读线索，不自动表示这些论文互相引用。

\paragraph{实际模型}全Maxwell十矩；非磁化三维Hammett-Perkins热流Fourier闭合，chi=sqrt(4/(9pi))；有限体积输运、局地隐式源与FFTW。

\paragraph{计算／实验细节}Fadeev岛、epsilon=.4, nb=.2, Lx×Lz=4piλ×2piλ, mi/me=25, vt, e/c=.35, Ti=Te, beta=1, ωpe/Ωce=2，10\%磁扰动；平均0-1.5tA。

\paragraph{核实的发现}非局地闭合得到1.20di片宽，PIC1.15di，Hall .18di；离子压力支撑区与速度、非旋转对称度改善；平均率λ\textasciicircum{}-0.45比局地\textasciicircum{}-0.2更接近PIC\textasciicircum{}-0.8。

\paragraph{限制与未完成问题}仍未复现PIC尺寸标度，离子非旋转对称度偏低；次级岛影响速率；电子各向异性偏弱；非磁化闭合不应直接外推强导引场。Langmuir阻尼测得.044ωpe，精确Vlasov .034ωpe，存在约29\%误差。

\paragraph{本报告判断}成功是特定离子结构改善与减少调参，不是所有动力学完全重现。

\section{Yoo2017Heating：2017年}\label{paper:Yoo2017Heating}

\textbf{Electron heating and energy inventory during asymmetric reconnection in a laboratory plasma}\citep{Yoo2017Heating}\

Jongsoo Yoo;  Byungkeun Na;  J. Jara-Almonte;  Masaaki Yamada;  Hantao Ji;  V. Roytershteyn;  M. R. Argall;  W. Fox;  Li-Jen Chen.

\textsf{核读范围：}所记录版本全部18个PDF物理页；正文、方法、结果、讨论、图表、参考文献及适用文内附录均核读。；\textsf{图表：}直接查看PDF物理页7, 8, 10, 12, 14, 15；Figs4, 5, 8, 10, 11;  Table1;  Eqs5–10；2026-11-13纠正定位：本地PDF图10、11在第15页，式(10)在第14页，图8在第12页。.

\href{run:papers/Yoo2017Heating.pdf}{打开本地PDF：\texttt{papers/Yoo2017Heating.pdf}}\quad\href{https://doi.org/10.1002/2017JA024152}{DOI／版本入口}

\paragraph{在报告中的研究脉络}\hyperref[story:energy]{电子独有重联与热化}；\hyperref[story:diagnostic]{实验与可区分诊断}。关联表示本报告的阅读线索，不自动表示这些论文互相引用。

\paragraph{实际模型}本文研究两侧密度不同的磁重联实验，并使用二维和三维全动理学粒子模拟解释电子加热及物种能量收支。实验重点测量低密度侧分离面附近的电场、温度和波动；模拟比较温度矩与速度分布形状的变化。因此它既约束空间非对称条件下的能量分配，也检验分布趋于麦克斯韦形态是否必然伴随额外净增能。

\paragraph{计算／实验细节}实验两侧密度比约为 8，重联磁场分别约为 145 与 168 高斯，残余导场小于重联场的一成，离子皮肤深度约 7.5 厘米。粒子模拟质量比为 900、密度比为 10，初始两物种等温，高密度侧热压与磁压比为 0.5。二维盒长为 10 乘 10 个离子皮肤深度、网格 1024 平方；三维第三方向长 3 个离子皮肤深度，增加 320 层网格。

\paragraph{核实的发现}低密度侧分离面具有增强的电子加热。图 10 中电子和离子沉积率分别为 0.35±0.07 与 0.44±0.09 兆瓦，占输入磁能约百分之二十五和三十一；两物种热能变化项均约 0.30 兆瓦。二维与三维模拟温度剖面相近，三维电子分布却更接近麦克斯韦形态，说明分布热化与净能量增长可以表现不同。

\paragraph{限制与未完成问题}实验波动与温升的相关系数约为 0.29，显著性水平约为 0.086，不能确立低混杂漂移波的强因果作用。模拟的电子温度、频率比和有限盒尺度与实验仍有差异。正文温度矩定义及表格密度单位存在排印疑点，应按标准定义和正文量级核对，不能据此反推代码错误。原始数据入口已记录，但本次未重新分析。

\paragraph{本报告判断}非对称实验下两物种沉积率相近，直接限制了对称实验约二比一结果的外推。三维波动可能促进分布散射而不显著增加温度，这提示流体闭合应分别接受能量预算和分布诊断检验。本文支持针对具体条件研究物种分配，但未建立跨装置、跨采样位置的普适比例；正式版 18 页及能量收支图已完整核读。

\section{E553：2018年}\label{paper:E553}

\textbf{Using the maximum entropy distribution to describe electrons in reconnecting current sheets}\citep{E553}\

Jonathan Ng;  Ammar H. Hakim;  Amitava Bhattacharjee.

\textsf{核读范围：}arXiv1808.07822v1全部6页，含最大熵闭合推导与全部重建结果；\textsf{图表：}图1–5，直接核验页3采样位置、页4图2–4的上游平顶、X点三角形、Grad负值区域.

\href{run:papers/E553.pdf}{打开本地PDF：\texttt{papers/E553.pdf}}\quad\href{https://doi.org/10.1063/1.5041758}{DOI／版本入口}

\paragraph{在报告中的研究脉络}\hyperref[story:energy]{电子独有重联与热化}。关联表示本报告的阅读线索，不自动表示这些论文互相引用。

\paragraph{实际模型}从PSC PIC数据取矩，最大熵f=exp(α·m)重建5/10/14/21矩分布，并对照Grad20；不是演化高阶矩的流体重联模拟。固定目标矩U时最小化h(α)-α·U，Hessian=<mm\textsuperscript{T}f>正定给对称双曲结构；14/21含|v|\textsuperscript{4}确保可积性，但并非所有物理可实现矩都在该指数族的参数域内。

\paragraph{计算／实验细节}反平行Harris PIC质量比25、ωpe/ωce2、Ti/Te5；tωci20，上游/X点/下游1.7di处边长2.5de采样盒，各约10000粒子，Newton迭代求α。分布以动量表示但使用非相对论p=mev闭合。

\paragraph{核实的发现}十矩捕获上游各向异性，十四矩的四次矩补平顶；X点及排出区须21矩才近似三角形，仍丢失双峰与多次穿越轨道细节。Grad20产生负分布，最大熵保持正性；增加矩数改善分布误差但计算代价大。

\paragraph{限制与未完成问题}未做21矩时演、稳定性或重联率验证；重建所用矩直接来自PIC，不能把拟合成功当预测成功。PIC噪声影响误差，Grad比较使用|f|进入Hellinger式但未说明重归一化，非严格同类概率距离；作者H定义相当常用平方距离。闭合无粒子捕获历史信息；速度积分加优化每时步昂贵。

\paragraph{本报告判断}第一章推导约束熵极值和矩系统双曲性，第四章明确高阶矩的可实现性、计算成本以及分布细节的不足，避免把闭合论文算作新流体模拟现象。

\section{L244：2018年}\label{paper:L244}

\textbf{A reduced Landau-gyrofluid model for magnetic reconnection driven by electron inertia}\citep{L244}\

E. Tassi;  D. Grasso;  D. Borgogno;  T. Passot;  P. L. Sulem.

\textsf{核读范围：}HAL02049658v1封面及正文1–26全读；物理页1–27；\textsf{图表：}式2.1–4.21、Hamiltonian/Casimir推导及图1–7；阻尼算子、图3–6视觉复核.

\href{run:papers/L244.pdf}{打开本地PDF：\texttt{papers/L244.pdf}}\quad\href{https://doi.org/10.1017/S002237781800051X}{DOI／版本入口}

\paragraph{在报告中的研究脉络}\hyperref[story:closure]{压力、热流与闭合}；\hyperref[story:energy]{电子独有重联与热化}。关联表示本报告的阅读线索，不自动表示这些论文互相引用。

\paragraph{实际模型}从Scott型gyrofluid按epsilon低频、delta\textsuperscript{2}=me/mi、beta\_e\textasciitilde{}delta\textsuperscript{2}约化为六场；含离子FLR、电子双温、电子惯性和非局地Landau闭合。实际重联只算冷离子、Ni=Tperpi=0、Tperpe恒零的子模型。

\paragraph{计算／实验细节}delta=.1, beta\_e=.05, lambda3；线性半盒Lx46, Ly24pi, Lz96pi；非线性ky1/12, kz1/48，Ly12pi, Lz48pi。显式AB3、y/z并行；非线性热流刚性使步长至多缩小100倍，改用单螺度坐标y*=y+(kz/ky)z的等价二维推进。全文未列具体格点数及绝对步长。

\paragraph{核实的发现}Landau算子用局地平行梯度和磁场方向二阶空间平均的Fourier分母，非简单固定k0扩散。阻尼降低平行温度扰动、线性增长率，高m斜模可稳定；同一gamma*t上岛面积曲线相似，意味着真实时间中发展变慢，而非未受影响。饱和岛大小近似不变，细密度/电流/温度结构变平滑。闭合去除自由能，未出现动理学参考中的重联—加热延迟。

\paragraph{限制与未完成问题}斜单螺度不是多螺度三维湍流；没有热离子完整六场非线性验证。无Landau时数值能损约总能量.2\%，却约为同期磁能释放20\%，不可忽略。未出现延迟可能是闭合过度耗散，作者提出高阶闭合问题。自由能损失解释为未解析相混合/电子加热，非直接分布函数测量。

\paragraph{本报告判断}给出可推导、可检验的闭合成功与失败并存例子；必须同时检验时序、饱和、能量预算，不能以一项吻合宣称动理学等价。

\section{L250：2018年}\label{paper:L250}

\textbf{On the Hall-mediated resistive tearing instability of highly elongated current sheets}\citep{L250}\

Grigory Vekstein.

\textsf{核读范围：}arXiv1811.11401全部18页，全部常ψ/非常ψ体制及两张参数分区图；\textsf{图表：}图1–2与式1–67；页10/13非常ψ公式及页17/18分区图直接视觉核验.

\href{run:papers/L250.pdf}{打开本地PDF：\texttt{papers/L250.pdf}}\quad\href{https://arxiv.org/abs/1811.11401}{DOI／版本入口}

\paragraph{在报告中的研究脉络}\hyperref[story:onset]{启动、爆发与磁岛}。关联表示本报告的阅读线索，不自动表示这些论文互相引用。

\paragraph{实际模型}电阻Hall MHD，冷离子而电子热压提供可压缩响应，无黏性、无电子惯性；Harris无力平衡By=B0tanh(x/L), Bz=B0sech(x/L)，二维线性渐近分析。此文L是平衡片厚尺度，不是沿片总长。

\paragraph{计算／实验细节}S=LVA/ηm\ensuremath{\gg}1, d=di/L≤1，k以1/L归一化，β=4πΓP0/B0\textsuperscript{2}，与常用8πP/B\textsuperscript{2}不同。Δ'L=2(1/k-k)，外层准静态要求γτA/k<1；极低β由惯性承担压缩的区域未纳入主分区。无数值计算/网格。

\paragraph{核实的发现}统一常ψ的MHD、压缩控制Hall、近EMHD三支和非常ψ的新增两支。非常ψ高β有γτA≈kd、δ/L≈S\textasciicircum{}-1/2(kd)\textasciicircum{}-1/2；低β有γτA≈β\textasciicircum{}1/2kd、δ再乘β\textasciicircum{}-1/4，增长式不显含S但断线仍依赖有限η。β≥1全谱最快模在d>S\textasciicircum{}-1/4后移向更长波：设Z=dS\textasciicircum{}1/4，kmax≈Z\textasciicircum{}-1/3S\textasciicircum{}-1/4、γmaxτA≈Z\textasciicircum{}2/3S\textasciicircum{}-1/2。

\paragraph{限制与未完成问题}这是渐近比例和体制图，未求数值系数/完整本征函数，也未验证非线性快率。固定k体制内的S独立不是对任意S取极限：S变化会跨过常ψ界线；高β最快增长仍含S\textasciicircum{}-1/3 d\textasciicircum{}2/3。须检验沿片长度容纳更长最强模、外层准静态和δ>de；d→∞不能简单外推，外层惯性会改Δ'。结论中的太阳/磁尾量级仅示例，不是观测事件验证。

\paragraph{本报告判断}可与L122比较说明为何保留Δ'、β、波长与外层匹配至关重要；适合第一章教学用的分支匹配和第三章‘Hall增加增长却减少初生岛数’条件化结论。

\section{L251：2018年}\label{paper:L251}

\textbf{On the role of system size in Hall MHD magnetic reconnection}\citep{L251}\

Christopher Bard;  J. Dorelli.

\textsf{核读范围：}arXiv1710.03612v1全部5页；\textsf{图表：}图1-7全部视觉核查.

\href{run:papers/L251.pdf}{打开本地PDF：\texttt{papers/L251.pdf}}\quad\href{https://doi.org/10.1063/1.5010785}{DOI／版本入口}

\paragraph{在报告中的研究脉络}\hyperref[story:global]{层长、驱动与全局完成}。关联表示本报告的阅读线索，不自动表示这些论文互相引用。

\paragraph{实际模型}二维可压缩Hall有限体积；不含显式电阻与电子惯性，以网格尺度控制数值磁耗散。

\paragraph{计算／实验细节}Fadeev平衡epsilon.4；Lx=4pi lambda、Ly=2pi lambda，x周期、y导体。lambda6.04-25.46di，网格.033-.2di，部分至.012di；最大95.37di用.1与.04di。初始T常数，p=.5(1-epsilon\textsuperscript{2})/den\textsuperscript{2}+.5p0，扰动磁场.1。

\paragraph{核实的发现}网格细化到约.04-.05di出现从长薄层到开放双尺度层的转变；中间细度可既有开放排气又有磁岛弹回。lambda14.48时.025di可消除弹回，但最大系统两种网格均弹回并产生次级岛；快峰值并不保证快速完成合并。

\paragraph{限制与未完成问题}细化同时改变数值电阻与可解析色散尺度，不是固定物理耗散下的收敛扫描。没有证明物理双稳态或普适.1极限；与PIC次级岛差别归因离子陀螺黏性只是待检验解释。文末速率与lambda正比表述和正文下降趋势不一致，应查图与具体平均窗。

\paragraph{本报告判断}重要证据是Hall模型也能出现弹回；模型差异和离散耗散差异不能混为一谈。

\section{L256：2018年}\label{paper:L256}

\textbf{Temperature gradient driven heat flux closure in fluid simulations of collisionless reconnection}\citep{L256}\

F. Allmann‐Rahn;  T. Trost;  Rainer Grauer.

\textsf{核读范围：}期刊PDF1-16全部；\textsf{图表：}图1-10及表1；图1、3-6、8-10与表1视觉复核.

\href{run:papers/L256.pdf}{打开本地PDF：\texttt{papers/L256.pdf}}\quad\href{https://doi.org/10.1017/s002237781800048x}{DOI／版本入口}

\paragraph{在报告中的研究脉络}\hyperref[story:closure]{压力、热流与闭合}。关联表示本报告的阅读线索，不自动表示这些论文互相引用。

\paragraph{实际模型}全Maxwell十矩、对压力偏离各向同性的部分作Laplacian闭合，不是直接扩散完整温度；以Vlasov热流高阶矩作对照。

\paragraph{计算／实验细节}GEM质量比25、256×128基准及1024×512加密；另流体质量比400/1836。k0s分别为GEM3/ds、100×50di WHBG的1/(3ds)、岛合并1/(2ds)；不是一组参数全场景。岛lambda5/10/15/25di，两套网格512×256到2048×1024。

\paragraph{核实的发现}原松弛闭合的热流散度形状和符号在某些区域错误；梯度闭合改善上游假热流与电子非对角压力。岛合并平均0-1.5tA率标度约lambda\textasciicircum{}(-.73)，出现反弹且无二次岛。表1lambda5至25平均率.180→.052；两套网格差小。GEM高分辨梯度流体后期斜率近Vlasov，但启动仍迟；高分辨局地松弛反而更慢。

\paragraph{限制与未完成问题}k0按与动理学一致调节，尺寸标度也依赖k0；不能宣称无参预测。实际质量比只展示流体结构，不是同质量比完整动理学验证。显式Laplacian子循环中不交换块边界、不更新速度密度，作者承认块边界误差。GEM温度条件缺常规1/2，岛平衡坐标与边界文字不一致，复现需输入核对。

\paragraph{本报告判断}重要改进是直接检查热流矩再验证非线性反弹；其成功同时有问题相关的校准与数值限制。

\section{Yamada2018：2018年}\label{paper:Yamada2018}

\textbf{The two-fluid dynamics and energetics of the asymmetric magnetic reconnection in laboratory and space plasmas}\citep{Yamada2018}\

Masaaki Yamada;  Li‐Jen Chen;  Jongsoo Yoo;  Shan Wang;  W. Fox;  Jonathan Jara-Almonte;  Hantao Ji;  W. Daughton;  A. Lê;  James L. Burch;  Barbara L. Giles;  M. Hesse;  Thomas Earle Moore;  R. B. Torbert.

\textsf{核读范围：}期刊PDF1-11全部；未核读补充电影；\textsf{图表：}图1-8与表1-2；PDF4、8-10页视觉核验.

\href{run:papers/Yamada2018.pdf}{打开本地PDF：\texttt{papers/Yamada2018.pdf}}\quad\href{https://doi.org/10.1038/s41467-018-07680-2}{DOI／版本入口}

\paragraph{在报告中的研究脉络}\hyperref[story:diagnostic]{实验与可区分诊断}。关联表示本报告的阅读线索，不自动表示这些论文互相引用。

\paragraph{实际模型}MRX氦放电与MMS2015-10-16观测，辅以VPIC全粒子模拟；不是双流体流体求解器。

\paragraph{计算／实验细节}密度比约10，无显著导引场；MRX电子流由Ui-J/(en)反演。VPIC40di×20di，4200×2100，mi/me100，omega\_pe/omega\_ce2，400粒子/格/物种约70亿总粒子；t=96/omega\_ci在2.5di×3di盒积分能量。

\paragraph{核实的发现}强不对称时Hall四极接近双极，电子与离子停滞点、X点分离；电子停滞点偏移5-10de，离子约0.3-0.5di。电子能量转移集中于电子停滞点且垂直电功主导；离子被偏向低密度侧的势阱加速，MMS有反向离子束而实验较热化。表2非对称MRX入射磁能归一化后出射磁能.44、离子.31、电子.25；PIC分别.43/.32/.25。

\paragraph{限制与未完成问题}MMS有限轨迹不能检验全局能量分配；电场滤除高于离子回旋频率成分，不能用此排除高频能量交换。作者未找到波幅与率的确定关系。MRX电子速度是推算；模拟二维且质量比100。方法段将X写周期、Z写导体/反射，与X为片法向及单tanh初值需澄清，复现前应查输入文件。

\paragraph{本报告判断}支持跨尺度两种粒子解耦形态及近X点能量转移图样；不能据此宣称Hall项本身耗散，或五矩/十矩闭合已被验证。

\section{Yoon2018Heating：2018年}\label{paper:Yoon2018Heating}

\textbf{Fast Ion Heating in Transient Collisionless Magnetic Reconnection via an Intrinsic Stochastic Mechanism}\citep{Yoon2018Heating}\

Young Dae Yoon;  Paul M. Bellan.

\textsf{核读范围：}所记录版本全部7个PDF物理页；正文、方法、结果、讨论、图表、参考文献及适用文内附录均核读。；\textsf{图表：}直接查看PDF物理页2, 3, 4, 5, 6；All Figs1–3;  Eqs2–17.

\href{run:papers/Yoon2018Heating.pdf}{打开本地PDF：\texttt{papers/Yoon2018Heating.pdf}}\quad\href{https://doi.org/10.3847/2041-8213/aaf0a3}{DOI／版本入口}

\paragraph{在报告中的研究脉络}\hyperref[story:energy]{电子独有重联与热化}。关联表示本报告的阅读线索，不自动表示这些论文互相引用。

\paragraph{实际模型}本文由二维电子磁流体模型生成重联电磁场，再在场中独立积分测试离子轨道。电子惯性保留，冷电子或正压条件使电子正则涡量冻结，而磁场仍可重联；离子背景在流体计算中固定。因此离子速度展宽来自附加的粒子轨道计算，并非一个自洽双温流体方程组直接演化出的离子与电子温度。

\paragraph{计算／实验细节}电子流体取恒密度和扰动哈里斯电流片，200 乘 600 个网格点覆盖 30 乘 120 个电子皮肤深度，演化至 580 个电子回旋逆频率。6000 个测试离子初始静止，比较冻结早期场与随时间增长的场。图示质量比为 10，另报告 20、50、100 的定性相似结果；有限初温例的热速为电子皮肤深度与回旋频率乘积的百分之一。

\paragraph{核实的发现}增长的电子流体场使测试离子轨道分离、相空间混合和流出速度方差增加，而冻结早期场主要形成有序运动。近似分析指出，入流方向电场变化尺度小于离子皮肤深度时可出现失磁化，出流方向也容易满足判据。用磁重联实验的场强和尺度估计得到约 14.6 的失磁化参数，与快速离子增能在量级上相容。

\paragraph{限制与未完成问题}测试离子不反过来改变电场，分布保持非麦克斯韦形态，也未计算碰撞熵产。电子非正压张量被舍弃，单凭准中性不能普遍忽略密度梯度；原文局地等离子体压强参数的表达也不能直接充当定量闭合。最终加热量、质量比与时间标度仍待研究，原始流体实现部分依赖前序方法论文，本次未重新运行。

\paragraph{本报告判断}这是流体生成场能够帮助研究离子局部增能的直接例子，但应明确其电子磁流体加测试粒子的混合方法身份。它说明霍尔尺度电势可能触发轨道失磁化和速度展宽，尚未解决电子温度、自洽离子反馈及不可逆热化比例。与实验的量级一致不能排除碰撞或波粒作用；本文全部 7 页及三幅图已核读。

\section{F659：2019年}\label{paper:F659}

\textbf{Three‐Dimensional Magnetic Reconnection With a Spatially Confined X‐Line Extent: Implications for Dipolarizing Flux Bundles and the Dawn‐Dusk Asymmetry}\citep{F659}\

Yi‐Hsin Liu;  T. C. Li;  M. Hesse;  W. J. Sun;  J. Liu;  J. Burch;  J. A. Slavin;  K. Huang.

\textsf{核读范围：}arXiv:1901.10195v1全部9页，含正文1–8页、参考文献8–9页；无文内附录。实际阅读版本是预印本，不冒充12页正式版。；\textsf{图表：}图2–9；直接查看PDF3–7页，比较不同色标的Bz与离子喷流、图8完整欧姆预算和电子漂移分解。.

\href{run:papers/F659.pdf}{打开本地PDF：\texttt{papers/F659.pdf}}\quad\href{https://doi.org/10.1029/2019ja026539}{DOI／版本入口}

\paragraph{在报告中的研究脉络}\hyperref[story:three]{三维传播与有限X线}。关联表示本报告的阅读线索，不自动表示这些论文互相引用。

\paragraph{实际模型}三维全PIC动理学对照；厚电流片夹住薄片以限制X线沿电流方向的延伸。Hall输运是解释机制，并非一项独立Hall流体模拟。

\paragraph{计算／实验细节}薄/厚半宽0.5/4di，过渡S=5di，nb=.3n0；薄片长度31、12、8.4、4di。域32×64×16di、768×1536×384格，mi/me=75、ωpe/Ωce=4。x/y周期，z导体/粒子反射；固定离子漂移以降低过渡处惯性失衡，厚区Ti/Te=5。率为对称中面重联磁通变化除薄片长度。

\paragraph{核实的发现}同一X线分为电子漂移侧活跃区与离子漂移侧受抑区；薄片长度小于约10di时率和喷流显著下降。电子沿电流反方向搬走Bz，削弱出流驱动和进一步变薄；受抑长度约为电子漂移速度乘变薄时间。长例tΩci=12局地非理想电场约.12，而沿线积分再除31di约.068；二者不是同一种率。活跃区主要由非对角电子压力散度支撑断线。

\paragraph{限制与未完成问题}10di是所用初值和参数下的尺度估计；没有本篇质量比扫描。初始状态非严格三维动理学平衡；反平行、有限周期盒与人为厚片限制外推。沿线对流惯性闭积分为零是对称/周期下的结果，密度加权压力梯度积分小不是一般恒等式。与早期双流体受抑区层形不同尚待解释；行星晨昏差异仅为机制候选。

\paragraph{本报告判断}连接GEM局地快率、早期有限长X线、传播和全局几何：Hall输运可帮助二维排气，也能在第三方向搬走维持某区域重联所需的磁通。不要把PIC中的压力破冻结改称Hall项直接断线。

\section{F743：2019年}\label{paper:F743}

\textbf{Kinetic Verification of the Stochastic Ion Heating Mechanism in Collisionless Magnetic Reconnection}\citep{F743}\

Young Dae Yoon;  Paul M. Bellan.

\textsf{核读范围：}所记录版本全部5个PDF物理页；正文、方法、结果、讨论、图表、参考文献及适用文内附录均核读。；\textsf{图表：}直接查看PDF物理页2, 3, 4；All Figs1–5;  Eqs1–6.

\href{run:papers/F743.pdf}{打开本地PDF：\texttt{papers/F743.pdf}}\quad\href{https://doi.org/10.3847/2041-8213/ab5b0a}{DOI／版本入口}

\paragraph{在报告中的研究脉络}\hyperref[story:energy]{电子独有重联与热化}。关联表示本报告的阅读线索，不自动表示这些论文互相引用。

\paragraph{实际模型}本文先分析哈里斯平衡中的电势及有限导场对离子失磁化区域的影响，再用自洽全动理学粒子模拟检验温度与判据的空间联系，并跟踪测试离子的轨道和伪磁矩。它是前序电子流体加测试粒子机制的动理学验证，保留完整粒子分布，不能计作霍尔流体独立求解分物种热方程的实例。

\paragraph{计算／实验细节}二维双哈里斯电流片采用周期边界，半宽等于离子皮肤深度，背景密度为峰值的 0.2，两物种初始等温，质量比为 100。1024 乘 2048 网格覆盖 40.96 乘 81.92 个离子皮肤深度，每格 100 至 600 个粒子，光速与阿尔芬速度比为 20。导场比扫描零、0.1、0.3、0.5；混合离子例含一成质量比 500 的重离子。

\paragraph{核实的发现}电场梯度失磁化判据满足的区域与离子中心二阶矩温度增强区相关，在所考察范围内二者随导场增加而减弱。较高质量电荷比离子更容易满足判据并升温；跟踪轨道显示相邻粒子分离及伪磁矩跳变。静态哈里斯解给出导场和电流片宽度的阈值趋势，但作者明确不把该阈值直接等同于演化重联的精确边界。

\paragraph{限制与未完成问题}图 5 仍有判据不能解释的额外平行温升，作者将其留作后续问题。有限时间内两条轨道分离和高温区对应不是完整熵产预算；本文未进行广泛三维、压强比、质量比及数值收敛扫描。电场梯度判据也没有给出电子能量变化，因而不能从这组模拟直接推出一般的离子与电子不可逆加热比例。

\paragraph{本报告判断}该研究加强了低至中等导场条件下随机离子增能机制的证据，并说明不同质量电荷比可造成选择性加热。然而机制的有效参数范围、其他并存增能过程及定量能量份额仍需检验。后续 L272 使用了这里的粒子证据，不能重复计作独立验证；本次完整核读正式版 5 页并直接检查全部五幅图和关键公式。

\section{L243：2019年}\label{paper:L243}

\textbf{A GPU-enabled implicit Finite Volume solver for the ideal two-fluid plasma model on unstructured grids}\citep{L243}\

Alonso Asensio,  Isaac;  Alvarez Laguna,  Alejandro;  Mohamed Hassanine Aissa;  Stefaan Poedts;  Nataly Ozak;  Andrea Lani.

\textsf{核读范围：}arXiv1809.02532v2全部23页，含代码清单1–4与完整性能表；正式期刊2019年CPC239卷16–32页；\textsf{图表：}图1–9、表1–11；直接核验页18图7/8及页20性能图9、表10/11.

\href{run:papers/L243.pdf}{打开本地PDF：\texttt{papers/L243.pdf}}\quad\href{https://doi.org/10.1016/j.cpc.2019.01.019}{DOI／版本入口}

\paragraph{在报告中的研究脉络}\hyperref[ch:numerics]{模型与数值可信度}。关联表示本报告的阅读线索，不自动表示这些论文互相引用。

\paragraph{实际模型}五矩理想双流体+Maxwell+双曲散度清理共18变量，COOLFluiD非结构有限体积；二阶最小二乘重构、Venkatakrishnan限制器，流体AUSM+up、场改进Rusanov，BDF2全隐式Newton、数值差分Jacobian、GMRES/Jacobi。

\paragraph{计算／实验细节}每GPU线程一个单元并计算通量/源项及Jacobian，CSR装配与PARALUTION避免主机往返；边界等仍在CPU。波测N64/128/256、CFL1/2/5；GEM双电流片全周期256\textsuperscript{2}/512\textsuperscript{2}、mi/me25、Te/Ti.2、γ5/3，tΩci25用10000步即dtΩci.0025，P100完整计算49小时16分。

\paragraph{核实的发现}GPU与同算法CPU磁通曲线吻合，空间二阶通过波误差测量。仅装配移至GPU约3–4倍，512\textsuperscript{2} P100总加速14.23倍；小问题PARALUTION可能反而变慢。大网格主要耗时仍在Jacobian装配，完整耦合矩阵显存限制某些组合。

\paragraph{限制与未完成问题}14.23倍相对本机单CPU核心，换更快K40主机CPU基准约8.15倍，不能解释为相对现代多核集群加速。PETSc3.9/PARALUTION1.1性能为当年版本；波及同代码CPU对照不证明压力闭合和电子扩散区正确。GEM部分尺度和时间步细节委托前作，本文非完整独立复现说明；无网格收敛的局部Ohm预算。

\paragraph{本报告判断}报告明确区分全隐式、矩阵装配、物理精度与GPU加速，并标注硬件和比较基准；不把性能论文作为新重联物理发现。

\section{L247：2019年}\label{paper:L247}

\textbf{High-order two-fluid plasma solver for direct numerical simulations of plasma flows with full transport phenomena}\citep{L247}\

Z. Li;  D. Livescu.

\textsf{核读范围：}PDF1–44: entire text and references including Appendix A plasma transport,  B quasineutrality,  C single-fluid limits.；\textsf{图表：}Direct visual Eq15–19 (p11),  identical-temperature restriction (p16),  Figs1, 3, 4, 7, 9, 10, 11 (pp18, 21, 25, 27, 35),  AppendixA EqsA1–A9, A15–A19 (pp31, 33). All other figures/captions read in text..

\href{run:papers/L247.pdf}{打开本地PDF：\texttt{papers/L247.pdf}}\quad\href{https://doi.org/10.1063/1.5082190}{DOI／版本入口}

\paragraph{在报告中的研究脉络}\hyperref[story:energy]{电子独有重联与热化}。关联表示本报告的阅读线索，不自动表示这些论文互相引用。

\paragraph{实际模型}构建准中性、忽略电子惯性与位移电流的Braginskii两流体方程，保留独立物种内能、磁场方向相关黏性和热流、摩擦、热力、电阻加热及温度交换；附录列出输运系数和单流体极限。但第IV节明确四项实际演示均强制Ti=Te，并删去独立离子能量演化，实际测试的是单温扩展Hall系统，不能当作两温分配验证。

\paragraph{计算／实验细节}波测试用384点；激波Rei为6.3×10\textsuperscript{3}和2.3×10\textsuperscript{4}。重联di/L=1、Rei约124–850、Rem约58–890，主例426/112；参考长度10 m、密度10\textsuperscript{15} m\textsuperscript{-3}。CFDNS用六阶紧致差分和物理输运；主例约256×512收敛，以768×1536为参考，观测阶数6.58。正文未给具体时间步，另有更高Re和激波网格扫描。

\paragraph{核实的发现}在所选物理输运系数下，波色散、激波和重联基准得到高阶网格收敛；局地共同温度图显示磁能转化及热扩散的重要性。激波误差阶数约6.08–6.14，重联磁通在所试较大黏性和磁雷诺数处趋于弱敏感。论文提供了完整物种热方程及物理正则化的实现路径，但没有给出独立Ti/Te、交换过程或加热比例的验证结果。

\paragraph{限制与未完成问题}最重要边界是所有验证算例固定Ti=Te，不能从完整方程形式推断已演示双温预测。经典碰撞输运在冠层扩散层可能不满足局地平衡排序；磁通及饱和时间随网格或Re收敛，也不能证明微观热化和物种分配正确。零电子惯性等约化排除相关电子尺度；原始输运推导、CFDNS源代码及独立双温测试未复核，也无实验温度比对照。

\paragraph{本报告判断}应将本篇用作能量方程、各向异性输运和数值收敛的方法证据，尤其支持区分物理输运与隐含数值耗散。作者关于高Re宏观结果可代表无碰撞重联的讨论，只由其实际磁通等指标检验，不能扩张为无碰撞加热闭合获验证。全文44页含附录均已读；正式出版年为2019，2018是预印本日期，两者不是独立文献。

\section{L262：2019年}\label{paper:L262}

\textbf{Fast Magnetic Reconnection: Secondary Tearing Instability and Role of the Hall Term}\citep{L262}\

E. Papini;  Simone Landi;  L. Del Zanna.

\textsf{核读范围：}arXiv1906.06779v2全部14页；\textsf{图表：}表1与图1-9正文、图注；图2-9及模型/算法页视觉核查.

\href{run:papers/L262.pdf}{打开本地PDF：\texttt{papers/L262.pdf}}\quad\href{https://doi.org/10.3847/1538-4357/ab4352}{DOI／版本入口}

\paragraph{在报告中的研究脉络}\hyperref[story:onset]{启动、爆发与磁岛}。关联表示本报告的阅读线索，不自动表示这些论文互相引用。

\paragraph{实际模型}二维可压缩电阻Hall MHD，Gamma5/3含欧姆热，无电子惯性与电子压力张量。

\paragraph{计算／实验细节}力自由B=(0, tanh(x/a), sech(x/a))，rho1、beta.5，a/L=S\textasciicircum{}(-1/3)。x±20a自由流出、y周期Ly=100pi a；10随机相位模，势种子1e-6。y Fourier、x四阶紧致差分、RK3、标量磁势保持无散；线性1024×128，非线性4096×512。全局S1e5-1e8；Hall非线性S1e6、di/原始内层=.2/1.4/3。

\paragraph{核实的发现}自发次级片厚度约初始的.1-.15，在形成阶段局地a*B*/eta近常数，近压强平衡。高局地S约>3e4时长宽比接近S*\textasciicircum{}(-1/3)，较低值更接近SP；次级增长2.3对初级.48，以同一全局时钟计。原始Hall弱到不改初级增长时，次级增长仍增17\%；较强一组近加倍。再强Hall时未见次级片；局地重标度后增长仍O(1)。

\paragraph{限制与未完成问题}gamma\_rec是平均对数磁通增长或Bx\_rms指数拟合，不是E/(BVA)；末期非线性喷流不能由线性增长系数推出。Hall算例的次级长宽比并不吻合Pucci2017拟合标度，原文明确承认；参数样本有限。强Hall次级片消失可能来自快速破裂、X几何，也可能分辨率不足，未唯一确定。原文最快波数正1/4及末段内层正1/2的印刷符号不能照抄；应为相应负幂。

\paragraph{本报告判断}最重要的是弱初级Hall能在递归次级层变强；它同时展示现象增长快与理论标度未验证可以并存。

\section{L263：2019年}\label{paper:L263}

\textbf{Fast Recursive Reconnection and the Hall Effect: Hall-MHD Simulations}\citep{L263}\

Chen Shi;  Anna Tenerani;  Marco Velli;  San Lu.

\textsf{核读范围：}arXiv1906.05961v2全文1-13页；版面自动日期2021不改期刊2019年；\textsf{图表：}图1-8正文图注；图2/4-8与模型页逐图视觉核验.

\href{run:papers/L263.pdf}{打开本地PDF：\texttt{papers/L263.pdf}}\quad\href{https://doi.org/10.3847/1538-4357/ab33ff}{DOI／版本入口}

\paragraph{在报告中的研究脉络}\hyperref[story:onset]{启动、爆发与磁岛}；\hyperref[story:energy]{电子独有重联与热化}。关联表示本报告的阅读线索，不自动表示这些论文互相引用。

\paragraph{实际模型}2.5D可压缩标量绝热Hall MHD，显式均匀电阻；无电子惯性、电子压力项、压力张量或欧姆热反馈。

\paragraph{计算／实验细节}双周期8×.8，双Harris片，rho∞=T=B=1，a0=.02；4096×2048、dx.00195/dy.000391，eta1e-5、S1e5。谱微分、显式RK3；相同随机相位1-128模种子；di=0/.001/.002/.003/.004/.005/.01七组。种子幅度、去混叠和过滤细节未列。

\paragraph{核实的发现}递归生岛、开放X点持续喷出次级岛、较强Hall开放无次级岛三种情形。Hall增大使启动提前、排气角约10至28度。固定归一化磁通斜率.032至.058；MHD晚期.034高于最弱Hall.032，不是Hall组在任意时段总更快。面内电流的电阻耗散占比由40\%趋55\%，平行占比约45\%；这是eta Jperp\textsuperscript{2}而非Hall项直接做功。中心线By谱由浅段约-.8及陡段-2转为近单幂-2.22；Ez谱不同且小尺度Ex可占主导。

\paragraph{限制与未完成问题}高阶最细片已到网格尺度；X点半长高di收敛到约.007仅3-4个dx，作者承认可能分辨率限制。没有独立网格/eta扫描，速率上限.06不宜普适化；谱是固定中心线和不同时间窗，不等于体积各向同性湍流谱。绝热压力不接收etaJ\textsuperscript{2}，所以“耗散”是磁能移除诊断，不是自洽粒子加热分配。波数用周期频率N/L约定，角波数需乘2pi；没有3D及压力各向异性。

\paragraph{本报告判断}补强Hall既能提前次级增长又能改变后期岛产生的证据；所谓Hall抑岛需明确阶段、几何、耗散和分辨率。

\section{L264：2019年}\label{paper:L264}

\textbf{Statistics of magnetic reconnection and turbulence in Hall-MHD and hybrid-PIC simulations}\citep{L264}\

Emanuele Papini;  Luca Franci;  Simone Landi;  Petr Hellinger;  Andrea Verdini;  Lorenzo Matteini.

\textsf{核读范围：}Nuovo Cimento42C, 23（2019）正式会议论文，物理页1–4全部；不是同组2018 arXiv1810.02210长文；\textsf{图表：}图1–2的统计、色标和谱全部视觉核验；页3.

\href{run:papers/L264.pdf}{打开本地PDF：\texttt{papers/L264.pdf}}\quad\href{https://doi.org/10.1393/ncc/i2019-19023-0}{DOI／版本入口}

\paragraph{在报告中的研究脉络}\hyperref[story:energy]{电子独有重联与热化}。关联表示本报告的阅读线索，不自动表示这些论文互相引用。

\paragraph{实际模型}二维可压电阻黏性Hall伪谱与CAMELIA动理学离子/无质量等温电子混合比较，另关Hall的MHD；面外均匀平均场、自由衰减Alfvén扰动。

\paragraph{计算／实验细节}256\textsuperscript{2}di方盒、2048\textsuperscript{2}网格、全周期，kinj di≈.2、初始δBrms/B0≈.24、beta2，Hall模型η=ν=.001，2/3滤波。成熟比较时刻MHD180、Hall150、混合200离子时间；详细推进、粒子数量和事件算法需另文。

\paragraph{核实的发现}Hall与混合磁谱在流体到亚离子范围接近，约-5/3接-3；离子速度谱仅大尺度相近，密度谱幅值因Hall绝热声速而偏移。平均重联时间MHD/Hall/混合约9.6/12.2/4.3离子时间，故磁谱相同不等于事件率相同。首批重联同时小尺度非线性时间骤减，随后20–30时间单位向较大尺度扩展，作者解释为磁岛合并。

\paragraph{限制与未完成问题}γrec/Ωi为频率型指标，本短文未完整定义事件追踪与误差，不能将图中.08直接当Ez/(BVA)。单一二维初值与参数，成熟比较时间不同；能量跨尺度跳转根据时序和谱代理推断，不是完整分尺度能流因果预算。离子加热与分布未被流体磁谱吻合验证。

\paragraph{本报告判断}保留其会议论文身份；价值在跨模型同时给谱与事件统计，揭示两类指标并不等价。与电子独有重联的谱相似问题互补，不作为新独立大规模实验。

\section{TenBarge2019MMS：2019年}\label{paper:TenBarge2019MMS}

\textbf{An Extended MHD Study of the 16 October 2015 MMS Diffusion Region Crossing}\citep{TenBarge2019MMS}\

J. M. TenBarge;  J. Ng;  J. Juno;  L. Wang;  A. H. Hakim;  A. Bhattacharjee.

\textsf{核读范围：}1-17（全部）；\textsf{图表：}PDF5、6、9、10、11页直接视觉核验式12、图1/6-10；17页全文及13图图注读完。正式HTML新增第14图未视觉核验。.

\href{run:papers/TenBarge2019MMS.pdf}{打开本地PDF：\texttt{papers/TenBarge2019MMS.pdf}}\quad\href{https://doi.org/10.1029/2019JA026731}{DOI／版本入口}

\paragraph{在报告中的研究脉络}\hyperref[story:energy]{电子独有重联与热化}。关联表示本报告的阅读线索，不自动表示这些论文互相引用。

\paragraph{实际模型}两物种全Maxwell五矩/十矩Gkeyll，局地Hammett–Perkins型压力松弛；直接有限矩流体证据。

\paragraph{计算／实验细节}mi/me100, c/vAs500；2D40.96×20.48di, 2048×1024；3D40.96×20.48×10.24di, 1024×512×256；周期双片；Bg/Bsx.099, Bmx/Bsx-1.7, Tem/Tes12.3, Tim/Tis5.63；k0s ds=.1，种子.1，磁噪声.0005。

\paragraph{核实的发现}图1直接给离子/电子热能和流能，十矩2D/3D全域预算相近、离子热增量较大；电子各向异性峰约2.5近MMS幅度但位置在磁鞘侧而观测在磁层侧；局地闭合抑制LHDI，剪切/kink仍出现。

\paragraph{限制与未完成问题}热能增量不是物种熵产Qi/Qe；五矩过快LHDI并明显总能损失，不可信作分配预测；未给网格/质量比独立收敛；nm/ns=16.7与低密度磁层文字冲突，l0=2π/k0≈d与k0d=.1冲突；不静默修正输入。

\paragraph{本报告判断}模型身份、温度/能量/熵产区别及跨文献判断见独立加热专题；不把本篇单独作为普适分配定律。

\section{F792：2020年}\label{paper:F792}

\textbf{Scaling of Magnetic Reconnection With a Limited X‐Line Extent}\citep{F792}\

Kai Huang;  Yi‐Hsin Liu;  Quanming Lu;  Michael Hesse.

\textsf{核读范围：}GRL47 e2020GL088147正式版全部10页（正文1–8、参考文献8–10）；从USTC作者公开目录取得。旧5页预印本保留，但本次判断以正式版为准。；\textsf{图表：}图2–5及式1–3；直接查看PDF4、6、7页，核对活跃长度定义、动量预算、梯形重叠模型及拟合参数。.

\href{run:papers/F792_published.pdf}{打开本地PDF：\texttt{papers/F792\_published.pdf}}\quad\href{https://doi.org/10.1029/2020gl088147}{DOI／版本入口}

\paragraph{在报告中的研究脉络}\hyperref[story:three]{三维传播与有限X线}。关联表示本报告的阅读线索，不自动表示这些论文互相引用。

\paragraph{实际模型}三维VPIC动理学参数扫描与简化几何模型；延续F659的有限长薄片装置，量化Hall磁通输运造成的受抑区。

\paragraph{计算／实验细节}32×64×16di域、640×1280×320格、每物种超过2.6×10\textasciicircum{}10粒子；mi/me=25、c=20VA、上游β=.3，δBz=.05B0。薄片长度40、30、20、15、10、9、7.5、5di；半宽/过渡同F659。活跃区定义为δ<2de, local；率按整个薄片长度或活跃长度分别归一。

\paragraph{核实的发现}有活跃区的算例按Lact归一后率约O(.1)，按整个X线长度归一则随其缩短下降。Jy与Bz沿电流方向错位，使其乘积和离子出流驱动力降低。梯形模型给短活跃区Vmax\ensuremath{\propto}Lact/(2Ltrans)；图5用R*=.15、Vm=.75VA、Ltrans=4di及由模拟拟合的Lsupp得到主要趋势。

\paragraph{限制与未完成问题}R与R*的比值本身来自定义，预测内容在R*近恒定和Lsupp模型，不能把恒等变换当独立理论验证。拟合使用本组数据，不是无参数预测。7.5、5di例没有活跃区，仍有种子引发的短暂率/喷流，零输出模型解释不了；9di例存在长瞬态。作者明确未严格建立最大喷流与总率的定量联系。导引场、高β、斜模与不同离子转向解释尚待检验。

\paragraph{本报告判断}将F659的现象发现推进到可检验标度，同时保留拟合、瞬态和诊断定义；与F659属于同一研究序列，不能充当两套独立跨模型复现。

\section{Kaptanoglu2020：2020年}\label{paper:Kaptanoglu2020}

\textbf{Two-temperature effects in Hall-MHD simulations of the HIT-SI experiment}\citep{Kaptanoglu2020}\

A. A. Kaptanoglu;  T. E. Benedett;  K. D. Morgan;  C. J. Hansen;  T. R. Jarboe.

\textsf{核读范围：}PDF 1–35: all body,  methods,  results,  conclusions,  references and appended repeated/enlarged figure pages;  no internal appendix.；\textsf{图表：}Direct PDF visual: Eq.1–13/Table I (pp3–4),  Figs.2–4 (pp7–9),  Figs.5–13/scan results (pp11–17);  overview of appended repeated images pp19–35. Key Ti/Te map and heating power plot inspected. Some noncentral scan figures reviewed at contact-sheet scale..

\href{run:papers/Kaptanoglu2020_accepted.pdf}{打开本地PDF：\texttt{papers/Kaptanoglu2020\_accepted.pdf}}\quad\href{https://doi.org/10.1063/5.0006311}{DOI／版本入口}

\paragraph{在报告中的研究脉络}\hyperref[story:energy]{电子独有重联与热化}；\hyperref[story:diagnostic]{实验与可区分诊断}。关联表示本报告的阅读线索，不自动表示这些论文互相引用。

\paragraph{实际模型}在PSI-Tet和NIMROD中采用三维两温Hall MHD，独立演化Ti、Te，欧姆定律保留Hall漂移及增强电子质量的简化电流时间导数惯性项。电子温度随电子流对流和压缩；电子得欧姆热、离子得黏性热，两者通过库仑交换及各向异性热传导联系。闭合仍是标量压力、Spitzer电阻和常数各向同性黏性，另有人工粒子扩散。

\paragraph{计算／实验细节}注入频率14.5–68.5 kHz，边界密度0.75×10\textsuperscript{19} m\textsuperscript{-3}、温度3 eV，黏性550 m\textsuperscript{2}/s、扩散250 m\textsuperscript{2}/s，有效质量比100。PSI-Tet用103887四面体、2.8 cm三阶单元，等效9 mm，Crank–Nicolson步长≤40 ns；NIMROD用1.8 cm四阶极向单元和11环向模。两代码注入器几何不同，另扫描壁温、密度、扩散和增强质量。

\paragraph{核实的发现}独立两温改变了单温模型的温度、密度与电流演化；较高驱动频率下离子明显热于电子。图4显示两物种相似的大尺度空间结构，但温度幅度不同；图3列各物种热源、交换和壁面热流，积分范围包括注入器。作者称与离子Doppler和初步Thomson数据定性一致，但明确未做严格定量验证。压缩功也不能全部视为不可逆热量。

\paragraph{限制与未完成问题}温度依赖壁面条件和人工扩散，尚未达到对扩散系数不敏感的收敛极限；真实驱动电路、中性/原子/辐射过程及更真实黏性仍缺失。已有勘误DOI10.1063/5.0032632经元数据确认，但更正全文未取得；本次只保证35页接收稿完整核读，不宣称掌握更正后的全部公式或精确比值。两个代码的差异也不全来自两温处理。

\paragraph{本报告判断}这是本分支取得的最直接装置尺度两温Hall实例，证明可同时计算独立温度、局地热图和带输运损失的物种预算。但结果受驱动、边界与耗散闭合约束，且部分功率来自注入器反转过程，并非仅重联X点。应采用定性和条件性结论；未核勘误、人工扩散敏感性及缺失电路模型使其不能支持参数无关的Qi/Qe或精密实验验证。

\section{L267：2020年}\label{paper:L267}

\textbf{Electron-Only Reconnection in Plasma Turbulence}\citep{L267}\

F. Califano;  Silvio Sergio Cerri;  M. Faganello;  D. Laveder;  Manuela Sisti;  Matthew W. Kunz.

\textsf{核读范围：}arXiv1810.03957v3全部12页；\textsf{图表：}所有图注与正文；图2、3、5、6、8视觉核对.

\href{run:papers/L267.pdf}{打开本地PDF：\texttt{papers/L267.pdf}}\quad\href{https://doi.org/10.3389/fphy.2020.00317}{DOI／版本入口}

\paragraph{在报告中的研究脉络}\hyperref[story:energy]{电子独有重联与热化}。关联表示本报告的阅读线索，不自动表示这些论文互相引用。

\paragraph{实际模型}二维三速度维混合Vlasov：动理学离子、等温惯性电子，非纯电子—离子双流体，作为比较证据。

\paragraph{计算／实验细节}20pi di方盒、1024\textsuperscript{2}，dx=.0614di约.74de；质量比144，beta\_i1、Te=Ti。离子速度51\textsuperscript{3}点覆盖±5vthi，存在高波数谱滤波。初始磁扰动rms.15；A注入.1<=kdi<=.6，B至.3；另.21幅值比较只文字报告。

\paragraph{核实的发现}小尺度注入A中重联主要为电子独有喷流，发展后仍无对应离子喷流；大尺度注入B产生离子参与的长重联层。两组能谱相似但重联类型不同；手选重联层A约3.5di、最长8-9di，B多数大于10di。电子广义磁通与普通磁通解耦；两组厚度都在约3de。

\paragraph{限制与未完成问题}自动统计7182/460个电流峰并非都在重联；人工确认重联层样本小。de只有约1.36格，作者也承认临界分辨率与滤波可能改变广义磁通；欧姆式省去de\textsuperscript{2} grad div E。等温电子不能预测电子加热，J·Eprime也包含体流能转换；未建立普适注入阈值。

\paragraph{本报告判断}提出湍流注入尺度选择电子独有重联的可检验机制；不是Hall MHD已准确替代动理学的证明。

\section{L268：2020年}\label{paper:L268}

\textbf{Oblique Tearing Mode Instability: Guide Field and Hall Effect}\citep{L268}\

Chen Shi;  Marco Velli;  Fulvia Pucci;  Anna Tenerani;  Maria Elena Innocenti.

\textsf{核读范围：}arXiv2007.00607v3 全部11页，含全部本征值扫描和讨论；\textsf{图表：}图1–7、式1–14；页8/9增长率、频率及本征函数图直接视觉核验.

\href{run:papers/L268.pdf}{打开本地PDF：\texttt{papers/L268.pdf}}\quad\href{https://doi.org/10.3847/1538-4357/abb6fa}{DOI／版本入口}

\paragraph{在报告中的研究脉络}\hyperref[story:three]{三维传播与有限X线}。关联表示本报告的阅读线索，不自动表示这些论文互相引用。

\paragraph{实际模型}均匀密度不可压缩电阻Hall MHD，无电子惯性/黏性；Harris反向场与常量导向场，压力平衡；旋转坐标后求复本征值。

\paragraph{计算／实验细节}S=aVA/ηm=10\textasciicircum{}4，di/a=0–1，Bg/B0=0–100。SciPy四阶残差控制配点法，y\ensuremath{\in}[-15, 15]，远场指数衰减；未公布固定网格数/容差及完整收敛扫描，忽略背景电阻衰减。

\paragraph{核实的发现}固定较大kx时强导向场可使斜模比平行模快约5\%，但完整(kx, kz)扫描的全局最快模仍kz=0，只要沿重联场方向可容纳足够长波。Hall斜模有实频率及传播；导向场使频率随kz非单调并抑制法向振荡。最大频率相速度随di近线性，最大频率与最大增长位置不同。

\paragraph{限制与未完成问题}仅线性本征问题；相速度不是非线性X线前沿速度。L224是惯性无力片固定波数，两文不能直接互相否定。沿片有限长度排除长波可能改变优势；本文未实际求解线绑定边界，无电子压强张量/热化。

\paragraph{本报告判断}用于区别固定波数斜模优势和全谱最快模。作者名核为Chen Shi（时辰）。

\section{Ng2020：2020年}\label{paper:Ng2020}

\textbf{An improved ten-moment closure for reconnection and instabilities}\citep{Ng2020}\

Jonathan Ng;  Ammar H. Hakim;  Liang Wang;  Amitava Bhattacharjee.

\textsf{核读范围：}1-8（作者稿全部，含附录A）；\textsf{图表：}图1-7、式4-6、附录A；图像复核待补；主要结果页已视觉复核.

\href{run:papers/Ng2020.pdf}{打开本地PDF：\texttt{papers/Ng2020.pdf}}\quad\href{https://doi.org/10.1063/5.0012067}{DOI／版本入口}

\paragraph{在报告中的研究脉络}\hyperref[story:global]{层长、驱动与全局完成}。关联表示本报告的阅读线索，不自动表示这些论文互相引用。

\paragraph{实际模型}十矩全Maxwell温度梯度热流，与局地各向同性化和非局地无磁化闭合比较；FV输运、局地隐式源、超时间步抛物子步。

\paragraph{计算／实验细节}Fadeev与Ng2015同系；lambda5di用640×320，同dx扫描系统尺寸，mi/me25, k0s=1/ds；导引场0-.4-1-2等。LHDI另例mi/me36, Ti/Te10, L=rhoi, vt, e=.06c, 盒6.4L×12.8L, 模数8。

\paragraph{核实的发现}平均率随岛尺寸下降可接近PIC约lambda\textasciicircum{}-0.8；岛合并最大率随Bg升到1先升后降，不能套用Harris片单调下降。离子宽层及非旋转对称度改善。相同k0=1/ds可同时激发重联与LHDI。

\paragraph{限制与未完成问题}LHDI增长率默认1.33Ωci、调k为ky后.86，动力学1.1，局地闭合.34，非局地.84；仍非完全吻合。电子各向异性空间结构仍不准、跨场热流过大；未研究LHDI对重联的非线性作用。附录指出零热流、无碰撞的对称稳态停滞点会驱动不物理Pxx趋零。

\paragraph{本报告判断}梯度闭合并非更高阶就自动正确；改善来自k\textsuperscript{2}/|k0|的尺度选择，仍有自由参数和磁化缺失。

\section{E552：2021年}\label{paper:E552}

\textbf{Fluid simulations of three-dimensional reconnection that capture the lower-hybrid drift instability}\citep{E552}\

F. Allmann‐Rahn;  S. Lautenbach;  Rainer Grauer;  R. D. Sydora.

\textsf{核读范围：}arXiv2008.06440v3 1-20全部，对应2021年期刊；\textsf{图表：}图1-9；闭合式、图4、6-9视觉复核.

\href{run:papers/E552.pdf}{打开本地PDF：\texttt{papers/E552.pdf}}\quad\href{https://doi.org/10.1017/s0022377820001683}{DOI／版本入口}

\paragraph{在报告中的研究脉络}\hyperref[story:closure]{压力、热流与闭合}。关联表示本报告的阅读线索，不自动表示这些论文互相引用。

\paragraph{实际模型}全Maxwell五矩与十矩比较；压力偏离各向同性的梯度闭合chi3, k0s1/ds；Vlasov用于二维电流片不稳定性对照。

\paragraph{计算／实验细节}CWENO+RK3，Maxwell FDTD，Laplacian二阶中心，每子循环交换压力边界，三维用8次；压力对角下限。三维100×25×6di、2400×600×150、mi/me100, c/VA30，Harris a.5, nb.2, Ti/Te5；磁通种子.1、电子密度噪1e-8。独立静电LHDI nb.001、Ti/Te10、a约.67、质量比36/100/250至流体1836。

\paragraph{核实的发现}静电LHDI质量比100时Vlasov/十矩梯度/五矩/十矩松弛增长率2.2/2.4/3.2/.67 Omega\_i；最快模随质量比变化，固定单一模比较会误导。电磁支Vlasov约.41、流体.21，均出现超过B0的磁扰动及强kink。三维较二维启动早、后期片变宽，出流湍动强而电子层相对适中；密度噪1e-10至1e-2使LHDI启动约30→4 Omega\_i\textasciicircum{}-1。

\paragraph{限制与未完成问题}高质量比无同参数Vlasov验证；电磁支结构、启动与增长仍不同，异常高饱和原因未解释。CSSI只是形态相似的候选，未确诊。三维率是z平均电场代理；电子de约2.4格，未给全三维网格收敛。对角下限不等于完整压力正定保证；印刷Lap约+k0\textsuperscript{2}应理解为负号Fourier符号。

\paragraph{本报告判断}该工作同时检验重联与不稳定性并强调噪声/启动竞争，优于仅比较峰值；但不能把3D形态相近称为完整动理学等价。

\section{Finelli2021：2021年}\label{paper:Finelli2021}

\textbf{Bridging hybrid- and full-kinetic models with Landau-fluid electrons. I. 2D magnetic reconnection}\citep{Finelli2021}\

F. Finelli;  S. S. Cerri;  F. Califano;  F. Pucci;  D. Laveder;  G. Lapenta;  T. Passot.

\textsf{核读范围：}1-17（全部）；\textsf{图表：}PDF3-4、9、11-14直接视觉核验式6-8、图3-8，全文图注均读。.

\href{run:papers/Finelli2021.pdf}{打开本地PDF：\texttt{papers/Finelli2021.pdf}}\quad\href{https://doi.org/10.1051/0004-6361/202140279}{DOI／版本入口}

\paragraph{在报告中的研究脉络}\hyperref[story:energy]{电子独有重联与热化}。关联表示本报告的阅读线索，不自动表示这些论文互相引用。

\paragraph{实际模型}混合Vlasov动理学质子+旋转对称各向异性Landau流体电子(HVLF)，对照等温HVM及iPIC；不是纯双流体。

\paragraph{计算／实验细节}24π×12πdi, 1024×512，质子51\textsuperscript{3}速度±6.4vthi；质量比100, Bg/B0.25, βp1, Te/Ti.25；Δx.074di=.74de，初始ρe.035di；Δt流体.005Ωp\textasciicircum{}-1且3电子子步，PIC.2ωpi\textasciicircum{}-1=.2Ωe\textasciicircum{}-1，每物种约7000ppc。

\paragraph{核实的发现}HVLF较等温电子更接近PIC的重联率、电子喷流、局地各向异性和对离子反馈；firehose/mirror的限制部分保留。

\paragraph{限制与未完成问题}没有电子FLR/回旋不稳定性和完整俘获；沿弯曲场线热流近似；初始ρe空间未充分解析，存在少量PIC数值冷却；原文把详细能量分配留待后续，无不可逆Qi/Qe/实验事件验证。

\paragraph{本报告判断}模型身份、温度/能量/熵产区别及跨文献判断见独立加热专题；不把本篇单独作为普适分配定律。

\section{L272：2021年}\label{paper:L272}

\textbf{How Hall electric fields intrinsically chaotize and heat ions during collisionless magnetic reconnection}\citep{L272}\

Young Dae Yoon;  Paul M. Bellan.

\textsf{核读范围：}Physics of Plasmas28 022113正式版10个PDF物理页（含出版封面），正文编号1–8、参考文献9页；无文内附录。Caltech作者库下载，替换主链接，旧16页稿保留。；\textsf{图表：}图1–7、式1–48；直接查看试粒子轨道/能量、Harris阈值、PIC场和温度图，核对式18–20的系数。.

\href{run:papers/L272_published.pdf}{打开本地PDF：\texttt{papers/L272\_published.pdf}}\quad\href{https://doi.org/10.1063/5.0040374}{DOI／版本入口}

\paragraph{在报告中的研究脉络}\hyperref[story:energy]{电子独有重联与热化}；\hyperref[story:diagnostic]{实验与可区分诊断}。关联表示本报告的阅读线索，不自动表示这些论文互相引用。

\paragraph{实际模型}电子正则涡量/低压或正压流体推导、给定场Boris试粒子和SMILEI全PIC对照；研究Hall电场的离子轨道效应，不是Hall流体独立计算的不可逆加热。

\paragraph{计算／实验细节}波场1000试粒子，ω/Ωc=.4、α=.5或1.5；倒抛物势10000试粒子，α=.4或1.04。PIC初始半厚di、mi/me=100、Bg/Brec=0, .1, .3, .5；域、网格与粒子数指向2019前序文献，本篇未列全。图6/7明确复用2019结果，非独立新一轮验证。

\paragraph{核实的发现}给定强电场梯度可使轨道间距与能量展宽增长；PIC中较大梯度指标区域与离子高温区大致对应，增加导引场使范围和程度减小，.5例强效应受抑。Harris近似给出与片厚及导场有关的阈值；高质量电荷比更易满足其指标。提供Hall几何—轨道—温度的可检验联系。

\paragraph{限制与未完成问题}J·(J×B)=0仍成立，分物种电功非零不等于总Hall耗散。图1/2的TE按文中定义是粒子能量分布标准差，非通常的中心速度二阶矩温度；线性倒抛物系统的指数分离也不足以单独证明混沌或熵增。原文由准中性直接采用不可压缩需另加假设；式18/19按其α定义少1/4，与式20不一致，报告从运动方程复算而不推断代码错误。强导场、高β和三维可由其他机制主导，作者也承认。

\paragraph{本报告判断}保留场结构与轨道失磁化的证据，降低其普适热化措辞的强度。与Yoo2013原始实验并读：实验确证面内场加速与下游升温，但未唯一判定微观热化机制；与E512并列区分速度空间展宽、压力内能和碰撞熵产。

\section{L275：2021年}\label{paper:L275}

\textbf{Scaling theory of three-dimensional magnetic reconnection spreading}\citep{L275}\

Milton Arencibia;  P. A. Cassak;  M. A. Shay;  E. R. Priest.

\textsf{核读范围：}arXiv2108.00294v1全部15页，含完整方法、讨论与参考文献；\textsf{图表：}图1-7与表I-III；公式1、2、5-8、17，图3-7及表I/III视觉核验.

\href{run:papers/L275.pdf}{打开本地PDF：\texttt{papers/L275.pdf}}\quad\href{https://doi.org/10.1063/5.0052189}{DOI／版本入口}

\paragraph{在报告中的研究脉络}\hyperref[story:three]{三维传播与有限X线}。关联表示本报告的阅读线索，不自动表示这些论文互相引用。

\paragraph{实际模型}三维F3D准中性冷电子Hall惯性系统与电阻MHD对照；离子绝热，初始电流全由电子承担。

\paragraph{计算／实验细节}102.4×51.2×256，dxdy.05, dz1，周期；质量比25，n1，Ti1-1.5；初始厚度1，另.5/2/3。导引场0或3；Hall电阻0，MHD电阻.004且片厚.32。D4xy1.6e-5；D4z无导引场.16、有导引场.64；dt.02/.01。磁扰动.005、z半长15（9/12/30检查）、噪声1e-5；重联代理为距X点±2di处平均|By|超过.04。

\paragraph{核实的发现}Faraday方程的Ex沿z梯度解释前方是否生成正确极性的X型By：反平行主要Hall输运，强导引场主要场线弯曲驱动的整体流。表I-III直接比较电场梯度/By变化及磁张力/流加速，量级与符号一致而非逐点精确。w0=2矩形域持续扩展、速度-.51VA；同厚方域只移动不伸长。压力种子磁声脉冲约1.75VA双向，重联-.92VA单向。

\paragraph{限制与未完成问题}研究主要核对早期By拓扑的传播，稳定快重联有延迟（例z=-45，扰动约t60到、快重联约t170起）；厚度均匀才可抵消统一延迟。方/矩形域差别归因自由磁能不足仍为单组假说，未做能量因果扫描。MHD反平行仅在t≤200无可分辨传播，理论仍有电阻扩散。v1式1/2/7符号与文内电子-z方向不一致；式17将两个y偏移平方放在同一指数，会变成中心峰乘exp(-16)，不是声称的两侧压力峰。需核对实际输入，不断言代码错误。

\paragraph{本报告判断}新意是把不同传播机制统一到可检验的电场梯度，并暴露域形状和启动种子问题；报告从正确SI方程重推符号，不照抄v1公式。

\section{E474：2022年}\label{paper:E474}

\textbf{Gyrofluid analysis of electron β e effects on collisionless reconnection}\citep{E474}\

Camille Granier;  D. Borgogno;  D. Grasso;  E. Tassi.

\textsf{核读范围：}arXiv2110.03052v1全部30页，含完整Hamiltonian约化和附录双内层匹配；期刊登记2022，所读预印本2021；\textsf{图表：}图1–13、表1及全部附录；页10/17/20/23直接核验修正色散、非线性双加速、能量和密度.

\href{run:papers/E474.pdf}{打开本地PDF：\texttt{papers/E474.pdf}}\quad\href{https://doi.org/10.1017/s0022377822000010}{DOI／版本入口}

\paragraph{在报告中的研究脉络}\hyperref[story:energy]{电子独有重联与热化}。关联表示本报告的阅读线索，不自动表示这些论文互相引用。

\paragraph{实际模型}有限βe、电子FLR及B平行扰动的双场Hamiltonian gyrofluid；冷离子主分析，Ni=Ui=0在Hamiltonian子代数层约化；quasistatic等温闭合，2D取自3D的形式极限。G10e(k)=exp(-k\textsuperscript{2}βe de\textsuperscript{2}/4)。

\paragraph{计算／实验细节}伪谱、三阶Adams–Bashforth、短尺度数值滤波，最高2880\textsuperscript{2}；Aeq=λsech\textsuperscript{2}(x/λ)的61项谱近似，非线性盒[-1.5π, 1.5π]×[-π, π]，Δ'=14.31，mi/me=100，de=.06–.1，βe=.1–1.5。线性另扫弱Δ'、βe至2及不同质量比。

\paragraph{核实的发现}βe→0常ψ双内层匹配新增γ\textsuperscript{2}修正：γ=2ky deρsΔ'/(πλ)+γ\textsuperscript{2}deπλ/(4kyρs\textsuperscript{2})，在ρe\ensuremath{\ll}γde/(kyρs)\ensuremath{\ll}ρs且常ψ成立时改善数值拟合。固定de, ρs升βe稳定；固定质量比/de升βe同时升ρs，可先促进再抑制。高β非线性出现减速间隔及第二次超指数阶段；电子密度有嵌套四极和较平缓梯度，正则不变量仍细丝化。

\paragraph{限制与未完成问题}有限βe不是单独电子温度开关；γ是对数通量增长率不是归一化重联电场。可靠能量段限定总能数值损失<1\%；X点移动引起部分增长率振荡。Epe=ρs\textsuperscript{2}Ne\textsuperscript{2}/2是密度扰动能，不能作真实温升。准静态闭合在严格kz=0不直接成立，作者取形式极限；冷离子/质量比/滤波使高β空间应用有限。预印本式3.1的ρe\textsuperscript{2}系数与G10e定义差2，应由指数直接展开得ρe\textsuperscript{2}/2。

\paragraph{本报告判断}第一章补充参数关联ρe=de√(βe/2)=ρs√(me/mi)，第三章按固定变量比较相反趋势，并将新修正只用于其常ψ有效域。

\section{F589：2022年}\label{paper:F589}

\textbf{First-principles theory of the rate of magnetic reconnection in magnetospheric and solar plasmas}\citep{F589}\

Yi‐Hsin Liu;  P. A. Cassak;  Xiaocan Li;  Michael Hesse;  Shan‐Chang Lin;  K. J. Genestreti.

\textsf{核读范围：}arXiv:2203.14268v1全部10页，含方法、讨论和参考文献；式1–18核验。；\textsf{图表：}图1–4及对应预算；直接查看PDF4、5、6页的相空间图、几何示意与公式。.

\href{run:papers/F589.pdf}{打开本地PDF：\texttt{papers/F589.pdf}}\quad\href{https://doi.org/10.1038/s42005-022-00854-x}{DOI／版本入口}

\paragraph{在报告中的研究脉络}\hyperref[ch:numerics]{模型与数值可信度}。关联表示本报告的阅读线索，不自动表示这些论文互相引用。

\paragraph{实际模型}低β反平行准稳态Hall能量输运及压力亏损理论，以二维全PIC检验；不是独立Hall MHD模拟。Hall项总功为零，但单物种仍可从Hall电场获能。

\paragraph{计算／实验细节}VPIC76.8di×38.4di，12288×6144格，约150亿宏粒子，mi/me400，Ti/Te1，背景β=.01，ωpe/Ωce4，初始半宽di，背景密度等于片峰密度；x周期，z导体/粒子反射；磁扰动.03Bx0；Ohm预算平均.085Ωci\textsuperscript{-1}。

\paragraph{核实的发现}Hall场使入流电磁能转向排气而不在X线堆积，压力亏损由磁张力补偿而开放排气。反向离子束建立Pizz，预测压力增量为离子到电子尺度磁压差的2/3。相似张角假设给Bxe\textsuperscript{2}/Bxi\textsuperscript{2}≈sqrt(me/mi)，基本模型率.172(400)、.157(1836)，引入外流热压修正约.075更接近该算例。

\paragraph{限制与未完成问题}稳定单X线、低β、反平行、近不可压、相似张角和Alfvén外流是关键假设；总Hall功为零不能单独保证压力亏损。基本率.172不是与模拟.075精确一致。页7写Li/dm≈Slope，按几何及式16到17应为其倒数。一般3D、强导引场、高β和全局尺度不由此证明。

\paragraph{本报告判断}第一章补充压力、张角和率的推导；与F600对照，区分局部准稳态机制与全局磁通处理率。

\section{L283：2023年}\label{paper:L283}

\textbf{Discovery of double Hall pattern associated with collisionless magnetic reconnection in dusty plasmas}\citep{L283}\

Shu-Di Yang;  Liang Wang;  Chuanfei Dong.

\textsf{核读范围：}arXiv:2206.08553v2接受稿6页全文核读，包括三物种方程、初值、磁场产生项、速率与参考文献；\textsf{图表：}图1—6及图注全部核对；放大物理页3的嵌套反号四极场和闭合电流图.

\href{run:papers/L283.pdf}{打开本地PDF：\texttt{papers/L283.pdf}}\quad\href{https://doi.org/10.1093/mnras/stad1453}{DOI／版本入口}

\paragraph{在报告中的研究脉络}\hyperref[story:diagnostic]{实验与可区分诊断}。关联表示本报告的阅读线索，不自动表示这些论文互相引用。

\paragraph{实际模型}二维PSC全动力学三物种PIC：电子、正离子、负尘埃；不是常规电子—离子Hall MHD数值结果。用分尺度的三流体动量解释内外两套Hall电流。

\paragraph{计算／实验细节}32di0×16di0，1000×500，x周期、z导体壁，Harris半厚di0、扰动0.15B0di0；mi/me=100、md/mi=16、Zd=1、nd/nb=5、nb/n0=0.1、Ti/Te=5、Td/Te=1、VAe0/c=0.5、Δtωpe0≈0.22；每种约7.5×10\textasciicircum{}7数值粒子。尘埃带电量和质量固定；为准中性增加了离子背景密度。

\paragraph{核实的发现}在清晰的电子、离子、尘埃尺度分离下，内侧传统Hall四极结构外出现极性相反的四极结构；后者较晚建立，内层随后收窄。尘埃流向片心聚集再随喷流排出，晚期含尘例重联进展和速率低于无尘例。外侧正离子/电子的有效正电流与较慢负尘埃形成反号闭合。

\paragraph{限制与未完成问题}单组缩减质量和电荷参数，未扫描真实尘埃质量/充电演化/碰撞/三维；加入尘埃同时改变离子密度与总质量，不能说严格只改变一个物理因素。尘埃聚集促成颗粒生长仅是设想，未模拟凝聚。原稿di0定义误排成等离子体频率量纲，扰动坐标也与xz几何不一致；不直接抄录。没有实验中已观测到本研究双Hall图样的证据。

\paragraph{本报告判断}纳入新的嵌套反号Hall现象，明确它是需要三物种扩展的边界证据，不能算作标准两物种Hall MHD已解释的现象。

\section{L294：2023年}\label{paper:L294}

\textbf{Three‐Dimensional Magnetic Reconnection Spreading in Current Sheets of Non‐Uniform Thickness}\citep{L294}\

Milton Arencibia;  P. A. Cassak;  Michael A. Shay;  Jiong Qiu;  S. M. Petrinec;  Haoming Liang.

\textsf{核读范围：}正式期刊纠正版全部15页，页15勘误仅涉及参考文献与图4图例；非仅arXiv摘要；\textsf{图表：}图1–5、表1、式1–21；图4–5及表1直接视觉检查，重算式16的磁尾时间.

\href{run:papers/L294.pdf}{打开本地PDF：\texttt{papers/L294.pdf}}\quad\href{https://doi.org/10.1029/2022ja030999}{DOI／版本入口}

\paragraph{在报告中的研究脉络}\hyperref[story:three]{三维传播与有限X线}。关联表示本报告的阅读线索，不自动表示这些论文互相引用。

\paragraph{实际模型}无导向场三维Hall加电子惯性F3D；冷电子、绝热离子，均匀密度与渐变片厚；研究沿初始电子载流方向的X线扩展。

\paragraph{计算／实验细节}盒102.4×51.2×256di，格0.05×0.05×1di，dt=0.02，质量比25；双周期Harris，全周期，只扰动上层，幅0.005、沿z半长15di，噪声10\textasciicircum{}-5。w1=0.75–2.5、w2=1.25–2等13组表列参数；变化宽度参数4并试8、12。D4xy=1.6×10\textasciicircum{}-5、D4z=0.16，文中称做过系数试验。起始快重联由两侧电子层边界平均|By|超过0.04判定。

\paragraph{核实的发现}薄到厚时By扰动保持窄束，仅内嵌磁场Bup≈B0w1/w参与，扩展速度vs≈VA di w1/w\textsuperscript{2}，带起始厚度记忆；厚到薄为VA di/w，无该记忆。拟合斜率0.919±0.082、截距0.044±0.041支持标度。w1=1, w2=2过渡耗时约20，与积分预测19接近；加倍/三倍过渡宽度时间相应增加。

\paragraph{限制与未完成问题}嵌入对vs的代换作者明确称半经验；要求渐变、By先传播后片塌缩，近似忽略ΔJ项。主要证据是场阈值传播，非完整三维电压诊断。无导向场、非对称、电子漂移波、全局曲率；不能直接定量解释所有耀斑/磁层顶。表1(0.75, 2)预测0.19实测0.26却写偏差-27.9\%，不应照录。按页11给出的di=720km、VA=1400km/s、Δz=133di、w2=8.9di、w1=0.89–3.5di、α=2代式16，本调研重算约10.0–23.5min；原文2.5–23.6min下限未复现。

\paragraph{本报告判断}可靠新现象是传播速度的路径记忆和方向不对称，不是快率统一公式。报告给出积分推导并标识应用数值复算差异，避免把作者定性观测联系升级为验证。

\section{Yoon2023PhaseSpace：2023年}\label{paper:Yoon2023PhaseSpace}

\textbf{Phase-space Analysis of Ordered and Disordered Nonthermal Ion Energization during Magnetic Reconnection}\citep{Yoon2023PhaseSpace}\

Young Dae Yoon;  Paul M. Bellan;  Gunsu S. Yun.

\textsf{核读范围：}所记录版本全部9个PDF物理页；正文、方法、结果、讨论、图表、参考文献及适用文内附录均核读。；\textsf{图表：}直接查看PDF物理页2, 3, 4, 5, 6, 7, 8；All Figs1–7;  Eqs1–5 and AppendixEqA1.

\href{run:papers/Yoon2023PhaseSpace.pdf}{打开本地PDF：\texttt{papers/Yoon2023PhaseSpace.pdf}}\quad\href{https://doi.org/10.3847/1538-4357/acf1f5}{DOI／版本入口}

\paragraph{在报告中的研究脉络}\hyperref[story:energy]{电子独有重联与热化}。关联表示本报告的阅读线索，不自动表示这些论文互相引用。

\paragraph{实际模型}本文采用二维空间、三维速度的全动理学粒子模拟，结合正则动量守恒和有效势分析离子相空间结构，再以蒙特卡罗测试粒子计算最大李雅普诺夫指数。温度由粒子分布的中心二阶矩得到，能够区分多束有序运动与速度分布展宽。它提供动理学机制诊断，并不是有限矩霍尔流体温度闭合的模拟。

\paragraph{计算／实验细节}反平行哈里斯片初始两物种等温，半宽等于离子皮肤深度，质量比 100，阿尔芬速度为光速的 0.025 倍。1024 乘 8192 网格覆盖 51.2 乘 409.6 个离子皮肤深度，每格 200 至 1000 个粒子；流出方向周期、入流方向开放。200 个测试离子在冻结的粒子模拟场中演化，以贝内廷算法计算最大指数，并舍弃最初 100 个逆离子等离子体频率的暂态。

\paragraph{核实的发现}较高温度矩可以由互相分离的冷束叠加造成，各束自身温度甚至低于上游。作者区分入流、拾取状、单阱与双阱四个运动阶段，用相空间分布解释各方向温度。单阱阶段测试粒子的平均最大李雅普诺夫指数为正，为混沌混合提供较直接证据；附录逐项积分电场和磁力贡献，说明不同阶段的加速来源。

\paragraph{限制与未完成问题}最大指数来自冻结场和有限时间；实时演化场只报告定性相似，没有展示相同校准的指数计算。二维反平行几何及缩小质量比限制外推范围。第 5 节明确把磁能转为离子能的精确份额、有序多束与随机展宽的相对份额及其时空变化列为未决问题，也没有给出两物种不可逆熵产生热的统一预算。

\paragraph{本报告判断}相较早期仅展示相近轨道分离，该文更清楚地区分有序增能、中心矩温升和混沌展宽，因此适合约束流体模型应保留哪些信息。不过其进展属于机制辨识，不能表述为已解决普适离子与电子加热比。论文全部 9 页、附录及七幅图均已核读；其中对 L272 的年份误写按正式发表信息纠正为 2021 年。

\section{Huang2024：2024年}\label{paper:Huang2024}

\textbf{Three-dimensional plasmoid-mediated reconnection and turbulence in Hall magnetohydrodynamics}\citep{Huang2024}\

Yi-Min Huang;  Amitava Bhattacharjee.

\textsf{核读范围：}1-11（arXiv v2全部）；\textsf{图表：}图1-9；图像复核待补；主要结果页已视觉复核.

\href{run:papers/Huang2024.pdf}{打开本地PDF：\texttt{papers/Huang2024.pdf}}\quad\href{https://doi.org/10.1063/5.0216561}{DOI／版本入口}

\paragraph{在报告中的研究脉络}\hyperref[story:three]{三维传播与有限X线}。关联表示本报告的阅读线索，不自动表示这些论文互相引用。

\paragraph{实际模型}三维可压Hall MHD，标量绝热压力、pe=pi、无电子惯性，电阻黏性超电阻；压力方程不含显式耗散加热。

\paragraph{计算／实验细节}立方盒、两磁通管合并，h=.01L, Bg约Brec, 初始pi=pe=1, rho=1, gamma=5/3；S=2e5, Pm=1, di/L=.002或.001；etaH=1e-13；网格2000×1000×2000（40亿格），中心dy=1e-4L；x/y导体滑移, z周期。

\paragraph{核实的发现}同参二维均趋单X线；三维L/di500也趋准二维，1000保持自生湍动与多重联点。单X线例通量率与磁能下降更快，Hall湍动态快于电阻MHD但慢于单X线。中面涡各向异性近尺度无关，远离中面大涡部分恢复GS标度。

\paragraph{限制与未完成问题}三维率取z平均场的二维X点磁通，是代理诊断，不是全部3D连通性。亚di尺度不足，k\textasciicircum{}-11/3只是参考线，未确证幂律。结构函数限于x-z切片，不能直接与全3D采样比较；仅两组di/固定S。

\paragraph{本报告判断}三维自由度可保留湍动但非必然加速。Petschek相似仅几何，不证明慢激波机制；闭合与真实电子物理仍未检验。

\section{E473：2025年}\label{paper:E473}

\textbf{GREENY: A Full-F 2D Gyrofluid Reconnection Code}\citep{E473}\

F. F. Locker;  Markus Held;  Tobias Stocker-Waldhuber;  Alexander Stürz;  M. Rinner;  A. Kendl.

\textsf{核读范围：}arXiv:2502.10219v1全部17页核读，包含Full-F/δF方程、算法1、制造解、守恒检查、全部算例和参考文献；独立补充文件和期刊修订全文尚未取得，不宣称核读这些材料；\textsf{图表：}图1—10和表1—3逐项核对；另放大检查页9制造解误差、页12形态比较、页13超黏性扫描、页14岛/涡示例.

\href{run:papers/E473.pdf}{打开本地PDF：\texttt{papers/E473.pdf}}\quad\href{https://doi.org/10.1016/j.cpc.2025.109807}{DOI／版本入口}

\paragraph{在报告中的研究脉络}\hyperref[story:onset]{启动、爆发与磁岛}。关联表示本报告的阅读线索，不自动表示这些论文互相引用。

\paragraph{实际模型}GREENY低β二维等温强导引场gyrofluid，演化全密度及电子/离子平行正则动量，保留离子FLR，忽略电子FLR与曲率。Full-F不表示任意磁扰动、任意β或完整速度分布。发表于CPC的版本另登记F802，SSRN为F803，属于同一研究的版本线索。

\paragraph{计算／实验细节}Arakawa括号、Karniadakis多步；起步用后向Euler初始化历史状态，FFT或有限差分空间算子。极化方程PCG/迭代Teague/SOR；该v1的Ampère SOR保留变密度但省离子gyroaverage，直接FFT保留FLR但冻结密度以避免卷积，不能同时宣称该实现完整Full-F与完整FLR。制造解SOR误差h\textsuperscript{2}，FFT/PCG/Teague示例h\textsuperscript{4}。典型1024\textsuperscript{2}、32点/ρs；β=10\textasciicircum{}-5—10\textasciicircum{}-2形态扫描，近边界Δ'box=0.15、β=0.005扫描超黏ν=10\textasciicircum{}-12—10\textasciicircum{}-6。

\paragraph{核实的发现}小β例更长的Y形片，较大β例更开放X形；暖离子电流驱动长宽比2的例出现双岛；同向平行电流吸引也可因共转涡形成的涡量层推迟合并。超黏性可改变前期多岛/合并和起始时刻：过强抹平扰动、过弱出现格点颗粒噪声；主峰部分重合不能替代全时序与能量收敛。所谓rate γ=dln|δA|/dt是增长率，不是归一重联电场。

\paragraph{限制与未完成问题}v1多处排版/说明矛盾：图6正文和图注左右对应相反；图7正文称主事件形状/起始无直接联系，图注和结论又强调超黏性影响，故只保留有条件的数值敏感性结论。能量式磁能写在种类求和内、势能系数等需从一致方程重新推导。制造解阶数不是全部非线性算例已收敛的证明。期刊修订可能更改求解器，本文限制只指实际核读v1。

\paragraph{本报告判断}纳入数值实现和Full-F适用边界；与L314串联时保留版本区别。按一项研究计数，不把预印本、SSRN和期刊版算三篇。

\section{F606：2025年}\label{paper:F606}

\textbf{Can plasmoid-mediated reconnection occur in collisionless systems?}\citep{F606}\

Keita Akutagawa;  Shinsuke Imada;  Munehito Shoda.

\textsf{核读范围：}arXiv:2507.11904v1全部13页，含附录A分辨率对照和全部图1–22。；\textsf{图表：}表I/II、图1–22；直接查看PDF12页图22的粗细网格差异。.

\href{run:papers/F606.pdf}{打开本地PDF：\texttt{papers/F606.pdf}}\quad\href{https://doi.org/10.1063/5.0274477}{DOI／版本入口}

\paragraph{在报告中的研究脉络}\hyperref[story:onset]{启动、爆发与磁岛}。关联表示本报告的阅读线索，不自动表示这些论文互相引用。

\paragraph{实际模型}KAMMUY 2.5D全PIC与独立电阻MHD对照；非Hall MHD参数扫描。PIC用Yee/Leapfrog及Langdon–Marder修正；MHD为HLLD、二阶MUSCL/RK和upwind CT。

\paragraph{计算／实验细节}PIC主扫描mi/me=1、2、4、8、16、32；δ=2di，100di×50di，另对4/16用200di×100di；Δ=.1de、dt=.05/ωpe，片内100ppc、背景20ppc，Ti/Te=1，ωpe/Ωce=1，背景β=.04。附录另有mi/me100但仅Δ=.5de的粗网格例。MHDη=.001、Δ=1/16，背景β=.25。

\paragraph{核实的发现}主扫描1–8形成次级磁岛，16/32未形成；磁岛出现不伴随PIC全域能量转化增强，MHD磁岛阶段则增强。部分移动次级X点处压力张量与惯性项相抵、电子非理想场很小。把16质量比的网格从.5de细化到.1de，原有磁岛消失。

\paragraph{限制与未完成问题}没有真实质量比1836的充分分辨模拟；不产生磁岛的外推依赖EDR长约20de和相关片宽约di，不能当作普遍定理。De>0不等同于不可逆耗散，局部X点非理想场/全域能量转化/重联率不应互换。PIC与MHD参数不同。原式4相对E*=E+ve×B的电子电荷符号有疑点，报告采用标准负号，不照抄。

\paragraph{本报告判断}纳入磁岛是否促进重联的条件性反例及分辨率检验；不把作者关于1836的推测改写成直接观测结果。

\section{F703：2025年}\label{paper:F703}

\textbf{Quadrupolar density structures in driven magnetic reconnection experiments with a guide field}\citep{F703}\

Thomas W O Varnish;  Joe M. Chen;  Simran Chowdhry;  Rishabh Datta;  George V. Dowhan;  Lansing S. Horan;  Nicholas Jordan;  E. R. Neill;  Akash P. Shah;  Roman V. Shapovalov;  Brendan J. Sporer;  R. D. McBride;  Jack Hare.

\textsf{核读范围：}arXiv:2412.02556v1全部12页，含附录A–D的零点校准、对称性、几何和视线积分模型；期刊年份2025。；\textsf{图表：}图1–9，关键对照为图4/5实测、图8电阻MHD、图9规定密度模型；直接查看PDF4/8/9页。.

\href{run:papers/F703.pdf}{打开本地PDF：\texttt{papers/F703.pdf}}\quad\href{https://doi.org/10.1063/5.0251581}{DOI／版本入口}

\paragraph{在报告中的研究脉络}\hyperref[story:diagnostic]{实验与可区分诊断}。关联表示本报告的阅读线索，不自动表示这些论文互相引用。

\paragraph{实际模型}MAIZE脉冲功率双倾斜碳丝阵实验；GORGON三维电阻MHD作为对照；Hall解释用既有二维结果启发的规定四极密度模型，不是本文完成的自洽三维Hall计算。

\paragraph{计算／实验细节}峰值约500kA、上升200ns，倾角0/22.5/45度给Bg/Brec=0/.4/1；阵直径10mm、高16mm、间距22mm；1064nm、2ns双视线干涉，di估约.5mm，观测跨20mm约40di，β约1。GORGON100微米、480×400×280格。

\paragraph{核实的发现}导引场增加时密度图出现四极增强/减弱区，230与325ns不同炮次有相似形态；单纯扭曲几何和电阻MHD投影无该四极结构，加入四极密度后能定性再现两视线图。

\paragraph{限制与未完成问题}没有直接测出层内Hall磁场、温度、完整流速或局部Ohm预算，二维视线积分不能唯一重建3D。GORGON密度高2–3倍原因未定。零点依赖烧蚀不随倾角变等假设；正文与附录对b=0零条纹视线的描述不一致。有Hall特征不等于无碰撞；密度团块没有磁诊断不能直接叫磁岛。

\paragraph{本报告判断}作为大范围、β约1的导引场密度形态证据纳入，并提出带合成诊断的三维Hall/双流体检验；不把几何拟合说成模型已验证。

\section{F741：2025年}\label{paper:F741}

\textbf{An analytical model of “Electron-Only” magnetic reconnection rates}\citep{F741}\

Yi‐Hsin Liu;  P. S. Pyakurel;  Xiaocan Li;  Michael Hesse;  Naoki Bessho;  K. J. Genestreti;  Shiva Bikram Thapa.

\textsf{核读范围：}本地10页作者稿全部核读：正文、逐步速率推导、Methods、参考文献和数据说明；核对期刊入口Communications Physics 8, 128(2025)，未把作者稿当作排版定稿；\textsf{图表：}图1—5及图注全部核对；特别放大页6图5的模型/PIC比较与局部归一化差异.

\href{run:papers/F741.pdf}{打开本地PDF：\texttt{papers/F741.pdf}}\quad\href{https://doi.org/10.1038/s42005-025-02034-z}{DOI／版本入口}

\paragraph{在报告中的研究脉络}\hyperref[story:energy]{电子独有重联与热化}。关联表示本报告的阅读线索，不自动表示这些论文互相引用。

\paragraph{实际模型}电子单独重联的准稳态多尺度解析模型，由二维P3D PIC验证；不是完整双流体数值计算。令离子电流远小于电子电流，把电子Lorentz力写成磁张力与磁压；面内Hall电场减速电子，下游离子反压通过经验函数加入。

\paragraph{计算／实验细节}真实质量比1836，Bg=-8Bx0，βi=3.54、βe=0.35（按总场），Ti/Te≈10，双Harris片半厚0.06di；域Lx/di=1.28、2.56、3.84、5.12、7.68，最小域Lz=2.56di，其余方形。周期边界，Δ=0.21de，Δt=2.5×10\textasciicircum{}-5 Ωci\textasciicircum{}-1，6000粒子/单元，c=300VAi0。作者稿关于ρi/di在讨论和方法分别写1.23、1.33，未取为精确校准值。

\paragraph{核实的发现}电子峰值喷流和决定磁通输运的远下游平台速度不同。Lx=2.56di例峰值约0.15VAe0，但平台仅1.7VAi0≈0.04VAe0。模型给U/VAi0=(di/δ0)√[(1-s\textsuperscript{2})(1-f)/(2+(de/δ0)\textsuperscript{2})]，f=sech(Lx/1.28di)，δ0≈Lx s/4；由连续性和压力亏损几何联立求s及R≈sU/VAi0。周期小系统反压把全局离子归一率约束到O(1)，去掉反压时作者给出约4.28的可能上限。

\paragraph{限制与未完成问题}f的函数形式和1.28di参数是启发式选择，几何L0≈Lx/4亦为近似；多数比较误差在约因子2内。图5f局部电子归一理论约0.4—0.5，而固定z=de测量约0.1；作者以局部Bx取点差一倍解释四倍率差，并非无参数精确吻合。4.28是开放系统预测而不是本文周期模拟测得。模型不处理显著离子电流、三维、强驱动和完整瞬态；‘无离子喷流’不等于离子压强毫无影响。

\paragraph{本报告判断}纳入第一章速率推导和第三章电子单独重联；明确区分EMHD近似、PIC证据、经验反压参数、峰值/输运速度及两种归一化。

\section{F788：2025年}\label{paper:F788}

\textbf{Magnetic reconnection,  plasmoids and numerical resolution}\citep{F788}\

García Morillo,  José María;  Alexandros Alexakis.

\textsf{核读范围：}先读arXiv:2406.08951v1全部8页，再下载并全文核对JFM1007 R3(2025)正式版全部10页；以正式版为判断依据。；\textsf{图表：}图1–5、表1，直接查看正式版PDF7页电流谱/重联率图；比较正式版新增讨论。.

\href{run:papers/F788_published.pdf}{打开本地PDF：\texttt{papers/F788\_published.pdf}}\quad\href{https://doi.org/10.1017/jfm.2025.109}{DOI／版本入口}

\paragraph{在报告中的研究脉络}\hyperref[story:onset]{启动、爆发与磁岛}。关联表示本报告的阅读线索，不自动表示这些论文互相引用。

\paragraph{实际模型}二维不可压电阻MHD Orszag–Tang涡；非Hall或双流体模拟，是数值误差对照。GHOST伪谱、RK4、2/3去混叠、双周期2π方域，ν=η。

\paragraph{计算／实验细节}初始磁能.5、动能.125，k<16随机相位扰动占总能量.25\%；N=1024–32768，最高充分分辨例η=5×10\textsuperscript{-6}、SL约5.4×10\textsuperscript{5}；比较均方电流峰值时刻t约1.9，率以区段平均入流速度/局部Bmax定义。

\paragraph{核实的发现}本初值与采样阶段，充分分辨时未见磁岛且率近SP标度；欠分辨时出现磁岛、虚假的近S无关快率和高波数谱堆积。正式版强调电流谱峰/截止值≥10仅是本设置经验判据。

\paragraph{限制与未完成问题}不能推出所有磁岛均为数值伪影；正式版明确有限片寿命、扰动衰减、对流流出与初值形状都可限制磁岛长到非线性幅度。不是静态无限长SP层线性理论的直接证伪。文中笼统称有限元指数收敛不适用于任意固定阶有限元。

\paragraph{本报告判断}补充局部电流谱和拓扑收敛要求；与Hall/双流体磁岛结果对照时保留模型差异。新版讨论比v1更审慎，索引改指正式版。

\section{L296：2025年}\label{paper:L296}

\textbf{A scalable multidimensional fully implicit solver for Hall magnetohydrodynamics}\citep{L296}\

Luis Chacón.

\textsf{核读范围：}2025 OSTI接收稿封面及正文/附录全部29页（正文1-28）；2024 arXiv v1全部25页亦读，用于版本核对；\textsf{图表：}图1-4、表1全部视觉核验；模型式5/6与附录A-C逐式检查.

\href{run:papers/L296_accepted.pdf}{打开本地PDF：\texttt{papers/L296\_accepted.pdf}}\quad\href{https://doi.org/10.1016/j.jcp.2025.113789}{DOI／版本入口}

\paragraph{在报告中的研究脉络}\hyperref[ch:numerics]{模型与数值可信度}。关联表示本报告的阅读线索，不自动表示这些论文互相引用。

\paragraph{实际模型}可压缩Hall流体，固定Ti/Te温比，p由加权速度vstar输运。接收稿改用de\textsuperscript{2} De(j/rho)惯性以匹配v≈vi下能量；不是精确有限质量双流体。全部实际测试de=0。

\paragraph{计算／实验细节}二阶单元中心有限体积、矢势兼容curl；波用Crank-Nicolson、GEM用BDF2。Newton-FGMRES与物理块Schur、多重网格V(3, 3)、阻尼Jacobi .7；EMHD内层GMRES最多10次或残差降至.1。GEM256\textsuperscript{2}拉伸网格，最小dx.02/dy.05，dt.05、di1、eta.005、nu\_i.05、nu\_e1e-4、kappa.02、温比5。弱扩展32\textsuperscript{2}单MPI至4096\textsuperscript{2}/16384MPI，dt.01、100步。

\paragraph{核实的发现}Hall电子块与离子流动块分离的预条件降低刚性求解迭代；波相位和HiFi GEM能量曲线吻合，多数非线性阶段接近，末期约1\%差。16384倍问题规模时总时间132至1130秒，2-3次Newton/外层GMRES每步；隐式步比显式CFL最高8400，实际速度最高75倍，粗网格反而慢25倍。

\paragraph{限制与未完成问题}弱扩展而非强扩展；nu\_e随细网格由1e-4降至6e-9，非固定物理参数收敛。性能段总时长1，早于GEM非线性启动约20。附录平滑分析仅均匀周期、无背景电子流、de0、刚性及足够网格黏性条件；不是一般异质三维证明。理论能量含电子应力而实现动量忽略它，GEM另以v替vstar匹配HiFi；形式定理与实际算例须分开。密度初值印刷为sech而非与tanh平衡相容的sech\textsuperscript{2}；热能定义多写1/2，不能机械照抄，未核代码不判实现错误。

\paragraph{本报告判断}解决大时间步求解成本的重要方法证据；没有验证有限de、全三维或任意高S新物理。预条件中的近似不能当成实际残差删除物理项；接受稿与v1必须区分。

\section{TS6FEC2025：2025年}\label{paper:TS6FEC2025}

\textbf{The scaling of the ion heating and electrostatic potential in spherical tokamak}\citep{TS6FEC2025}\

T. Ahmadi;  Y. Ono;  H. Tanaka;  H. Tanabe.

\textsf{核读范围：}PDF1–8 all manuscript,  methods,  results,  summary,  references;  no appendix.；\textsf{图表：}All Figs1–7 visually inspected on pp2–6 and Eq1 p2 / heating definition p7 checked directly. Key graph axes/offsets read from image..

\href{run:papers/TS6FEC2025.pdf}{打开本地PDF：\texttt{papers/TS6FEC2025.pdf}}\quad\href{https://conferences.iaea.org/event/392/papers/35991/}{DOI／版本入口}

\paragraph{在报告中的研究脉络}\hyperref[story:energy]{电子独有重联与热化}。关联表示本报告的阅读线索，不自动表示这些论文互相引用。

\paragraph{实际模型}2025年IAEA会议预稿结合TS-6实验、解析标度与VPIC动理学模拟。导言提及双流体Hall MHD，但正文没有独立Hall方程、闭合、网格或结果，不能计为Hall加热模拟。实验Ti由中性Doppler层析及离子—中性松弛假设推断，电势由浮动电势和电子温度反演；PIC采用约化质量比，需要与实际诊断量仔细区分。

\paragraph{计算／实验细节}实验重联场约0.014–0.03 T，典型导向场0.15–0.16 T，场强比约4–9，并比较氢与氩。PIC质量比200、电子等离子体与回旋频率比2、初始Ti=Te、每格每物种200粒子，演化约50个离子回旋时间。x周期，y/z导电且反射粒子；网格/步长文字缺关系符号，未报告实际总格数及质量比、粒子数收敛。

\paragraph{核实的发现}作者报告四极电势与下游Doppler温度及软X射线的时空关联，提出电势差随重联场和导向场乘积变化，以及导向场调节离子增能的趋势。所谓实验与模拟约20\%以内的差别指PIC电势幅度，既不是Hall结果，也不是离子/电子热分配误差。这些是已核读会议稿的主张，尚不能据此独立确认普适不可逆热化机制或加热比例。

\paragraph{限制与未完成问题}无完整物种热预算；PDF7的PIC温升写成mi[〈v\textsuperscript{2}〉-〈v\textsuperscript{2}〉\textsubscript{0}]/2，未明确扣均流，可能混入有序能量。浮动电势转换假设不充分，图3电势拟合有427 V截距，图4误差较大，图2色标与温升叙述也需澄清。时间关联不唯一确定因果，缺原始反演、合成探针和数值收敛；未确认正式期刊版本，不接受其普适或反应堆外推结论。

\paragraph{本报告判断}该8页会议稿适合作为近期实验与PIC线索，以及“必须核全文才能确认模型身份”的边界例。即使标题、导言或摘要提Hall，只有实际报告的模拟和诊断才构成证据；电势相符也不等于热量、熵增或Ti/Te相符。应保留作者主张与综述判断的距离，继续核查诊断反演和扣除均流后的内能，而不能据此宣称Hall加热比已获验证。

\section{Closure2026：2026年}\label{paper:Closure2026}

\textbf{Capturing Secondary Kinetic Instabilities in Three‐Dimensional Dayside Reconnection Using an Improved Gradient‐Based Closure}\citep{Closure2026}\

K. Bradshaw;  A. H. Hakim;  James Juno;  J. Pawlak;  Jason M. TenBarge;  A. Bhattacharjee.

\textsf{核读范围：}1-23（arXiv2602.16960v1全部，含附录A）；\textsf{图表：}图1-11、附录A1及式A1-A12；图像复核待补；主要结果页已视觉复核.

\href{run:papers/Closure2026.pdf}{打开本地PDF：\texttt{papers/Closure2026.pdf}}\quad\href{https://doi.org/10.1029/2026ja035250}{DOI／版本入口}

\paragraph{在报告中的研究脉络}\hyperref[story:closure]{压力、热流与闭合}。关联表示本报告的阅读线索，不自动表示这些论文互相引用。

\paragraph{实际模型}Gkeyll三维非对称十矩；温度梯度闭合与局地闭合对照；顶点对称热扩散和alpha=.75限制器。

\paragraph{计算／实验细节}盒40.96×20.48×10.24di，全周期双电流片，1152×576×288，mi/me100, c/vA200, Bg/Bsx=.099；nm/ns=.062, Bxm/Bxs=-1.696, Tem/Tes12.3, Tim/Tis5.63；k0αdα约1，磁噪声.0005，psi种子.1。

\paragraph{核实的发现}梯度闭合出现约5Ωci\textasciicircum{}-1内LHDI、约20后kink与通量绳/次级岛，局地同参不出现；三维率先峰后降，压力张量支撑X点非理想场。梯度闭合成本约3倍。

\paragraph{限制与未完成问题}温度各向异性峰值侧别与MMS/PIC不符，非旋转对称度与各向异性区域过宽，作者认为背景跨场扩散产生不物理各向异性。独立2D Vlasov对照kink仅数个离子周期出现，流体需数十周期，且遗漏小尺度模。正文未给质量比/网格独立收敛序列。

\paragraph{本报告判断}证明某些先前受抑制的横向不稳定性可在十矩中出现，不是准确恢复全部三维动理学。建议的磁场方向热流及弱各向同性化仍为未来改进。

\section{F600：2026年}\label{paper:F600}

\textbf{System size dependence of collisionless reconnection rate}\citep{F600}\

Yi-Min Huang;  Naoki Bessho;  Li-Jen Chen;  Judith T. Karpen;  Amitava Bhattacharjee.

\textsf{核读范围：}arXiv:2604.18787v1全部6页；正文、方法、结果、讨论及参考文献。正式PRL入口另存，未把期刊版当作已读版本。；\textsf{图表：}图1–4；直接查看PDF3页图1平均区间及5页图3/4。.

\href{run:papers/F600.pdf}{打开本地PDF：\texttt{papers/F600.pdf}}\quad\href{https://doi.org/10.1103/xpsr-4hqn}{DOI／版本入口}

\paragraph{在报告中的研究脉络}\hyperref[story:global]{层长、驱动与全局完成}。关联表示本报告的阅读线索，不自动表示这些论文互相引用。

\paragraph{实际模型}2D Harris片，WarpX全PIC与HMHD双温Hall流体对照；流体用超电阻破坏冻结，电子和离子分别推进绝热温度。

\paragraph{计算／实验细节}Lx=25.6ℓ、Lz=12.8ℓ、δ=.5ℓ，固定δ/L；背景密度.2n0、Te/Ti=.2、质量比25。PIC ℓ/di=1–30，20格/di，每物种平均250粒/格，ωpe/ωce=2；Hall到100，di\textsuperscript{3}VA/ηH=10\textsuperscript{4}、diVA/ν=100，中心Δx=.1di、Δz=.02di，上下对称半域；离子初始噪声10\textsuperscript{-3}VA。

\paragraph{核实的发现}保持宏观形状并以初始B0VA归一化、相同重联磁通比较，尺寸增大使全局率降低和排气张角缩小。尺寸增30倍PIC平均率降约2.6倍，增100倍Hall降约3.4倍。Hall小尺寸近ℓ\textasciicircum{}(-1/3)，大尺寸变缓；PIC持续产磁岛，Hall进入单X点；Hall波列未见于PIC。

\paragraph{限制与未完成问题}不是任意初始条件下的普适幂律；大尺度是否平台未定，质量比25、二维、闭合和超电阻限制外推。PIC ℓ=10重复例率差12\%。v1页4写平均磁通.05–.3，但图1阴影明确约.5–3；精确平均区间需作者澄清，不能默默纠正文献。

\paragraph{本报告判断}补入全局重联率的尺度依赖争议；要求同时报告δ/L、归一化、比较磁通阶段及统计区间。Hall和PIC趋势相同不等于细结构相同。

\section{F635：2026年}\label{paper:F635}

\textbf{Counter-streaming heat-flux closure for electron-only collisionless magnetic reconnection}\citep{F635}\

Madox C. McGrae-Menge;  Jacob R. Pierce;  Maria Almanza;  Alexander Velberg;  Nathaniel Barbour;  W. Dorland;  Nuno Loureiro;  Frederico Fiúza;  E. P. Alves.

\textsf{核读范围：}arXiv2607.07106v1全部10页；公开三个输入文件已下载核对关键设置；\textsf{图表：}图1-6全部视觉核查；特别核对图4最佳时间指标、图5独立色标.

\href{run:papers/F635.pdf}{打开本地PDF：\texttt{papers/F635.pdf}}\quad\href{https://doi.org/10.48550/arxiv.2607.07106}{DOI／版本入口}

\paragraph{在报告中的研究脉络}\hyperref[story:closure]{压力、热流与闭合}。关联表示本报告的阅读线索，不自动表示这些论文互相引用。

\paragraph{实际模型}2026预印本：固定离子背景，两个电子十矩流体；对照单电子十矩及OSIRIS PIC。不是通常电子—离子双流体。

\paragraph{计算／实验细节}48de×64de双周期，480×640；力自由双层半厚de，Brec=.5、Bg/Brec=.1、beta\_rec=.1。PIC每格625粒子、二阶形函数、dt=.066667/omega\_pe；流体Gkeyll分维有限体积CFL2/3、k=.1/de各向同性松弛。两电子群按片上下来源标记，tanh分配宽de、密度下限.03n0；初始速度相同。Zenodo输入中plasmaBetae0=.3等旧声明不参与实际温度初始化；实际vth=.111803与Brec=.5给beta=.1。

\paragraph{核实的发现}两群相对漂移与不同压力可重建总三阶中心矩的大部分结构，改善单流体的密度空洞与局部压力堆积；热流符号/形状更接近PIC。在示例中单群内部热流对局部Qxxx贡献约4\%；但双流体重建热流幅度仍偏高。

\paragraph{限制与未完成问题}图4是每个PIC参数算例在时间内取最大相似度，非双流体全参数性能或时间平均；图5流体色标不同，不能直接认定幅度一致。高beta时人为分群梯度引发自由膨胀，晚期同群再度多束化；三维、动态分裂和离子耦合均未验证。156倍是按粒子相空间变量数与20矩的表示压缩估算，不是实测运行加速。仍有弱松弛而非严格Q=0。

\paragraph{本报告判断}有价值的新方向是利用来源分群减少每群的闭合负担；不能推成单电子流体在数学上不可能表达这些矩，或热流闭合已解决。

\section{F655：2026年}\label{paper:F655}

\textbf{Secondary Electron-Only Reconnection Driven by Large Scale Ion-Coupled Reconnection and Electron Kelvin-Helmholtz Instabilities in Hybrid Simulations of Solar Wind Turbulence}\citep{F655}\

Joaquín Espinoza-Troni;  Giuseppe Arrò;  F. Califano;  Julia E. Stawarz;  Pablo S. Moya.

\textsf{核读范围：}arXiv:2605.11139v1全部17页核读，包括主文、附录A入流归一化、附录B早期聚类和参考文献；作为2026预印本证据；\textsf{图表：}图1—11及图注全部核对；放大物理页5的次级电子喷流、页8的剪切阈值和涡链、页11的聚类分布.

\href{run:papers/F655.pdf}{打开本地PDF：\texttt{papers/F655.pdf}}\quad\href{https://doi.org/10.48550/arxiv.2605.11139}{DOI／版本入口}

\paragraph{在报告中的研究脉络}\hyperref[story:energy]{电子独有重联与热化}。关联表示本报告的阅读线索，不自动表示这些论文互相引用。

\paragraph{实际模型}二维空间、三维速度的Eulerian hybrid Vlasov-Maxwell：动力学离子、有限惯性等温流体电子，准中性、含Hall与惯性电场椭圆求解；没有电子压力张量或电子温度演化。

\paragraph{计算／实验细节}周期域100πdi见方、3072\textsuperscript{2}空间网格、51\textsuperscript{3}离子速度网格[-5vthi, 5vthi]；质量比100、Ti/Te=1、βi=1、δBrms/B0≈0.28，初始波数k/k0=1—6，约50di注入尺度的自由衰减湍流，dt=0.01Ωi\textasciicircum{}-1，格点尺度电磁滤波。由这些数直接算Δx/de≈1.023，电子惯性长度仅约一个网格，本文未给独立网格/滤波收敛组。

\paragraph{核实的发现}展示两条生成次级电子单独重联的路径：大尺度离子耦合重联喷流挤压/合并磁岛；电子剪切形成Kelvin–Helmholtz相容涡链再扭曲磁场。次级例电子喷流约5VAi, in、没有相关离子喷流，局部De峰比主重联例高约十倍；这不是总加热或总重联率十倍。剪切例kde≈0.14、波长4.4di，速度跃迁大于基于whistler恢复力的估算阈值。HDBSCAN晚期516片中约69\%厚度<di且长度<5di，算法本身没有判别这些片是否实际重联。

\paragraph{限制与未完成问题}预印本、二维、单次湍流实现、有限注入尺度、约一格/de和滤波均限制电子尺度定量结论。69\%是电流片比例而非电子单独重联事件发生率；50.2\%片的|平均De|超过全域标准差也不等于不可逆耗散比例。X点速度被取为当地E×B/B\textsuperscript{2}，非理想区此等同不普遍成立，应与磁通鞍点轨迹核对。页6把VAe/VAi质量比平方根写反，附录A的主事件时间360.7与正文320.7不一致。EKHI用饱和波长回估阈值，支持相容性但非完整线性本征模证明。

\paragraph{本报告判断}纳入两条重要的跨尺度生成路径，但标为待收敛和三维验证的候选机制；不把作者关于耗散/太阳风普遍性的较强表述直接升级为结论。

\section{L312：2026年}\label{paper:L312}

\textbf{Cross-scale observations of expanding bifurcated current sheets in fast magnetic reconnection}\citep{L312}\

Jongsoo Yoo;  Masaaki Yamada;  Tongnyeol Rhee;  Hantao Ji;  Ari Le;  Tai Phan;  Sung Hyun Son;  A.B. Robbins.

\textsf{核读范围：}2026-11-04 Article in Press全文1-15页（含Methods）；下载49页透明同行评议文件，但不计为已逐页核读；\textsf{图表：}图1-9全部图注；图3/4/6-9、表1及式1-9视觉/数值核验.

\href{run:papers/L312.pdf}{打开本地PDF：\texttt{papers/L312.pdf}}\quad\href{https://doi.org/10.1038/s41467-026-79220-2}{DOI／版本入口}

\paragraph{在报告中的研究脉络}\hyperref[ch:numerics]{模型与数值可信度}；\hyperref[story:diagnostic]{实验与可区分诊断}。关联表示本报告的阅读线索，不自动表示这些论文互相引用。

\paragraph{实际模型}MRX实验、MMS/PSP事件与VPIC全动理学对照；使用双流体电子动量分析，不是Hall MHD数值计算。

\paragraph{计算／实验细节}MRX氢气3.5mTorr，320-340μs，约250个有清晰EBCS的重复放电平均，di2.5cm；高分辨一维探针z6di。MMS2017-06-19、PSP2021-04-29各一例。小PIC20×20di/1280\textsuperscript{2}、mi/me1836、Ti/Te5；大PIC200×40di/4032×800、mi/me100、Ti/Te1、初始半厚5di，开放流出，均ωpe/ωce2、2000粒子/格。MRX几何PIC30×60di、200×448、质量比25、ωpe/ωce1、200粒子/物种/格、dtωpe.24。

\paragraph{核实的发现}沿分离线延伸的双峰电流、电势井、离子加速和排气密度增强在三环境出现；MRX两侧片厚拟合.29/.70di，MMS排气约20di，PSP约4000di且边界约20di。MRX电势估算12/6V与测量11/5V近；MMS估算6/5kV与测量5/7kV量级近；PSP电场不确定无法检验。低beta .012大岛破坏长程边界，.12/.36局部延伸，1.2大域保持至边界。MMS扣除上游温度后的焓流占优，PSP体流能流占优。

\paragraph{限制与未完成问题}主数值为PIC，不能证明任意Hall/五矩闭合均能重现。实验按EBCS明显事件挑选，两个空间个例及假定Ti=Te的实验beta不构成普适阈值统计。漂移电流吻合只在强Ex分离线时窗；下游距离由几何长宽比推断；1D磁梯度诊断在多维区失效。式4的Delta B\textsuperscript{2}应由磁压差积分理解；原文式1交叉项符号需连同坐标手性核查。PIC大/小盒同时改变质量比和温比，未给独立网格/质量比收敛；代码公开但自定义分析需索取。2026-11-10访问其Figshare API返回404，不能声称已取得源数据。

\paragraph{本报告判断}新实验对照目标是分离线电流层的宏观延伸和能量分配差异；势井标度可用Hall/双流体检验，低杂波、束分布和准确加热仍需张量/动理学及3D。

\section{L313：2026年}\label{paper:L313}

\textbf{Electron inertia and magnetic reconnection}\citep{L313}\

Allen H. Boozer.

\textsf{核读范围：}arXiv2509.14400v5（2026年1月）全部7页及参考文献；\textsf{图表：}全部式1–37；图1及页2/3/5公式直接视觉核验.

\href{run:papers/L313.pdf}{打开本地PDF：\texttt{papers/L313.pdf}}\quad\href{https://doi.org/10.1063/5.0304789}{DOI／版本入口}

\paragraph{在报告中的研究脉络}\hyperref[story:energy]{电子独有重联与热化}。关联表示本报告的阅读线索，不自动表示这些论文互相引用。

\paragraph{实际模型}离子静止、j=-en ve的电子惯性理论，略去电阻、电子压强和黏性；定义B*=B+curl(me j/(ne\textsuperscript{2}))并证明理想输运。一般离子流动不能只靠选取参考系消除。

\paragraph{计算／实验细节}无新数值/实验数据。恒密度、忽略位移电流时B*=(1-de\textsuperscript{2}\ensuremath{\nabla}\textsuperscript{2})B；量级例n=10\textasciicircum{}20 m\textasciicircum{}-3、de约0.5 cm、qR约20 m。

\paragraph{核实的发现}B可以重联而电子正则涡量B*=B-(me/e)curl ve满足∂tB*=curl(ve×B*)。给出三维直接推导；讨论三维混沌理想拉伸下弱非理想性的对数时间敏感性。

\paragraph{限制与未完成问题}要求光滑电子流及所设闭合；图1借用既有理想混沌流结果。三维并非必然混沌，10–100个对流时间不是普适率定律。导引中心估计不能推广到非磁化电子薄层。页2引言Voigt符号与最终式35不一致；式34位移电流项漏c\textasciicircum{}-2，35恢复。报告直接取旋度推出电子冻结，不沿用式30的速度表述。

\paragraph{本报告判断}第一章区分磁场和正则涡量拓扑、可逆惯性与不可逆加热。属理论澄清，不能当作普适快重联机制的数值验证。

\section{L314：2026年}\label{paper:L314}

\textbf{Gyrofluid magnetic reconnection beyond the δF approximation}\citep{L314}\

F. F. Locker;  M. Rinner;  Markus Held;  A. Kendl.

\textsf{核读范围：}arXiv:2603.06234v2全部22页，含附录A内外层推导；正式2026期刊另有入口。；\textsf{图表：}图1–13、表I–III；直接查看PDF6页Δ'算例，14/15页de扫描与17页长宽比扫描。.

\href{run:papers/L314.pdf}{打开本地PDF：\texttt{papers/L314.pdf}}\quad\href{https://doi.org/10.1063/5.0332416}{DOI／版本入口}

\paragraph{在报告中的研究脉络}\hyperref[story:onset]{启动、爆发与磁岛}。关联表示本报告的阅读线索，不自动表示这些论文互相引用。

\paragraph{实际模型}GREENY二维低β、等温Full-F陀螺流体，演化全密度和正则平行动量；保留离子FLR和非Boussinesq任意波长极化，电子FLR略去，磁扰动仍小；不是全f动理学或任意β模型。

\paragraph{计算／实验细节}方域32ρs、1024\textsuperscript{2}，Arakawa+四阶超扩散，ν10\textsuperscript{-6}–10\textsuperscript{-8}并与迎风对照；例Fig4 β=.005、τi=1、m=1；mi/me1836，长宽比到16；初始电子/离子载流分别比较。de扫描通过β等参数变化，不能把全部范围视为同一低β物理区间。

\paragraph{核实的发现}任意波长极化相对LWL在展示算例中较快；线性算子非正规，存在暂态放大可能；非线性快速增长同时出现谱展宽，不能等同线性暂态。离子初始载流时率通常比电子载流慢约10倍；长域多X点处于不同重联阶段；τi主要改变局部形态，未见主磁岛数随τi系统增加。

\paragraph{限制与未完成问题}非正规性分析使用简化平衡和选定变量范数，不能用条件数或1/ε直接当物理能量增长量。v2式42后k=1却给Δ'=14.3108，与2(1/k-k)=0矛盾；有符号μe为负，de公式应取其绝对值。图11/12轴标签写de\textsuperscript{2}，图注却写de且行列方向相反，参数解释不唯一；若按极小de且固定真实质量比理解，会超出低β排序。未经原始输入澄清不作定量跨区间比较。引言将F788欠分辨增加磁岛表述为抑制磁岛，方向相反。附录存在积分上下限排印疑点，不能当作无误教材照抄。

\paragraph{本报告判断}纳入Full-F极化、暂态放大、多X点不同步等候选新现象，同时单列归一化、排序和算式审计；第一章教学推导采用自洽范数与符号，不盲抄原式。

\Needspace{10\baselineskip}

\section{未完成核读的候选：不据此作内容结论}

\small\begin{longtable}{P{24mm}P{97mm}P{27mm}}

\toprule 文献键／年份 & 题名与原始入口 & 当前状态\\\midrule\endfirsthead

\toprule 文献键／年份 & 题名与原始入口 & 当前状态\\\midrule\endhead

E407 / 1973 & \href{http://www.osti.gov/scitech/biblio/4575517-tearing-mode-instabilties-rotating-plasma-hall-effect}{Tearing mode instabilties of a rotating plasma with Hall effect.} & 未获匹配全文\\

L001 / 1978 & \href{https://doi.org/10.1088/0029-5515/18/8/003}{Drift-tearing modes in a tokamak plasma} & 未获匹配全文\\

E409 / 1981 & \href{https://doi.org/10.1017/s0022377800022996}{Hall current effects on tearing modes in rotating reverse field plasmas} & 未获匹配全文\\

L002 / 1983 & \href{https://doi.org/10.1029/gl010i006p00475}{Hall current effect on tearing mode instability} & 未获匹配全文\\

E391 / 1983 & \href{https://ntrs.nasa.gov/search.jsp?R=19830053518}{Hall current effect on tearing mode instability. [possible cause of magnetic field reconnection in space]} & 未获匹配全文\\

L003 / 1988 & \href{https://doi.org/10.1016/0010-4655(88)90222-6}{On the theory of steady state reconnection} & 未获匹配全文\\

F781 / 1989 & \href{https://www.osti.gov/servlets/purl/7161998}{Direct implicit simulation of collisionless magnetic reconnection} & 有PDF，待核读\\

E544 / 1989 & \href{https://ui.adsabs.harvard.edu/abs/1989ESASP.285...55H/abstract}{Magnetic reconnection simulation using the 2.5 D electromagnetic direct implicit code "AVANTI".} & 未获匹配全文\\

E499 / 1990 & \href{https://doi.org/10.1088/0029-5515/30/12/008}{Sawtooth reconnection} & 未获匹配全文\\

L004 / 1990 & \href{https://doi.org/10.1029/gm058p0063}{Steady magnetic field reconnection} & 未获匹配全文\\

E516 / 1991 & \href{https://doi.org/10.1103/physrevlett.66.425}{Collisionless m=1 tearing mode} & 未获匹配全文\\

E501 / 1991 & \href{https://doi.org/10.1103/physrevlett.66.1458}{Collisionless reconnection and the sawtooth crash} & 未获匹配全文\\

E458 / 1991 & \href{https://doi.org/10.1063/1.859804}{Electron inertia and lower-hybrid fluctuation-driven modified low-m tearing modes} & 未获匹配全文\\

E528 / 1991 & \href{https://doi.org/10.1029/91gl00235}{Sheared field tearing mode instability and creation of flux ropes in the Earth magnetotail} & 未获匹配全文\\

L005 / 1992 & \href{https://doi.org/10.1063/1.860467}{Magnetic reconnection in electron magnetohydrodynamics} & 未获匹配全文\\

E403 / 1993 & \href{https://doi.org/10.1063/1.860880}{Modification of short scale-length tearing modes by the Hall field} & 未获匹配全文\\

Waelbroeck1993 / 1993 & \href{https://doi.org/10.1103/physrevlett.70.3259}{Onset of the sawtooth crash} & 未获匹配全文\\

E510 / 1993 & \href{https://doi.org/10.1088/0741-3335/35/6/009}{Pressure gradient turbulent transport and collisionless reconnection} & 未获匹配全文\\

F719 / 1993 & \href{https://doi.org/10.1063/1.860735}{The theory of the early nonlinear stage of m =1 reconnection in tokamaks*} & 未获匹配全文\\

E489 / 1994 & \href{https://doi.org/10.1029/94ja01767}{Observation of anomalous slow‐mode shock and reconnection layer in the dayside magnetopause} & 未获匹配全文\\

Biskamp1995 / 1995 & \href{https://doi.org/10.1103/physrevlett.75.3850}{Ion-Controlled Collisionless Magnetic Reconnection} & 未获匹配全文\\

E511 / 1995 & \href{https://doi.org/10.1063/1.871192}{Macro-particle simulations of collisionless magnetic reconnection} & 未获匹配全文\\

E538 / 1995 & \href{https://doi.org/10.1063/1.871088}{Nonlinear dynamics of the m=1 kink-tearing instability in a modified magnetohydrodynamic model} & 未获匹配全文\\

F768 / 1995 & \href{http://hdl.handle.net/10068/262417}{Whistler dominated quasi-collisionless magnetic reconnection} & 未获匹配全文\\

F782 / 1996 & \href{http://www.osti.gov/scitech/biblio/489441-magnetic-field-reconnection-collisionless-plasmas}{Magnetic field reconnection in collisionless plasmas} & 未获匹配全文\\

E531 / 1996 & \href{http://www.osti.gov/scitech/biblio/489565-magnetic-reconnection-collisionless-regimes-reversed-shear}{Magnetic reconnection in collisionless regimes with reversed shear} & 未获匹配全文\\

F769 / 1996 & \href{http://hdl.handle.net/11568/54239}{Reconnection of magnetic field lines near critical points} & 未获匹配全文\\

L008 / 1997 & \href{https://doi.org/10.1029/97gl52961}{Breakup of the electron current layer during 3‐D collisionless magnetic reconnection} & 未获匹配全文\\

L009 / 1997 & \href{https://doi.org/10.1063/1.872338}{Collisional and collisionless magnetic reconnection} & 未获匹配全文\\

F710 / 1997 & \href{https://doi.org/10.1063/1.872215}{Comment on “Collisionless m=1 reconnection in tokamaks” [Phys. Plasmas 3,  2411 (1996)]} & 未获匹配全文\\

E506 / 1997 & \href{https://doi.org/10.1063/1.872088}{Particle simulation study of collisionless driven reconnection in a sheared magnetic field} & 未获匹配全文\\

L010 / 1997 & \href{https://doi.org/10.1016/s0079-1946(97)00206-1}{Three-dimensional HALL-MHD simulations of magnetopause reconnection} & 未获匹配全文\\

L012 / 1997 & \href{https://doi.org/10.1029/96ja03634}{Undriven magnetic reconnection in magnetohydrodynamics and Hall magnetohydrodynamics} & 未获匹配全文\\

E325 / 1998 & \href{http://ui.adsabs.harvard.edu/abs/1998APS..DPP.K6S43M/abstract}{Collisionless Hall Magnetohydrodynamic Reconnection} & 未获匹配全文\\

E518 / 1998 & \href{https://doi.org/10.1063/1.873046}{Fast reconnection due to localized anomalous resistivity} & 未获匹配全文\\

L013 / 1998 & \href{https://doi.org/10.1063/1.873005}{Forced magnetic field line reconnection in electron magnetohydrodynamics} & 未获匹配全文\\

F706 / 1998 & \href{https://doi.org/10.1017/s0022377898006655}{Hamiltonian vortices and reconnection in a magnetized plasma} & 未获匹配全文\\

E345 / 1998 & \href{http://ui.adsabs.harvard.edu/abs/1998PlPhR..24..929Z/abstract}{Influence of the Hall effect on the reconnection process in helical geometry} & 未获匹配全文\\

L014 / 1998 & \href{https://doi.org/10.1103/physrevlett.80.4430}{Invariants and Geometric Structures in Nonlinear Hamiltonian Magnetic Reconnection} & 未获匹配全文\\

E505 / 1998 & \href{https://doi.org/10.1029/97ja02699}{Kinetic quasi‐viscous and bulk flow inertia effects in collisionless magnetotail reconnection} & 未获匹配全文\\

F676 / 1998 & \href{https://works.swarthmore.edu/theses/680}{Magnetic Reconnection Studies on SSX} & 未获匹配全文\\

E481 / 1998 & \href{https://doi.org/10.1029/98gl52510}{On the kinetic structure of the magnetotail reconnection layer} & 未获匹配全文\\

E525 / 1998 & \href{https://doi.org/10.48550/arxiv.astro-ph/9804166}{Reconnection via the Tearing Instability} & 有PDF，待核读\\

E502 / 1998 & \href{https://doi.org/10.1029/97ja03528}{Structure of the dissipation region during collisionless magnetic reconnection} & 未获匹配全文\\

F779 / 1999 & \href{https://doi.org/10.1017/s0022377899007497}{Collisionless magnetic reconnection in the presence of an external driving flow} & 未获匹配全文\\

L015 / 1999 & \href{https://doi.org/10.1088/0256-307x/16/11/028}{Difference Between the Ion Gyroviscous Effect and Hall Magnetohydrodynamic Effect in the Magnetic Reconnection} & 未获匹配全文\\

L016 / 1999 & \href{https://doi.org/10.1029/1998ja900094}{Impulsive reconnection dynamics in collisionless laboratory and space plasmas} & 未获匹配全文\\

L017 / 1999 & \href{https://doi.org/10.1088/0256-307x/16/6/026}{Modified Magnetohydrodynamic Model of Magnetic Reconnection} & 未获匹配全文\\

E428 / 1999 & \href{https://ui.adsabs.harvard.edu/abs/1999APS..DPP.CP104B/abstract}{Numerical Study of Two-Fluid Effects on Magnetic Reconnection in Merging Spheromaks} & 未获匹配全文\\

E418 / 1999 & \href{https://doi.org/10.1029/1999gl003627}{Sudden disruption of a thin current sheet in collisionless Hall magnetohydrodynamics due to secondary tearing and coalescence instabilities} & 未获匹配全文\\

F647 / 1999 & \href{https://doi.org/10.1063/1.873436}{The diffusion region in collisionless magnetic reconnection} & 未获匹配全文\\

E500 / 1999 & \href{https://doi.org/10.1029/1999gl900481}{The scaling of collisionless,  magnetic reconnection for large systems} & 未获匹配全文\\

E515 / 2000 & \href{https://doi.org/10.1063/1.873822}{Collisionless reconnection in instabilities and due to external forcing} & 未获匹配全文\\

L018 / 2000 & \href{https://doi.org/10.1029/1999ja000357}{Collisionless reconnection: Effects of Hall current and electron pressure gradient} & 未获匹配全文\\

E337 / 2000 & \href{http://ui.adsabs.harvard.edu/abs/2000APS..DPPNP1059Y/abstract}{Comparison of hybrid and Hall-MHD simulations of reconnection} & 未获匹配全文\\

F764 / 2000 & \href{https://doi.org/10.1063/1.874076}{Fast collisionless reconnection in the whistler frequency range} & 未获匹配全文\\

L019 / 2000 & \href{https://doi.org/10.1029/1999gl003688}{Field‐aligned currents generated in magnetotail reconnection: 3D Hall‐MHD simulations} & 未获匹配全文\\

F648 / 2001 & \href{https://doi.org/10.1029/1999ja001002}{Collisionless magnetic reconnection: Electron processes and transport modeling} & 未获匹配全文\\

F695 / 2001 & \href{https://doi.org/10.1086/319012}{Exact Solutions for Two‐dimensional Steady State Magnetic Reconnection in Partially Ionized Plasmas} & 未获匹配全文\\

Otto2001 / 2001 & \href{https://doi.org/10.1029/1999ja001005}{Geospace Environment Modeling (GEM) magnetic reconnection challenge: MHD and Hall MHD—constant and current dependent resistivity models} & 未获匹配全文\\

BirnHesse2001 / 2001 & \href{https://doi.org/10.1029/1999ja001001}{Geospace Environment Modeling (GEM) magnetic reconnection challenge: Resistive tearing,  anisotropic pressure and Hall effects} & 未获匹配全文\\

L024 / 2001 & \href{https://doi.org/10.1029/1999ja001006}{Geospace Environment Modeling magnetic reconnection challenge: Simulations with a full particle electromagnetic code} & 未获匹配全文\\

E379 / 2001 & \href{http://t2r2.star.titech.ac.jp/cgi-bin/publicationinfo.cgi?q_publication_content_number=CTT100469369}{Geotail observations of the Hall current system: Evidence of magnetic reconnection at the kinetic level} & 未获匹配全文\\

L026 / 2001 & \href{https://doi.org/10.1029/2001ja900038}{Geotail observations of the Hall current system: Evidence of magnetic reconnection in the magnetotail} & 未获匹配全文\\

F758 / 2001 & \href{https://doi.org/10.1029/2000ja000184}{Global hybrid simulation of the dayside reconnection layer and associated field‐aligned currents} & 未获匹配全文\\

E394 / 2001 & \href{http://ui.adsabs.harvard.edu/abs/2001AGUSM..SM31A06F/abstract}{Hall MHD computations of magnetic reconnection with a fully implicit 2 1/2 D nonlinear code} & 未获匹配全文\\

L027 / 2001 & \href{https://doi.org/10.1029/2000ja000290}{Hall effects on the generation of field‐aligned currents in three‐dimensional magnetic reconnection} & 未获匹配全文\\

Ma2001 / 2001 & \href{https://doi.org/10.1029/1999ja001004}{Hall magnetohydrodynamic reconnection: The Geospace Environment Modeling challenge} & 未获匹配全文\\

L030 / 2001 & \href{https://doi.org/10.1029/2000ja000398}{Hybrid and Hall‐MHD simulations of collisionless reconnection: Dynamics of the electron pressure tensor} & 未获匹配全文\\

E346 / 2001 & \href{http://ui.adsabs.harvard.edu/abs/2001AGUFMSM41B0811N/abstract}{Magnetic Reconnection in Kelvin-Helmholtz Vortices: Hall Dynamics} & 未获匹配全文\\

E329 / 2001 & \href{https://doi.org/10.1029/2000ja900151}{Magnetopause reconnection induced by magnetosheath Hall‐MHD fluctuations} & 有PDF，待核读\\

E376 / 2001 & \href{https://doi.org/10.1029/2000gl012797}{Particle Hall‐MHD simulation of collisionless reconnection: Ion gyro‐radius correction} & 未获匹配全文\\

Grasso2001 / 2001 & \href{https://doi.org/10.1103/physrevlett.86.5051}{Phase Mixing and Island Saturation in Hamiltonian Reconnection} & 未获匹配全文\\

L032 / 2001 & \href{https://doi.org/10.1063/1.1353803}{Recent developments in collisionless reconnection theory: Applications to laboratory and space plasmas} & 未获匹配全文\\

F646 / 2001 & \href{https://doi.org/10.1103/physrevlett.87.265003}{Scaling of Collisionless Forced Reconnection} & 未获匹配全文\\

E368 / 2001 & \href{https://ui.adsabs.harvard.edu/abs/2001sps..proc..140F/abstract}{Two topics pointing to the importance of resolving Hall effects in reconnection jets} & 未获匹配全文\\

E360 / 2002 & \href{http://ui.adsabs.harvard.edu/abs/2002AGUFMSM61A0471H/abstract}{3D Hall MHD Reconnection Dynamics in a Strongly Sheared System} & 未获匹配全文\\

F765 / 2002 & \href{https://doi.org/10.1063/1.1459455}{A new laboratory experiment on magnetic reconnection} & 未获匹配全文\\

E356 / 2002 & \href{https://ui.adsabs.harvard.edu/abs/2002AGUFMSM61A0472B/abstract}{Collisionless Hall MHD Reconnection Dynamics: Is the Nonlinear Reconnection Rate Independent of the Mechanism that Breaks Field Lines?} & 未获匹配全文\\

F714 / 2002 & \href{https://doi.org/10.1103/physrevlett.89.055002}{Driven Reconnection about a Magnetic X-Line with Strong Guide Component} & 未获匹配全文\\

F631 / 2002 & \href{https://doi.org/10.1063/1.1494433}{Experimental study of lower-hybrid drift turbulence in a reconnecting current sheet} & 有PDF，待核读\\

L034 / 2002 & \href{https://doi.org/10.1029/2001ja000126}{Fingerprints of collisionless reconnection at the separator,  I,  Ambipolar‐Hall signatures} & 未获匹配全文\\

F598 / 2002 & \href{https://doi.org/10.1029/2002gl015314}{GEM reconnection challenge: Implicit kinetic simulations with the physical mass ratio} & 未获匹配全文\\

E364 / 2002 & \href{https://ui.adsabs.harvard.edu/abs/2002AGUFMSM61A0480K/abstract}{Hall MHD and hybrid simulations of the magnetic reconnection: A comparative study} & 未获匹配全文\\

E354 / 2002 & \href{https://doi.org/10.1063/1.1474424}{Hybrid and Hall-magnetohydrodynamics simulations of collisionless reconnection: Effect of the ion pressure tensor and particle Hall-magnetohydrodynamics} & 未获匹配全文\\

E524 / 2002 & \href{https://doi.org/10.1063/1.1435002}{Kinetic regimes of high frequency magnetic reconnection in a neutral sheet configuration} & 未获匹配全文\\

E417 / 2002 & \href{http://ui.adsabs.harvard.edu/abs/2002APS..DPPCP1016C/abstract}{Linear and nonlinear regimes of the 2-D Kelvin-Helmholtz/Tearing instability in Hall MHD.} & 未获匹配全文\\

E341 / 2002 & \href{http://ui.adsabs.harvard.edu/abs/2002AGUFMSM61A0444O/abstract}{Magnetic Reconnection in the Presence of Flow Shear and Hall Physics} & 未获匹配全文\\

E367 / 2002 & \href{https://ui.adsabs.harvard.edu/abs/2002APS..DPPUI1005Y/abstract}{Plasma Pressure Tensor Effects on Reconnection: Hybrid and Hall-MHD Simulations} & 未获匹配全文\\

E410 / 2002 & \href{https://ui.adsabs.harvard.edu/abs/2002APS..DPPCP1015M/abstract}{Quasilinear Hall Dynamo Theory For Two-Fluid Tearing Instability} & 未获匹配全文\\

E503 / 2002 & \href{https://doi.org/10.1088/0741-3335/44/12b/327}{Recent advances in collisionless magnetic reconnection} & 未获匹配全文\\

F567 / 2002 & \href{https://doi.org/10.1029/2001ja009148}{Reconnection onset in the tail of Earth's magnetosphere} & 未获匹配全文\\

E371 / 2002 & \href{http://ui.adsabs.harvard.edu/abs/2002AGUFMSM61A0467D/abstract}{The Role of Hall Electric Fields in Two-dimensional Forced Magnetic Reconnection} & 未获匹配全文\\

Huba2002 / 2002 & \href{https://doi.org/10.1063/1.1514970}{Three-dimensional Hall magnetic reconnection} & 未获匹配全文\\

L037 / 2003 & \href{https://doi.org/10.1134/1.1592555}{A hall dynamo effect driven by two-fluid tearing instability} & 未获匹配全文\\

E542 / 2003 & \href{http://citeseerx.ist.psu.edu/viewdoc/summary?doi=10.1.1.959.2557}{Driven Collisionless Reconnection about a Magnetic X-line in the Absence of a Magnetic Guide Field} & 未获匹配全文\\

E384 / 2003 & \href{http://ui.adsabs.harvard.edu/abs/2003EAEJA.....7302D/abstract}{Effects of Hall Electric Fields on the Saturation of Flux Pile-up Magnetic Reconnection} & 未获匹配全文\\

E519 / 2003 & \href{https://doi.org/10.1088/0741-3335/45/12a/020}{Electron kinetics in collisionless magnetic reconnection} & 未获匹配全文\\

E455 / 2003 & \href{https://doi.org/10.1063/1.1615242}{Electron-inertia effects on driven magnetic field reconnection} & 未获匹配全文\\

F730 / 2003 & \href{https://doi.org/10.1063/1.1566026}{Global extended magnetohydrodynamic studies of fast magnetic reconnection} & 未获匹配全文\\

E401 / 2003 & \href{https://doi.org/10.1016/s0375-9601(02)01807-8}{Hall MHD effects on the 2D Kelvin–Helmholtz/tearing instability} & 未获匹配全文\\

E351 / 2003 & \href{http://ui.adsabs.harvard.edu/abs/2003APS..DPPGP1117S/abstract}{Hall Mediated Reconnection in Systems with Wide Current Sheets} & 未获匹配全文\\

L038 / 2003 & \href{https://doi.org/10.1063/1.1590980}{Hall current effects in dynamic magnetic reconnection solutions} & 有PDF，待核读\\

L039 / 2003 & \href{https://doi.org/10.1063/1.1591182}{Magnetohydrodynamic modeling of two-dimensional reconnection in the Magnetic Reconnection Experiment} & 未获匹配全文\\

F673 / 2003 & \href{https://doi.org/10.1103/physrevlett.90.035003}{Measurement of the Current Sheet during Magnetic Reconnection in a Toroidal Plasma} & 未获匹配全文\\

L040 / 2003 & \href{https://doi.org/10.1063/1.1554739}{Numerical modeling of linear drift-tearing mode stability} & 未获匹配全文\\

L041 / 2003 & \href{https://doi.org/10.1063/1.1559971}{Plasma pressure tensor effects on reconnection: Hybrid and Hall-magnetohydrodynamics simulations} & 未获匹配全文\\

L042 / 2003 & \href{https://doi.org/10.1029/2002ja009699}{Signatures of collisionless magnetic reconnection} & 未获匹配全文\\

L043 / 2003 & \href{https://doi.org/10.1029/2003ja009900}{Structure of the Hall current system in the vicinity of the magnetic reconnection site} & 未获匹配全文\\

L044 / 2003 & \href{https://doi.org/10.1063/1.1578998}{Three-dimensional electron magnetohydrodynamic reconnection. I. Fields,  currents,  and flows} & 未获匹配全文\\

L045 / 2003 & \href{https://doi.org/10.1063/1.1578999}{Three-dimensional electron magnetohydrodynamic reconnection. II. Tilt and precession of a field-reversed configuration} & 未获匹配全文\\

L046 / 2003 & \href{https://doi.org/10.1063/1.1579000}{Three-dimensional electron magnetohydrodynamic reconnection. III. Energy conversion and electron heating} & 未获匹配全文\\

L047 / 2003 & \href{https://doi.org/10.1063/1.1579001}{Three-dimensional electron magnetohydrodynamic reconnection. IV. Instabilities,  fluctuations,  and emissions} & 未获匹配全文\\

E484 / 2003 & \href{https://doi.org/10.1029/2001ja009180}{Whistler‐mediated magnetic reconnection in large systems: Magnetic flux pileup and the formation of thin current sheets} & 未获匹配全文\\

E382 / 2004 & \href{https://ui.adsabs.harvard.edu/abs/2004APS..DPPBP1073B/abstract}{Collisionless Reconnection Dynamics Using PIC and Hall MHD Simulations: A Comparative Study} & 未获匹配全文\\

L048 / 2004 & \href{https://doi.org/10.1063/1.1791640}{Collisionless magnetic reconnection with arbitrary guide field} & 未获匹配全文\\

F720 / 2004 & \href{https://doi.org/10.1063/1.1796691}{Driven reconnection about a magnetic X point in weakly collisional plasmas} & 未获匹配全文\\

E338 / 2004 & \href{https://doi.org/10.1088/0256-307x/21/7/059}{Effects of Hall Current in the Driven Reconnection with Various Scales} & 未获匹配全文\\

E456 / 2004 & \href{http://ui.adsabs.harvard.edu/abs/2004cosp...35.1242M/abstract}{Electron inertia modifications to X-point reconnection in solar flares} & 未获匹配全文\\

E532 / 2004 & \href{https://doi.org/10.1086/425646}{Exact Steady State Reconnection Solutions in Weakly Collisional Plasmas} & 未获匹配全文\\

L050 / 2004 & \href{https://doi.org/10.1063/1.1826094}{Exact models for Hall current reconnection with axial guide fields} & 有PDF，待核读\\

Huba2004 / 2004 & \href{https://doi.org/10.1103/physrevlett.93.175003}{Hall Magnetic Reconnection Rate} & 未获匹配全文\\

L052 / 2004 & \href{https://doi.org/10.1029/2003gl018689}{Hall and Turbulence Effects on Magnetic Reconnection} & 未获匹配全文\\

E400 / 2004 & \href{https://doi.org/10.1017/s0022377804002880}{Hall dynamics and resistive tearing instability} & 未获匹配全文\\

HubaGuide2004 / 2004 & \href{https://doi.org/10.1063/1.1834592}{Hall magnetic reconnection: Guide field dependence} & 未获匹配全文\\

F560 / 2004 & \href{https://doi.org/10.1146/annurev.astro.42.053102.134039}{Impulsive Magnetic Reconnection in the Earth's Magnetotail and the Solar Corona} & 未获匹配全文\\

L054 / 2004 & \href{https://doi.org/10.5194/angeo-22-935-2004}{Influence of the Hall term on KH instability and reconnection inside KH vortices} & 有PDF，待核读\\

E429 / 2004 & \href{http://ui.adsabs.harvard.edu/abs/2004AGUFMSM51D..07Y/abstract}{Laboratory Studies of the Physics of Two-Fluids MHD Physics for Magnetic Reconnection} & 未获匹配全文\\

L055 / 2004 & \href{https://doi.org/10.1103/physrevlett.93.045002}{Measurement of the Hall Dynamo Effect during Magnetic Reconnection in a High-Temperature Plasma} & 未获匹配全文\\

L056 / 2004 & \href{https://doi.org/10.1029/2004ja010478}{On magnetic reconnection regimes and associated three‐dimensional asymmetries: Hybrid,  Hall‐less hybrid,  and Hall‐MHD simulations} & 未获匹配全文\\

E480 / 2004 & \href{https://doi.org/10.1029/2004gl019553}{On the generation and structure of the quadrupole magnetic field in the reconnection process: Comparative simulation study} & 未获匹配全文\\

F666 / 2004 & \href{https://doi.org/10.1103/physrevlett.93.105001}{Structure of the Magnetic Reconnection Diffusion Region from Four-Spacecraft Observations} & 有PDF，待核读\\

E324 / 2004 & \href{http://ui.adsabs.harvard.edu/abs/2004APS..DPPBP1080H/abstract}{The Hall Magnetic Reconnection Rate} & 未获匹配全文\\

E353 / 2004 & \href{https://ui.adsabs.harvard.edu/abs/2004AGUSMSM43A..04H/abstract}{Time Scales for Energy Release in Hall Magnetic Reconnection} & 未获匹配全文\\

L060 / 2004 & \href{https://doi.org/10.1063/1.1776565}{Turbulence in low-β reconnection} & 未获匹配全文\\

L061 / 2004 & \href{https://doi.org/10.1063/1.1773778}{Two-fluid tearing instability in force-free magnetic configuration} & 未获匹配全文\\

E520 / 2004 & \href{https://doi.org/10.1063/1.1758718}{Two-scale structure of the current layer controlled by meandering motion during steady-state collisionless driven reconnection} & 未获匹配全文\\

L062 / 2005 & \href{https://doi.org/10.1016/j.physd.2005.06.003}{Bifurcation analysis of magnetic reconnection in Hall-MHD-systems} & 未获匹配全文\\

L064 / 2005 & \href{https://doi.org/10.1063/1.1872893}{Current singularities: Drivers of impulsive reconnection} & 未获匹配全文\\

E397 / 2005 & \href{https://terpconnect.umd.edu/~drake/news/TriggeringSolarMayhem.pdf}{Drake and his colleagues confirmed their theoretical predictions with supercomputer simulations. Researchers haven't been able to detect Hall reconnection in the sun because the key effects occur in such small regions,  but satellites have spotted it in the earth's magnetosphere.} & 排除：非论文\\

F783 / 2005 & \href{https://doi.org/10.1029/2005ja011229}{Electron surfing acceleration in magnetic reconnection} & 有PDF，待核读\\

F572 / 2005 & \href{https://doi.org/10.1029/2004ja010948}{Externally driven magnetic reconnection in the presence of a normal magnetic field} & 未获匹配全文\\

L066 / 2005 & \href{https://doi.org/10.1029/2004gl022058}{Forced magnetic reconnection} & 有PDF，待核读\\

E508 / 2005 & \href{https://doi.org/10.1029/2004gl021245}{Generalized Ohm's law in a 3‐D reconnection experiment} & 未获匹配全文\\

L067 / 2005 & \href{https://doi.org/10.1063/1.1870003}{Hall effect and fine structures in magnetic reconnection with high plasma β} & 未获匹配全文\\

L068 / 2005 & \href{https://doi.org/10.1029/2004ja010675}{Hall effect in incompressible magnetic reconnection} & 未获匹配全文\\

L069 / 2005 & \href{https://doi.org/10.1016/j.jastp.2005.03.013}{Hall effect on magnetic reconnection at the Earth's magnetopause} & 未获匹配全文\\

F678 / 2005 & \href{https://arxiv.org/pdf/physics/0502001}{Hysteresis and the Onset of Fast Magnetic Reconnection} & 重复：Cassak2005\\

F563 / 2005 & \href{https://doi.org/10.1029/2004ja010751}{Kinetic theory of collisionless tearing at the magnetopause} & 未获匹配全文\\

E427 / 2005 & \href{http://ui.adsabs.harvard.edu/abs/2005AGUFMSM11B..02Y/abstract}{Laboratory Studies of Magnetic Reconnection and the Physics of Two-Fluids MHD} & 未获匹配全文\\

L070 / 2005 & \href{https://doi.org/10.1063/1.1900094}{Linear theory for fast collisionless magnetic reconnection in the lower-hybrid frequency range} & 未获匹配全文\\

F739 / 2005 & \href{https://doi.org/10.1063/1.1883184}{Magnetic reconnection on the ion-skin-depth scale in the dusty magnetotail of a comet} & 未获匹配全文\\

F689 / 2005 & \href{https://doi.org/10.1016/j.asr.2005.01.057}{Magnetic reconnection within MHD-scale Kelvin–Helmholtz vortices triggered by electron inertial effects} & 未获匹配全文\\

L071 / 2005 & \href{https://doi.org/10.1029/2005gl023362}{Magnetic reconnection within rolled‐up MHD‐scale Kelvin‐Helmholtz vortices: Two‐fluid simulations including finite electron inertial effects} & 未获匹配全文\\

E446 / 2005 & \href{https://doi.org/10.1134/1.2031627}{Modeling of the tearing instability in unreduced two-fluid magnetohydrodynamics} & 未获匹配全文\\

L072 / 2005 & \href{https://doi.org/10.1088/0029-5515/45/11/006}{Nonlinear behaviour of collisionless double tearing mode induced by electron inertia} & 未获匹配全文\\

F650 / 2005 & \href{http://lucid.igpp.ucla.edu/lessons/ess265/2005/7th-ISSS/CONTENTS/DATA_PDF/P-1-29.PDF}{Roles of Initial Current Carriers in Three-dimensional Magnetic Reconnection} & 未获匹配全文\\

L073 / 2005 & \href{https://doi.org/10.1007/s10598-005-0026-5}{Simulation of Space Plasma Allowing for the Hall Effect: Magnetic Reconnection} & 未获匹配全文\\

E439 / 2005 & \href{http://absimage.aps.org/image/DPP05/MWS_DPP05-2005-001148.pdf}{Study of Two-Fluid MHD Physics and the Effects of Boundary Conditions on Magnetic Reconnection in MRX} & 未获匹配全文\\

E430 / 2005 & \href{https://ui.adsabs.harvard.edu/abs/2005APS..DPPLP1096Y/abstract}{Study of Two-fluid Effects during Magnetic Reconnection in a Laboratory Experiment} & 未获匹配全文\\

L074 / 2005 & \href{https://doi.org/10.1103/physrevlett.95.045003}{Suppression of Hall-Term Effects by Gyroviscous Cancellation in Steady Collisionless Magnetic Reconnection} & 未获匹配全文\\

L075 / 2005 & \href{https://doi.org/10.1016/j.asr.2005.06.054}{The role of Hall currents on incompressible magnetic reconnection} & 未获匹配全文\\

L077 / 2005 & \href{https://doi.org/10.1029/2005ja011228}{The “Newton Challenge”: Kinetic aspects of forced magnetic reconnection} & 未获匹配全文\\

L078 / 2005 & \href{https://doi.org/10.1017/s0022377805003922}{Theoretical and numerical simulation study of the structure of the Hall electric field in the vicinity of the magnetic reconnection site} & 未获匹配全文\\

E436 / 2005 & \href{http://absimage.aps.org/image/DPP05/MWS_DPP05-2005-000619.pdf}{Two-fluid effects on 3D reconnection in the SSX experiment with comparisons to space data} & 未获匹配全文\\

F565 / 2006 & \href{https://doi.org/10.1063/1.2375039}{A generalized Petschek magnetic reconnection rate} & 未获匹配全文\\

E377 / 2006 & \href{https://ui.adsabs.harvard.edu/abs/2006cosp...36.3552L/abstract}{A numerical study of low frequency wave in Hall MHD reconnection with various plasma} & 未获匹配全文\\

L080 / 2006 & \href{http://hdl.handle.net/1903/4139}{Catastrophe Model for the Onset of Fast Magnetic Reconnection} & 有PDF，待核读\\

Knoll2006 / 2006 & \href{https://doi.org/10.1103/physrevlett.96.135001}{Coalescence of Magnetic Islands in the Low-Resistivity,  Hall-MHD Regime} & 未获匹配全文\\

F780 / 2006 & \href{http://library.tue.nl/csp/dare/LinkToRepository.csp?recordnumber=686527}{Contour dynamics modelling of collisionless magnetic reconnection} & 未获匹配全文\\

L082 / 2006 & \href{https://doi.org/10.1103/physrevlett.97.135002}{Coupling between Global Geometry and the Local Hall Effect Leading to Reconnection-Layer Symmetry Breaking} & 未获匹配全文\\

L083 / 2006 & \href{https://doi.org/10.1029/2005ja011536}{Density depletion and Hall effect in magnetic reconnection} & 未获匹配全文\\

L084 / 2006 & \href{https://doi.org/10.1142/s0217984906011621}{ELECTRON PARALLEL COMPRESSIBILITY IN THE NONLINEAR DEVELOPMENT OF TWO-DIMENSIONAL COLLISIONLESS MAGNETOHYDRODYNAMIC RECONNECTION} & 未获匹配全文\\

E425 / 2006 & \href{http://absimage.aps.org/image/DPP06/MWS_DPP06-2006-000627.pdf}{Experimental Study of Two-fluid Effects on Magnetic Reconnection} & 未获匹配全文\\

Yamada2006 / 2006 & \href{https://doi.org/10.1063/1.2203950}{Experimental study of two-fluid effects on magnetic reconnection in a laboratory plasma with variable collisionality} & 未获匹配全文\\

L086 / 2006 & \href{https://doi.org/10.1063/1.2212397}{Forced Hall magnetic reconnection: Parametric variation of the “Newton Challenge”} & 未获匹配全文\\

F738 / 2006 & \href{https://doi.org/10.1103/physrevlett.97.085003}{Forced Magnetic Reconnection and Field Penetration of an Externally Applied Rotating Helical Magnetic Field in the TEXTOR Tokamak} & 未获匹配全文\\

E326 / 2006 & \href{http://ui.adsabs.harvard.edu/abs/2006AGUFMSM41A1443H/abstract}{Forced Reconnection in the Hall Limit} & 未获匹配全文\\

E322 / 2006 & \href{http://absimage.aps.org/image/DPP06/MWS_DPP06-2006-001101.pdf}{Hall Reconnection in MST} & 未获匹配全文\\

L087 / 2006 & \href{https://doi.org/10.1063/1.2398933}{Hall assisted forced magnetic reconnection} & 未获匹配全文\\

L088 / 2006 & \href{https://doi.org/10.1051/0004-6361:20054391}{Hall effects on dynamic magnetic reconnection at an X-type neutral point} & 未获匹配全文\\

F786 / 2006 & \href{http://absimage.aps.org/image/DPP05/MWS_DPP05-2005-020014.pdf}{How Does Collisionless Magnetic Reconnection Work in the Presence of a Guide Magnetic Field} & 未获匹配全文\\

F581 / 2006 & \href{https://doi.org/10.1063/1.2347101}{Kinetic Vlasov simulations of collisionless magnetic reconnection} & 有PDF，待核读\\

F670 / 2006 & \href{https://doi.org/10.1029/2006gl026318}{Magnetic reconnection induced by weak Kelvin‐Helmholtz instability and the formation of the low‐latitude boundary layer} & 未获匹配全文\\

F778 / 2006 & \href{https://ntrs.nasa.gov/search.jsp?R=20060013113}{Particle Demagnetization in Collisionless Magnetic Reconnection} & 未获匹配全文\\

F575 / 2006 & \href{https://doi.org/10.1029/2005ja011521}{Reproducing the bipolar magnetic signature at the jet leading edge by three‐dimensional reconnection with nonzero guide field} & 未获匹配全文\\

E358 / 2006 & \href{https://doi.org/10.1063/1.2218815}{Scalings of steady state Hall magnetohydrodynamic reconnection in high-beta plasmas} & 未获匹配全文\\

E339 / 2006 & \href{http://absimage.aps.org/image/DPP06/MWS_DPP06-2006-000316.pdf}{Stability of a resistive Sweet-Parker reconnection layer and onset of fast reconnection in Hall MHD.} & 未获匹配全文\\

E363 / 2006 & \href{http://ui.adsabs.harvard.edu/abs/2006AGUFMSM41A1454D/abstract}{Structure of the electromagnetic field in Hall MHD turbulence and relation to reconnection} & 未获匹配全文\\

L090 / 2006 & \href{https://doi.org/10.15017/26818}{Study of Thermal Transport Effects on Drift-Tearing Mode} & 未获匹配全文\\

E365 / 2006 & \href{http://absimage.aps.org/image/DPP06/MWS_DPP06-2006-000815.pdf}{The Hall Effect and Magnetic Fluctuations during Fast Reconnection in MRX} & 未获匹配全文\\

L091 / 2006 & \href{https://doi.org/10.1063/1.2363353}{The Hall dynamo effect and nonlinear mode coupling during sawtooth magnetic reconnection} & 未获匹配全文\\

L092 / 2006 & \href{https://doi.org/10.1016/j.asr.2004.11.045}{The role of the Hall effect in collisionless magnetic reconnection} & 未获匹配全文\\

F604 / 2006 & \href{https://doi.org/10.1063/1.2220534}{Time evolution of the electron diffusion region and the reconnection rate in fully kinetic and large system} & 未获匹配全文\\

L093 / 2006 & \href{https://doi.org/10.1063/1.2180729}{Two fluid effects on three-dimensional reconnection in the Swarthmore Spheromak Experiment with comparisons to space data} & 未获匹配全文\\

E383 / 2007 & \href{http://ui.adsabs.harvard.edu/abs/2007AGUFMSH44A1721C/abstract}{A First-Principles Analytical Theory for 2D Magnetic Reconnection in Electron and Hall Magnetohydrodynamics} & 未获匹配全文\\

E359 / 2007 & \href{https://doi.org/10.1016/j.chinastron.2007.10.004}{A Numerical Study of Low-frequency Waves in Hall MHD Reconnection} & 未获匹配全文\\

E385 / 2007 & \href{http://absimage.aps.org/image/DPP07/MWS_DPP07-2007-000437.pdf}{A first-principles analytical theory for 2D magnetic reconnection in electron and Hall MHD.} & 未获匹配全文\\

L095 / 2007 & \href{https://doi.org/10.1063/1.2734948}{Catastrophic onset of fast magnetic reconnection with a guide field} & 未获匹配全文\\

L096 / 2007 & \href{https://doi.org/10.1017/s0022377807006757}{Collapse of the neutral current sheet and reconnection at micro-scales} & 有PDF，待核读\\

L097 / 2007 & \href{https://doi.org/10.1017/s0022377806006313}{Collisionless magnetic reconnection dynamics with electron inertia and parallel electron compressibility} & 未获匹配全文\\

E527 / 2007 & \href{https://doi.org/10.1088/0741-3335/49/5/008}{Comparison between resistive and collisionless double tearing modes for nearby resonant surfaces} & 有PDF，待核读\\

F800 / 2007 & \href{https://doi.org/10.48550/arxiv.0705.3186}{Dissipation of magnetic fields in neutron star crusts due to development of a tearing mode} & 有PDF，待核读\\

L099 / 2007 & \href{https://doi.org/10.5281/zenodo.47873}{Global Two-Fluid Simulations Of Magnetic Reconnection} & 有PDF，待核读\\

L100 / 2007 & \href{https://doi.org/10.1063/1.2774003}{Gyrokinetic simulations of collisionless magnetic reconnection} & 未获匹配全文\\

L101 / 2007 & \href{https://doi.org/10.1016/j.asr.2007.09.035}{Hall MHD reconnection with an initial guide field By0} & 未获匹配全文\\

L102 / 2007 & \href{https://doi.org/10.1063/1.2820904}{Is the “out-of-plane” magnetic perturbation always a quadrupole in the Hall-mediated magnetic reconnection?} & 未获匹配全文\\

E486 / 2007 & \href{https://doi.org/10.1103/physrevlett.99.175002}{Magnetic Reconnection Rate in Space Plasmas: A Fractal Approach} & 未获匹配全文\\

F618 / 2007 & \href{https://doi.org/10.1063/1.2800854}{Magnetic reconnection during collisionless,  stressed,  X-point collapse using particle-in-cell simulation} & 有PDF，待核读\\

E333 / 2007 & \href{http://absimage.aps.org/image/DPP07/MWS_DPP07-2007-001493.pdf}{Measurements of Linear and Nonlinear Hall Reconnection} & 未获匹配全文\\

L103 / 2007 & \href{https://doi.org/10.1029/2006ja012158}{Multi‐point observations of the Hall electromagnetic field and secondary island formation during magnetic reconnection} & 有PDF，待核读\\

E451 / 2007 & \href{http://absimage.aps.org/image/DPP07/MWS_DPP07-2007-001372.pdf}{Numerical Studies of Linear Two Fluid Tearing Modes in Slab and Cylindrical Geometries} & 未获匹配全文\\

F561 / 2007 & \href{https://doi.org/10.11728/cjss2007.02.096}{Numerical Study of Low Frequency Wave in Hall MHD Reconnection With Various Plasma β} & 未获匹配全文\\

E347 / 2007 & \href{https://ui.adsabs.harvard.edu/abs/2007AGUSMSM54A..05V/abstract}{Occurrence and signatures of the Hall-mediated regimes of magnetic reconnection} & 未获匹配全文\\

L104 / 2007 & \href{https://doi.org/10.1063/1.2759893}{On the two-fluid modification of the resistive tearing instability} & 未获匹配全文\\

L105 / 2007 & \href{https://doi.org/10.1143/jpsj.77.014501}{Poloidal Flow Generated by Drift-Tearing Mode} & 未获匹配全文\\

L106 / 2007 & \href{https://doi.org/10.15017/26835}{Poloidal Flow Generated by Tearing Mode} & 有PDF，待核读\\

F662 / 2007 & \href{https://doi.org/10.1063/1.2740595}{Progress in understanding magnetic reconnection in laboratory and space astrophysical plasmas} & 未获匹配全文\\

L108 / 2007 & \href{https://doi.org/10.1103/physrevlett.99.235001}{Steady-State Properties of Driven Magnetic Reconnection in 2D Electron Magnetohydrodynamics} & 未获匹配全文\\

E521 / 2007 & \href{https://doi.org/10.1103/physrevlett.98.265002}{Temperature Gradients in Fast Collisionless Magnetic Reconnection} & 未获匹配全文\\

L109 / 2007 & \href{https://doi.org/10.1063/1.2741241}{The bimodal distribution structure of electron density brought about by the Hall effect in the electron diffusion region during magnetic reconnection} & 未获匹配全文\\

L110 / 2007 & \href{https://hdl.handle.net/10289/2593}{The influence of Hall currents,  plasma viscosity and electron inertia on magnetic reconnection solutions} & 有PDF，待核读\\

L111 / 2007 & \href{https://doi.org/10.1143/jpsj.76.064501}{Thermal Transport Effects on Drift-Tearing Mode} & 未获匹配全文\\

E453 / 2007 & \href{http://absimage.aps.org/image/DPP07/MWS_DPP07-2007-001398.pdf}{Transition From Two-Fluid to Resistive MHD Tearing Instability in Slab and Cylindrical Geometries} & 未获匹配全文\\

L112 / 2007 & \href{https://doi.org/10.1002/cjg2.1069}{Whistle and Hall Magnetic Field Associated with Near‐Earth Magnetotail Reconnection Observed by Cluster} & 未获匹配全文\\

E404 / 2008 & \href{https://doi.org/10.1088/1009-0630/10/4/02}{A Simulation Study of Hall Effect on Double Tearing Modes} & 未获匹配全文\\

F722 / 2008 & \href{https://doi.org/10.1063/1.3029737}{A new fast reconnection model in a collisionless regime} & 有PDF，待核读\\

E366 / 2008 & \href{http://absimage.aps.org/image/DPP08/MWS_DPP08-2008-000188.pdf}{A quantitative,  comprehensive analytical model for ``fast'' magnetic reconnection in Hall MHD} & 未获匹配全文\\

L114 / 2008 & \href{https://doi.org/10.5281/zenodo.3783590}{COMPUTATIONAL STUDY OF THE INTERNAL KINK MODE EVOLUTION AND ASSOCIATED MAGNETIC RECONNECTION PHENOMENA} & 有PDF，待核读\\

L115 / 2008 & \href{https://doi.org/10.1063/1.2917916}{Calculation of resistive magnetohydrodynamics and two-fluid tearing modes by example of reversed-field-pinch-like plasma} & 有PDF，待核读\\

L116 / 2008 & \href{https://doi.org/10.1051/0004-6361:200809624}{Chromospheric magnetic reconnection: Two-fluid simulations of coalescing current loops} & 有PDF，待核读\\

F613 / 2008 & \href{https://doi.org/10.1051/0004-6361:200809624}{Chromospheric magnetic reconnection: two-fluid simulations of coalescing current loops} & 重复：L116\\

E340 / 2008 & \href{https://doi.org/10.1029/2007ja012769}{Cluster observations of collisionless Hall reconnection at high‐latitude magnetopause} & 未获匹配全文\\

L117 / 2008 & \href{https://doi.org/10.1088/0741-3335/50/3/035005}{Collisionless tearing modes in cylindrical geometry} & 未获匹配全文\\

F591 / 2008 & \href{https://doi.org/10.1086/588088}{Development of a Turbulent Outflow During Electron‐Positron Magnetic Reconnection} & 有PDF，待核读\\

L118 / 2008 & \href{https://doi.org/10.1063/1.2969439}{Effect of zonal flow caused by microturbulence on the double tearing mode} & 未获匹配全文\\

L120 / 2008 & \href{https://doi.org/10.1088/0256-307x/25/8/055}{Generation of Electric Field and Net Charge in Hall Reconnection} & 未获匹配全文\\

F692 / 2008 & \href{http://citeseerx.ist.psu.edu/viewdoc/summary?doi=10.1.1.601.8469}{Generation of Suprathermal Electrons during Magnetic Reconnection} & 未获匹配全文\\

Dorfman2008 / 2008 & \href{https://doi.org/10.1063/1.2904600}{Global axisymmetric simulations of two-fluid reconnection in an experimentally relevant geometry} & 未获匹配全文\\

L124 / 2008 & \href{https://doi.org/10.1088/0741-3335/51/1/015009}{Linear tearing mode stability equations for a low collisionality toroidal plasma} & 未获匹配全文\\

L125 / 2008 & \href{https://doi.org/10.1063/1.2876622}{Magnetic reconnection with electron viscosity in electron magnetohydrodynamics} & 未获匹配全文\\

L126 / 2008 & \href{https://doi.org/10.1063/1.2895388}{Magnetic reconnection with pressure gradient effect in compressible electron magnetohydrodynamics} & 未获匹配全文\\

F577 / 2008 & \href{https://doi.org/10.48550/arxiv.0812.4579}{Model of Reconnection of Weakly Stochastic Magnetic Field and its Testing} & 有PDF，待核读\\

L128 / 2008 & \href{https://doi.org/10.1063/1.2980286}{Nonlinear dynamics of rotating drift-tearing modes in tokamak plasmas} & 未获匹配全文\\

E450 / 2008 & \href{http://absimage.aps.org/image/DPP08/MWS_DPP08-2008-000738.pdf}{Numerical Studies of Nonlinear Two Fluid Tearing Modes in Cylindrical RFPs} & 未获匹配全文\\

F585 / 2008 & \href{https://doi.org/10.1029/2008ja013535}{Observations and simulations of asymmetric magnetic field reconnection} & 未获匹配全文\\

E335 / 2008 & \href{https://ui.adsabs.harvard.edu/abs/2008AGUFMSM23C..05C/abstract}{On the properties and limitations of magnetic reconnection in Hall MHD} & 未获匹配全文\\

E447 / 2008 & \href{http://absimage.aps.org/image/DPP08/MWS_DPP08-2008-001107.pdf}{Parametric analysis of the two-fluid,  resistive tearing instability} & 未获匹配全文\\

F748 / 2008 & \href{https://doi.org/10.1063/1.2999532}{Physics of collisionless reconnection in a stressed X-point collapse} & 有PDF，待核读\\

Simakov2008 / 2008 & \href{https://doi.org/10.1103/physrevlett.101.105003}{Quantitative,  Comprehensive,  Analytical Model for Magnetic Reconnection in Hall Magnetohydrodynamics} & 未获匹配全文\\

L131 / 2008 & \href{https://doi.org/10.1029/2007gl033033}{THEMIS observations of modified Hall fields in asymmetric magnetic field reconnection} & 未获匹配全文\\

L132 / 2008 & \href{https://doi.org/10.1029/2007ja012852}{The 2.5‐D analytical model of steady‐state Hall magnetic reconnection} & 未获匹配全文\\

L133 / 2008 & \href{https://doi.org/10.1063/1.2901194}{The Hall fields and fast magnetic reconnection} & 未获匹配全文\\

E355 / 2008 & \href{http://ui.adsabs.harvard.edu/abs/2008AGUSMSM24A..07M/abstract}{The Importance of the Hall Effect in Magnetic Reconnection: Comparisons of Hybrid and Hall- less Hybrid Simulations} & 未获匹配全文\\

L134 / 2008 & \href{https://doi.org/10.1103/physreve.77.026713}{Three-dimensional lattice Boltzmann model for magnetic reconnection} & 未获匹配全文\\

L135 / 2008 & \href{https://doi.org/10.11728/cjss2008.02.114}{Whistler Wave and Hall Effect in Magnetic Reconnection} & 有PDF，待核读\\

L136 / 2009 & \href{https://doi.org/10.1090/psapm/067.2/2605283}{A numerical study of magnetic reconnection: a central scheme for Hall MHD} & 未获匹配全文\\

L138 / 2009 & \href{https://doi.org/10.1088/1367-2630/11/6/063008}{Being on time in magnetic reconnection} & 未获匹配全文\\

F731 / 2009 & \href{https://doi.org/10.1103/physrevlett.102.245003}{Collisionless Magnetic Reconnection via Alfvén Eigenmodes} & 有PDF，待核读\\

L139 / 2009 & \href{https://doi.org/10.1090/psapm/067.2/2605264}{Collisionless magnetic reconnection in a five-moment two-fluid electron-positron plasma} & 有PDF，待核读\\

E491 / 2009 & \href{https://doi.org/10.5194/angeo-27-905-2009}{Collisionless magnetic reconnection: analytical model and PIC simulation comparison} & 有PDF，待核读\\

E373 / 2009 & \href{http://absimage.aps.org/image/DPP09/MWS_DPP09-2009-001646.pdf}{Current sheet extension and reconnection scaling in collisionless,  hyperresistive,  Hall MHD} & 未获匹配全文\\

E545 / 2009 & \href{http://id.nii.ac.jp/1013/00001414/}{Electron Dynamics in Steady Collisionless Driven Reconnection} & 未获匹配全文\\

Jain2009 / 2009 & \href{https://doi.org/10.1063/1.3134045}{Electron scale structures in collisionless magnetic reconnection} & 未获匹配全文\\

L140 / 2009 & \href{https://doi.org/10.48550/arxiv.0905.4765}{Extension of the Electron Dissipation Region in Collisionless Hall MHD Reconnection} & 重复：Sullivan2009\\

F775 / 2009 & \href{https://doi.org/10.1063/1.3106685}{Fast magnetic reconnection in a kinked current sheet} & 未获匹配全文\\

L142 / 2009 & \href{https://doi.org/10.1029/2009ja014544}{In situ observations of reconnection Hall magnetic fields at Mars: Evidence for ion diffusion region encounters} & 未获匹配全文\\

F661 / 2009 & \href{https://doi.org/10.1146/annurev-astro-082708-101726}{Magnetic Reconnection in Astrophysical and Laboratory Plasmas} & 未获匹配全文\\

L143 / 2009 & \href{https://doi.org/10.1063/1.3122050}{Magnetic reconnection with pressure tensor in electron magnetohydrodynamics} & 未获匹配全文\\

E327 / 2009 & \href{http://absimage.aps.org/image/DPP09/MWS_DPP09-2009-020006.pdf}{Measurements of Hall Reconnection in MST} & 未获匹配全文\\

E479 / 2009 & \href{https://doi.org/10.1088/0004-637x/700/1/63}{NUMERICAL TESTS OF FAST RECONNECTION IN WEAKLY STOCHASTIC MAGNETIC FIELDS} & 有PDF，待核读\\

E399 / 2009 & \href{https://doi.org/10.1063/1.3276534}{Nonlinear evolution of double tearing mode in Hall magnetohydrodynamics} & 未获匹配全文\\

F602 / 2009 & \href{https://doi.org/10.1029/2008ja013427}{Observation of waves near lower hybrid frequency in the reconnection region with thin current sheet} & 未获匹配全文\\

F767 / 2009 & \href{https://doi.org/10.1063/1.3266419}{On the role of a nonscalar electron pressure in the collisionless magnetic reconnection} & 未获匹配全文\\

L146 / 2009 & \href{https://doi.org/10.1088/0741-3335/51/5/055018}{Parametric analysis of the two-fluid tearing instability} & 未获匹配全文\\

L147 / 2009 & \href{https://doi.org/10.1063/1.3079075}{Particle-in-cell simulations of tearing modes in reversed-field-pinch-like plasma} & 未获匹配全文\\

L148 / 2009 & \href{https://doi.org/10.1029/2009ja014257}{Plasmoid‐like structures in multiple X line Hall MHD reconnection} & 未获匹配全文\\

L150 / 2009 & \href{https://doi.org/10.1088/0004-637x/705/1/907}{RELATIVISTIC TWO-FLUID SIMULATIONS OF GUIDE FIELD MAGNETIC RECONNECTION} & 有PDF，待核读\\

E342 / 2009 & \href{https://doi.org/10.1029/2009gl037239}{Reconnection Hall current system observed in the magnetotail and in the ionosphere} & 未获匹配全文\\

F573 / 2009 & \href{https://doi.org/10.1063/1.3155453}{Reconnection in semicollisional,  low-β plasmas} & 未获匹配全文\\

E332 / 2009 & \href{http://absimage.aps.org/image/DPP09/MWS_DPP09-2009-000900.pdf}{Reconnection properties in Hall MHD Plasmas} & 未获匹配全文\\

F609 / 2009 & \href{https://doi.org/10.5281/zenodo.47876}{Simulation And Analysis Of Magnetic Reconnection In An Experimental Geometry} & 有PDF，待核读\\

L151 / 2009 & \href{https://doi.org/10.5281/zenodo.16803}{Simulation and Analysis of Magnetic Reconnection in a Laboratory Plasma Astrophysics Experiment} & 有PDF，待核读\\

L152 / 2009 & \href{https://doi.org/10.1063/1.3086867}{Structure of the dissipation region in fluid simulations of asymmetric magnetic reconnection} & 未获匹配全文\\

L153 / 2009 & \href{https://doi.org/10.1088/0004-637x/696/2/1385}{TWO-FLUID MAGNETOHYDRODYNAMIC SIMULATIONS OF RELATIVISTIC MAGNETIC RECONNECTION} & 有PDF，待核读\\

E462 / 2009 & \href{https://doi.org/10.1063/1.3078470}{Tearing mode with guide field gradient in electron magnetohydrodynamics} & 未获匹配全文\\

E423 / 2009 & \href{https://arxiv.org/abs/0902.2074v1}{Two-Fluid MHD Simulations of Relativistic Magnetic Reconnection} & 重复：L153\\

L155 / 2009 & \href{https://doi.org/10.1063/1.3068470}{Two-fluid regimes of the resistive and collisionless tearing instability} & 未获匹配全文\\

F776 / 2009 & \href{https://ntrs.nasa.gov/search.jsp?R=20090012450}{What Turns Reconnection On and Off} & 未获匹配全文\\

E433 / 2010 & \href{http://ui.adsabs.harvard.edu/abs/2010AGUFMSM51B1813J/abstract}{A Two Fluid Code to Study Cross Scale Coupling in Collisionless Magnetic Reconnection} & 未获匹配全文\\

L158 / 2010 & \href{https://doi.org/10.1103/physrevlett.104.205001}{Asymmetry of the Ion Diffusion Region Hall Electric and Magnetic Fields during Guide Field Reconnection: Observations and Comparison with Simulations} & 有PDF，待核读\\

F756 / 2010 & \href{https://doi.org/10.1029/2009ja014650}{Cluster observations of bidirectional beams caused by electron trapping during antiparallel reconnection} & 未获匹配全文\\

E513 / 2010 & \href{https://doi.org/10.1063/1.3430640}{Collisionless magnetic reconnection in the presence of a sheared velocity field} & 未获匹配全文\\

E372 / 2010 & \href{http://absimage.aps.org/image/DPP10/MWS_DPP10-2010-001669.pdf}{Diamagnetic Effects on Asymmetric Reconnection: A Comparative PIC and Hall MHD Study} & 未获匹配全文\\

E440 / 2010 & \href{http://absimage.aps.org/image/DPP10/MWS_DPP10-2010-000576.pdf}{Fast 2D magnetic reconnection: relating anomalous resistivity and two-fluid MHD models} & 未获匹配全文\\

F663 / 2010 & \href{https://doi.org/10.48550/arxiv.1005.3760}{Fast and slow two-fluid magnetic reconnection} & 重复：L160\\

L161 / 2010 & \href{https://doi.org/10.1088/0029-5515/50/2/025014}{Linear and nonlinear stability of drift-tearing mode} & 未获匹配全文\\

F734 / 2010 & \href{https://doi.org/10.1080/10420150903516569}{Local modeling of collisionless magnetic reconnection in the dayside magnetosphere} & 未获匹配全文\\

F660 / 2010 & \href{https://doi.org/10.1103/revmodphys.82.603}{Magnetic reconnection} & 未获匹配全文\\

L162 / 2010 & \href{https://doi.org/10.1029/2009gl041941}{Magnitude of the Hall fields during magnetic reconnection} & 未获匹配全文\\

E334 / 2010 & \href{https://doi.org/10.1063/1.3521336}{Measurements of impulsive reconnection driven by nonlinear Hall dynamics} & 未获匹配全文\\

E357 / 2010 & \href{http://absimage.aps.org/image/DPP10/MWS_DPP10-2010-000419.pdf}{Measurements of nonlinear Hall-driven reconnection in the reversed field pinch} & 未获匹配全文\\

E330 / 2010 & \href{http://absimage.aps.org/image/DPP10/MWS_DPP10-2010-001121.pdf}{On the question of hysteresis in Hall MHD Reconnection} & 未获匹配全文\\

Sullivan2010 / 2010 & \href{https://doi.org/10.1063/1.3518752}{On the question of hysteresis in Hall magnetohydrodynamic reconnection} & 未获匹配全文\\

E424 / 2010 & \href{https://ui.adsabs.harvard.edu/abs/2010APS..DPPCP9139J/abstract}{Pressure anisotropy in two-fluid simulations of magnetic reconnection} & 未获匹配全文\\

L165 / 2010 & \href{https://doi.org/10.48550/arxiv.1010.0222}{Simulation of Fast Magnetic Reconnection using a Two-Fluid Model of Collisionless Pair Plasma without Anomalous Resistivity} & 有PDF，待核读\\

F586 / 2010 & \href{https://doi.org/10.1017/s1743921311006508}{Study of magnetic reconnection in collisional and collionless plasmas in Magnetic Reconnection Experiment (MRX)} & 有PDF，待核读\\

E415 / 2010 & \href{http://absimage.aps.org/image/DPP10/MWS_DPP10-2010-001620.pdf}{Tearing Instabilities in the Reversed-Field Pinch in the Presence of Hall Currents and Pressure Gradients} & 未获匹配全文\\

F745 / 2010 & \href{https://doi.org/10.5194/angeo-28-1327-2010}{The Alfvén edge in asymmetric reconnection} & 有PDF，待核读\\

E464 / 2010 & \href{https://doi.org/10.1063/1.3328819}{The nonlinear evolution of tearing mode with electron viscosity in electron magnetohydrodynamics} & 未获匹配全文\\

E393 / 2010 & \href{https://ui.adsabs.harvard.edu/abs/2010APS..DPPCP9128E/abstract}{The role of kinetic effects and parallel electric fields within the Hall current system of collisionless reconnection} & 未获匹配全文\\

E421 / 2010 & \href{https://research.manchester.ac.uk/en/studentTheses/d99f1bda-cb55-47e2-8d82-ebe240e98456}{Two-Fluid Models of Magnetic Reconnection} & 未获匹配全文\\

E529 / 2010 & \href{http://edoc.mpg.de/472531}{Two-dimensional and three-dimensional collisionless reconnection} & 未获匹配全文\\

E426 / 2010 & \href{https://ui.adsabs.harvard.edu/abs/2010AAS...21545326Z/abstract}{Two-fluid Simulations of Relativistic Magnetic Reconnection in Pair Plasmas} & 未获匹配全文\\

F634 / 2010 & \href{https://doi.org/10.1029/2010ja015335}{Wave properties in the magnetic reconnection diffusion region with high β: Application of the k‐filtering method to Cluster multispacecraft data} & 未获匹配全文\\

E370 / 2011 & \href{http://ui.adsabs.harvard.edu/abs/2011AGUFMSM21A1983M/abstract}{3D Dynamics of Magnetopause Reconnection Using Hall-MHD Global Simulations} & 未获匹配全文\\

L167 / 2011 & \href{https://doi.org/10.1029/2010ja015942}{A 2.5-D electron Hall-MHD analytical model of steady state Hall magnetic reconnection in a compressible plasma} & 未获匹配全文\\

E442 / 2011 & \href{http://absimage.aps.org/image/DPP11/MWS_DPP11-2011-000787.pdf}{A Two-Fluid Reconnection Code,  Implementing the Electron Pressure Tensor With New Anisotropic Equations of State} & 未获匹配全文\\

L168 / 2011 & \href{https://doi.org/10.1063/1.3568828}{A comparison of fluid and kinetic models of steady magnetic reconnection} & 未获匹配全文\\

E416 / 2011 & \href{https://doi.org/10.1140/epjd/e2011-20106-4}{A study of shear flows induced by nonlinear evolution of double tearing modes in Hall magnetohydrodynamics} & 未获匹配全文\\

L169 / 2011 & \href{https://doi.org/10.48550/arxiv.1103.1076}{Collisionless magnetic reconnection: Flux quanta,  field lines,  `composite electrons' -- Is the quantum-Hall effect involved in its micro-scale physics?} & 有PDF，待核读\\

E412 / 2011 & \href{https://ui.adsabs.harvard.edu/abs/2011APS..DPPNP9019A/abstract}{Diamagnetic Effects on Double Tearing Modes in Hall-MHD and PIC Simulations} & 未获匹配全文\\

F777 / 2011 & \href{https://doi.org/10.1063/1.3642609}{Dissipation mechanism in 3D magnetic reconnection} & 未获匹配全文\\

L170 / 2011 & \href{https://doi.org/10.5194/angeo-29-679-2011}{Downward auroral currents from the reconnection Hall-region} & 有PDF，待核读\\

F588 / 2011 & \href{https://doi.org/10.1007/s11434-011-4440-0}{Electron dynamics in collisionless magnetic reconnection} & 有PDF，待核读\\

L171 / 2011 & \href{https://doi.org/10.1029/2011ja016991}{Evidence of the origin of the Hall magnetic field for reconnection: Hall MHD reconstruction results from Cluster observations} & 未获匹配全文\\

F562 / 2011 & \href{https://doi.org/10.1088/0004-637x/743/1/51}{FAST MAGNETIC RECONNECTION AND SPONTANEOUS STOCHASTICITY} & 有PDF，待核读\\

F584 / 2011 & \href{https://doi.org/10.1103/physrevlett.106.215003}{Fast Magnetic Reconnection in Laser-Produced Plasma Bubbles} & 有PDF，待核读\\

F762 / 2011 & \href{https://doi.org/10.1088/0029-5515/51/3/033003}{Fast linear growth of collisionless double tearing modes in a cylindrical plasma} & 未获匹配全文\\

F749 / 2011 & \href{https://doi.org/10.1063/1.3659035}{Gyrokinetic simulations of the tearing instability} & 有PDF，待核读\\

L172 / 2011 & \href{https://doi.org/10.1029/2010ja015992}{Hall MHD simulations of collisionless multiple X line reconnection} & 未获匹配全文\\

E328 / 2011 & \href{http://ui.adsabs.harvard.edu/abs/2012decs.confE...2L/abstract}{Hall Reconnection in Partially Ionized Plasmas in the Magnetic Reconnection Experiment} & 未获匹配全文\\

F574 / 2011 & \href{https://doi.org/10.1063/1.3643792}{Influences of sub-Alfvénic shear flows on nonlinear evolution of magnetic reconnection in compressible plasmas} & 未获匹配全文\\

F630 / 2011 & \href{https://doi.org/10.1016/j.crhy.2010.11.004}{Ion acceleration in antiparallel collisionless magnetic reconnection: Kinetic and fluid aspects} & 有PDF，待核读\\

E507 / 2011 & \href{https://doi.org/10.1029/gm090p0155}{Magnetic Reconnection: A Kinetic Treatment} & 未获匹配全文\\

E495 / 2011 & \href{http://hdl.handle.net/1903/12005}{Magnetic islands produced by reconnection in large current layers: A statistical approach to modeling at global scales} & 有PDF，待核读\\

E478 / 2011 & \href{https://doi.org/10.5194/npg-18-675-2011}{Magnetic reconnection as an element of turbulence} & 有PDF，待核读\\

F680 / 2011 & \href{https://doi.org/10.1063/1.3658034}{Mechanisms of impulsive magnetic reconnection: Global and local aspects} & 未获匹配全文\\

F789 / 2011 & \href{https://doi.org/10.1088/0004-637x/739/2/72}{ONSET OF FAST MAGNETIC RECONNECTION IN PARTIALLY IONIZED GASES} & 有PDF，待核读\\

F583 / 2011 & \href{https://doi.org/10.1063/1.3609771}{Scaling of the magnetic reconnection rate with symmetric shear flow} & 未获匹配全文\\

E406 / 2011 & \href{http://ui.adsabs.harvard.edu/abs/2011AGUFMSM23B2046M/abstract}{Secondary Tearing Mode in Hall MHD and Its Application} & 未获匹配全文\\

L176 / 2011 & \href{https://doi.org/10.1029/156gm18}{The Hall Current System for Magnetic Reconnection in the Magnetotail} & 未获匹配全文\\

F619 / 2011 & \href{http://hdl.handle.net/1903/11865}{The Kinetic Structure of Collisionless Slow Shocks and Reconnection Exhausts} & 有PDF，待核读\\

L177 / 2011 & \href{https://doi.org/10.1063/1.3662430}{The inner structure of collisionless magnetic reconnection: The electron-frame dissipation measure and Hall fields} & 有PDF，待核读\\

L178 / 2011 & \href{https://doi.org/10.1063/1.3602859}{Theory and simulations of the scaling of magnetic reconnection with symmetric shear flow} & 未获匹配全文\\

L179 / 2011 & \href{https://doi.org/10.5194/npg-18-871-2011}{Three-dimensional magnetic reconnection through a moving magnetic null} & 有PDF，待核读\\

F599 / 2012 & \href{https://doi.org/10.11728/cjss2012.06.778}{Characteristics of Reconnection Diffusion Region in the Solar Wind} & 有PDF，待核读\\

F569 / 2012 & \href{https://doi.org/10.5194/npg-19-145-2012}{Collisionless magnetic reconnection in a plasmoid chain} & 有PDF，待核读\\

Ohia2012 / 2012 & \href{https://doi.org/10.1103/PhysRevLett.109.115004}{Demonstration of Anisotropic Fluid Closure Capturing the Kinetic Structure of Magnetic Reconnection} & 未获匹配全文\\

E494 / 2012 & \href{http://en.cnki.com.cn/Article_en/CJFDTOTAL-NCDL201203012.htm}{Distribution and evolution of plasma structure in region near X-line of magnetic reconnection} & 未获匹配全文\\

E386 / 2012 & \href{https://doi.org/10.1063/1.3697561}{Effects of out-of-plane shear flows on fast reconnection in a two-dimensional Hall magnetohydrodynamics model} & 未获匹配全文\\

F590 / 2012 & \href{https://doi.org/10.1029/2012ja017896}{Electron dynamics in the reconnection ion diffusion region} & 未获匹配全文\\

F727 / 2012 & \href{https://doi.org/10.1029/2012ja017742}{Experimental evaluation of common spacecraft data analysis techniques for reconnection region analysis in a laboratory plasma} & 未获匹配全文\\

E348 / 2012 & \href{http://absimage.aps.org/image/DPP12/MWS_DPP12-2012-000125.pdf}{Guide Field Effects on Hall Reconnection in a Laboratory Plasma} & 未获匹配全文\\

E331 / 2012 & \href{https://ui.adsabs.harvard.edu/abs/2012cosp...39..641G/abstract}{Hall electric fields during reconnection} & 未获匹配全文\\

L181 / 2012 & \href{https://doi.org/10.1029/2011ja017006}{Hall magnetohydrodynamic effects for three‐dimensional magnetic reconnection with finite width along the direction of the current} & 未获匹配全文\\

L182 / 2012 & \href{https://doi.org/10.1063/1.4754000}{Linear stability of the tearing mode with two-fluid and curvature effects in tokamaks} & 未获匹配全文\\

F594 / 2012 & \href{https://doi.org/10.1088/0004-637x/757/2/173}{MAGNETIC RECONNECTION,  HELICITY DYNAMICS,  AND HYPER-DIFFUSION} & 未获匹配全文\\

E483 / 2012 & \href{https://doi.org/10.1063/1.3694119}{Magnetic reconnection in high-energy-density laser-produced plasmas} & 未获匹配全文\\

E526 / 2012 & \href{https://doi.org/10.48550/arxiv.1210.0630}{Nonlinear Acceleration Mechanism of Collisionless Magnetic Reconnection} & 有PDF，待核读\\

L183 / 2012 & \href{https://doi.org/10.1051/0004-6361/201219098}{Nonlinear wave propagation and reconnection at magnetic X-points in the Hall MHD regime} & 有PDF，待核读\\

L186 / 2012 & \href{https://doi.org/10.1063/1.4754151}{Reconnection events in two-dimensional Hall magnetohydrodynamic turbulence} & 未获匹配全文\\

F601 / 2012 & \href{https://doi.org/10.1029/2011ja017464}{Scaling of the inner electron diffusion region in collisionless magnetic reconnection} & 未获匹配全文\\

L187 / 2012 & \href{https://doi.org/10.1088/0029-5515/52/6/063020}{Seed island formation by forced magnetic reconnection} & 未获匹配全文\\

E460 / 2012 & \href{http://absimage.aps.org/image/TSF12/MWS_TSF12-2012-000189.pdf}{Simulations of Experiments on Electron Magnetohydrodynamic Reconnection in a Field Reversed Configuration} & 未获匹配全文\\

E402 / 2012 & \href{http://absimage.aps.org/image/APR12/MWS_APR12-2012-000859.pdf}{Supersonic regime of the Hall-MHD tearing instability} & 未获匹配全文\\

E405 / 2012 & \href{https://doi.org/10.1063/1.4739787}{Supersonic regime of the Hall-magnetohydrodynamics resistive tearing instability} & 未获匹配全文\\

L188 / 2012 & \href{https://doi.org/10.1029/2012gl054574}{Three‐dimensional,  impulsive magnetic reconnection in a laboratory plasma} & 未获匹配全文\\

E395 / 2013 & \href{https://ui.adsabs.harvard.edu/abs/2013shin.confE..78D/abstract}{3-D Simulations of Turbulent and Non-Turbulent Hall Magnetic Reconnection Via the Four-Field Equations} & 未获匹配全文\\

F664 / 2013 & \href{https://doi.org/10.1063/1.4811092}{A review of pressure anisotropy caused by electron trapping in collisionless plasma,  and its implications for magnetic reconnection} & 未获匹配全文\\

F687 / 2013 & \href{https://doi.org/10.1088/1749-4699/6/1/014002}{Application of PDSLin to the magnetic reconnection problem} & 未获匹配全文\\

F773 / 2013 & \href{https://doi.org/10.1103/physrevlett.110.265004}{Bifurcated Structure of the Electron Diffusion Region in Three-Dimensional Magnetic Reconnection} & 未获匹配全文\\

F795 / 2013 & \href{https://doi.org/10.1002/jgra.50403}{Cluster observations near reconnection X lines in Earth's magnetotail current sheet} & 未获匹配全文\\

F667 / 2013 & \href{https://doi.org/10.3389/fphy.2013.00031}{Collisionless magnetic reconnection in space plasmas} & 有PDF，待核读\\

E419 / 2013 & \href{https://doi.org/10.1063/1.4789449}{Comment on “Supersonic regime of the Hall-magnetohydrodynamics resistive tearing instability” [Phys. Plasmas 19,  072519 (2012)]} & 未获匹配全文\\

E468 / 2013 & \href{http://ui.adsabs.harvard.edu/abs/2013AGUFMSH13A2039W/abstract}{Forced Magnetic Reconnection at an X-point: Particle-In-Cell and Ten-Moment Extended MHD Simulations} & 未获匹配全文\\

F709 / 2013 & \href{https://doi.org/10.1002/jgra.50136}{Formation of a transient front structure near reconnection point in 3‐D PIC simulations} & 未获匹配全文\\

F801 / 2013 & \href{https://doi.org/10.18725/oparu-2600}{Gyrokinetic simulation of microtearing turbulence} & 未获匹配全文\\

E414 / 2013 & \href{http://absimage.aps.org/image/DPP13/MWS_DPP13-2013-001107.pdf}{Hall-MHD resistive tearing modes in finite-aspect-ratio cylindrical geometry} & 未获匹配全文\\

F564 / 2013 & \href{https://doi.org/10.1002/jgra.50442}{How to identify reconnecting current sheets in incompressible Hall MHD turbulence} & 未获匹配全文\\

F677 / 2013 & \href{https://doi.org/10.1103/physrevlett.110.205002}{Identification of Intermittent Multifractal Turbulence in Fully Kinetic Simulations of Magnetic Reconnection} & 有PDF，待核读\\

F736 / 2013 & \href{https://doi.org/10.1007/s11467-013-0369-x}{Low frequency Whistler waves excited in fast magnetic reconnection processes} & 未获匹配全文\\

E490 / 2013 & \href{https://research.manchester.ac.uk/en/studentTheses/26c3b17d-87ca-4d98-b5b5-3a3d78e0dd03}{Magnetic reconnection and particle acceleration in semi-collisional plasmas} & 未获匹配全文\\

E352 / 2013 & \href{https://ui.adsabs.harvard.edu/abs/2013AGUFMSH51B2092G/abstract}{Magnetic reconnection in Hall-MHD including electron inertia} & 未获匹配全文\\

F596 / 2013 & \href{https://doi.org/10.1002/jgra.50457}{Mechanisms of field‐aligned current formation in magnetic reconnection} & 未获匹配全文\\

F787 / 2013 & \href{https://doi.org/10.1103/physrevlett.111.135001}{New Electric Field in Asymmetric Magnetic Reconnection} & 有PDF，待核读\\

F785 / 2013 & \href{https://doi.org/10.1063/1.4811120}{On phase diagrams of magnetic reconnection} & 未获匹配全文\\

E485 / 2013 & \href{https://doi.org/10.1063/1.4829039}{On the structure of guide magnetic field in the inertia-driven magnetic reconnection with the presence of shear flow} & 未获匹配全文\\

L192 / 2013 & \href{https://doi.org/10.1093/mnras/stt1155}{Onset of Hall-mediated magnetic reconnection in weakly ionized astrophysical plasmas} & 未获匹配全文\\

F582 / 2013 & \href{https://doi.org/10.1063/1.4802941}{Plasmoid instability in high-Lundquist-number magnetic reconnection} & 有PDF，待核读\\

F688 / 2013 & \href{https://doi.org/10.1063/1.4825132}{Propagation speed of rotation signals for field lines undergoing magnetic reconnection} & 未获匹配全文\\

F579 / 2013 & \href{https://doi.org/10.1093/mnras/stt167}{Reconnection-driven plasmoids in blazars: fast flares on a slow envelope} & 有PDF，待核读\\

F665 / 2013 & \href{https://doi.org/10.1103/physrevlett.110.135004}{Regimes of the Electron Diffusion Region in Magnetic Reconnection} & 未获匹配全文\\

E420 / 2013 & \href{https://doi.org/10.1063/1.4789450}{Response to “Comment on ‘Supersonic regime of the Hall-magnetohydrodynamics resistive tearing instability’” [Phys. Plasmas 20,  014703 (2013)]} & 未获匹配全文\\

E517 / 2013 & \href{https://doi.org/10.1063/1.4826201}{Reversible collisionless magnetic reconnection} & 未获匹配全文\\

E413 / 2013 & \href{https://doi.org/10.1016/j.physleta.2013.07.014}{The effects of a shear flow on nonlinear evolutions of double tearing mode in Hall magnetohydrodynamics} & 未获匹配全文\\

L194 / 2013 & \href{https://doi.org/10.1002/jgra.50417}{The influence of out‐of‐plane shear flow on Hall magnetic reconnection and FTE generation} & 未获匹配全文\\

L195 / 2013 & \href{https://doi.org/10.1017/s0022377812001201}{The scaling of reconnection rate in the two-fluid model of a collisionless forced magnetic reconnection} & 未获匹配全文\\

L196 / 2013 & \href{https://lirias.kuleuven.be/handle/123456789/463177}{Two-fluid 20-moment simulation of fast magnetic reconnection} & 海报／幻灯片\\

L198 / 2014 & \href{https://doi.org/10.33915/etd.6627}{A Study of Magnetic Reconnection: From 2D Energy Release to 3D Spreading and Localization} & 未获匹配全文\\

F605 / 2014 & \href{https://doi.org/10.1063/1.4875739}{Asymmetric evolution of magnetic reconnection in collisionless accretion disk} & 有PDF，待核读\\

L199 / 2014 & \href{https://doi.org/10.1002/2014ja020190}{Bursty magnetic reconnection under slow shock‐generated whistler waves} & 未获匹配全文\\

E435 / 2014 & \href{https://agu.confex.com/agu/fm14/webprogram/Paper27518.html}{Comparison of Two-Fluid and Kinetic Simulations of Magnetic Reconnection in Line-Tied Systems} & 未获匹配全文\\

L200 / 2014 & \href{https://doi.org/10.1063/1.4902106}{Conversion of the dominantly ideal perturbations into a tearing mode after a sawtooth crash} & 未获匹配全文\\

E454 / 2014 & \href{http://absimage.aps.org/image/DPP14/MWS_DPP14-2014-000895.pdf}{Effect of parallel heat flux in the gyroviscous tensor on tearing mode instability in two-fluid magnetohydrodynamic model} & 未获匹配全文\\

F610 / 2014 & \href{https://scholars.unh.edu/dissertation/2170}{Estimating proximity to the asymmetric reconnection X-line} & 未获匹配全文\\

F702 / 2014 & \href{https://doi.org/10.1088/0741-3335/57/1/014017}{Fast sawtooth reconnection at realistic Lundquist numbers} & 未获匹配全文\\

L203 / 2014 & \href{https://doi.org/10.48550/arxiv.1409.6985}{Gaussian-Moment Relaxation Closures for Verifiable Numerical Simulation of Fast Magnetic Reconnection in Plasma} & 有PDF，待核读\\

L205 / 2014 & \href{https://doi.org/10.1002/2013ja019224}{Interaction of magnetic reconnection and Kelvin‐Helmholtz modes for large magnetic shear: 1. Kelvin‐Helmholtz trigger} & 未获匹配全文\\

L206 / 2014 & \href{https://doi.org/10.1002/2013ja019225}{Interaction of magnetic reconnection and Kelvin‐Helmholtz modes for large magnetic shear: 2. Reconnection trigger} & 未获匹配全文\\

F725 / 2014 & \href{https://doi.org/10.3389/fphy.2014.00076}{Lessons on collisionless reconnection from quantum fluids} & 有PDF，待核读\\

F797 / 2014 & \href{https://doi.org/10.1103/physrevlett.112.105004}{Magnetic Reconnection in Plasma under Inertial Confinement Fusion Conditions Driven by Heat Flux Effects in Ohm’s Law} & 有PDF，待核读\\

E476 / 2014 & \href{http://hdl.handle.net/11583/2536090}{Magnetic reconnection in high-temperature plasmas: from fluid to gyrofluid modeling} & 未获匹配全文\\

L207 / 2014 & \href{https://doi.org/10.1063/1.4863498}{Nonlinear dynamics of the tearing mode with two-fluid and curvature effects in tokamaks} & 未获匹配全文\\

L208 / 2014 & \href{https://doi.org/10.1088/0029-5515/54/12/123018}{Nonlinear mutual destabilization of the tearing mode and ion temperature gradient mode} & 未获匹配全文\\

F742 / 2014 & \href{https://doi.org/10.1051/0004-6361/201424083}{Relativistic magnetic reconnection in collisionless ion-electron plasmas explored with particle-in-cell simulations} & 有PDF，待核读\\

E349 / 2014 & \href{http://absimage.aps.org/image/DPP14/MWS_DPP14-2014-001290.pdf}{Scaling of Hall Reconnection in the Coalescence Problem in Large Systems} & 未获匹配全文\\

F653 / 2014 & \href{https://doi.org/10.1186/bf03353282}{Strong electron heating and non-Maxwellian behavior in magnetic reconnection} & 有PDF，待核读\\

F713 / 2014 & \href{https://doi.org/10.1186/bf03353267}{Study of local reconnection physics in a laboratory plasma} & 有PDF，待核读\\

E444 / 2014 & \href{http://citeseerx.ist.psu.edu/viewdoc/summary?doi=10.1.1.717.5563}{Studying Incomplete Sawtooth Reconnection in ASDEX Upgrade with the Non-linear Two-fluid 3D MHD Code M3D-C1} & 未获匹配全文\\

E523 / 2014 & \href{https://doi.org/10.1186/bf03353257}{The origins of electrical resistivity in magnetic reconnection: Studies by 2D and 3D macro particle simulations} & 有PDF，待核读\\

E493 / 2014 & \href{https://doi.org/10.1002/2014gl060440}{The scale of the magnetotail reconnecting current sheet in the presence of O+} & 未获匹配全文\\

F770 / 2014 & \href{https://doi.org/10.1063/1.4885636}{Three dimensional instabilities of an electron scale current sheet in collisionless magnetic reconnection} & 未获匹配全文\\

L209 / 2014 & \href{https://doi.org/10.1103/physreve.89.031101}{Three-dimensional fast magnetic reconnection driven by relativistic ultraintense femtosecond lasers} & 未获匹配全文\\

E434 / 2014 & \href{http://hdl.handle.net/1721.1/95854}{Two-fluid simulations of magnetic reconnection with a kinetic closure for the electron pressure anisotropy} & 未获匹配全文\\

F737 / 2015 & \href{https://doi.org/10.1088/0004-637x/805/2/134}{ASYMMETRIC MAGNETIC RECONNECTION IN WEAKLY IONIZED CHROMOSPHERIC PLASMAS} & 有PDF，待核读\\

L211 / 2015 & \href{https://doi.org/10.1002/2015ja021489}{Asymmetric magnetic reconnection with a flow shear and applications to the magnetopause} & 有PDF，待核读\\

E431 / 2015 & \href{https://nifs-repository.repo.nii.ac.jp/records/7192}{Basic Study on the High-Resolution Two-Fluid Simulation of Magnetic Reconnection} & 未获匹配全文\\

F757 / 2015 & \href{https://doi.org/10.1002/2015gl066820}{Development of bifurcated current sheets in solar wind reconnection exhausts} & 未获匹配全文\\

F750 / 2015 & \href{https://doi.org/10.1038/srep08080}{Direct evidence for kinetic effects associated with solar wind reconnection} & 有PDF，待核读\\

E522 / 2015 & \href{https://doi.org/10.1063/1.4933120}{Evolution of electron current sheets in collisionless magnetic reconnection} & 未获匹配全文\\

E457 / 2015 & \href{http://ui.adsabs.harvard.edu/abs/2015EGUGA..17.7964P/abstract}{Fast magnetic reconnection in thin current sheets: effects of different current profiles and electron inertia in Ohm's law.} & 未获匹配全文\\

F603 / 2015 & \href{https://doi.org/10.1063/1.4932332}{Frozen flux violation,  electron demagnetization and magnetic reconnection} & 未获匹配全文\\

F568 / 2015 & \href{https://doi.org/10.1002/2014ja020520}{Generation of Alfvénic waves and turbulence in reconnection jets} & 有PDF，待核读\\

L218 / 2015 & \href{https://doi.org/10.1088/0256-307x/32/9/095202}{How Does the Guide Field Affect the Asymmetry of Hall Magnetic and Electric Fields in Fast Magnetic Reconnection?} & 未获匹配全文\\

F571 / 2015 & \href{https://doi.org/10.1088/0741-3335/58/1/014021}{Magnetic reconnection: from the Sweet–Parker model to stochastic plasmoid chains} & 有PDF，待核读\\

E438 / 2015 & \href{http://absimage.aps.org/image/DPP15/MWS_DPP15-2015-001543.pdf}{Measuring Properties of Magnetic Reconnection in Nonlinear Resistive and Two-Fluid Toroidal Simulations of Sawteeth} & 未获匹配全文\\

L219 / 2015 & \href{https://doi.org/10.1088/0029-5515/55/5/053016}{Modelling and analytic studies of sheared flow effects on tearing modes} & 未获匹配全文\\

L220 / 2015 & \href{https://doi.org/10.1002/2015gl065129}{Modification of the Hall physics in magnetic reconnection due to cold ions at the Earth's magnetopause} & 未获匹配全文\\

E497 / 2015 & \href{http://absimage.aps.org/image/DPP15/MWS_DPP15-2015-001927.pdf}{Recent Developments in Reconnection Theory: the Plasmoid Instability,  Self-Generated Turbulence,  and Implications for Laboratory and Space Plasmas} & 未获匹配全文\\

E343 / 2015 & \href{https://doi.org/10.1063/1.4922057}{Reconnection dynamics with secondary tearing instability in compressible Hall plasmas} & 未获匹配全文\\

Ohia2015 / 2015 & \href{https://doi.org/10.1002/2015GL067117}{Scaling laws for magnetic reconnection,  set by regulation of the electron pressure anisotropy to the firehose threshold} & 未获匹配全文\\

F685 / 2015 & \href{https://doi.org/10.1103/physrevlett.114.205004}{Slowing of Magnetic Reconnection Concurrent with Weakening Plasma Inflows and Increasing Collisionality in Strongly Driven Laser-Plasma Experiments} & 有PDF，待核读\\

L222 / 2015 & \href{https://doi.org/10.33915/etd.5178}{Theory and Simulations of Incomplete Reconnection During Sawteeth Due to Diamagnetic Effects} & 未获匹配全文\\

F637 / 2015 & \href{https://doi.org/10.1615/int.j.uncertaintyquantification.2015006492}{UNCERTAINTY QUANTIFICATION OF THE GEM CHALLENGE MAGNETIC RECONNECTION PROBLEM USING THE MULTILEVEL MONTE CARLO METHOD} & 未获匹配全文\\

F684 / 2016 & \href{https://doi.org/10.1002/2016gl069877}{Cold ion demagnetization near the X‐line of magnetic reconnection} & 有PDF，待核读\\

E422 / 2016 & \href{http://ui.adsabs.harvard.edu/abs/2016cosp...41E.697G/abstract}{Effect of electron inertia in two-fluid magnetic reconnection} & 未获匹配全文\\

E477 / 2016 & \href{http://absimage.aps.org/image/DPP16/MWS_DPP16-2016-000949.pdf}{Energy principle for a gyrofluid model and variational description of collisionless tearing modes} & 未获匹配全文\\

F642 / 2016 & \href{https://doi.org/10.1063/1.4944390}{Extreme ultra-violet burst,  particle heating,  and whistler wave emission in fast magnetic reconnection induced by kink-driven Rayleigh-Taylor instability} & 有PDF，待核读\\

F729 / 2016 & \href{http://citeseerx.ist.psu.edu/viewdoc/summary?doi=10.1.1.743.7627}{Fluid and Magnetofluid Modeling of Relativistic Magnetic Reconnection} & 未获匹配全文\\

E396 / 2016 & \href{https://ui.adsabs.harvard.edu/abs/2016EGUGA..18.9105L/abstract}{Hall-MHD simulations of the magnetosphere-northward solar wind interface : the Kelvin-Helmholtz instability as an entry mechanism for the solar wind through mixing and reconnections} & 未获匹配全文\\

F791 / 2016 & \href{https://doi.org/10.1002/2015sw001313}{Inside the Black Box: Magnetic Reconnection and the Magnetospheric Multiscale Mission} & 未获匹配全文\\

E387 / 2016 & \href{https://doi.org/10.1063/1.4940322}{Ion and electron dynamics generating the Hall current in the exhaust far downstream of the reconnection x-line} & 未获匹配全文\\

L226 / 2016 & \href{https://doi.org/10.1002/2016gl071044}{Kinetic Alfvén wave explanation of the Hall fields in magnetic reconnection} & 有PDF，待核读\\

E556 / 2016 & \href{https://doi.org/10.1088/1361-6587/59/2/025007}{Local properties of magnetic reconnection in nonlinear resistive- and extended-magnetohydrodynamic toroidal simulations of the sawtooth crash} & 有PDF，待核读\\

E375 / 2016 & \href{https://ui.adsabs.harvard.edu/abs/2016EGUGA..1813235A/abstract}{Magnetic reconnection at the magnetopause: Low-energy ions and modification of the Hall physics} & 未获匹配全文\\

F683 / 2016 & \href{https://doi.org/10.1002/2015ja022005}{Magnetic reconnection in Saturn's magnetotail: A comprehensive magnetic field survey} & 未获匹配全文\\

E509 / 2016 & \href{https://doi.org/10.1103/physrevlett.117.015001}{Magnetospheric Multiscale Observations of the Electron Diffusion Region of Large Guide Field Magnetic Reconnection} & 未获匹配全文\\

F669 / 2016 & \href{https://doi.org/10.1002/2016gl068783}{Magnetospheric Multiscale observations of magnetic reconnection associated with Kelvin‐Helmholtz waves} & 未获匹配全文\\

L228 / 2016 & \href{https://doi.org/10.1103/physrevlett.117.185102}{Observations of Hall Reconnection Physics Far Downstream of the X Line} & 有PDF，待核读\\

F728 / 2016 & \href{https://doi.org/10.1063/1.4961431}{On Lorentz invariants in relativistic magnetic reconnection} & 未获匹配全文\\

F649 / 2016 & \href{https://doi.org/10.1098/rspa.2016.0479}{Perspectives on magnetic reconnection} & 有PDF，待核读\\

L229 / 2016 & \href{https://doi.org/10.1093/mnras/stw384}{Relativistic tearing and drift-kink instabilities in two-fluid simulations} & 有PDF，待核读\\

F587 / 2016 & \href{https://doi.org/10.1093/mnras/stw2309}{Resistive tearing instability in electron MHD: application to neutron star crusts} & 有PDF，待核读\\

E388 / 2016 & \href{http://ui.adsabs.harvard.edu/abs/2016AGUFMSM13A2184T/abstract}{Structures of the Hall magnetic field in the dayside magnetic reconnection inferred from Geotail data} & 未获匹配全文\\

L231 / 2016 & \href{https://doi.org/10.1063/1.4943053}{Tearing and surface preserving electron magnetohydrodynamic modes in a current layer} & 有PDF，待核读\\

E361 / 2016 & \href{https://ui.adsabs.harvard.edu/abs/2016AAS...22714103M/abstract}{The plasmoid instability and Hall effect during chromospheric magnetic reconnection} & 未获匹配全文\\

E378 / 2016 & \href{https://confit.atlas.jp/guide/event/jpgu2016/subject/PEM17-11/detail}{The role of Hall magnetic field in large-scale magnetic reconnection dynamics and structure} & 未获匹配全文\\

E443 / 2016 & \href{http://absimage.aps.org/image/DPP16/MWS_DPP16-2016-000133.pdf}{Three-dimensional two-fluid investigation of 3D-localized magnetic reconnection and its relation to whistler waves} & 未获匹配全文\\

L232 / 2016 & \href{https://doi.org/10.1002/2016ja022650}{Two types of whistler waves in the hall reconnection region} & 未获匹配全文\\

F626 / 2016 & \href{https://doi.org/10.1063/1.4948721}{Understanding the dynamics and energetics of magnetic reconnection in a laboratory plasma: Review of recent progress on selected fronts} & 有PDF，待核读\\

E514 / 2016 & \href{https://doi.org/10.1017/s002237781600088x}{‘Ideally’ unstable current sheets and the triggering of fast magnetic reconnection} & 有PDF，待核读\\

E437 / 2017 & \href{https://doi.org/10.1063/1.4982812}{A generalized two-fluid picture of non-driven collisionless reconnection and its relation to whistler waves} & 有PDF，待核读\\

E380 / 2017 & \href{https://confit.atlas.jp/guide/event/jpguagu2017/subject/PEM13-03/detail}{Cold ion heating in the vicinity of the Hall field region in dayside magnetic reconnection} & 未获匹配全文\\

E504 / 2017 & \href{https://doi.org/10.3847/1538-4357/aa9754}{Collisionless Reconnection in Magnetohydrodynamic and Kinetic Turbulence} & 有PDF，待核读\\

F668 / 2017 & \href{https://doi.org/10.1017/s0022377817000812}{Disruption of Alfvénic turbulence by magnetic reconnection in a collisionless plasma} & 有PDF，待核读\\

L233 / 2017 & \href{https://doi.org/10.1002/2017ja024233}{Double‐peaked core field of flux ropes during magnetic reconnection} & 未获匹配全文\\

L234 / 2017 & \href{https://doi.org/10.3847/1538-4357/aa7554}{Effect of Radiation on Chromospheric Magnetic Reconnection: Reactive and Collisional Multi-fluid Simulations} & 未获匹配全文\\

F763 / 2017 & \href{https://escholarship.org/uc/item/0kf95194}{Electron and Ion Acceleration Associated with Magnetotail Reconnection} & 未获匹配全文\\

E390 / 2017 & \href{https://ui.adsabs.harvard.edu/abs/2019EGUGA..21.6241W/abstract}{Electron scale quadrants of the Hall quadrupolar structure observed during asymmetric magnetic reconnection} & 未获匹配全文\\

L235 / 2017 & \href{https://doi.org/10.1103/physrevlett.118.175101}{Electron-Scale Quadrants of the Hall Magnetic Field Observed by the Magnetospheric Multiscale spacecraft during Asymmetric Reconnection} & 有PDF，待核读\\

F732 / 2017 & \href{https://doi.org/10.1103/physrevlett.118.125002}{Experimental Verification of the Role of Electron Pressure in Fast Magnetic Reconnection with a Guide Field} & 未获匹配全文\\

F566 / 2017 & \href{https://doi.org/10.1017/s0022377817000782}{Forced magnetic reconnection} & 有PDF，待核读\\

L237 / 2017 & \href{https://doi.org/10.1088/2058-6272/aa7c17}{Generation of Alfvén wave energy during magnetic reconnection in Hall MHD} & 未获匹配全文\\

E398 / 2017 & \href{https://doi.org/10.1063/1.5004430}{Hall effect on tearing mode instabilities in tokamak} & 未获匹配全文\\

F652 / 2017 & \href{https://doi.org/10.1063/1.5000071}{How anomalous resistivity accelerates magnetic reconnection} & 有PDF，待核读\\

E362 / 2017 & \href{http://ui.adsabs.harvard.edu/abs/2017AGUFMSM13B2376D/abstract}{Kinetic Alfven wave explanation of the Hall signals in magnetic reconnection} & 未获匹配全文\\

E541 / 2017 & \href{https://doi.org/10.1063/1.5007860}{Kinetic-scale flux rope reconnection in periodic and line-tied geometries} & 未获匹配全文\\

E547 / 2017 & \href{https://doi.org/10.1088/2058-6272/aa5801}{Linear tearing modes in an electron-positron plasma} & 未获匹配全文\\

F681 / 2017 & \href{https://doi.org/10.48550/arxiv.1703.01238}{Magnetic Reconnection in Turbulent Diluted Plasmas} & 有PDF，待核读\\

F686 / 2017 & \href{https://doi.org/10.1088/1361-6587/aa6803}{Magnetic reconnection in the high-energy density regime} & 未获匹配全文\\

F721 / 2017 & \href{https://doi.org/10.1088/1741-4326/aa6c70}{Magneto-thermal reconnection processes,  related mode momentum and formation of high energy particle populations} & 有PDF，待核读\\

E488 / 2017 & \href{https://doi.org/10.1063/1.4974758}{On the influence of environmental parameters on mixing and reconnection caused by the Kelvin-Helmholtz instability at the magnetopause} & 未获匹配全文\\

E392 / 2017 & \href{http://urlib.net/sid.inpe.br/mtc-m21b/2017/12.21.17.31}{Particle-in-cell simulations of asymmetric guide-field reconnection: quadrupolar structure of Hall magnetic field} & 海报／幻灯片\\

F705 / 2017 & \href{https://doi.org/10.1002/2016ja023666}{Quadrupolar pattern of the asymmetric guide‐field reconnection} & 未获匹配全文\\

L239 / 2017 & \href{https://doi.org/10.1002/2016ja023620}{Quantitative analysis of a Hall system in the exhaust of asymmetric magnetic reconnection} & 有PDF，待核读\\

E543 / 2017 & \href{https://doi.org/10.1063/1.4973328}{Secondary fast reconnecting instability in the sawtooth crash} & 有PDF，待核读\\

E336 / 2017 & \href{http://ui.adsabs.harvard.edu/abs/2017EGUGA..1915103I/abstract}{Simulations of Hall reconnection in partially ionized plasmas} & 未获匹配全文\\

E408 / 2017 & \href{http://absimage.aps.org/image/DPP17/MWS_DPP17-2017-000073.pdf}{Tearing mode dynamics and sawtooth oscillation in Hall-MHD} & 未获匹配全文\\

F790 / 2017 & \href{https://doi.org/10.13016/m24q7qq78}{The Effects of Turbulence on Magnetic Reconnection at the Magnetopause} & 有PDF，待核读\\

E381 / 2017 & \href{https://ui.adsabs.harvard.edu/abs/2017shin.confE..42C/abstract}{Tripolar Hall Magnetic Field Signatures in Magnetic Reconnection in the Solar Wind and Magnetopause} & 未获匹配全文\\

F608 / 2017 & \href{https://doi.org/10.1063/1.4975086}{Turbulent transport in 2D collisionless guide field reconnection} & 有PDF，待核读\\

E441 / 2017 & \href{https://ui.adsabs.harvard.edu/abs/2017AGUFMSM11C2328N/abstract}{Two-fluid studies including pressure tensor effects of current sheet instabilities and reconnection in the magnetotail} & 未获匹配全文\\

L242 / 2017 & \href{https://doi.org/10.1063/1.4986116}{Two-fluid tearing mode instability in cylindrical geometry} & 未获匹配全文\\

F612 / 2018 & \href{https://doi.org/10.1002/2017gl076330}{An Electron‐Scale Current Sheet Without Bursty Reconnection Signatures Observed in the Near‐Earth Tail} & 未获匹配全文\\

L245 / 2018 & \href{https://doi.org/10.1063/1.5016345}{An intuitive two-fluid picture of spontaneous 2D collisionless magnetic reconnection and whistler wave generation} & 有PDF，待核读\\

E445 / 2018 & \href{https://ui.adsabs.harvard.edu/abs/2018cosp...42E3772Y/abstract}{An intuitive two-fluid picture of spontaneous collisionless magnetic reconnection and whistler wave generation involving frozen-in canonical vorticity} & 未获匹配全文\\

L246 / 2018 & \href{https://doi.org/10.1088/1742-6596/1031/1/012020}{Fast magnetic reconnection: The "ideal" tearing instability in classic,  Hall,  and relativistic plasmas} & 有PDF，待核读\\

F715 / 2018 & \href{https://doi.org/10.1063/1.5021492}{General connected and reconnected fields in plasmas} & 未获匹配全文\\

E369 / 2018 & \href{https://ui.adsabs.harvard.edu/abs/2018AGUFMSM13B2865L/abstract}{Hall electric field in pre-reconnection magnetotail thin current sheet} & 未获匹配全文\\

F793 / 2018 & \href{https://doi.org/10.1002/2017ja025019}{MMS Observation of Asymmetric Reconnection Supported by 3‐D Electron Pressure Divergence} & 有PDF，待核读\\

F675 / 2018 & \href{https://doi.org/10.3847/1538-4357/aa9edb}{Magnetic Reconnection in Strongly Magnetized Regions of the Low Solar Chromosphere} & 有PDF，待核读\\

F671 / 2018 & \href{https://doi.org/10.1038/s41467-018-07415-3}{Magnetic reconnection driven by electron dynamics} & 有PDF，待核读\\

L248 / 2018 & \href{http://paduaresearch.cab.unipd.it/10907/1/tesi_Luigi_Cordaro.pdf}{Magnetic reconnection in fusion plasmas} & 未获匹配全文\\

F624 / 2018 & \href{https://doi.org/10.1029/2018ja025626}{North‐South Asymmetric Kelvin‐Helmholtz Instability and Induced Reconnection at the Earth's Magnetospheric Flanks} & 有PDF，待核读\\

L249 / 2018 & \href{https://doi.org/10.1063/1.5009389}{Numerical analysis of two-fluid tearing mode instability in a finite aspect ratio cylinder} & 未获匹配全文\\

E374 / 2018 & \href{http://ui.adsabs.harvard.edu/abs/2018shin.confE.240S/abstract}{Onset and nonlinear evolution of fast reconnection: Lundquist number and Hall effects} & 未获匹配全文\\

L252 / 2018 & \href{https://doi.org/10.1063/1.5030439}{Quadrupolar and hexapolar Hall magnetic field during asymmetric magnetic reconnection without a guide field} & 未获匹配全文\\

E344 / 2018 & \href{http://ui.adsabs.harvard.edu/abs/2018shin.confE.188M/abstract}{Recombination and the Hall Effect in Simulations of Chromospheric Magnetic Reconnection} & 未获匹配全文\\

L253 / 2018 & \href{https://doi.org/10.1029/2017ja024729}{Reconnection in the Martian Magnetotail: Hall‐ MHD With Embedded Particle‐in‐Cell Simulations} & 未获匹配全文\\

L254 / 2018 & \href{https://doi.org/10.1029/2018ja026051}{Reconstruction of the Electron Diffusion Region of Magnetotail Reconnection Seen by the MMS Spacecraft on 11 July 2017} & 未获匹配全文\\

L255 / 2018 & \href{https://doi.org/10.1029/2018ja025251}{Structures of Hall Fields in Asymmetric Magnetic Reconnection} & 未获匹配全文\\

E487 / 2018 & \href{https://doi.org/10.1063/1.5037774}{Study of magnetic reconnection in large-scale magnetic island coalescence via spatially coupled MHD and PIC simulations} & 有PDF，待核读\\

E452 / 2018 & \href{https://meetings.aps.org/Meeting/DPP18/Session/CO5.1}{Two-fluid edge-plasma SOL turbulent transport driven by interchange/tearing instabilities} & 未获匹配全文\\

L258 / 2018 & \href{https://doi.org/10.1029/2018gl078614}{Upward Frequency Sweeping Due to the Nonlinear Resonance of Multiwave Coupling and Its Significant Contribution to the Hall Effect at the Reconnection Point} & 未获匹配全文\\

L259 / 2019 & \href{https://doi.org/10.3847/1538-4357/ab14ef}{A Parametric Study of the Structure of Hall Magnetic Field Based on Kinetic Simulations. I. Anti-parallel Magnetic Reconnection in an Asymmetric Current Sheet} & 未获匹配全文\\

L260 / 2019 & \href{https://doi.org/10.3847/1538-4357/ab35d9}{A Parametric Study of the Structure of Hall Magnetic Field Based on Kinetic Simulations. II. Asymmetric Magnetic Reconnection with a Guide Field} & 未获匹配全文\\

F694 / 2019 & \href{https://doi.org/10.1063/1.5086170}{Asymmetric and endogenous magnetic reconnection processes in quasi-collisionless plasmas} & 有PDF，待核读\\

F774 / 2019 & \href{https://doi.org/10.3389/fspas.2019.00070}{Collisionless Magnetic Reconnection and Waves: Progress Review} & 有PDF，待核读\\

F658 / 2019 & \href{https://doi.org/10.1029/2018ja026451}{Dissipation of Earthward Propagating Flux Rope Through Re‐reconnection with Geomagnetic Field: An MMS Case Study} & 有PDF，待核读\\

E449 / 2019 & \href{http://ui.adsabs.harvard.edu/abs/2019APS..DPPP10083G/abstract}{Effect of Two-Fluid Equilibrium Flow on Tearing Linear Stability} & 未获匹配全文\\

L261 / 2019 & \href{https://doi.org/10.1088/0256-307x/36/1/015201}{Electron Magnetohydrodynamics Magnetic Reconnection Experiment on Keda Linear Magnetized Plasma Device} & 未获匹配全文\\

E548 / 2019 & \href{https://dataspace.princeton.edu/handle/88435/dsp015d86p299f}{Fluid closures for the modelling of reconnection and instabilities in magnetotail current sheets} & 未获匹配全文\\

F682 / 2019 & \href{https://doi.org/10.1103/physrevd.99.063017}{Gravitational electromotive force in magnetic reconnection around Schwarzschild black holes} & 有PDF，待核读\\

F632 / 2019 & \href{https://doi.org/10.1063/1.5109368}{Inner and outer electron diffusion region of antiparallel collisionless reconnection: Density dependence} & 有PDF，待核读\\

F645 / 2019 & \href{https://doi.org/10.1029/2019gl083782}{Magnetospheric MultiscaleMission Observations of Reconnecting Electric Fields in the Magnetotail on Kinetic Scales} & 有PDF，待核读\\

F623 / 2019 & \href{https://doi.org/10.3847/2041-8213/ab4b5a}{Mechanism of Reconnection on Kinetic Scales Based on Magnetospheric Multiscale Mission Observations} & 未获匹配全文\\

F754 / 2019 & \href{https://doi.org/10.1063/1.5086402}{Numerical study of tearing mode seeding in tokamak X-point plasma} & 有PDF，待核读\\

F726 / 2019 & \href{https://doi.org/10.3847/1538-4357/ab29f6}{Relativistic Tearing Mode in Pair Plasmas and Application to Magnetic Giant Flares} & 未获匹配全文\\

F693 / 2019 & \href{https://doi.org/10.1029/2019gl083282}{Satellite Data‐Based 3‐D Simulation of Kelvin‐Helmholtz Instability and Induced Magnetic Reconnection at the Earth's Magnetopause} & 有PDF，待核读\\

E533 / 2019 & \href{https://doi.org/10.1063/1.5122225}{The electron canonical battery effect in magnetic reconnection: Completion of the electron canonical vorticity framework} & 有PDF，待核读\\

F638 / 2020 & \href{https://doi.org/10.1029/2019ja027410}{Comparative Analysis of the Vlasiator Simulations and MMS Observations of Multiple X‐Line Reconnection and Flux Transfer Events} & 未获匹配全文\\

L266 / 2020 & \href{https://doi.org/10.1088/1741-4326/ab9564}{Effect of bootstrap current on the onset of neoclassical tearing mode} & 未获匹配全文\\

F723 / 2020 & \href{https://doi.org/10.1063/5.0015240}{Grad–Shafranov reconstruction of the magnetic configuration in the reconnection X-point vicinity in compressible plasma} & 未获匹配全文\\

F747 / 2020 & \href{https://doi.org/10.1029/2019ja027490}{MESSENGER Observations of Mercury's Nightside Magnetosphere Under Extreme Solar Wind Conditions: Reconnection‐Generated Structures and Steady Convection} & 未获匹配全文\\

E482 / 2020 & \href{https://escholarship.org/uc/item/0nv873w9}{Magnetic Reconnection and Turbulence in the Inner Heliosphere} & 未获匹配全文\\

E539 / 2020 & \href{https://doi.org/10.1088/1361-6587/abcab3}{Magnetic reconnection in 3D fusion devices: non-linear reduced equations and linear current-driven instabilities} & 未获匹配全文\\

F570 / 2020 & \href{https://doi.org/10.1002/essoar.10504213.1}{Magnetotail Reconnection Asymmetries in a Small,  Earth-Like Magnetosphere} & 未获匹配全文\\

F707 / 2020 & \href{https://doi.org/10.48550/arxiv.2009.08779}{Major Scientific Challenges and Opportunities in Understanding Magnetic Reconnection and Related Explosive Phenomena in Solar and Heliospheric Plasmas} & 有PDF，待核读\\

E540 / 2020 & \href{https://doi.org/10.1063/5.0022133}{Multiparametric study of tearing modes in thin current sheets} & 有PDF，待核读\\

F657 / 2020 & \href{https://doi.org/10.1029/2020ja027985}{Multiscale Coupling During Magnetopause Reconnection: Interface Between the Electron and Ion Diffusion Regions} & 有PDF，待核读\\

F752 / 2020 & \href{https://doi.org/10.3847/2041-8213/ab8e31}{Observation of Energy Conversion Near the X-line in Asymmetric Guide-field Reconnection} & 未获匹配全文\\

F622 / 2020 & \href{https://doi.org/10.1029/2020gl088761}{Physical Implication of Two Types of Reconnection Electron Diffusion Regions With and Without Ion‐Coupling in the Magnetotail Current Sheet} & 有PDF，待核读\\

F580 / 2020 & \href{https://doi.org/10.1063/1.5128376}{Reconnection from a turbulence perspective} & 有PDF，待核读\\

E323 / 2020 & \href{http://absimage.aps.org/image/DPP20/MWS_DPP20-2020-001941.pdf}{Reconnection in Hall-MHD} & 未获匹配全文\\

E389 / 2020 & \href{https://doi.org/10.1029/2020ja028162}{Reconnection‐Driven Dynamics at Ganymede's Upstream Magnetosphere: 3‐D Global Hall MHD and MHD‐EPIC Simulations} & 未获匹配全文\\

E551 / 2020 & \href{https://doi.org/10.1029/2020ja027879}{Reconstruction of Electron and Ion Distribution Functions in a Magnetotail Reconnection Diffusion Region} & 有PDF，待核读\\

E463 / 2020 & \href{https://doi.org/10.1063/5.0022332}{Tearing mode analysis in electron magnetohydrodynamics with pressure gradient} & 未获匹配全文\\

F805 / 2020 & \href{https://etheses.whiterose.ac.uk/28202/1/Geng_PhD_Thesis.pdf}{The physics of a collisionless micro-scale tearing mode} & 有PDF，待核读\\

E350 / 2020 & \href{http://absimage.aps.org/image/DPP20/MWS_DPP20-2020-001547.pdf}{Three-dimensional plasmoid-mediated reconnection in Hall magnetohydrodynamics} & 未获匹配全文\\

L269 / 2020 & \href{https://doi.org/10.3847/1538-4357/ab7551}{Using a Higher-order Numerical Scheme to Study the Hall Magnetic Reconnection} & 未获匹配全文\\

F651 / 2021 & \href{https://doi.org/10.48550/arxiv.2112.03316}{A novel kinetic mechanism for the onset of fast magnetic reconnection and plasmoid instabilities in collisionless plasmas} & 有PDF，待核读\\

L270 / 2021 & \href{https://doi.org/10.1029/2021ja029550}{Electrodynamic Contributions to the Hall‐ and Parallel Electric Fields in Collisionless Magnetic Reconnection} & 未获匹配全文\\

L271 / 2021 & \href{https://doi.org/10.1103/physreve.104.065206}{Energy transfer in reconnection and turbulence} & 有PDF，待核读\\

E411 / 2021 & \href{https://doi.org/10.1017/s0022377821000726}{Hall instability: origin,  properties and asymptotic theory for its tearing mode} & 有PDF，待核读\\

F760 / 2021 & \href{https://doi.org/10.1063/5.0037227}{Impact of electron temperature anisotropy on the collisionless tearing mode instability in the presence of a strong guide field} & 有PDF，待核读\\

L273 / 2021 & \href{https://doi.org/10.1088/1741-4326/ac3365}{Influence of Hall effect and toroidal flow on the plasmoid formation and incomplete reconnection in a low resistivity plasma in tokamak} & 未获匹配全文\\

F700 / 2021 & \href{https://doi.org/10.1088/1361-6587/ac43f0}{Kelvin-Helmholtz instability and induced magnetic reconnection at the Earth’s magnetopause: a 3D simulation based on satellite data} & 有PDF，待核读\\

L274 / 2021 & \href{https://doi.org/10.3847/1538-4357/ac31b1}{MMS Observation of the Hall Field in an Asymmetric Magnetic Reconnection with Guide Field} & 未获匹配全文\\

F643 / 2021 & \href{https://doi.org/10.3389/fspas.2021.760885}{MMS Observations of Double Mid-Latitude Reconnection Ion Beams in the Early Non-Linear Phase of the Kelvin-Helmholtz Instability} & 有PDF，待核读\\

F698 / 2021 & \href{https://doi.org/10.1029/2021ja029415}{Magnetic Reconnection Within the Boundary Layer of a Magnetic Cloud in the Solar Wind} & 有PDF，待核读\\

F595 / 2021 & \href{https://doi.org/10.48550/arxiv.2104.12013}{Magnetic reconnection as an energy cascade process} & 重复：L271\\

E492 / 2021 & \href{https://doi.org/10.5194/angeo-39-991-2021}{Magnetotail reconnection asymmetries in an ion-scale,  Earth-like magnetosphere} & 有PDF，待核读\\

F576 / 2021 & \href{https://doi.org/10.1088/2058-6272/abe038}{Nonlinear evolution and secondary island formation of the double tearing mode in a hybrid simulation} & 未获匹配全文\\

E465 / 2021 & \href{https://doi.org/10.1063/5.0066269}{Numerical studies on electron magnetohydrodynamics tearing mode instability} & 未获匹配全文\\

E496 / 2021 & \href{https://doi.org/10.1002/essoar.10508377.1}{Reconnection as an Energy Cascade} & 未获匹配全文\\

F593 / 2021 & \href{https://doi.org/10.1103/physrevx.11.041005}{Reconnection-Controlled Decay of Magnetohydrodynamic Turbulence and the Role of Invariants} & 有PDF，待核读\\

F628 / 2021 & \href{https://doi.org/10.1029/2021gl094201}{Shear Alfvén Waves Driven by Magnetic Reconnection as an Energy Source for the Aurora Borealis} & 有PDF，待核读\\

F656 / 2021 & \href{https://doi.org/10.3847/1538-4357/abcf29}{The Formation of Electron Outflow Jets with Power-law Energy Distribution in Guide-field Magnetic Reconnection} & 未获匹配全文\\

F701 / 2021 & \href{https://doi.org/10.1088/1741-4326/ac077d}{W7-X and the sawtooth instability: towards realistic simulations of current-driven magnetic reconnection} & 有PDF，待核读\\

L276 / 2022 & \href{https://doi.org/10.2307/j.ctv1w7v22j.15}{12 Cross-discipline study of the two-fluid dynamics of magnetic reconnection in laboratory and magnetopause plasmas} & 未获匹配全文\\

E432 / 2022 & \href{https://doi.org/10.2307/j.ctv1w7v22j.8}{5 Development of two-fluid theory for reconnection coordinated with key observations} & 未获匹配全文\\

F796 / 2022 & \href{https://doi.org/10.3847/1538-4357/ac7da6}{Anisotropic Electron Heating in Turbulence-driven Magnetic Reconnection in the Near-Sun Solar Wind} & 有PDF，待核读\\

F712 / 2022 & \href{https://doi.org/10.3847/1538-4357/ac8667}{Energy Transport during 3D Small-scale Reconnection Driven by Anisotropic Plasma Turbulence} & 有PDF，待核读\\

L277 / 2022 & \href{https://doi.org/10.3847/1538-4357/aca27f}{Evidence of the Hall Effect Providing the Core Field inside a Magnetic Flux Rope in an Antiparallel Magnetic Reconnection} & 未获匹配全文\\

F740 / 2022 & \href{https://doi.org/10.1088/1741-4326/ac984f}{Fast magnetic reconnection and driven plasma rotation in reversed central magnetic shear configuration} & 未获匹配全文\\

F746 / 2022 & \href{https://doi.org/10.1063/5.0082331}{Influence of anomalous perpendicular transport on linear tearing mode dynamics in tokamak plasmas} & 有PDF，待核读\\

F614 / 2022 & \href{https://doi.org/10.1007/s41614-022-00107-y}{Kinetic Alfven wave (KAW) eigenmode in magnetosphere magnetic reconnection} & 有PDF，待核读\\

L278 / 2022 & \href{https://doi.org/10.1088/1741-4326/ac6617}{Magnetic reconnection rate during sawtooth crashes in ASDEX Upgrade} & 未获匹配全文\\

E537 / 2022 & \href{https://doi.org/10.1017/s0022377822001088}{Microscopic scales of linear tearing modes: a tutorial on boundary layer theory for magnetic reconnection} & 有PDF，待核读\\

L279 / 2022 & \href{https://doi.org/10.33915/etd.11608}{Scaling Theory of 3D Magnetic Reconnection X-Line Spreading} & 未获匹配全文\\

L280 / 2022 & \href{https://doi.org/10.1029/2022gl099638}{Simulation of Magnetospheric Sawtooth Oscillations: The Role of Kinetic Reconnection in the Magnetotail} & 未获匹配全文\\

L281 / 2022 & \href{https://doi.org/10.1088/1741-4326/ac8f59}{Stabilization of double tearing mode growth by resonant magnetic perturbations} & 未获匹配全文\\

L282 / 2022 & \href{https://doi.org/10.1029/2022ja030686}{The Helicity Sign of Flux Transfer Event Flux Ropes and Its Relationship to the Guide Field and Hall Physics in Magnetic Reconnection at the Magnetopause} & 有PDF，待核读\\

E461 / 2023 & \href{https://doi.org/10.1063/5.0155211}{Asymptotic scalings of fluid,  incompressible “electron-only” reconnection instabilities: Electron-magnetohydrodynamics tearing modes} & 有PDF，待核读\\

E534 / 2023 & \href{https://doi.org/10.1029/2022ja031209}{Characteristics of Magnetic Reconnection With a Finite X‐Line Length} & 未获匹配全文\\

F620 / 2023 & \href{https://doi.org/10.1029/2023gl104733}{Does Magnetic Reconnection Occur in the Near Lunar Surface Environment?} & 有PDF，待核读\\

L284 / 2023 & \href{https://doi.org/10.1063/5.0150929}{Effect of a guide field on the turbulence like properties of magnetic reconnection} & 有PDF，待核读\\

L285 / 2023 & \href{https://doi.org/10.3389/fspas.2023.1069888}{Electron magnetohydrodynamics Grad–Shafranov reconstruction of the magnetic reconnection electron diffusion region} & 有PDF，待核读\\

F636 / 2023 & \href{https://doi.org/10.48550/arxiv.2305.05701}{Experimental mitigation of fast magnetic reconnection in multiple interacting laser-produced plasmas} & 有PDF，待核读\\

L286 / 2023 & \href{https://doi.org/10.3847/1538-4357/aca760}{Hall Fields and Current Systems of Magnetic Reconnection under Asymmetric Conditions} & 未获匹配全文\\

L287 / 2023 & \href{https://doi.org/10.1088/1402-4896/acd3bb}{Hall effect on the magnetic reconnections during the evolution of a three-dimensional magnetic flux rope} & 有PDF，待核读\\

L288 / 2023 & \href{https://noesis.uis.edu.co/handle/20.500.14071/12490}{Impact of Hall effect and ambipolar diffusion on the magnetic reconnection rate and energy distribution in the Earth’s magnetotail} & 未获匹配全文\\

E549 / 2023 & \href{https://doi.org/10.1029/2022ja031134}{MESSENGER Observations of Reconnection in Mercury's Magnetotail Under Strong IMF Forcing} & 未获匹配全文\\

L289 / 2023 & \href{https://doi.org/10.1063/5.0169243}{Magnetic reconnection during sawteeth crashes} & 未获匹配全文\\

F641 / 2023 & \href{https://doi.org/10.3847/25c2cfeb.e22a8d1f}{Major Scientific Challenges and Opportunities in Understanding Magnetic Reconnection and Related Explosive Phenomena in Heliophysics and Beyond} & 未获匹配全文\\

L290 / 2023 & \href{https://doi.org/10.3847/1538-4357/acbfb1}{Multifluid Simulations of Upper-chromospheric Magnetic Reconnection with Helium–Hydrogen Mixture} & 有PDF，待核读\\

F753 / 2023 & \href{https://doi.org/10.1063/5.0106785}{Non-ideal MHD growth of current interchange tearing modes at plasma edge and response to externally-imposed flow} & 有PDF，待核读\\

F640 / 2023 & \href{https://doi.org/10.1063/5.0159182}{Numerical investigation of the 2/1 double tearing mode in EAST with the CLT code} & 未获匹配全文\\

L291 / 2023 & \href{https://doi.org/10.3847/1538-4357/acbdf6}{Observation of the Hall Magnetic Reconnection As Close As 56 Solar Radii from the Sun} & 未获匹配全文\\

E530 / 2023 & \href{https://doi.org/10.1063/5.0174167}{Post-disruption reconnection event driven by a runaway current} & 有PDF，待核读\\

L292 / 2023 & \href{https://doi.org/10.48550/arxiv.2307.13268}{Preferential acceleration of heavy ions in magnetic reconnection: Hybrid-kinetic simulations with electron inertia} & 有PDF，待核读\\

L293 / 2023 & \href{https://doi.org/10.1063/5.0137123}{Role of edge-localized neoclassical tearing modes in quiescent H-mode plasmas in the DIII-D tokamak} & 有PDF，待核读\\

F592 / 2023 & \href{https://doi.org/10.1103/physrevlett.130.205201}{Subion-Scale Turbulence Driven by Magnetic Reconnection} & 有PDF，待核读\\

L295 / 2023 & \href{https://doi.org/10.1051/0004-6361/202346659}{Two-fluid reconnection jets in a gravitationally stratified atmosphere} & 有PDF，待核读\\

F699 / 2023 & \href{https://doi.org/10.48550/arxiv.2307.05736}{Two-fluid reconnection jets in a gravitationally stratified atmosphere} & 重复：L295\\

E466 / 2024 & \href{https://doi.org/10.1063/5.0205061}{A numerical study of electron-magnetohydrodynamics tearing modes in parameter ranges of experimental interest} & 有PDF，待核读\\

F733 / 2024 & \href{https://doi.org/10.1063/5.0226026}{An examination of nonlinear collisionless magnetic reconnection through eigenmode decomposition} & 未获匹配全文\\

F716 / 2024 & \href{https://doi.org/10.3389/fphy.2024.1439949}{Analytical solution of steady reconnection outflows in a time-varying three-dimensional reconnection model with generalized spatiotemporal distributions} & 有PDF，待核读\\

F629 / 2024 & \href{https://doi.org/10.1063/5.0191674}{Comparative simulations of Kelvin–Helmholtz induced magnetic reconnection at the Earth's magnetospheric flanks} & 有PDF，待核读\\

F633 / 2024 & \href{https://doi.org/10.1029/2024gl108270}{Counter‐Helical Magnetic Flux Ropes From Magnetic Reconnections in Space Plasmas} & 未获匹配全文\\

F784 / 2024 & \href{https://doi.org/10.3847/1538-4357/ad7f54}{Current Sheet Broadening due to Turbulence in Three-dimensional Collisionless Reconnection} & 未获匹配全文\\

F654 / 2024 & \href{https://doi.org/10.1186/s40623-024-02003-w}{Duskward displacement of plasmoids and reconnection in the near-Earth magnetotail} & 有PDF，待核读\\

F771 / 2024 & \href{https://doi.org/10.3847/1538-4357/ad6a19}{Electron-only Reconnection and Ion Heating in 3D3V Hybrid-Vlasov Plasma Turbulence} & 有PDF，待核读\\

L297 / 2024 & \href{https://doi.org/10.1016/j.fpp.2024.100055}{Gyrofluid simulations of turbulence and reconnection in space plasmas} & 有PDF，待核读\\

L298 / 2024 & \href{https://doi.org/10.1088/1361-6587/ad2b8f}{Hall physics during magnetic reconnection with collision effect} & 未获匹配全文\\

E546 / 2024 & \href{https://doi.org/10.1063/5.0185897}{Influence of ion-to-electron temperature ratio on tearing instability and resulting subion-scale turbulence in a low- β e collisionless plasma} & 有PDF，待核读\\

E467 / 2024 & \href{https://doi.org/10.1088/2058-6272/ada421}{Instabilities of ideal magnetohydrodynamics mode and neoclassical tearing mode stabilization by electron cyclotron current drive for EHL-2 spherical torus} & 未获匹配全文\\

F607 / 2024 & \href{https://doi.org/10.3847/2041-8213/ad2774}{Intermittent Electron-only Reconnection at Lunar Mini-magnetospheres} & 有PDF，待核读\\

F597 / 2024 & \href{https://doi.org/10.1029/2024ja032669}{Kinetic Signatures,  Dawn‐Dusk Asymmetries,  and Flux Transfer Events Associated With Mercury's Dayside Magnetopause Reconnection From 3D MHD‐AEPIC Simulations} & 未获匹配全文\\

F794 / 2024 & \href{https://doi.org/10.3847/1538-4357/ad4fff}{Laboratory Study of Magnetic Reconnection in Lunar-relevant Mini-magnetospheres} & 有PDF，待核读\\

F711 / 2024 & \href{https://doi.org/10.1063/5.0169397}{Localization of beam generated whistler wave and turbulence generation in reconnection region of magnetopause} & 未获匹配全文\\

F799 / 2024 & \href{https://doi.org/10.1063/5.0229910}{Neural network based fast prediction of double tearing modes in advanced tokamak plasmas} & 未获匹配全文\\

L299 / 2024 & \href{https://doi.org/10.1063/5.0213566}{Numerical analysis of three-dimensional magnetopause-like reconnection properties by Hall MHD simulation for SPERF-AREX} & 未获匹配全文\\

E498 / 2024 & \href{https://doi.org/10.1088/1361-6587/ad4e66}{Numerical study of the influence of electron inertial effects and electron dynamics on tearing mode instability} & 未获匹配全文\\

E535 / 2024 & \href{https://doi.org/10.1063/5.0186879}{Numerical study of the suppression of magnetic reconnection onset with injected plasma} & 未获匹配全文\\

F772 / 2024 & \href{https://doi.org/10.1051/0004-6361/202348573}{Preferential acceleration of heavy ions in magnetic reconnection: Hybrid-kinetic simulations with electron inertia} & 未获匹配全文\\

L300 / 2024 & \href{https://doi.org/10.1088/1361-6587/ada1f9}{Stability of the experimental and numerical spheromaks with shear toroidal flow induced by magnetic reconnection} & 未获匹配全文\\

F615 / 2024 & \href{https://doi.org/10.58088/x3fm-h327}{Structure of the diffusion region in three dimensional magnetic reconnection} & 有PDF，待核读\\

L301 / 2024 & \href{https://doi.org/10.1016/j.physleta.2024.129897}{Systematic 2.5D resistive MHD simulations with ambipolar diffusion and Hall effect for fast magnetic reconnection} & 有PDF，待核读\\

E470 / 2025 & \href{https://resolver.obvsg.at/urn:nbn:at:at-ubi:1-175860}{2D Full-F Gyrofluid Magnetic Reconnection} & 未获匹配全文\\

E536 / 2025 & \href{https://doi.org/10.1063/5.0243286}{Comparative study of magnetic reconnection rates in ASDEX Upgrade and EAST tokamaks} & 未获匹配全文\\

F690 / 2025 & \href{https://doi.org/10.3847/1538-4357/ae03c5}{Determining the 3D Dynamics of Solar Flare Magnetic Reconnection} & 有PDF，待核读\\

F621 / 2025 & \href{https://doi.org/10.3847/1538-4357/adfedd}{Effect of Initial Current Sheet and Stratification on Reconnection and Coronal Heating} & 未获匹配全文\\

L303 / 2025 & \href{https://doi.org/10.1029/2025ja034289}{Effective Resistivity for Magnetohydrodynamic Simulation of Collisionless Magnetic Reconnection} & 有PDF，待核读\\

F751 / 2025 & \href{https://doi.org/10.1029/2024gl113104}{Evidence of a Magnetic Null in Electron‐Only Reconnection} & 未获匹配全文\\

L304 / 2025 & \href{https://doi.org/10.3847/1538-4357/ae0ffc}{Fundamental Understanding of Magnetic Reconnection via Spiral CT Scan. V. Ion Demagnetization,  Electron Demagnetization,  and Hall Effect} & 未获匹配全文\\

F802 / 2025 & \href{https://doi.org/10.1016/j.cpc.2025.109807}{GREENY: A full-F 2D gyrofluid reconnection code} & 重复：E473\\

F803 / 2025 & \href{https://doi.org/10.2139/ssrn.5201132}{Greeny: A Full-F 2dgyrofluid Reconnection Code} & 重复：E473\\

L305 / 2025 & \href{https://doi.org/10.1063/5.0259029}{Impact of tearing instability on the global quadrupole Hall magnetic field in an anti-parallel magnetic reconnection} & 未获匹配全文\\

F708 / 2025 & \href{https://doi.org/10.1093/pasj/psaf123}{Influence of kinetic effects in large-scale magnetic reconnection with multi-hierarchy simulation code KAMMUY} & 有PDF，待核读\\

L306 / 2025 & \href{https://doi.org/10.1029/2024gl114151}{Laboratory Observations of Hall Magnetic Field in Electron‐Only Magnetic Reconnection With a Guide Field} & 未获匹配全文\\

F724 / 2025 & \href{https://doi.org/10.1029/2025gl117351}{Laboratory Realization of Magnetopause‐Like Reconnection in Space Plasma Environment Research Facility (SPERF)} & 未获匹配全文\\

L307 / 2025 & \href{https://doi.org/10.1029/2025ja033875}{Magnetic Reconnection Signatures in Near‐Earth Magnetotail Dipolarization Processes: Hall Electrons} & 未获匹配全文\\

F759 / 2025 & \href{https://doi.org/10.1029/2024gl113889}{New Insights on the High Reconnection Rate and the Diminishment of Ion Outflow} & 有PDF，待核读\\

E550 / 2025 & \href{https://doi.org/10.1029/2025gl117612}{Observation of a Knotted Electron Diffusion Region in Earth's Magnetotail Reconnection} & 有PDF，待核读\\

F616 / 2025 & \href{https://doi.org/10.1029/2024sw004057}{Simulating Extreme Space Weather With Kinetic Magnetotail Reconnection} & 未获匹配全文\\

F559 / 2025 & \href{https://doi.org/10.1007/s11214-025-01210-5}{Simulation Models for Exploring Magnetic Reconnection} & 有PDF，待核读\\

F735 / 2025 & \href{https://doi.org/10.1080/10420150.2025.2475389}{Study of driven magnetic reconnection with a multiscale analysis} & 未获匹配全文\\

F627 / 2025 & \href{https://doi.org/10.3847/1538-4357/ae2014}{The Enhancement of Ion Heating in Kinetic,  Antiparallel Reconnection in the Presence of a Flow Shear} & 有PDF，待核读\\

Huang2026 / 2025 & \href{https://doi.org/10.1080/10420150.2025.2598789}{The morphology of large-scale magnetic reconnection: the role of plasmoid instability and Hall effects} & 未获匹配全文\\

F704 / 2025 & \href{https://doi.org/10.1088/1361-6587/adfe19}{Turbulence excited by the electron beam and evolution of current sheets near magnetopause magnetic reconnection regions} & 未获匹配全文\\

F578 / 2025 & \href{https://doi.org/10.3847/1538-4357/ae046e}{Unveiling a New β Scaling of the Tearing Instability in Weakly Collisional Plasmas} & 有PDF，待核读\\

F639 / 2026 & \href{https://doi.org/10.1063/5.0323293}{A numerical experiment on oscillatory magnetic reconnection in a laboratory plasma system driven by alternating currents} & 有PDF，待核读\\

L309 / 2026 & \href{https://arxiv.org/abs/2607.14286}{A structure-preserving Numerical Method for the Compressible Resistive-Hall-MHD System} & 有PDF，待核读\\

Teh2026 / 2026 & \href{https://doi.org/10.1029/2026ja035335}{Asymmetric Hall Magnetic Field Structures in Oppositely Directed Jets During Guide‐Field Magnetopause Reconnection} & 未获匹配全文\\

F804 / 2026 & \href{https://dspace.mit.edu/handle/1721.1/169938}{Collisionless Plasma Physics below the Ion Scales: Instabilities,  Reconnection,  and Turbulence} & 未获匹配全文\\

F696 / 2026 & \href{https://doi.org/10.1088/2058-6272/ae78b7}{Effects of internal current-sheet structure on magnetic reconnection evolution} & 未获匹配全文\\

F744 / 2026 & \href{https://doi.org/10.1063/5.0347826}{Ion dynamics in the ion diffusion region of symmetric magnetic reconnection} & 未获匹配全文\\

F718 / 2026 & \href{https://doi.org/10.1103/9l36-74my}{Laboratory Observation of Transition from Collisional Slow to Collisionless Fast Reconnection} & 未获匹配全文\\

L315 / 2026 & \href{https://doi.org/10.1007/s11207-026-02714-1}{Magnetic Reconnection Process in Partially Ionized Fluids: Insights from the Solar Chromosphere} & 有PDF，待核读\\

F691 / 2026 & \href{https://doi.org/10.21203/rs.3.rs-8443380/v1}{Magnetic Reconnections Stochastic-Adaptive Control for Plasma Fusion} & 有PDF，待核读\\

Arshad2026 / 2026 & \href{https://doi.org/10.1029/2026ja035547}{Modeling Kinetic Reconnection Onset in Global Magnetospheric Simulations} & 未获匹配全文\\

L317 / 2026 & \href{https://doi.org/10.1088/1741-4326/ae4e49}{Numerical modeling of linear and nonlinear evolution of tearing modes and resistive-kink modes with runaway electrons} & 未获匹配全文\\

L318 / 2026 & \href{https://doi.org/10.13140/rg.2.2.22493.17125}{Plasmoid-Mediated 2D Magnetic Reconnection in Partially Ionized Plasmas} & 有PDF，待核读\\

F617 / 2026 & \href{https://doi.org/10.1007/s11038-026-09597-z}{Recent advances in magnetic reconnection in Mars’ space environment} & 有PDF，待核读\\

L319 / 2026 & \href{https://doi.org/10.48550/arxiv.2610.23494}{Reconnection in 2D Inertial Magnetohydrodynamics} & 有PDF，待核读\\

F755 / 2026 & \href{https://doi.org/10.1088/1741-4326/ae3332}{Synergetic effects of E × B shear flow and diamagnetic flow on boundary-perturbation-driven neoclassical tearing modes} & 未获匹配全文\\

F761 / 2026 & \href{https://doi.org/10.1063/5.0336253}{Tearing and Kelvin–Helmholtz dynamics in fully kinetic particle-in-cell simulations of electron-scale current sheets} & 有PDF，待核读\\

L320 / 2026 & \href{https://doi.org/10.1088/1361-6587/ae7156}{Tearing stability prediction combining toroidal calculations with a two-fluid slab layer} & 有PDF，待核读\\

E555 / 2026 & \href{https://doi.org/10.1088/1402-4896/ae52bc}{The behavior of electron flow during asymmetric magnetic reconnection} & 未获匹配全文\\

L321 / 2026 & \href{https://doi.org/10.1140/epjd/s10053-025-01116-5}{The impact of localized out-of-plane antisymmetric flows and Hall effect on the magnetic reconnection in a compressible plasma} & 未获匹配全文\\

E554 / 2026 & \href{https://doi.org/10.3847/1538-4357/ae91d4}{Transition of Magnetic Reconnection Regimes in Partially Ionized Plasmas} & 有PDF，待核读\\

\bottomrule\end{longtable}\normalsize